\documentclass[10pt, letterpaper]{article}
\usepackage[nomarginpar, margin=2.75cm, tmargin=3cm, bmargin=1.75cm]{geometry}
\usepackage[        
    colorlinks=true,      
    linkcolor=black,      
    urlcolor=blue,       
    citecolor=black       
]{hyperref}
\usepackage[italian]{babel}
\usepackage{template}
\usepackage{float}
\usepackage{graphicx}
\usepackage{caption}
\usepackage{lastpage} 
\renewcommand{\arraystretch}{1.05} % migliora la leggibilità
\renewcommand{\contentsname}{Indice}
\renewcommand{\listfigurename}{Elenco dei grafici}
\renewcommand{\listtablename}{Elenco delle tabelle}
\usepackage{fancyhdr}
%Comandi per livello di sottosezioni = 3
\setcounter{tocdepth}{4}
\setcounter{secnumdepth}{4}
\newcommand{\trisubsection}[1]{\paragraph{#1}\mbox{}\\}

\pagestyle{fancy}
\fancyhf{}
\fancyhead[L]{SnakeByte}
\fancyhead[R]{Piano di qualifica}
\fancyfoot[C]{Pagina \thepage\ di \pageref{LastPage}}

\begin{document}
\setcounter{tocdepth}{5}
\setcounter{secnumdepth}{5}

\begin{titlepage}
    \begin{center}
        \begin{center}
            
\includegraphics[width=0.6\textwidth]{./img/logo.pdf}
        \end{center}
        \vspace{4cm}
        \huge\textbf{Piano di qualifica}\par
        \vspace{2cm}
        \large \textbf{SnakeByte} (Gruppo 1):\\
        \large Valeria Baleanu, Leonardo Pellizzon, Filippo Venzo, Giuseppe De Fina, \\
        Francesco Pasqual, Christian Libralato, Luca Granziero \\
        (2109911, 2111006, 2113705, 2113187, 2103119, 2101047, 2075512)
        \vfill
        \small
        \begin{center}
            \begin{tabular}{|c|c|c|c|}
                \hline
                \multicolumn{4}{|c|}{\textbf{Informazioni documento}}                                          \\
                \hline
                \rowcolor{lightgray} \textbf{Versione} & \textbf{Data} & \textbf{Stato} & \textbf{Destinatari} \\
                \hline
                1.0.3                                  & 01/04/2026    & Verificato     &
                \begin{tabular}[c]{@{}l@{}}
                    \textbf{Interni:} SnakeByte \\
                    \textbf{Esterni:} prof. Vardanega Tullio, prof. Cardin Riccardo, Vimar
                \end{tabular}                      \\
                \hline
            \end{tabular}
        \end{center}
        \vfill
        \large Contatti: snakebyteteam@gmail.com
    \end{center}
\end{titlepage}

\begin{center}
    \begin{tabularx}{\textwidth}{|c|c|c|c|c|X|}
        \hline
        \multicolumn{6}{|c|}{\textbf{Registro delle modifiche}}                                                                                                                                                                    \\
        \hline
        \rowcolor{lightgray} \textbf{Versione} & \textbf{Data} & \textbf{Autore} & \textbf{Verificatore}    & \textbf{Approvatore} & \textbf{Descrizione}                                                                          \\
        \hline 1.0.3                           & 01/04/2026    & L. Granziero    & G. De Fina               & -                    & Aggiornamento allo sprint 10                                                                  \\
        \hline
        1.0.2                                  & 13/03/2026    & F. Venzo        & C. Libralato             & -                    & Aggiornamento allo sprint 9                                                                   \\
        \hline
        1.0.1                                  & 8/02/2026     & L. Pellizzon    & F. Pasqual               & -                    & Aggiornamento allo sprint 8                                                                   \\
        \hline
        1.0.0                                  & 13/02/2026    & -               & -                        & F. Venzo             & Approvazione                                                                                  \\
        \hline
        0.2.2                                  & 10/02/2026    & F. Venzo        & F. Pasqual, L. Granziero & -                    & Aggiunto elenco delle tabelle                                                                 \\
        \hline
        0.2.1                                  & 04/02/2026    & F. Venzo        & V. Baleanu               & -                    & Aggiunta descrizione ai grafici, campagna di test di unità e campagna di test di integrazione \\
        \hline
        0.2.0                                  & 28/12/2025    & V. Baleanu      & F. Pasqual               & -                    & Aggiunta tabelle delle soglie per le metriche di processo e di prodotto                       \\
        \hline
        0.1.1                                  & 27/12/2025    & V. Baleanu      & F. Pasqual               & -                    & Aggiunta Introduzione                                                                         \\
        \hline
        0.1.0                                  & 22/12/2025    & V. Baleanu      & F. Pasqual               & -                    & Prima Stesura                                                                                 \\
        \hline
    \end{tabularx}
\end{center}

\newpage

\tableofcontents

\newpage
\listoffigures
\newpage
\listoftables
\newpage

\section{Introduzione}
\subsection{Finalità del documento}
Il presente documento ha lo scopo di definire e documentare gli obiettivi di qualità
che il gruppo \textit{SnakeByte} si impegna a perseguire durante l'intero ciclo di vita del
progetto. In particolare, il Piano di qualifica${_G}$ comprende i seguenti elementi:
\begin{itemize}
    \item \textbf{Piano della Qualità}: tutte le attività volte a definire gli obiettivi
          qualitativi e a stabilire i processi e le risorse necessarie per il loro conseguimento;
    \item \textbf{Controllo di Qualità}: le attività finalizzate ad accertare la conformità dei processi
          e dei prodotti agli standard, ai requisiti e alle metriche definite;
    \item \textbf{Miglioramento continuo}: le azioni intraprese per analizzare i risultati
          del controllo della qualità e per introdurre interventi correttivi e migliorativi
          sui processi e sul prodotto.
\end{itemize}
Il documento si rivolge sia ai membri del gruppo, come guida per il mantenimento
degli standard qualitativi, sia ai committenti e al proponente, come dimostrazione
dell'impegno del gruppo verso la qualità del prodotto finale.
\subsection{Glossario}
Il documento cita alcuni termini la cui definizione può risultare ambigua. Per questo è possibile consultare
il glossario${_G}$ il quale contiene le definizioni di tali espressioni, che saranno marcate da una lettera G a
pedice.
\subsection{Riferimenti}
\subsubsection{Riferimenti Normativi}
\begin{itemize}
    \item \textbf{Norme di Progetto}:
          \\ \url{https://snakebyteteam.github.io/3-PB/Documenti_interni/Norme_di_progetto/Norme_di_Progetto_v2.0.0.pdf}
          \\ (consultato il 05/04/2026).
    \item \textbf{Vimar View4Life Capitolato di Ingegneria del Software Università di Padova 2025-2026}:
          \\ \url{https://www.math.unipd.it/~tullio/IS-1/2025/Progetto/C9.pdf}
          \\ (consultato il 05/04/2026).
\end{itemize}
\subsubsection{Riferimenti Informativi}
\begin{itemize}
    \item \textbf{Glossario}:
          \\ \url{https://snakebyteteam.github.io/3-PB/Documenti_interni/Glossario/Glossario_v2.0.0.pdf}
          \\ (consultato il 05/04/2026).
    \item \textbf{Standard ISO 9000:2015 - Quality management systems. Fundamentals and vocabulary}:
          \\ \url{https://www.iso.org/standard/45481.html}
          \\ (consultato il 05/04/2026).
    \item \textbf{Standard ISO/IEC 12207:1995}:
          \\ \url{https://www.math.unipd.it/~tullio/IS-1/2009/Approfondimenti/ISO_12207-1995.pdf}
          \\ (consultato il 05/04/2026).
    \item \textbf{Standard ISO/IEC 9126:2001}:
          \\ \url{https://it.wikipedia.org/wiki/ISO/IEC_9126}
          \\ (consultato il 05/04/2026).
\end{itemize}

\section{Qualità di processo}
La qualità di processo${_G}$ ha lo scopo di garantire che le attività di sviluppo
del software siano pianificate, eseguite, monitorate e migliorate in modo
sistematico, al fine di ridurre il rischio di errori o incongruenze
durante l'intero ciclo di vita del progetto.

Per la definizione e la classificazione dei processi di qualità, il progetto
fa riferimento allo standard ISO/IEC 12207:1995, che descrive i processi
del ciclo di vita del software e li suddivide nelle seguenti categorie:
\begin{itemize}
    \item \textbf{Processi primari}: sono direttamente coinvolti nella
          realizzazione, fornitura e mantenimento del prodotto \textit{software};
    \item \textbf{Processi di supporto}: affiancano i processi primari fornendo
          attività di supporto come la documentazione e la verifica;
    \item \textbf{Processi organizzativi}: forniscono le risorse necessarie
          all'esecuzione dei processi primari e di supporto, favorendo il
          miglioramento continuo dell'organizzazione.
\end{itemize}
\subsection{Processi primari}
\subsubsection{Fornitura}
Nella tabella che segue il termine $BAC$ (\textit{Budget At Completion})
rappresenta il budget totale previsto inizialmente per il completamento
del progetto.
\begin{table}[H]
    \sogliatabella{
        \sogliariga{MPC1}{Planned Value}{$\geq 0$}{$\leq BAC$}
        \sogliariga{MPC2}{Earned Value}{$\geq 0$}{$\leq EAC$}
        \sogliariga{MPC3}{Actual Cost}{$\geq 0$}{$\leq EAC$}
        \sogliariga{MPC4}{Cost Performance Index}{$\geq 0.80$}{$1.00$}
        \sogliariga{MPC5}{Schedule Performance Index}{$\geq 0.80$}{$1.00$}
        \sogliariga{MPC6}{Estimate At Completion (EAC)}{$\leq 120\% \cdot BAC$}{$\leq BAC$}
        \sogliariga{MPC7}{Estimate To Complete}{$\geq 0$}{$\leq BAC$}
    }
    \caption{Tabella dei valori soglia per metriche di fornitura}
\end{table}

\begin{table}[H]
    \subsubsection{Sviluppo}
    \sogliatabella{
        \sogliariga{MPC8}{Requirements Stability Index}{$\geq 80\%$}{$100\%$}
    }
    \caption{Tabella dei valori soglia per metriche di sviluppo}
\end{table}


\begin{table}[H]
    \subsection{Processi di supporto}
    \subsubsection{Documentazione}
    \sogliatabella{
        \sogliariga{MPC9}{Indice di Gulpease}{$\geq 60$}{$\geq 80$}
        \sogliariga{MPC10}{Indice di Frammentazione (IF)}{$5\% \leq IF \leq 20\%$}{$10\% \leq IF \leq 15\%$}
    }
    \caption{Tabella dei valori soglia per metriche di documentazione}
\end{table}

\begin{table}[H]
    \subsubsection{Verifica}
    \sogliatabella{
        \sogliariga{MPC11}{Test Success Rate}{$\geq 80\%$}{$100\%$}
    }
    \caption{Tabella dei valori soglia per metriche di verifica}
\end{table}

\begin{table}[H]
    \subsubsection{Gestione della qualità}
    \sogliatabella{
        \sogliariga{MPC12}{Percentuale Metriche Soddisfatte}{$\geq 80\%$}{$100\%$}
    }
    \caption{Tabella dei valori soglia per metriche di gestione della qualità}
\end{table}

\begin{table}[H]
    \subsection{Processi organizzativi}
    \subsubsection{Gestione dei processi}
    \sogliatabella{
        \sogliariga{MPC13}{Percentuale Rischi Inattesi}{$\leq 20\%$}{$0\%$}
        \sogliariga{MPC14}{Labor Efficiency}{$\geq 70\%$}{$\geq 100\%$}
    }
    \caption{Tabella dei valori soglia per metriche di gestione dei processi}
\end{table}

\section{Qualità di prodotto}
La qualità di prodotto$_G$ è definita secondo un modello basato sullo standard
ISO/IEC 9126, che struttura la qualità in caratteristiche e sottocaratteristiche
misurabili attraverso metriche specifiche.
\\Per il prodotto \textit{software} sviluppato, il gruppo \textit{SnakeByte}
ha individuato le seguenti caratteristiche qualitative da monitorare e garantire:
\begin{itemize}
    \item \textbf{Funzionalità}: capacità di fornire funzioni che soddisfino
          i requisiti definiti;
    \item \textbf{Affidabilità}: capacità di mantenere il livello di prestazione
          richiesto quando utilizzato in condizioni specifiche;
    \item \textbf{Usabilità}: facilità con cui gli utenti possono apprendere,
          comprendere e utilizzare il prodotto;
    \item \textbf{Efficienza}: capacità di fornire prestazioni adeguate in
          relazione alle risorse utilizzate;
    \item \textbf{Manutenibilità}: facilità con cui il prodotto può essere
          modificato, corretto o esteso, favorendo interventi di manutenzione;
    \item \textbf{Portabilità}: capacità di essere trasferito e adattato a
          diversi ambienti operativi, piattaforme o configurazioni hardware e software.
\end{itemize}

\subsection{Funzionalità}
\begin{table}[H]
    \sogliatabella{
        \sogliariga{MPD1}{Percentuale Requisiti Obbligatori Soddisfatti}{$100\%$}{$100\%$}
        \sogliariga{MPD2}{Percentuale Requisiti Opzionali Soddisfatti}{$0\%$}{$100\%$}
        \sogliariga{MPD3}{Percentuale Requisiti Desiderabili Soddisfatti}{$0\%$}{$100\%$}
    }
    \caption{Tabella dei valori soglia per metriche di funzionalità}
\end{table}

\begin{table}[H]
    \subsection{Affidabilità}
    \sogliatabella{
        \sogliariga{MPD4}{Line Coverage}{$\geq 75\%$}{$100\%$}
        \sogliariga{MPD5}{Branch Coverage}{$\geq 75\%$}{$100\%$}
        \sogliariga{MPD6}{Statement Coverage}{$\geq 75\%$}{$100\%$}
    }
    \caption{Tabella dei valori soglia per metriche di affidabilità}
\end{table}

\begin{table}[H]
    \subsection{Usabilità}
    \sogliatabella{
        \sogliariga{MPD7}{Profondità Massima di Navigazione }{$\leq 7$ click}{$\leq 3$ click}
    }
    \caption{Tabella dei valori soglia per metriche di usabilità}
\end{table}

\begin{table}[H]
    \subsection{Efficienza}
    \sogliatabella{
        \sogliariga{MPD8}{Tempo Medio di Risposta}{$\leq 5$ secondi}{$\leq 1$ secondo}
    }
    \caption{Tabella dei valori soglia per metriche di efficienza}
\end{table}

\begin{table}[H]
    \subsection{Manutenibilità}
    \sogliatabella{
        \sogliariga{MPD9}{Complessità Ciclomatica}{$\leq 20$}{$\leq 10$}
    }
    \caption{Tabella dei valori soglia per metriche di manutenibilità}
\end{table}

\begin{table}[H]
    \subsection{Portabilità}
    \sogliatabella{
        \sogliariga{MPD10}{Compatibilità Cross-Browser}{$\geq 80\%$}{100\%}
    }
    \caption{Tabella dei valori soglia per metriche di portabilità}
\end{table}

\section{Strategie di Testing}
\subsection{Test di sistema}

\addTestS{Verificare che un Utente non autenticato possa effettuare l'autenticazione con il Sistema}{RF1-OB}{S}
\addTestS{Verificare che l'Utente non autenticato, durante il processo di autenticazione, possa inserire il proprio username}{RF2-OB}{S}
\addTestS{Verificare che l'Utente non autenticato, durante il processo di autenticazione, possa inserire la propria password}{RF3-OB}{S}
\addTestS{Verificare che l'Utente non autenticato riceva un messaggio errore a seguito di un tentativo di autenticazione fallito}{RF4-OB}{S}
\addTestS{Verificare che l'Utente autenticato possa impostare una nuova password}{RF5-OB}{S}
\addTestS{Verificare che l'Utente non registrato possa essere autenticato con la password temporanea fornita dall'amministratore e impostare una nuova password}{RF6-OB}{S}
\addTestS{Verificare che l'Utente autenticato con la password temporanea possa impostare una nuova password}{RF7-OB}{S}
\addTestS{Verificare che l'Utente non registrato riceva un errore a seguito di un tentativo di autenticazione con password temporanea fallito}{RF8-OB RF9-OB}{S}
\addTestS{Verificare che l'Utente autenticato con la password temporanea riceva un errore nel caso la nuova password sia uguale alla password temporanea}{RF10-OB}{S}
\addTestS{Verificare che l'Utente autenticato con la password temporanea riceva un errore nel caso la nuova password non rispetti i criteri richiesti}{RF11-OB}{S}
\addTestS{Verificare che l'amministratore possa visualizzare se un account MyVimar è configurato nel Sistema}{RF12-OB}{S}
\addTestS{Verificare che l'amministratore possa collegare un account MyVimar al Sistema}{RF13-OB}{S}
\addTestS{Verificare che l'amministratore possa rimuovere un account MyVimar dal Sistema}{RF14-OB}{S}
\addTestS{Verificare che l'amministratore possa visualizzare l'elenco degli utenti del Sistema}{RF15-DE}{S}
\addTestS{Verificare che l'amministratore, visualizzando l'elenco degli utenti del Sistema, possa visualizzare un utente del Sistema nel dettaglio}{RF16-OB}{S}
\addTestS{Verificare che l'amministratore, visualizzando un utente del Sistema nel dettaglio, ne visualizzi il nome}{RF17-OB}{S}
\addTestS{Verificare che l'amministratore, visualizzando un utente del Sistema nel dettaglio, ne visualizzi il cognome}{RF18-OB}{S}
\addTestS{Verificare che l'amministratore, visualizzando un utente del Sistema nel dettaglio, ne visualizzi l'username}{RF19-OB}{S}
\addTestS{Verificare che l'amministratore possa creare un utente Operatore Sanitario}{RF20-DE}{S}
\addTestS{Verificare che l'amministratore, creando un utente Operatore Sanitario, possa inserirne il nome}{RF21-DE}{S}
\addTestS{Verificare che l'amministratore, creando un utente Operatore Sanitario, possa inserirne il cognome}{RF22-DE}{S}
\addTestS{Verificare che l'amministratore, creando un utente Operatore Sanitario, possa inserire un username}{RF23-DE}{S}
\addTestS{Verificare che l'amministratore, creando un utente Operatore Sanitario, possa generare una password temporanea}{RF24-DE}{S}
\addTestS{Verificare che l'amministratore riceva un errore se l'username utilizzato per la creazione di un utente Operatore Sanitario è già in uso}{RF25-DE}{S}
\addTestS{Verificare che l'amministratore possa eliminare un utente Operatore Sanitario}{RF26-DE}{S}
\addTestS{Verificare che l'amministratore possa visualizzare tutti i reparti del Sistema}{RF27-OB}{S}
\addTestS{Verificare che l'amministratore possa visualizzare un reparto del Sistema nel dettaglio}{RF28-OB}{S}
\addTestS{Verificare che l'amministratore, visualizzando un reparto del Sistema nel dettaglio, possa visualizzarne il nome}{RF29-OB}{S}
\addTestS{Verificare che l'amministratore possa creare un nuovo reparto}{RF30-OB}{S}
\addTestS{Verificare che l'amministratore, creando un nuovo reparto, possa inserirne il nome}{RF31-OB}{S}
\addTestS{Verificare che l'amministratore riceva un errore se il nome inserito per creare il nuovo reparto è già in uso}{RF32-OB}{S}
\addTestS{Verificare che l'amministratore possa modificare il nome di un reparto del Sistema}{RF33-OB}{S}
\addTestS{Verificare che l'amministratore possa eliminare un reparto del Sistema}{RF34-OB}{S}
\addTestS{Verificare che l'amministratore possa assegnare un reparto ad un utente Operatore Sanitario}{RF35-OB}{S}
\addTestS{Verificare che l'amministratore, assegnando un reparto ad un utente Operatore Sanitario, selezioni un utente Operatore Sanitario}{RF36-OB}{S}
\addTestS{Verificare che l'amministratore, assegnando un reparto ad un utente Operatore Sanitario, selezioni un reparto del Sistema}{RF37-OB}{S}
\addTestS{Verificare che l'amministratore possa rimuovere un utente Operatore Sanitario dall'assegnazione ad un reparto}{RF38-OB}{S}
\addTestS{Verificare che l'amministratore possa assegnare un reparto del Sistema ad un appartamento}{RF39-OB}{S}
\addTestS{Verificare che l'amministratore, assegnando un reparto del Sistema ad un appartamento, selezioni un appartamento}{RF40-OB}{S}
\addTestS{Verificare che l'amministratore, assegnando un reparto del Sistema ad un appartamento, selezioni un reparto del Sistema}{RF41-OB}{S}
\addTestS{Verificare che l'amministratore possa rimuovere un appartamento dall'assegnazione ad un reparto}{RF42-OB}{S}
\addTestS{Verificare che l'utente possa visualizzare una dashboard che rappresenta lo stato del sistema}{RF43-OB}{S}
\addTestS{Verificare che l'utente, visualizzando la dashboard, possa visualizzare un modulo di gestione allarmi}{RF44-OB}{S}
\addTestS{Verificare che l'utente, visualizzando il modulo di gestione allarmi, possa visualizzare ogni singolo allarme}{RF45-OB}{S}
\addTestS{Verificare che l'utente, visualizzando un singolo allarme, possa visualizzare il segnale di pericolo}{RF46-OB}{S}
\addTestS{Verificare che l'utente, visualizzando un singolo allarme, possa visualizzarne il nome}{RF47-OB}{S}
\addTestS{Verificare che l'utente, visualizzando un singolo allarme, possa visualizzarne il tempo trascorso dallo scatto}{RF48-OB}{S}
\addTestS{Verificare che l'utente, visualizzando la dashboard, possa visualizzare un modulo di statistiche allarmi}{RF49-OB}{S}
\addTestS{Verificare che l'utente, visualizzando il modulo di statistiche allarmi, possa visualizzare il numero di allarmi risolti}{RF50-OB}{S}
\addTestS{Verificare che l'utente, visualizzando il modulo di statistiche allarmi, possa visualizzare il numero di allarmi attivi}{RF51-OB}{S}
\addTestS{Verificare che l'utente, visualizzando la dashboard, possa visualizzare un modulo di informazioni Utente}{RF52-OB}{S}
\addTestS{Verificare che l'utente, visualizzando il modulo di informazioni utente, possa visualizzare il nome Utente}{RF53-OB}{S}
\addTestS{Verificare che l'utente, visualizzando il modulo di informazioni utente, possa visualizzare il cognome Utente}{RF54-OB}{S}
\addTestS{Verificare che l'utente, visualizzando la dashboard, possa visualizzare un modulo di analisi clima}{RF55-OB}{S}
\addTestS{Verificare che l'utente, visualizzando la dashboard, possa visualizzare un modulo di analisi consumi}{RF56-OB}{S}
\addTestS{Verificare che l'utente, visualizzando la dashboard, possa visualizzare un modulo di analisi presenze}{RF57-OB}{S}
\addTestS{Verificare che l'amministratore possa aggiungere alla visualizzazione della dashboard il modulo di statistiche allarmi}{RF58-DE}{NI}
\addTestS{Verificare che l'amministratore possa rimuovere alla visualizzazione della dashboard il modulo di statistiche allarmi}{RF59-DE}{NI}
\addTestS{Verificare che l'amministratore possa aggiungere alla visualizzazione della dashboard il modulo di informazioni Utente}{RF60-DE}{NI}
\addTestS{Verificare che l'amministratore possa rimuovere alla visualizzazione della dashboard il modulo di informazioni Utente}{RF61-DE}{NI}
\addTestS{Verificare che l'amministratore possa aggiungere alla visualizzazione della dashboard il modulo di analisi clima}{RF62-DE}{NI}
\addTestS{Verificare che l'amministratore possa rimuovere alla visualizzazione della dashboard il modulo di analisi clima}{RF63-DE}{NI}
\addTestS{Verificare che l'amministratore possa aggiungere alla visualizzazione della dashboard il modulo di analisi consumi}{RF64-DE}{NI}
\addTestS{Verificare che l'amministratore possa rimuovere alla visualizzazione della dashboard il modulo di analisi consumi}{RF65-DE}{NI}
\addTestS{Verificare che l'utente possa risolvere un allarme}{RF66-OB}{S}
\addTestS{Verificare che l'utente possa visualizzare le analytics}{RF67-OB}{S}
\addTestS{Verificare che l'utente, visualizzando le analytics, possa visualizzare l'elenco dei suggerimenti risparmio energetico}{RF68-OB}{S}
\addTestS{Verificare che l'utente, visualizzando l'elenco dei suggerimenti risparmio energetico, possa visualizzare un suggerimento risparmio energetico}{RF69-OB}{S}
\addTestS{Verificare che l'utente, visualizzando le analytics, possa visualizzare il grafico dedicato al consumo energetico}{RF70-OB}{S}
\addTestS{Verificare che l'utente, visualizzando le analytics, possa visualizzare il grafico dedicato alle anomalie dell'impianto}{RF71-OB}{S}
\addTestS{Verificare che l'utente, visualizzando le analytics, possa visualizzare il grafico dedicato al rilevamento presenze}{RF72-OB}{S}
\addTestS{Verificare che l'utente, visualizzando le analytics, possa visualizzare il grafico dedicato alla presenza prolungata nello stesso ambiente}{RF73-OB}{S}
\addTestS{Verificare che l'utente, visualizzando le analytics, possa visualizzare il grafico dedicato alle variazioni di temperatura}{RF74-OB}{S}
\addTestS{Verificare che l'utente, visualizzando le analytics, possa visualizzare il grafico dedicato agli allarmi inviati e risolti}{RF75-OB}{S}
\addTestS{Verificare che l'utente, visualizzando le analytics, possa visualizzare il grafico dedicato alla frequenza degli allarmi}{RF76-OP}{S}
\addTestS{Verificare che l'utente, visualizzando le analytics, possa visualizzare il grafico dedicato alla frequenza delle cadute}{RF77-OP}{S}
\addTestS{Verificare che l'utente possa visualizzare un appartamento del Sistema}{RF78-OB}{S}
\addTestS{Verificare che l'utente, visualizzando un appartamento del Sistema, possa visualizzarne il nome}{RF79-OB}{S}
\addTestS{Verificare che l'utente, visualizzando un appartamento del Sistema, possa visualizzarne la mappa degli allarmi}{RF80-OP}{S}
\addTestS{Verificare che l'utente, visualizzando un appartamento del Sistema, possa visualizzarne le stanze}{RF81-OB}{S}
\addTestS{Verificare che l'utente possa visualizzare una stanza di un appartamento nel dettaglio}{RF82-OB}{S}
\addTestS{Verificare che l'amministratore, visualizzando una stanza nel dettaglio, possa visualizzarne il nome}{RF83-OB}{S}
\addTestS{Verificare che l'amministratore, visualizzando una stanza nel dettaglio, possa visualizzarne l'elenco dei dispositivi}{RF84-OB}{S}
\addTestS{Verificare che l'amministratore, visualizzando l'elenco dei dispositivi di una stanza, possa visualizzare un singolo dispositivo di tipo:
    \begin{itemize}
        \item termostato;
        \item sensore di caduta;
        \item sensore di presenza;
        \item punto luce;
        \item pulsante di allarme;
        \item porta di ingresso;
        \item tapparella.
    \end{itemize}
}{RF85-OB}{S}
\addTestS{Verificare che l'amministratore, visualizzando un dispositivo, possa visualizzarne il nome}{RF86-OB}{S}
\addTestS{Verificare che l'amministratore, visualizzando un dispositivo, possa visualizzarne lo stato}{RF87-OB}{S}
\addTestS{Verificare che l'amministratore, visualizzando un dispositivo, possa visualizzarne le azioni eseguibili}{RF88-OP}{S}
\addTestS{Verificare che l'amministratore possa abilitare un appartamento nel Sistema}{RF89-OB}{S}
\addTestS{Verificare che l'amministratore possa disabilitare un appartamento nel Sistema}{RF90-OB}{S}
\addTestS{Verificare che l'amministratore possa creare un allarme nel Sistema}{RF91-OB}{S}
\addTestS{Verificare che l'amministratore, nella creazione di un allarme, possa selezionare un appartamento}{RF92-OB}{S}
\addTestS{Verificare che l'amministratore, nella creazione di un allarme, possa selezionare un sensore}{RF93-OB}{S}
\addTestS{Verificare che l'amministratore, nella creazione di un allarme, possa selezionare un livello di priorità tra:
    \begin{itemize}
        \item priorità 1 (bianco);
        \item priorità 2 (verde);
        \item priorità 3 (arancione);
        \item priorità 4 (rossa)
    \end{itemize}
}{RF94-OB}{S}
\addTestS{Verificare che l'amministratore, nella creazione di un allarme, possa selezionare una soglia di scatto}{RF95-OB}{S}
\addTestS{Verificare che l'amministratore, nella creazione di un allarme, possa selezionare un orario di attivazione}{RF96-OB}{S}
\addTestS{Verificare che l'amministratore, nella creazione di un allarme, possa selezionare un orario di disattivazione}{RF97-OB}{S}
\addTestS{Verificare che l'amministratore riceva un errore se durante la creazione di un allarme non ha selezionato alcun sensore}{RF98-OB}{S}
\addTestS{Verificare che l'amministratore riceva un errore se durante la creazione di un allarme non ha selezionato alcun livello di priorità}{RF99-OB}{S}
\addTestS{Verificare che l'amministratore riceva un errore se durante la creazione di un allarme non ha selezionato alcuna soglia di scatto}{RF100-OB}{S}
\addTestS{Verificare che l'amministratore possa modificare la priorità di un allarme del Sistema}{RF101-OB}{S}
\addTestS{Verificare che l'amministratore possa modificare la soglia di scatto di un allarme del Sistema}{RF102-OB}{S}
\addTestS{Verificare che l'amministratore possa modificare l'orario di attivazione di un allarme del Sistema}{RF103-OB}{S}
\addTestS{Verificare che l'amministratore possa modificare l'orario di disattivazione di un allarme del Sistema}{RF104-OB}{S}
\addTestS{Verificare che l'amministratore possa abilitare un allarme nel Sistema}{RF105-OB}{S}
\addTestS{Verificare che l'amministratore possa disabilitare un allarme nel Sistema}{RF106-OB}{S}
\addTestS{Verificare che l'amministratore possa eliminare un allarme dal Sistema}{RF107-OB}{S}
\addTestS{Verificare che l'amministratore possa visualizzare le notifiche inviate dal Sistema}{RF108-OP}{S}
\addTestS{Verificare che l'amministratore, visualizzando le notifiche inviate dal Sistema, possa visualizzare una notifica nel dettaglio}{RF109-OP}{S}
\addTestS{Verificare che l'amministratore, visualizzando una notifica nel dettaglio possa visualizzarne il titolo}{RF110-OP}{S}
\addTestS{Verificare che l'amministratore, visualizzando una notifica nel dettaglio possa visualizzarne il tempo trascorso dall'invio}{RF111-OP}{S}
\makeTSTable


\subsection{Test di accettazione}
\addTestA{Verificare che il prodotto dia la possibilità all'utente non autenticato di essere riconosciuto tramite username e password}{S}
\addTestA{Verificare che il prodotto dia la possibilità all'utente non registrato di autenticarsi tramite username e password temporanea, forniti dall'amministratore}{S}
\addTestA{Verificare che il prodotto dia la possibilità, all'utente autenticato tramite password temporanea, di cambiare password}{S}
\addTestA{Verificare che il prodotto dia la possibilità all'utente amministratore di gestire l'account MyVimar collegato all'impianto}{S}
\addTestA{Verificare che il prodotto dia la possibilità all'utente amministratore di visualizzare l'elenco degli utenti}{S}
\addTestA{Verificare che il prodotto dia la possibilità all'utente amministratore di creare utenti Operatore Sanitario}{S}
\addTestA{Verificare che il prodotto dia la possibilità all'utente amministratore eliminare utenti Operatore Sanitario}{S}
\addTestA{Verificare che il prodotto dia la possibilità all'utente amministratore di generare una password temporanea per gli utenti Operatore Sanitario}{S}
\addTestA{Verificare che il prodotto dia la possibilità all'utente di visualizzare la lista dei reparti}{S}
\addTestA{Verificare che il prodotto dia la possibilità all'amministratore di creare un reparto}{S}
\addTestA{Verificare che il prodotto dia la possibilità all'amministratore di eliminare un reparto}{S}
\addTestA{Verificare che il prodotto dia la possibilità all'amministratore di assegnare un reparto ad un utente Operatore Sanitario}{S}
\addTestA{Verificare che il prodotto dia la possibilità all'amministratore di rimuovere l'assegnazione di un reparto ad un utente Operatore Sanitario}{S}
\addTestA{Verificare che il prodotto dia la possibilità all'amministratore di assegnare un reparto ad un appartamento}{S}
\addTestA{Verificare che il prodotto dia la possibilità all'amministratore di rimuovere l'assegnazione di un reparto ad un appartamento}{S}
\addTestA{Verificare che il prodotto dia la possibilità all'utente di visualizzare la dashboard}{S}
\addTestA{Verificare che il prodotto dia la possibilità all'utente di visualizzare il modulo gestione allarmi nella dashboard}{S}
\addTestA{Verificare che il prodotto dia la possibilità all'utente di visualizzare il modulo statistiche allarmi nella dashboard}{S}
\addTestA{Verificare che il prodotto dia la possibilità all'amministratore di aggiungere il modulo informazioni utente alla dashboard}{NI}
\addTestA{Verificare che il prodotto dia la possibilità all'amministratore di aggiungere il modulo analisi clima alla dashboard}{NI}
\addTestA{Verificare che il prodotto dia la possibilità all'amministratore di aggiungere il modulo analisi consumi alla dashboard}{NI}
\addTestA{Verificare che il prodotto dia la possibilità all'amministratore di rimuovere il modulo informazioni utente dalla dashboard}{NI}
\addTestA{Verificare che il prodotto dia la possibilità all'amministratore di rimuovere il modulo analisi clima dalla dashboard}{NI}
\addTestA{Verificare che il prodotto dia la possibilità all'amministratore di rimuovere il modulo analisi consumi dalla dashboard}{NI}
\addTestA{Verificare che il prodotto dia la possibilità all'utente di visualizzare il modulo informazioni utente nella dashboard}{S}
\addTestA{Verificare che il prodotto dia la possibilità all'utente di visualizzare il modulo analisi clima nella dashboard}{S}
\addTestA{Verificare che il prodotto dia la possibilità all'utente di visualizzare il modulo analisi consumi nella dashboard}{S}
\addTestA{Verificare che il prodotto dia la possibilità all'utente di risolvere un allarme}{S}
\addTestA{Verificare che il prodotto dia la possibilità all'utente di visualizzare le analytics}{S}
\addTestA{Verificare che il prodotto dia la possibilità all'utente di visualizzare i consigli di risparmio energetico}{S}
\addTestA{Verificare che il prodotto dia la possibilità all'utente di visualizzare il grafico dedicato al consumo energetico}{S}
\addTestA{Verificare che il prodotto dia la possibilità all'utente di visualizzare il grafico dedicato alle anomalie dell'impianto}{S}
\addTestA{Verificare che il prodotto dia la possibilità all'utente di visualizzare il grafico dedicato al rilevamento presenze}{S}
\addTestA{Verificare che il prodotto dia la possibilità all'utente di visualizzare il grafico dedicato alla presenza prolungata nello stesso ambiente}{S}
\addTestA{Verificare che il prodotto dia la possibilità all'utente di visualizzare il grafico dedicato alle variazioni di temperatura}{S}
\addTestA{Verificare che il prodotto dia la possibilità all'utente di visualizzare il grafico dedicato agli allarmi inviati e risolti}{S}
\addTestA{Verificare che il prodotto dia la possibilità all'utente di visualizzare il grafico dedicato alla frequenza degli allarmi}{S}
\addTestA{Verificare che il prodotto dia la possibilità all'utente di visualizzare il grafico dedicato alla frequenza delle cadute}{S}
\addTestA{Verificare che il prodotto dia la possibilità all'utente di visualizzare un appartamento}{S}
\addTestA{Verificare che il prodotto dia la possibilità all'utente di visualizzare la mappa degli allarmi di un appartamento}{S}
\addTestA{Verificare che il prodotto dia la possibilità all'utente di visualizzare l'elenco di stanze di un appartamento}{S}
\addTestA{Verificare che il prodotto dia la possibilità all'utente di visualizzare l'elenco dei dispositivi di una stanza}{S}
\addTestA{Verificare che il prodotto dia la possibilità all'amministratore di abilitare un appartamento}{S}
\addTestA{Verificare che il prodotto dia la possibilità l'amministratore di disabilitare un appartamento}{S}
\addTestA{Verificare che il prodotto dia la possibilità all'amministratore di creare un allarme}{S}
\addTestA{Verificare che il prodotto dia la possibilità all'amministratore di modificare un allarme}{S}
\addTestA{Verificare che il prodotto dia la possibilità all'amministratore di abilitare un allarme}{S}
\addTestA{Verificare che il prodotto dia la possibilità all'amministratore di disabilitare un allarme}{S}
\addTestA{Verificare che il prodotto dia la possibilità all'utente di ricevere notifiche dal Sistema}{S}
\addTestA{Verificare che il prodotto dia la possibilità all'utente di visualizzare le notifiche ricevute dal Sistema}{S}

\makeTATable

\subsection{Test di integrazione}
\addTestI{Contesto \textit{BreadcrumbService}; caso 'breadcrumb construction e navigation tracking': verifica che breadcrumbs\$ emette array corretto, re-emette on NavigationEnd, gestisce routes multi-level ignoring empty segments and building proper paths.}{Esito atteso: breadcrumbs per routes annidate costruiti correttamente, no double-slash, multi-segment paths concatenati.}{S}
\addTestI{Contesto \textit{InternalAuthService}; consolidated summary covering login flows, session management, refresh token handling, concurrent restore, logout and user stream: verifica che il servizio gestisce login POST con credenziali, esegue first-login con temp password, aggiorna token via refresh, ripristina sessione dai cookie, evita richieste refresh duplicate con operazioni concorrenti, resetta stato al logout, espone stream utente. Tutti i flussi richiedono withCredentials true e metodo POST.}{Esito atteso: login imposta session.username e token; first-login accetta temp password e nuova password; refresh aggiorna accessToken; restore concorrente single-request; logout pulisce token/role; currentUser inizialmente null; tutti gli endpoint con POST e withCredentials true; logout chiama /auth/logout.}{S}
\addTestI{Contesto \textit{AnalyticsComponent}; consolidated summary covering component creation, apartment loading, chart fetching and metric retrieval: verifica che il componente sia valorizzato, carica il primo appartamento di default invocando getAllApartments e getAnalytics('1'), aggiorna grafici al cambio appartamento con getAnalytics('2'), resituisce undefined se metrica non trovata, trova la metrica corretta con getChartByMetric('plant-anomalies').}{Esito atteso: component valorizzato; mockAnalyticsService.\textit{getAllApartments} invocato; mockAnalyticsService.\textit{getAnalytics} invocato con '1' al caricamento e con '2' al cambio; getChartByMetric('plant-anomalies') restituisce metrica con metric:'plant-anomalies'; getChartByMetric('unknown') restituisce undefined; feature tests RF61-RF62-RF63 verificano stesso flusso con stub API.}{S}
\addTestI{Contesto \textit{MainLayoutComponent}; caso 'carica navItems, gestisce MyVimar e logout': verifica che mockNavService.\textit{getNavItems} viene invocato e component.navItems è caricato; mostra warning MyVimar per admin; logout reindirizza a /auth/login.}{Esito atteso: mockNavService.\textit{getNavItems} viene invocato; component.navItems.length è maggiore di 0; il valore di component.showVimarAssociationWarning è appropriato; navigateSpy viene invocato con ['/auth/login'] al logout.}{S}
\addTestI{Contesto \textit{MainLayoutComponent}; caso 'reindirizza al login senza sessione e apre pannello notifiche': verifica che senza sessione utente reindirizza a /auth/login; le nuove notifiche realtime aprono il pannello.}{Esito atteso: component.navItems coincide con []; navigateSpy viene invocato con ['/auth/login']; il valore di component.isNotificationPanelOpen diventa true al ricevimento notifica.}{S}
\addTestI{Contesto MainLayout feature integration; caso 'RF46-RF47-RF48: carica navItems da ruolo utente, gestisce login redirect e notifiche realtime': verifica che navServiceStub.\textit{getNavItems} viene invocato con UserRole.AMMINISTRATORE, mostra navItems.length > 0 e warning MyVimar, reindirizza a /auth/login senza sessione, apre pannello notifiche su evento realtime.}{Esito atteso: navServiceStub.\textit{getNavItems} invocato con admin role; component.navItems > 0; showVimarAssociationWarning true per admin; navigateSpy con ['/auth/login'] senza sessione; isNotificationPanelOpen toggled true su notifica realtime.}{S}
\addTestI{Contesto device-taxonomy; consolidated summary covering device type resolution, endpoint labeling and dynamic discovery: verifica che il servizio risolve type thermostat esplicitamente e da endpoint sfeType ambiguo, inferisce entrance door da nome italiano, scenario devices da nome, fallback unknown; getEndpointLabel restituisce label leggibile per SFE\_Cmd\_* noti, fallback humanized per sconosciuti, tracks runtime discovery in getDiscoveredUnmappedSfeCodes, restituisce label dispositivo leggibile.}{Esito atteso: thermostat resolved explicit e inferito; entrance\_door da nome italiano; alarm\_button da scenario; unknown fallback; getEndpointLabel('SFE\_Cmd\_ChangeOverMode')='Comando cambio modalita HVAC'; unknown SFE code fallback umanizzato; getDeviceTypeLabel(PRESENCE\_SENSOR)='Sensore presenza'; unknown type fallback.}{S}
\addTestI{Contesto \textit{ElapsedTimePipe}; consolidated summary covering time formatting for past and future dates in compact format: verifica che la pipe ritorna valore originale per date non valide, mostra 'tra poco' per timestamp futuri, restituisce formato compatto secondi/minuti/ore/giorni per intervalli passati con suffisso 'fa'.}{Esito atteso: input invalido ritorna originale; timestamp futuro='tra poco'; 30s fa < 60s; 1m fa 60-3600s; 1h fa 3600-86400s; 1g fa >= 86400s; 3h fa e 1h 12m fa con ore/minuti; 5g fa per giorni.}{S}
\addTestI{Contesto notification-title util integration; consolidated summary covering fallback handling, priority mapping, enum validation and truncation: verifica che la funzione restituisce fallback per nomi non-string o vuoti, normalizza spazi, mappa priorita WHITE/GREEN/ORANGE/RED a simboli unicode, accetta priorita enum string case-insensitive, restituisce fallback per priorita sconosciute, tronca e appende ellipsis se eccede maxLength.}{Esito atteso: fallback per undefined/null/empty/invalid names; priority mapping WHITE=i, GREEN=bullet, ORANGE=!, RED=triangle; enum string e numeric supported; truncation con ellipsis; no error on small maxLength.}{S}
\addTestI{Contesto \textit{NotificationApiService}; consolidated summary for notification history fetching covering managed/unmanaged endpoints, cold observable behavior, session validation: verifica che il servizio chiama GET /managed e /unmanaged endpoints per storia notifiche, mappa payload verso NotificationEvent[], emette una volta e completa (non buffered), cold observable produce coppie di request per subscriber, restituisce [] se no sessione attiva senza HTTP calls.}{Esito atteso: managed/unmanaged GET invocati; result mappato; emissions=1 completed=true; cold observable behavior confirm; no session=no HTTP.}{S}
\addTestI{Contesto \textit{AlarmApiService}; consolidated summary for comprehensive CRUD and event management covering rules, alarms, resolution tracking: verifica che il servizio supporta GET /alarm-rules e GET /alarm-rules/:id, POST /alarm-rules, PUT /alarm-rules/:id, DELETE /alarm-rules/:id, GET /alarm-events/unmanaged e GET /alarm-events/managed per allarmi attivi/risolti con pagination (userId/limit/offset), GET /alarm-events/:id per dettaglio, PATCH /alarm-events/resolve per gestione. Input validation strict su numeric params, charset speciali handled in ID encoding, empty list handling e HTTP error propagation.}{Esito atteso: tutti gli endpoint supportati; GET/POST/PUT/DELETE/PATCH corretti; numeric validation userId/limit/offset; special chars in ID; empty list no error; HTTP errors propagated.}{S}
\addTestI{Contesto \textit{AlarmManagementRefreshService}; consolidated summary for hot observable refresh request management: verifica che il servizio espone observable hot per requestRefresh(), emette su richiesta esplicita, non bufferizza refresh precedenti per nuovi subscriber (hot behavior), refreshPromise resettato post-resolve.}{Esito atteso: hot observable behavior; refreshPromise cleared after resolve; single request per subscriber; no previous emissions buffering.}{S}
\addTestI{Contesto \textit{AlarmStateService}; consolidated summary for realtime alarm and notification state management covering initialization, add/update/remove operations, snapshot handling, unread counts: verifica che il servizio inizializza stato vuoto (alarms=[], notifications=[]), aggiunge/aggiorna allarmi via onAlarmTriggered preservando ID, rimuove via onAlarmResolved, antepone notifiche (più recente first), deduplica per notificationId, removeNotification e clearNotifications, setActiveAlarms imposta snapshot, emette conteggi activeAlarmsCount/unreadNotificationsCount.}{Esito atteso: init empty state; add/update preserve, no duplication; remove target only; notifications prepended; deduplicate by ID; snapshot handling merge; counters emitted correctly.}{S}
\addTestI{Contesto \textit{EventSubscriptionService}; consolidated summary part 1 - socket initialization and connection management: verifica che il servizio crea connessione socket via socketIoFactory su API\_BASE\_URL normalizzato (rtrim '/api' + '/ws'), risolve namespace ws correttamente, aggiorna stato ConnectionStatus su connect/reconnect/disconnect, riesegue join room precedenti al reconnect, supporta room join/leave con deduplicazione.}{Esito atteso: socket created on init; API\_BASE\_URL normalized; ws namespace resolved; connection status updated; rejoin on reconnect; dedup join/leave.}{S}
\addTestI{Contesto \textit{EventSubscriptionService}; consolidated summary part 2 - realtime events bootstrap and error handling: verifica che il servizio bootstrap room da plant/all al login se cache ward vuoto, evita bootstrap ripetuto stessa sessione, non interrompe realtime se bootstrap fallisce (fallback), riallinea room subscription senza page reload, ignora malformed events, disconnette socket alla destroy, su logout emette leave room.}{Esito atteso: bootstrap plant/all on login; single bootstrap per session; no error break; refresh ward subscription; malformed ignored; cleanup on destroy; logout emits leave.}{S}
\addTestI{Contesto \textit{EventSubscriptionService}; caso 'crea la connessione socket in initialize usando API\_BASE\_URL normalizzato': verifica che socketIoFactoryMock viene invocato con 'http://api.example.local/ws'.}{Esito atteso: socketIoFactoryMock viene invocato con 'http://api.example.local/ws'.}{S}
\addTestI{Contesto \textit{EventSubscriptionService}; caso 'risolve il namespace websocket ws quando API\_BASE\_URL termina con api': verifica che socketIoFactoryMock viene invocato con 'http://localhost/ws'.}{Esito atteso: socketIoFactoryMock viene invocato con 'http://localhost/ws'.}{S}
\addTestI{Contesto \textit{EventSubscriptionService}; caso 'aggiorna lo stato connessione su connect, reconnect attempt e disconnect': verifica che statuses coincide con [ ConnectionStatus.DISCONNECTED, ConnectionStatus.CONNECTED, ConnectionStatus.RECONNECTING, Co....}{Esito atteso: statuses coincide con [ ConnectionStatus.DISCONNECTED, ConnectionStatus.CONNECTED, ConnectionStatus.RECONNECTING, Co....}{S}
\addTestI{Contesto \textit{EventSubscriptionService}; caso 'esegue join e leave delle room con deduplicazione': verifica che fakeSocket.emit viene invocato con 'join-ward' e fakeSocket.emit viene invocato con 'leave-ward'.}{Esito atteso: fakeSocket.emit viene invocato con 'join-ward'; fakeSocket.emit viene invocato con 'leave-ward'; il valore di joinWardCalls.length è 2; il valore di leaveWardCalls.length è 1.}{S}
\addTestI{Contesto \textit{EventSubscriptionService}; caso 'riesegue il join delle room già sottoscritte al connect': verifica che il valore di joinCallsBeforeReconnect è 1 e il valore di joinCallsAfterReconnect.length è 2.}{Esito atteso: il valore di joinCallsBeforeReconnect è 1; il valore di joinCallsAfterReconnect.length è 2; joinCallsAfterReconnect.at(-1) coincide con ['join-ward', 'ward-a'].}{S}
\addTestI{Contesto \textit{EventSubscriptionService}; caso 'bootstrap delle room da plant all al login quando cache ward e vuota': verifica che httpClientSpy.get viene invocato con 'http://api.example.local/plant/all' e joinWardCalls coincide con [ ['join-ward', '21'], ['join-ward', '22'], ].}{Esito atteso: httpClientSpy.get viene invocato con 'http://api.example.local/plant/all'; joinWardCalls coincide con [ ['join-ward', '21'], ['join-ward', '22'], ].}{S}
\addTestI{Contesto \textit{EventSubscriptionService}; caso 'non interrompe il realtime quando bootstrap plant all fallisce': verifica che httpClientSpy.get viene invocato 1 volte e alarmStateSpy.\textit{onAlarmTriggered} viene invocato con expect.objectContaining(\{ id: 'active-alarm-bootstrap-fallback', \}).}{Esito atteso: httpClientSpy.get viene invocato 1 volte; alarmStateSpy.\textit{onAlarmTriggered} viene invocato con expect.objectContaining(\{ id: 'active-alarm-bootstrap-fallback', \}); alarmStateSpy.\textit{onNotificationReceived} viene invocato con expect.objectContaining(\{ notificationId: 'alarm-triggered-active-alarm-bootstrap-fallback', \}).}{S}
\addTestI{Contesto \textit{EventSubscriptionService}; caso 'evita bootstrap ripetuto per lo stesso utente nella stessa sessione': verifica che httpClientSpy.get viene invocato 1 volte.}{Esito atteso: httpClientSpy.get viene invocato 1 volte.}{S}
\addTestI{Contesto \textit{EventSubscriptionService}; caso '\textit{refreshWardRoomSubscription} riallinea le room joinate senza reload pagina': verifica che httpClientSpy.get viene invocato 2 volte e fakeSocket.emit viene invocato con 'join-ward'.}{Esito atteso: httpClientSpy.get viene invocato 2 volte; fakeSocket.emit viene invocato con 'join-ward'; fakeSocket.emit viene invocato con 'leave-ward'.}{S}
\addTestI{Contesto \textit{EventSubscriptionService}; caso 'ignora eventi raw malformati e payload backend non validi': verifica che alarmStateSpy.\textit{onAlarmTriggered} non viene invocato e alarmStateSpy.\textit{onAlarmResolved} non viene invocato.}{Esito atteso: alarmStateSpy.\textit{onAlarmTriggered} non viene invocato; alarmStateSpy.\textit{onAlarmResolved} non viene invocato; alarmStateSpy.\textit{onNotificationReceived} non viene invocato.}{S}
\addTestI{Contesto \textit{EventSubscriptionService}; caso 'usa location.origin quando API\_BASE\_URL non è configurato': verifica che socketIoFactoryMock viene invocato con `\$\{globalThis.location.origin\}/ws`.}{Esito atteso: socketIoFactoryMock viene invocato con `\$\{globalThis.location.origin\}/ws`.}{S}
\addTestI{Contesto \textit{EventSubscriptionService}; caso 'bootstrap accetta risposta wrappata in data[]': verifica che joinWardCalls coincide con [ ['join-ward', '31'], ['join-ward', '32'], ].}{Esito atteso: joinWardCalls coincide con [ ['join-ward', '31'], ['join-ward', '32'], ].}{S}
\addTestI{Contesto \textit{EventSubscriptionService}; caso 'bootstrap accetta risposta wrappata in plants[]': verifica che joinWardCalls coincide con [ ['join-ward', '41'], ['join-ward', '42'], ].}{Esito atteso: joinWardCalls coincide con [ ['join-ward', '41'], ['join-ward', '42'], ].}{S}
\addTestI{Contesto \textit{EventSubscriptionService}; caso 'su logout emette leave delle room joinate': verifica che fakeSocket.emit viene invocato con 'leave-ward'.}{Esito atteso: fakeSocket.emit viene invocato con 'leave-ward'.}{S}
\addTestI{Contesto \textit{EventSubscriptionService}; caso 'disconnette il socket e rimuove i listener alla distruzione': verifica che fakeSocket.\textit{removeAllListeners} viene invocato e fakeSocket.disconnect viene invocato.}{Esito atteso: fakeSocket.\textit{removeAllListeners} viene invocato; fakeSocket.disconnect viene invocato.}{S}
\addTestI{Contesto \textit{RealtimeAlarmEventNormalizerService}; consolidated summary for backend payload normalization: verifica che parseBackendTriggeredPayload estrae alarmEventId/wardId/alarmRuleId da payload nativo, accetta wardId numerico espresso come stringa, parseBackendResolvedPayload estrae alarmEventId e wardId opzionale.}{Esito atteso: triggered payload parsed correctly; wardId numerico string accepted; resolved payload extracted; optional wardId handled.}{S}
\addTestI{Contesto \textit{WardRealtimeCacheService}; consolidated summary for per-user ward cache management: verifica che mergeWardIds deduplica e normalizza valori stringa/numero, mantiene cache separata per utente, clearWardIds rimuove solo utente target, recupera cache vuota per JSON non valido.}{Esito atteso: dedup and normalize string/number; separate cache per user; selective clear; fallback empty cache.}{S}
\addTestI{Contesto \textit{WardSocketRoomCoordinatorService}; consolidated summary for room subscription lifecycle management: verifica che activateUser carica wardIds da cache e propone join, requestJoinRoom deduplica e persiste su cache, deactivateUser rimuove room runtime e pulisce cache, switch utente restituisce stanza da lasciare e cache del nuovo utente.}{Esito atteso: activate loads and proposes; join dedupped and cached; deactivate cleans; switch manages transitions.}{S}
\addTestI{Contesto \textit{AlarmConfigFormFieldOrchestratorHelper}; consolidated summary for form field orchestration in edit/create mode: verifica che applyModeState in edit disabilita campi non modificabili (name/plantId/deviceId/datapointId), resolveDeviceSelection ritorna datapoint leggibili solo e auto-selezione se singolo, resolveDatapointSelection normalizza id e restituisce datapoint atteso.}{Esito atteso: edit mode disabled fields correct; readable datapoints only; auto-select on single; ID normalized.}{S}
\addTestI{Contesto \textit{AlarmConfigFormValidationHelper}; consolidated summary for form validation in create and edit modes with enum and numeric datapoint types: verifica che create mode con enum forza operatore uguale e valida threshold su enum, con numerico applica validazione numerica pattern, edit mode applica stesso vincolo operatore enum, validazione numerica, senza metadati mantiene solo operatore corrente.}{Esito atteso: create enum forces = operator, numeric pattern validation; edit mode same; unsupported operator error.}{S}
\addTestI{Contesto \textit{AlarmRuleFormMapper}; consolidated summary for form value mapping from API model to form representation: verifica che il mapper mappa operatori >= come GREATER\_THAN\_OR\_EQUAL e <= come LESS\_THAN\_OR\_EQUAL, fallisce su operatori sconosciuti.}{Esito atteso: >= mapped GREATER\_THAN\_OR\_EQUAL; <= mapped LESS\_THAN\_OR\_EQUAL; unknown operators throw.}{S}
\addTestI{Contesto \textit{AlarmRuleRequestMapper}; consolidated summary for request payload serialization covering create, update and toggle operations: verifica che toUpdateRequest serializza operatore >= senza degradarlo, toCreateRequest serializza <= and normalizza soglia booleana in on/off, fallisce se thresholdValue vuoto o thresholdOperator null, toToggleRequest fallisce se name vuoto e normalizza operatori con padding dal backend.}{Esito atteso: >= preserved in serialize; <= preserved; boolean <-> on/off normalization; empty/null validation; padding normalization.}{S}
\addTestI{Contesto \textit{AlarmConfigFormStateService}; caso '\textit{ensurePlantsLoaded} in create mode carica e deduplica gli impianti': verifica che service.plants().map((plant) => plant.id) coincide con ['plant-1', 'plant-2', 'plant-3'] e service.plantsLoadError() è null.}{Esito atteso: service.plants().map((plant) => plant.id) coincide con ['plant-1', 'plant-2', 'plant-3']; service.plantsLoadError() è null.}{S}
\addTestI{Contesto \textit{AlarmConfigFormStateService}; caso '\textit{ensurePlantsLoaded} in edit mode non esegue chiamate': verifica che wardApiStub.\textit{getAvailablePlants} non viene invocato e wardApiStub.\textit{getWards} non viene invocato.}{Esito atteso: wardApiStub.\textit{getAvailablePlants} non viene invocato; wardApiStub.\textit{getWards} non viene invocato.}{S}
\addTestI{Contesto \textit{AlarmConfigFormStateService}; caso '\textit{loadDeviceOptionsByPlant} ritorna opzioni dispositivi e pulisce loading state': verifica che apartmentApiStub.loadApartmentViewForPlantId viene invocato con 'plant-1' e result coincide con options.}{Esito atteso: apartmentApiStub.loadApartmentViewForPlantId viene invocato con 'plant-1'; result coincide con options; il valore di service.isDevicesLoading() è false; service.devicesLoadError() è null.}{S}
\addTestI{Contesto \textit{AlarmConfigFormStateService}; caso '\textit{loadDeviceOptionsByPlant} in errore restituisce lista vuota e imposta messaggio': verifica che result coincide con [] e il valore di service.isDevicesLoading() è false.}{Esito atteso: result coincide con []; il valore di service.isDevicesLoading() è false; il valore di service.devicesLoadError() è 'Errore durante il caricamento dei dispositivi.'.}{S}
\addTestI{Contesto \textit{AlarmConfigFormStateService}; caso '\textit{resolveDatapointForEdit} usa il catalogo piante e restituisce il datapoint trovato': verifica che apartmentApiStub.\textit{getAllPlants} viene invocato 1 volte e datapointExtractionStub.findDatapointByDeviceAndDatapointId viene invocato con [\{ id: 'plant-1', name: 'Apt', rooms: [] \}].}{Esito atteso: apartmentApiStub.\textit{getAllPlants} viene invocato 1 volte; datapointExtractionStub.findDatapointByDeviceAndDatapointId viene invocato con [\{ id: 'plant-1', name: 'Apt', rooms: [] \}]; result coincide con expectedDatapoint.}{S}
\addTestI{Contesto \textit{AlarmConfigFormStateService}; caso '\textit{resolveDatapointForEdit} ritorna null se il catalogo piante fallisce': verifica che result è null.}{Esito atteso: result è null.}{S}
\addTestI{Contesto \textit{AlarmConfigStateService}; caso 'inizializza stato con alarms vuoti e nessun errore': verifica che await firstValueFrom(service.alarms\$) coincide con [] e await firstValueFrom(service.error\$) è null.}{Esito atteso: await firstValueFrom(service.alarms\$) coincide con []; await firstValueFrom(service.error\$) è null.}{S}
\addTestI{Contesto \textit{AlarmConfigStateService}; caso '\textit{loadAlarmRules} aggiorna stato locale quando API risponde con successo': verifica che apiStub.\textit{getAlarmRules} viene invocato 1 volte e alarms coincide con [alarmA, alarmB].}{Esito atteso: apiStub.\textit{getAlarmRules} viene invocato 1 volte; alarms coincide con [alarmA, alarmB]; alarms ha lunghezza 2; await firstValueFrom(service.error\$) è null.}{S}
\addTestI{Contesto \textit{AlarmConfigStateService}; caso '\textit{loadAlarmRules} espone messaggio errore e pulisce error precedente al retry': verifica che il valore di await firstValueFrom(service.error\$) è 'Errore durante il caricamento degli allarmi.' e await firstValueFrom(service.error\$) è null.}{Esito atteso: il valore di await firstValueFrom(service.error\$) è 'Errore durante il caricamento degli allarmi.'; await firstValueFrom(service.error\$) è null; await firstValueFrom(service.alarms\$) coincide con [alarmA].}{S}
\addTestI{Contesto \textit{AlarmConfigStateService}; caso 'getAlarmRuleById inoltra il valore di API in caso di successo': verifica che apiStub.\textit{getAlarmRule} viene invocato con 'alarm-1' e result coincide con alarmA.}{Esito atteso: apiStub.\textit{getAlarmRule} viene invocato con 'alarm-1'; result coincide con alarmA.}{S}
\addTestI{Contesto \textit{AlarmConfigStateService}; caso 'getAlarmRuleById in errore imposta messaggio e completa senza emissioni': verifica che il valore di emitted è false e il valore di await firstValueFrom(service.error\$) è 'Errore durante il recupero dell\textbackslash\{\}'allarme.'.}{Esito atteso: il valore di emitted è false; il valore di await firstValueFrom(service.error\$) è 'Errore durante il recupero dell\textbackslash\{\}'allarme.'.}{S}
\addTestI{Contesto \textit{AlarmConfigStateService}; caso '\textit{createAlarmRule} con form invalido non chiama API e imposta errore di validazione': verifica che il valore di emitted è false e apiStub.\textit{createAlarmRule} non viene invocato.}{Esito atteso: il valore di emitted è false; apiStub.\textit{createAlarmRule} non viene invocato; il valore di await firstValueFrom(service.error\$) è 'Dati del form non validi per la creazione dell\textbackslash\{\}'allarme.'.}{S}
\addTestI{Contesto \textit{AlarmConfigStateService}; caso '\textit{createAlarmRule} in errore API imposta messaggio e non altera lo stato locale': verifica che il valore di emitted è false e il valore di await firstValueFrom(service.error\$) è 'Errore durante la creazione dell\textbackslash\{\}'allarme.'.}{Esito atteso: il valore di emitted è false; il valore di await firstValueFrom(service.error\$) è 'Errore durante la creazione dell\textbackslash\{\}'allarme.'; await firstValueFrom(service.alarms\$) coincide con [alarmA]; \textit{refreshServiceStub}.\textit{requestRefresh} non viene invocato.}{S}
\addTestI{Contesto \textit{AlarmConfigStateService}; caso '\textit{createAlarmRule} pulisce l'errore precedente quando una creazione successiva va a buon fine': verifica che il valore di await firstValueFrom(service.error\$) è 'Dati del form non validi per la creazione dell\textbackslash\{\}'allarme.' e await firstValueFrom(service.error\$) è null.}{Esito atteso: il valore di await firstValueFrom(service.error\$) è 'Dati del form non validi per la creazione dell\textbackslash\{\}'allarme.'; await firstValueFrom(service.error\$) è null; await firstValueFrom(service.alarms\$) coincide con [alarmB].}{S}
\addTestI{Contesto \textit{AlarmConfigStateService}; caso '\textit{updateAlarmRule} mantiene il nome originale nel payload update e sostituisce l'allarme nello stato': verifica che apiStub.\textit{updateAlarmRule} viene invocato con 'alarm-1' e await firstValueFrom(service.alarms\$) coincide con [updatedAlarm, alarmB].}{Esito atteso: apiStub.\textit{updateAlarmRule} viene invocato con 'alarm-1'; await firstValueFrom(service.alarms\$) coincide con [updatedAlarm, alarmB]; \textit{refreshServiceStub}.\textit{requestRefresh} viene invocato 1 volte.}{S}
\addTestI{Contesto \textit{AlarmConfigStateService}; caso '\textit{updateAlarmRule} con form invalido non chiama API e imposta errore': verifica che il valore di emitted è false e apiStub.\textit{updateAlarmRule} non viene invocato.}{Esito atteso: il valore di emitted è false; apiStub.\textit{updateAlarmRule} non viene invocato; il valore di await firstValueFrom(service.error\$) è 'Dati del form non validi per l\textbackslash\{\}'aggiornamento dell\textbackslash\{\}'allarme.'.}{S}
\addTestI{Contesto \textit{AlarmConfigStateService}; caso '\textit{updateAlarmRule} in errore API imposta messaggio e mantiene lo stato invariato': verifica che il valore di emitted è false e il valore di await firstValueFrom(service.error\$) è 'Errore durante l\textbackslash\{\}'aggiornamento dell\textbackslash\{\}'allarme.'.}{Esito atteso: il valore di emitted è false; il valore di await firstValueFrom(service.error\$) è 'Errore durante l\textbackslash\{\}'aggiornamento dell\textbackslash\{\}'allarme.'; await firstValueFrom(service.alarms\$) coincide con [alarmA, alarmB]; \textit{refreshServiceStub}.\textit{requestRefresh} non viene invocato.}{S}
\addTestI{Contesto \textit{AlarmConfigStateService}; caso '\textit{updateAlarmRule} senza allarme locale usa il nome del form nel payload': verifica che apiStub.\textit{updateAlarmRule} viene invocato con 'missing-id'.}{Esito atteso: apiStub.\textit{updateAlarmRule} viene invocato con 'missing-id'.}{S}
\addTestI{Contesto \textit{AlarmConfigStateService}; caso '\textit{updateAlarmRule} sostituisce la riga locale anche se il backend ritorna un nuovo id': verifica che alarms coincide con [updatedWithNewId, alarmB] e il valore di alarms.some((alarm) => alarm.id === 'alarm-1') è false.}{Esito atteso: alarms coincide con [updatedWithNewId, alarmB]; il valore di alarms.some((alarm) => alarm.id === 'alarm-1') è false; il valore di alarms.some((alarm) => alarm.id === 'alarm-1-v2') è true; \textit{refreshServiceStub}.\textit{requestRefresh} viene invocato 1 volte.}{S}
\addTestI{Contesto \textit{AlarmConfigStateService}; caso '\textit{toggleEnabled} aggiorna enabled costruendo payload dai dati locali': verifica che apiStub.\textit{updateAlarmRule} viene invocato con 'alarm-1' e await firstValueFrom(service.alarms\$) coincide con [toggledAlarm].}{Esito atteso: apiStub.\textit{updateAlarmRule} viene invocato con 'alarm-1'; await firstValueFrom(service.alarms\$) coincide con [toggledAlarm]; \textit{refreshServiceStub}.\textit{requestRefresh} viene invocato 1 volte.}{S}
\addTestI{Contesto \textit{AlarmConfigStateService}; caso '\textit{toggleEnabled} su allarme assente non chiama API e imposta errore': verifica che il valore di emitted è false e apiStub.\textit{updateAlarmRule} non viene invocato.}{Esito atteso: il valore di emitted è false; apiStub.\textit{updateAlarmRule} non viene invocato; il valore di await firstValueFrom(service.error\$) è 'Allarme non trovato nello stato locale.'; \textit{refreshServiceStub}.\textit{requestRefresh} non viene invocato.}{S}
\addTestI{Contesto \textit{AlarmConfigStateService}; caso '\textit{toggleEnabled} in errore API imposta messaggio e mantiene lo stato invariato': verifica che il valore di emitted è false e il valore di await firstValueFrom(service.error\$) è 'Errore durante la modifica dello stato dell\textbackslash\{\}'allarme.'.}{Esito atteso: il valore di emitted è false; il valore di await firstValueFrom(service.error\$) è 'Errore durante la modifica dello stato dell\textbackslash\{\}'allarme.'; await firstValueFrom(service.alarms\$) coincide con [alarmA]; \textit{refreshServiceStub}.\textit{requestRefresh} non viene invocato.}{S}
\addTestI{Contesto \textit{AlarmConfigStateService}; caso '\textit{toggleEnabled} aggiorna subito lo stato locale anche se il backend ritorna un nuovo id': verifica che alarms coincide con [toggledWithNewId, alarmB] e il valore di alarms.some((alarm) => alarm.id === 'alarm-1') è false.}{Esito atteso: alarms coincide con [toggledWithNewId, alarmB]; il valore di alarms.some((alarm) => alarm.id === 'alarm-1') è false; il valore di alarms.some((alarm) => alarm.id === 'alarm-1-v2') è true.}{S}
\addTestI{Contesto \textit{AlarmConfigStateService}; caso '\textit{deleteAlarmRule} rimuove l'allarme locale quando API conferma': verifica che il valore di emitted è true e apiStub.\textit{deleteAlarmRule} viene invocato con 'alarm-1'.}{Esito atteso: il valore di emitted è true; apiStub.\textit{deleteAlarmRule} viene invocato con 'alarm-1'; await firstValueFrom(service.alarms\$) coincide con [alarmB]; \textit{refreshServiceStub}.\textit{requestRefresh} viene invocato 1 volte.}{S}
\addTestI{Contesto \textit{AlarmConfigStateService}; caso '\textit{deleteAlarmRule} in errore API imposta messaggio e mantiene lo stato invariato': verifica che il valore di emitted è false e il valore di await firstValueFrom(service.error\$) è 'Errore durante l\textbackslash\{\}'eliminazione dell\textbackslash\{\}'allarme.'.}{Esito atteso: il valore di emitted è false; il valore di await firstValueFrom(service.error\$) è 'Errore durante l\textbackslash\{\}'eliminazione dell\textbackslash\{\}'allarme.'; await firstValueFrom(service.alarms\$) coincide con [alarmA, alarmB]; \textit{refreshServiceStub}.\textit{requestRefresh} non viene invocato.}{S}
\addTestI{Contesto \textit{AlarmConfigStateService}; caso '\textit{deleteAlarmRule} pulisce errore precedente quando un delete successivo riesce': verifica che il valore di await firstValueFrom(service.error\$) è 'Errore durante l\textbackslash\{\}'eliminazione dell\textbackslash\{\}'allarme.' e await firstValueFrom(service.error\$) è null.}{Esito atteso: il valore di await firstValueFrom(service.error\$) è 'Errore durante l\textbackslash\{\}'eliminazione dell\textbackslash\{\}'allarme.'; await firstValueFrom(service.error\$) è null; await firstValueFrom(service.alarms\$) coincide con [alarmB]; \textit{refreshServiceStub}.\textit{requestRefresh} viene invocato 1 volte.}{S}
\addTestI{Contesto \textit{DeviceDatapointExtractionService}; caso '\textit{extractDeviceOptions} crea label stanza-dispositivo e normalizza enum duplicati': verifica che options ha lunghezza 1 e il valore di options[0].id è 'dev-1'.}{Esito atteso: options ha lunghezza 1; il valore di options[0].id è 'dev-1'; il valore di options[0].label è 'Soggiorno - Sensore UWB'; options[0].datapoints[0].enum coincide con ['Off', 'On'].}{S}
\addTestI{Contesto \textit{DeviceDatapointExtractionService}; caso '\textit{findReadableDatapoints} ritorna solo datapoint leggibili': verifica che datapoints.map((datapoint) => datapoint.id) coincide con ['dp-1', 'dp-2'] e il valore di datapoints.every((datapoint) => datapoint.readable) è true.}{Esito atteso: datapoints.map((datapoint) => datapoint.id) coincide con ['dp-1', 'dp-2']; il valore di datapoints.every((datapoint) => datapoint.readable) è true.}{S}
\addTestI{Contesto \textit{DeviceDatapointExtractionService}; caso '\textit{findReadableDatapoints} ritorna lista vuota se il device non esiste': verifica che datapoints coincide con [].}{Esito atteso: datapoints coincide con [].}{S}
\addTestI{Contesto \textit{DeviceDatapointExtractionService}; caso '\textit{getAllowedOperators} limita a uguale quando il datapoint ha enum': verifica che service.\textit{getAllowedOperators}(enumDatapoint) coincide con [ThresholdOperator.EQUAL\_TO] e service.\textit{getAllowedOperators}(numericDatapoint) coincide con [ ThresholdOperator.GREATER\_THAN, ThresholdOperator.GREATER\_THAN\_OR\_EQUAL, ThresholdOperator.L....}{Esito atteso: service.\textit{getAllowedOperators}(enumDatapoint) coincide con [ThresholdOperator.EQUAL\_TO]; service.\textit{getAllowedOperators}(numericDatapoint) coincide con [ ThresholdOperator.GREATER\_THAN, ThresholdOperator.GREATER\_THAN\_OR\_EQUAL, ThresholdOperator.L....}{S}
\addTestI{Contesto \textit{DeviceDatapointExtractionService}; caso 'isThresholdValueValid valida enum case-insensitive e soglie numeriche': verifica che il valore di service.isThresholdValueValid(enumDatapoint, 'on') è true e il valore di service.isThresholdValueValid(enumDatapoint, 'Unknown') è false.}{Esito atteso: il valore di service.isThresholdValueValid(enumDatapoint, 'on') è true; il valore di service.isThresholdValueValid(enumDatapoint, 'Unknown') è false; il valore di service.isThresholdValueValid(numericDatapoint, '10.5') è true; il valore di service.isThresholdValueValid(numericDatapoint, 'abc') è false.}{S}
\addTestI{Contesto \textit{DeviceDatapointExtractionService}; caso 'findDatapointByDeviceAndDatapointId trova e normalizza il datapoint dal catalogo plant': verifica che datapoint non è null e il valore di datapoint?.id è 'dp-lookup'.}{Esito atteso: datapoint non è null; il valore di datapoint?.id è 'dp-lookup'; datapoint?.enum coincide con ['Off', 'On'].}{S}
\addTestI{Contesto \textit{DeviceDatapointExtractionService}; caso 'findDatapointByDeviceAndDatapointId ritorna null se non trova corrispondenze': verifica che datapoint è null.}{Esito atteso: datapoint è null.}{S}
\addTestI{Contesto AlarmHistory feature integration; caso 'NO-RF-TAG inizializza pagina e service': verifica che component è valorizzato e alarmHistoryServiceStub.initialize viene invocato 1 volte.}{Esito atteso: component è valorizzato; alarmHistoryServiceStub.initialize viene invocato 1 volte.}{S}
\addTestI{Contesto AlarmHistory feature integration; caso 'NO-RF-TAG ordina righe per chiusura desc e applica fallback handler': verifica che il valore di rows[0].id è 'a-2' e il valore di rows[0].handlerUsername è 'sconosciuto'.}{Esito atteso: il valore di rows[0].id è 'a-2'; il valore di rows[0].handlerUsername è 'sconosciuto'; il valore di component.currentPage() è 2.}{S}
\addTestI{Contesto AlarmHistory feature integration; caso 'NO-RF-TAG delega paginazione next/previous al service': verifica che alarmHistoryServiceStub.nextPage viene invocato 1 volte e alarmHistoryServiceStub.previousPage viene invocato 1 volte.}{Esito atteso: alarmHistoryServiceStub.nextPage viene invocato 1 volte; alarmHistoryServiceStub.previousPage viene invocato 1 volte.}{S}
\addTestI{Contesto \textit{AlarmHistoryService}; caso 'initialize carica la prima pagina con offset 0': verifica che alarmApiStub.\textit{getResolvedAlarms} viene invocato con 99.}{Esito atteso: alarmApiStub.\textit{getResolvedAlarms} viene invocato con 99.}{S}
\addTestI{Contesto \textit{AlarmHistoryService}; caso 'initialize con lista vuota non imposta canGoNext': verifica che il valore di vm.alarms.length è 0 e vm.canGoNext è falsy.}{Esito atteso: il valore di vm.alarms.length è 0; vm.canGoNext è falsy.}{S}
\addTestI{Contesto \textit{AlarmHistoryService}; caso 'initialize con pageLimit+1 elementi imposta canGoNext e fa lo slice': verifica che vm.canGoNext è valorizzato e il valore di vm.alarms.length è 6.}{Esito atteso: vm.canGoNext è valorizzato; il valore di vm.alarms.length è 6.}{S}
\addTestI{Contesto \textit{AlarmHistoryService}; caso 'initialize con userId non numerico non chiama API': verifica che alarmApiStub.\textit{getResolvedAlarms} non viene invocato.}{Esito atteso: alarmApiStub.\textit{getResolvedAlarms} non viene invocato.}{S}
\addTestI{Contesto \textit{AlarmHistoryService}; caso 'initialize non propaga eccezioni al chiamante': verifica che () => service.initialize() non genera l'errore atteso.}{Esito atteso: () => service.initialize() non genera l'errore atteso.}{S}
\addTestI{Contesto \textit{AlarmHistoryService}; caso 'nextPage avanza alla pagina successiva': verifica che il valore di vm.currentPage è 2 e il valore di vm.pageOffset è 6.}{Esito atteso: il valore di vm.currentPage è 2; il valore di vm.pageOffset è 6; vm.canGoPrevious è valorizzato; vm.alarms.map((a) => a.id) coincide con ['resolved-10'].}{S}
\addTestI{Contesto \textit{AlarmHistoryService}; caso 'nextPage non fa nulla se canGoNext è false': verifica che alarmApiStub.\textit{getResolvedAlarms} viene invocato 1 volte.}{Esito atteso: alarmApiStub.\textit{getResolvedAlarms} viene invocato 1 volte.}{S}
\addTestI{Contesto \textit{AlarmHistoryService}; caso 'pagina successiva con meno di pageLimit elementi imposta canGoNext a false': verifica che vm.canGoNext è falsy e vm.alarms.map((a) => a.id) coincide con ['resolved-10', 'resolved-11'].}{Esito atteso: vm.canGoNext è falsy; vm.alarms.map((a) => a.id) coincide con ['resolved-10', 'resolved-11'].}{S}
\addTestI{Contesto \textit{AlarmHistoryService}; caso 'previousPage torna alla pagina precedente': verifica che il valore di vm.currentPage è 1 e il valore di vm.pageOffset è 0.}{Esito atteso: il valore di vm.currentPage è 1; il valore di vm.pageOffset è 0; vm.canGoPrevious è falsy; l'asserzione alarmApiStub.\textit{getResolvedAlarms} \textit{toHaveBeenNthCalledWith} 1 risulta soddisfatta.}{S}
\addTestI{Nel contesto \textit{AlarmHistoryService}, il caso 'previousPage non fa nulla se l'offset è già a zero' verifica che alarmApiStub.\textit{getResolvedAlarms} viene invocato 1 volte.}{Esito atteso: alarmApiStub.\textit{getResolvedAlarms} viene invocato 1 volte.}{S}
\addTestI{Contesto \textit{AlarmManagementService}; consolidated summary for RF44 initialization and pagination: verifica che initialize non carica automaticamente, admin role mostra tutti allarmi senza filtri, OSS usa endpoint unmanaged con userId, error handling con fallback messaggi, userId non-numeric blocks API call, nextPage incrementa offset con limit 6 quando pagina piena, previousPage decrementa fino a 0, pagina con <6 elementi imposta canGoNext=false, nextPage/previousPage non agiscono oltre limiti.}{Esito atteso: lazy init on demand; admin no-filter GET; OSS unmanaged endpoint; errors mapped to vm state; userId validation; pagination with limit 6; offset tracking; flag management.}{S}
\addTestI{Contesto \textit{AlarmManagementService}; consolidated summary for RF44 resolveAlarm and state management: verifica che resolveAlarm segue success flow (resolving true/false), VM reflects state transitions before/after resolve, local rows cleared e backend snapshot reload, automatic page down on empty non-first page, concurrent resolve requests tracked to completion, resolveError cleared on new resolve attempt, error path with network timeout, userId validation blocks resolve API.}{Esito atteso: resolving state toggled; VM aligned to backend post-resolve; pagination auto-adjust; concurrent requests joined; error clear on retry; timeout handled; userId validated.}{S}
\addTestI{Contesto \textit{AnalyticsApiService}; caso 'should be created': verifica che service è valorizzato.}{Esito atteso: service è valorizzato.}{S}
\addTestI{Contesto \textit{AnalyticsApiService} - \textit{getAllApartments}; caso 'should fetch all apartments and map them correctly': verifica che il valore di apartments.length è 2 e apartments[0] coincide con \{ id: '1', name: 'Apartment 1' \}.}{Esito atteso: il valore di apartments.length è 2; apartments[0] coincide con \{ id: '1', name: 'Apartment 1' \}; l'asserzione apartments[0] \textit{toHaveProperty} 'extraField' risulta non soddisfatta; il valore di req.request.method è 'GET'.}{S}
\addTestI{Contesto \textit{AnalyticsApiService} - \textit{getAnalytics}; caso 'should fetch analytics and map the complex DTO structure': verifica che il valore di result.apartmentId è apartmentId e il valore di result.analyticsInfo.length è 1.}{Esito atteso: il valore di result.apartmentId è apartmentId; il valore di result.analyticsInfo.length è 1; il valore di info.title è 'Test Chart'; info.suggestions.messages coincide con ['Suggestion 1'].}{S}
\addTestI{Contesto \textit{ApartmentApiService}; caso 'RF83-OBL abilita un appartamento nel Sistema': verifica che il valore di request.request.method è 'PATCH' e request.request.body coincide con \{\}.}{Esito atteso: il valore di request.request.method è 'PATCH'; request.request.body coincide con \{\}.}{S}
\addTestI{Contesto \textit{ApartmentApiService}; caso 'RF84-OBL disabilita un appartamento nel Sistema': verifica che il valore di request.request.method è 'PATCH' e request.request.body coincide con \{\}.}{Esito atteso: il valore di request.request.method è 'PATCH'; request.request.body coincide con \{\}.}{S}
\addTestI{Contesto \textit{ApartmentApiService}; caso 'RF72-OBL \textit{getAvailableApartments} ordina per nome e deduplica plant': verifica che il valore di request.request.method è 'GET' e apartments coincide con [ \{ id: 'p1', name: 'Alfa' \}, \{ id: 'p2', name: 'Zeta' \}, ].}{Esito atteso: il valore di request.request.method è 'GET'; apartments coincide con [ \{ id: 'p1', name: 'Alfa' \}, \{ id: 'p2', name: 'Zeta' \}, ].}{S}
\addTestI{Contesto \textit{ApartmentApiService}; caso 'RF72-OBL \textit{getAllPlants} supporta risposta wrappata in data[]': verifica che plants ha lunghezza 2 e il valore di plants[1].name è 'p10'.}{Esito atteso: plants ha lunghezza 2; il valore di plants[1].name è 'p10'.}{S}
\addTestI{Contesto \textit{ApartmentApiService}; caso 'RF72-OBL \textit{getCurrentApartment} usa activePlantId salvato e mappa apartment view': verifica che il valore di apartment.id è 'p2' e apartment.rooms ha lunghezza 1.}{Esito atteso: il valore di apartment.id è 'p2'; apartment.rooms ha lunghezza 1; apartment.rooms[0].devices ha lunghezza 1; il valore di apartment.rooms[0].devices[0].id è 'd-main'.}{S}
\addTestI{Contesto \textit{ApartmentApiService}; caso 'RF72-OBL \textit{getCurrentApartment} fallisce se non esistono plant disponibili': verifica che dopo rigetto, apartmentPromise genera l'errore atteso.}{Esito atteso: dopo rigetto, apartmentPromise genera l'errore atteso.}{S}
\addTestI{Contesto \textit{ApartmentApiService}; caso 'RF72-OBL setActivePlantId non emette evento se plant invariato': verifica che dispatchSpy non viene invocato e il valore di globalThis.localStorage.\textit{getItem}('activePlantId') è 'p1'.}{Esito atteso: dispatchSpy non viene invocato; il valore di globalThis.localStorage.\textit{getItem}('activePlantId') è 'p1'.}{S}
\addTestI{Contesto \textit{ApartmentApiService}; caso 'RF72-OBL \textit{getAllPlants} sincronizza wardIds nel realtime cache con utente loggato': verifica che wardRealtimeCacheStub.mergeWardIds viene invocato con 'u-1'.}{Esito atteso: wardRealtimeCacheStub.mergeWardIds viene invocato con 'u-1'.}{S}
\addTestI{Contesto \textit{NavService}; caso 'viene creato': verifica che service è valorizzato.}{Esito atteso: service è valorizzato.}{S}
\addTestI{Contesto \textit{NavService}; caso 'restituisce tutti i NavItems per AMMINISTRATORE': verifica che items.length è maggiore di 0 e il valore di items.some(i => i.label === "Dashboard") è true.}{Esito atteso: items.length è maggiore di 0; il valore di items.some(i => i.label === "Dashboard") è true; il valore di items.some(i => i.route.includes('dashboard')) è true; il valore di items.some(i => i.label === "Configurazione Allarmi") è true.}{S}
\addTestI{Contesto \textit{NavService}; caso 'restituisce solo i NavItems filtrati per OPERATORE SANITARIO': verifica che OSSItems.length è maggiore di 0 e il valore di roleReqItems.length è 0.}{Esito atteso: OSSItems.length è maggiore di 0; il valore di roleReqItems.length è 0; il valore di items.some(i => i.label === "Dashboard") è true; il valore di items.some(i => i.route.includes('dashboard')) è true.}{S}
\addTestI{Contesto \textit{DeviceApiService}; consolidated summary for RF79-RF81 covering device data retrieval, filtering, deduplication and write operations: verifica che getWritableEndpointRows GET /plant/{id}, mappa datapoint per stanza, POST /api/device con payload {datapointId, value}, propaga errori write, resolveActivePlantId GET /plant/all se mancante, esclude energy-measure devices, deduplica per nome mantenendo priorità, getCurrentValuePointsByDeviceIds GET valori, continua merge se un device call fallisce, error handling per timeout.}{Esito atteso: room mapping by id; datapoint array; POST write correct; error propagation; plant resolution; device filtering; deduplication priority; value map retrieval; resilience on partial failures.}{S}
\addTestI{Contesto \textit{NotificationService}; caso 'espone vm iniziale vuoto quando storico e stream in-session sono vuoti': verifica che vm coincide con \{ notifications: [], unreadCount: 0, \}.}{Esito atteso: vm coincide con \{ notifications: [], unreadCount: 0, \}.}{S}
\addTestI{Contesto \textit{NotificationService}; caso 'combina storico e in-session, deduplica per notificationId e ordina per sentAt desc': verifica che il valore di vm.unreadCount è 2 e vm.notifications coincide con [sharedInSession, inSessionOnly, historyOlder].}{Esito atteso: il valore di vm.unreadCount è 2; vm.notifications coincide con [sharedInSession, inSessionOnly, historyOlder]; vm.notifications ha lunghezza 3.}{S}
\addTestI{Contesto \textit{NotificationService}; caso 'in errore sullo storico applica fallback a [] e mantiene i dati in-session': verifica che vm.notifications coincide con [inSessionOnly] e il valore di vm.unreadCount è 1.}{Esito atteso: vm.notifications coincide con [inSessionOnly]; il valore di vm.unreadCount è 1.}{S}
\addTestI{Contesto \textit{NotificationService}; caso 'propaga nuovi unreadCount senza alterare la lista notifiche': verifica che vm.notifications coincide con [historyOlder] e il valore di vm.unreadCount è 5.}{Esito atteso: vm.notifications coincide con [historyOlder]; il valore di vm.unreadCount è 5.}{S}
\addTestI{Contesto \textit{NotificationService}; caso 'aggiorna il vm quando arrivano nuove notifiche in-session dopo la prima emissione': verifica che vm.notifications coincide con [inSessionOnly, historyOlder] e il valore di vm.unreadCount è 0.}{Esito atteso: vm.notifications coincide con [inSessionOnly, historyOlder]; il valore di vm.unreadCount è 0.}{S}
\addTestI{Contesto \textit{NotificationService}; caso 'deduplica tenendo la versione in-session quando notificationId è presente anche nello storico': verifica che vm.notifications coincide con [sharedInSession] e vm.notifications ha lunghezza 1.}{Esito atteso: vm.notifications coincide con [sharedInSession]; vm.notifications ha lunghezza 1.}{S}
\addTestI{Contesto \textit{NotificationService}; caso 'gestisce timestamp non validi senza lanciare errori e li ordina in fondo': verifica che vm.notifications coincide con [historyOlder, invalidTimestampNotification].}{Esito atteso: vm.notifications coincide con [historyOlder, invalidTimestampNotification].}{S}
\addTestI{Contesto \textit{NotificationService}; caso 'riusa la stessa pipeline vm con più subscriber senza richiamare di nuovo lo storico': verifica che vmA coincide con vmB e notificationApiStub.\textit{getNotificationsHistory} viene invocato 1 volte.}{Esito atteso: vmA coincide con vmB; notificationApiStub.\textit{getNotificationsHistory} viene invocato 1 volte.}{S}
\addTestI{Contesto \textit{NotificationService}; caso 'mantiene uno snapshot coerente per subscriber tardivi senza ricaricare lo storico': verifica che vmLateSubscriber.notifications coincide con [inSessionOnly, historyOlder] e notificationApiStub.\textit{getNotificationsHistory} viene invocato 1 volte.}{Esito atteso: vmLateSubscriber.notifications coincide con [inSessionOnly, historyOlder]; notificationApiStub.\textit{getNotificationsHistory} viene invocato 1 volte.}{S}
\addTestI{Contesto \textit{NotificationService}; caso 'non muta gli array sorgente ricevuti da API e \textit{AlarmStateService}': verifica che historicList coincide con [sharedHistoric, historyOlder] e inSessionList coincide con [sharedInSession, inSessionOnly].}{Esito atteso: historicList coincide con [sharedHistoric, historyOlder]; inSessionList coincide con [sharedInSession, inSessionOnly].}{S}
\addTestI{Contesto \textit{NotificationService}; caso 'deduplica anche duplicati interni alla stessa sorgente mantenendo l ultimo elemento': verifica che vm.notifications coincide con [duplicatedSessionB, duplicatedHistoryB] e vm.notifications ha lunghezza 2.}{Esito atteso: vm.notifications coincide con [duplicatedSessionB, duplicatedHistoryB]; vm.notifications ha lunghezza 2.}{S}
\addTestI{Contesto \textit{NotificationService}; caso 'mantiene coerenza: unreadCount segue lo stream dedicato anche quando la deduplica riduce la lista': verifica che vm.notifications ha lunghezza 1 e il valore di vm.unreadCount è 5.}{Esito atteso: vm.notifications ha lunghezza 1; il valore di vm.unreadCount è 5.}{S}
\addTestI{Contesto \textit{NotificationService}; caso 'filtra le notifiche dismissate localmente dall archivio unificato': verifica che vm.notifications.map((notification) => notification.notificationId) coincide con [ 'session-1', 'shared-id', ].}{Esito atteso: vm.notifications.map((notification) => notification.notificationId) coincide con [ 'session-1', 'shared-id', ].}{S}
\addTestI{Contesto \textit{NotificationService}; caso '\textit{removeNotification} delega dismiss + remove allo stato realtime': verifica che alarmStateStub.\textit{dismissNotification} viene invocato con 'n-100' e alarmStateStub.\textit{removeNotification} viene invocato con 'n-100'.}{Esito atteso: alarmStateStub.\textit{dismissNotification} viene invocato con 'n-100'; alarmStateStub.\textit{removeNotification} viene invocato con 'n-100'.}{S}
\addTestI{Contesto \textit{NotificationService}; caso '\textit{clearAllNotifications} delega dismiss multiplo + clear realtime': verifica che alarmStateStub.\textit{dismissNotifications} viene invocato con [ historyOlder.notificationId, inSessionOnly.notificationId, ] e alarmStateStub.\textit{clearNotifications} viene invocato 1 volte.}{Esito atteso: alarmStateStub.\textit{dismissNotifications} viene invocato con [ historyOlder.notificationId, inSessionOnly.notificationId, ]; alarmStateStub.\textit{clearNotifications} viene invocato 1 volte.}{S}
\addTestI{Contesto \textit{UserApiService}; consolidated summary for user CRUD and password decoding: verifica che getUsers GET /users, createUser POST /users con dto e decodifica tempPassword base64 al result, deleteUser DELETE /users/:id con id numerico, decodeTempPassword decodifica/mantiene base64 valida, empty string preserve, invalid base64 mantiene originale, atob exception fallback, whitespace trim per decode, spazi marginali gestiti.}{Esito atteso: GET users list; POST create returns decoded tempPassword; DELETE removes; base64 decoding works; invalid inputs preserved; exception handling fallback; whitespace handling correct.}{S}
\addTestI{Contesto \textit{AssignmentOperationsService}; caso '\textit{assignOperator} esegue operazione e poi ricarica i wards': verifica che result è undefined e apiStub.\textit{assignOperatorToWard} viene invocato con 1.}{Esito atteso: result è undefined; apiStub.\textit{assignOperatorToWard} viene invocato con 1; l'asserzione apiStub.\textit{getWards} \textit{toHaveBeenCalledOnce} risulta soddisfatta; apiStub.getPlantsByWardId viene invocato con 1.}{S}
\addTestI{Contesto \textit{AssignmentOperationsService}; caso '\textit{removeOperator} in errore espone errore testo da \textit{HttpErrorResponse}.error': verifica che il valore di emitted è false e storeStub.\textit{setError} viene invocato con 'Operatore non assegnato'.}{Esito atteso: il valore di emitted è false; storeStub.\textit{setError} viene invocato con 'Operatore non assegnato'; storeStub.\textit{setError} viene invocato 1 volte; apiStub.\textit{getWards} non viene invocato.}{S}
\addTestI{Contesto \textit{AssignmentOperationsService}; caso '\textit{getAvailablePlantsForWard} filtra i plant già assegnati in qualunque ward': verifica che result coincide con [\{ id: '203', name: 'App. 203' \}] e l'asserzione apiStub.\textit{getAvailablePlants} \textit{toHaveBeenCalledOnce} risulta soddisfatta.}{Esito atteso: result coincide con [\{ id: '203', name: 'App. 203' \}]; l'asserzione apiStub.\textit{getAvailablePlants} \textit{toHaveBeenCalledOnce} risulta soddisfatta.}{S}
\addTestI{Contesto \textit{AssignmentOperationsService}; caso 'dopo rimozione un appartamento torna disponibile': verifica che result coincide con [\{ id: '202', name: 'App. 202' \}] e result coincide con [ \{ id: '201', name: 'App. 201' \}, \{ id: '202', name: 'App. 202' \}, ].}{Esito atteso: result coincide con [\{ id: '202', name: 'App. 202' \}]; result coincide con [ \{ id: '201', name: 'App. 201' \}, \{ id: '202', name: 'App. 202' \}, ].}{S}
\addTestI{Contesto \textit{AssignmentOperationsService}; caso 'dopo eliminazione reparto, gli appartamenti del reparto eliminato tornano disponibili': verifica che result coincide con [\{ id: '203', name: 'App. 203' \}] e result coincide con [ \{ id: '201', name: 'App. 201' \}, \{ id: '203', name: 'App. 203' \}, ].}{Esito atteso: result coincide con [\{ id: '203', name: 'App. 203' \}]; result coincide con [ \{ id: '201', name: 'App. 201' \}, \{ id: '203', name: 'App. 203' \}, ].}{S}
\addTestI{Contesto \textit{AssignmentOperationsService}; caso '\textit{assignPlant} in errore usa \textit{HttpErrorResponse}.message quando non c'è error string': verifica che messageArg.length è maggiore di 0 e apiStub.\textit{getWards} non viene invocato.}{Esito atteso: messageArg.length è maggiore di 0; apiStub.\textit{getWards} non viene invocato; storeStub.\textit{setWards} non viene invocato.}{S}
\addTestI{Contesto \textit{AssignmentOperationsService}; caso '\textit{removePlant} in errore non-http usa fallback message': verifica che storeStub.\textit{setError} viene invocato con 'Operazione di assegnazione non riuscita.' e storeStub.\textit{setError} viene invocato 1 volte.}{Esito atteso: storeStub.\textit{setError} viene invocato con 'Operazione di assegnazione non riuscita.'; storeStub.\textit{setError} viene invocato 1 volte; apiStub.\textit{getWards} non viene invocato.}{S}
\addTestI{Contesto \textit{WardApiService}; consolidated summary for RF73-RF76 covering ward and relationship management endpoints: verifica che getWards GET /api/wards, getOperatorsByWardId GET /api/wards-users-relationships/:wardId, getAvailableOperators GET /users, getPlantsByWardId GET /api/wards-plants-relationships/:wardId, getAvailablePlants GET /api/plant/all, createWard POST /api/wards con dto, updateWard PUT /api/wards/:id, deleteWard DELETE /api/wards/:id, assignOperatorToWard POST /api/wards-users-relationships con wardId userId, removeOperatorFromWard DELETE /api/wards-users-relationships/:wardId/:userId, assignPlantToWard POST /api/wards-plants-relationships, removePlantFromWard DELETE /api/wards-plants-relationships/:plantId, mapping compatibility.}{Esito atteso: GET /wards lista; GET operators by ward; GET available operators; GET plants by ward; GET all available plants; POST create returns ward; PUT update preserves data; DELETE removes ward; POST relationship created; DELETE relationship removed; assignment endpoints work; backward compatibility maintained.}{S}
\addTestI{Contesto \textit{WardOperationsService}; caso '\textit{loadWards} aggiorna store e termina senza valore': verifica che result è undefined e l'asserzione apiStub.\textit{getWards} \textit{toHaveBeenCalledOnce} risulta soddisfatta.}{Esito atteso: result è undefined; l'asserzione apiStub.\textit{getWards} \textit{toHaveBeenCalledOnce} risulta soddisfatta; apiStub.getPlantsByWardId viene invocato con 1; apiStub.getOperatorsByWardId viene invocato con 1.}{S}
\addTestI{Contesto \textit{WardOperationsService}; caso '\textit{loadWards} con lista vuota non richiama le relationship API': verifica che l'asserzione apiStub.\textit{getWards} \textit{toHaveBeenCalledOnce} risulta soddisfatta e apiStub.getPlantsByWardId non viene invocato.}{Esito atteso: l'asserzione apiStub.\textit{getWards} \textit{toHaveBeenCalledOnce} risulta soddisfatta; apiStub.getPlantsByWardId non viene invocato; apiStub.getOperatorsByWardId non viene invocato; storeStub.\textit{setWards} viene invocato con [].}{S}
\addTestI{Contesto \textit{WardOperationsService}; caso '\textit{createWard} con 409 imposta messaggio funzionale': verifica che il valore di emitted è false e storeStub.addWard non viene invocato.}{Esito atteso: il valore di emitted è false; storeStub.addWard non viene invocato; storeStub.\textit{setError} viene invocato con 'Esiste già un reparto con questo nome.'; storeStub.\textit{setError} viene invocato 1 volte.}{S}
\addTestI{Contesto \textit{WardOperationsService}; caso '\textit{createWard} normalizza payload minimale mantenendo il modello Ward coerente': verifica che storeStub.addWard viene invocato con \{ id: 11, name: 'Oncologia', apartments: [], operators: [], \} e storeStub.\textit{setLoading} viene invocato con false.}{Esito atteso: storeStub.addWard viene invocato con \{ id: 11, name: 'Oncologia', apartments: [], operators: [], \}; storeStub.\textit{setLoading} viene invocato con false.}{S}
\addTestI{Contesto \textit{WardOperationsService}; caso '\textit{createWard} preserva apartments/operators dallo snapshot se il backend non li restituisce': verifica che storeStub.addWard viene invocato con \{ id: 1, name: 'Cardiologia nuova', apartments: ward.apartments, operators: ward.operators, \}.}{Esito atteso: storeStub.addWard viene invocato con \{ id: 1, name: 'Cardiologia nuova', apartments: ward.apartments, operators: ward.operators, \}.}{S}
\addTestI{Contesto \textit{WardOperationsService}; caso '\textit{updateWard} sostituisce il ward nel store': verifica che apiStub.\textit{updateWard} viene invocato con 1 e storeStub.\textit{replaceWard} viene invocato con updatedWard.}{Esito atteso: apiStub.\textit{updateWard} viene invocato con 1; storeStub.\textit{replaceWard} viene invocato con updatedWard; storeStub.\textit{replaceWard} viene invocato 1 volte; storeStub.\textit{setLoading} viene invocato con false.}{S}
\addTestI{Contesto \textit{WardOperationsService}; caso '\textit{updateWard} preserva apartments/operators quando backend restituisce solo id e name': verifica che storeStub.\textit{replaceWard} viene invocato con \{ id: 1, name: 'Cardiologia X', apartments: ward.apartments, operators: ward.operators, \} e storeStub.\textit{setLoading} viene invocato con false.}{Esito atteso: storeStub.\textit{replaceWard} viene invocato con \{ id: 1, name: 'Cardiologia X', apartments: ward.apartments, operators: ward.operators, \}; storeStub.\textit{setLoading} viene invocato con false.}{S}
\addTestI{Contesto \textit{WardOperationsService}; caso '\textit{updateWard} gestisce errore HTTP con messaggio testuale del server': verifica che storeStub.\textit{setError} viene invocato con 'update failed' e storeStub.\textit{replaceWard} non viene invocato.}{Esito atteso: storeStub.\textit{setError} viene invocato con 'update failed'; storeStub.\textit{replaceWard} non viene invocato.}{S}
\addTestI{Contesto \textit{WardOperationsService}; caso '\textit{deleteWard} rimuove il ward dal store': verifica che apiStub.\textit{deleteWard} viene invocato con 1 e storeStub.\textit{removeWard} viene invocato con 1.}{Esito atteso: apiStub.\textit{deleteWard} viene invocato con 1; storeStub.\textit{removeWard} viene invocato con 1; storeStub.\textit{removeWard} viene invocato 1 volte; storeStub.\textit{setLoading} viene invocato con false.}{S}
\addTestI{Contesto \textit{WardOperationsService}; caso '\textit{deleteWard} usa il fallback generico quando l'errore non è strutturato': verifica che storeStub.\textit{setError} viene invocato con 'Operazione sui reparti non riuscita.' e storeStub.\textit{removeWard} non viene invocato.}{Esito atteso: storeStub.\textit{setError} viene invocato con 'Operazione sui reparti non riuscita.'; storeStub.\textit{removeWard} non viene invocato.}{S}
\addTestI{Contesto \textit{WardOperationsService}; caso '\textit{createWard} usa il fallback generico quando l'errore non è strutturato': verifica che storeStub.\textit{setError} viene invocato con 'Operazione sui reparti non riuscita.'.}{Esito atteso: storeStub.\textit{setError} viene invocato con 'Operazione sui reparti non riuscita.'.}{S}
\addTestI{Contesto \textit{WardOperationsService}; caso 'in caso di errore generico usa il fallback message': verifica che storeStub.\textit{setError} viene invocato con 'Operazione sui reparti non riuscita.' e storeStub.\textit{setError} viene invocato 1 volte.}{Esito atteso: storeStub.\textit{setError} viene invocato con 'Operazione sui reparti non riuscita.'; storeStub.\textit{setError} viene invocato 1 volte; storeStub.\textit{setWards} non viene invocato.}{S}
\addTestI{Contesto \textit{WardOperationsService}; caso '\textit{loadWards} normalizza apartments dal dto': verifica che storeStub.\textit{setWards} viene invocato con [ \{ ...hydratedWard, apartments: [\{ id: '101', name: 'App. 101' \}], \}, ].}{Esito atteso: storeStub.\textit{setWards} viene invocato con [ \{ ...hydratedWard, apartments: [\{ id: '101', name: 'App. 101' \}], \}, ].}{S}
\addTestI{Contesto \textit{WardStore}; caso 'inizializza stato vuoto': verifica che await firstValueFrom(store.wards\$) coincide con [] e il valore di await firstValueFrom(store.isLoading\$) è false.}{Esito atteso: await firstValueFrom(store.wards\$) coincide con []; il valore di await firstValueFrom(store.isLoading\$) è false; await firstValueFrom(store.error\$) è null.}{S}
\addTestI{Contesto \textit{WardStore}; caso '\textit{setWards} sovrascrive i wards': verifica che wards coincide con [wardA, wardB] e wards ha lunghezza 2.}{Esito atteso: wards coincide con [wardA, wardB]; wards ha lunghezza 2.}{S}
\addTestI{Contesto \textit{WardStore}; caso 'addWard aggiunge il reparto e resetta error/loading': verifica che wards coincide con [wardA] e wards ha lunghezza 1.}{Esito atteso: wards coincide con [wardA]; wards ha lunghezza 1; il valore di await firstValueFrom(store.isLoading\$) è false; await firstValueFrom(store.error\$) è null.}{S}
\addTestI{Contesto \textit{WardStore}; caso '\textit{replaceWard} aggiorna il reparto target': verifica che wards coincide con [wardA, renamed] e wards ha lunghezza 2.}{Esito atteso: wards coincide con [wardA, renamed]; wards ha lunghezza 2.}{S}
\addTestI{Contesto \textit{WardStore}; caso '\textit{removeWard} elimina il reparto target': verifica che wards coincide con [wardB] e wards ha lunghezza 1.}{Esito atteso: wards coincide con [wardB]; wards ha lunghezza 1.}{S}
\addTestI{Contesto \textit{WardStore}; caso '\textit{patchPlant} aggiorna solo l appartamento richiesto': verifica che wards ha lunghezza 2 e wards[0].apartments[0] coincide con wardA.apartments[0].}{Esito atteso: wards ha lunghezza 2; wards[0].apartments[0] coincide con wardA.apartments[0]; wards[1].apartments[0] coincide con \{ id: '102', name: 'App. 102A' \}.}{S}
\addTestI{Contesto \textit{AlarmToggleSwitchComponent}; caso 'mostra lo stato iniziale non attivo': verifica che il valore di button.\textit{getAttribute}('aria-pressed') è 'false' e component.\textit{trackClass}() contiene 'border-slate-500'.}{Esito atteso: il valore di button.\textit{getAttribute}('aria-pressed') è 'false'; component.\textit{trackClass}() contiene 'border-slate-500'; component.thumbClass() contiene 'translate-x-0.5'.}{S}
\addTestI{Contesto \textit{AlarmToggleSwitchComponent}; caso 'mostra lo stato attivo e la classe corretta': verifica che il valore di button.\textit{getAttribute}('aria-pressed') è 'true' e component.\textit{trackClass}() contiene 'border-yellow-700'.}{Esito atteso: il valore di button.\textit{getAttribute}('aria-pressed') è 'true'; component.\textit{trackClass}() contiene 'border-yellow-700'; component.thumbClass() contiene 'translate-x-4'.}{S}
\addTestI{Contesto \textit{AlarmToggleSwitchComponent}; caso 'emette il nuovo valore quando viene togglato': verifica che emitSpy viene invocato con true e emitSpy viene invocato 1 volte.}{Esito atteso: emitSpy viene invocato con true; emitSpy viene invocato 1 volte.}{S}
\addTestI{Contesto \textit{AlarmToggleSwitchComponent}; caso 'non emette quando e disabilitato': verifica che emitSpy non viene invocato e il valore di (fixture.nativeElement.querySelector('button') as HTMLButtonElement).disabled è true.}{Esito atteso: emitSpy non viene invocato; il valore di (fixture.nativeElement.querySelector('button') as HTMLButtonElement).disabled è true.}{S}
\addTestI{Contesto \textit{AlarmToggleSwitchComponent}; caso 'usa aria-label custom quando fornita': verifica che il valore di (fixture.nativeElement.querySelector('button') as HTMLButtonElement).\textit{getAttribute}('aria-label') è 'Attiva allarme'.}{Esito atteso: il valore di (fixture.nativeElement.querySelector('button') as HTMLButtonElement).\textit{getAttribute}('aria-label') è 'Attiva allarme'.}{S}
\addTestI{Contesto \textit{AlarmConfigFormComponent}; caso 'crea il componente': verifica che component è valorizzato.}{Esito atteso: component è valorizzato.}{S}
\addTestI{Contesto \textit{AlarmConfigFormComponent}; caso '\textit{buildForm} applica i validatori required ai campi richiesti': verifica che il valore di component.form.controls.plantId.invalid è true e il valore di component.form.controls.deviceId.invalid è true.}{Esito atteso: il valore di component.form.controls.plantId.invalid è true; il valore di component.form.controls.deviceId.invalid è true; il valore di component.form.controls.datapointId.invalid è true; il valore di component.form.controls.priority.invalid è true.}{S}
\addTestI{Contesto \textit{AlarmConfigFormComponent}; caso 'in create mode mantiene il dispositivo bloccato finché non viene selezionato un plant': verifica che il valore di component.form.controls.deviceId.disabled è true e apartmentApiStub.loadApartmentViewForPlantId viene invocato con 'plant-1'.}{Esito atteso: il valore di component.form.controls.deviceId.disabled è true; apartmentApiStub.loadApartmentViewForPlantId viene invocato con 'plant-1'; il valore di component.form.controls.deviceId.enabled è true; il valore di component.deviceOptions()[0]?.id è 'sensor-1'.}{S}
\addTestI{Contesto \textit{AlarmConfigFormComponent}; caso 'mostra solo datapoint leggibili dopo la selezione dispositivo': verifica che component.datapointOptions().map((datapoint) => datapoint.id) coincide con [ 'dp-readable-1', 'dp-readable-2', ] e il valore di component.datapointOptions().every((datapoint) => datapoint.readable) è true.}{Esito atteso: component.datapointOptions().map((datapoint) => datapoint.id) coincide con [ 'dp-readable-1', 'dp-readable-2', ]; il valore di component.datapointOptions().every((datapoint) => datapoint.readable) è true.}{S}
\addTestI{Contesto \textit{AlarmConfigFormComponent}; caso 'in create mode non mostra il campo posizione read-only': verifica che positionInput è null.}{Esito atteso: positionInput è null.}{S}
\addTestI{Contesto \textit{AlarmConfigFormComponent}; caso 'in edit mode mostra il campo posizione read-only valorizzato': verifica che positionInput non è null e il valore di positionInput?.value è 'Appartamento 2 - Ingresso - Sensore porta'.}{Esito atteso: positionInput non è null; il valore di positionInput?.value è 'Appartamento 2 - Ingresso - Sensore porta'.}{S}
\addTestI{Contesto \textit{AlarmConfigFormComponent}; caso 'in edit mode non mostra il campo datapoint nel form': verifica che datapointInput è null.}{Esito atteso: datapointInput è null.}{S}
\addTestI{Contesto \textit{AlarmConfigFormComponent}; caso 'in edit mode mostra thresholdOperator come campo readonly': verifica che thresholdOperatorReadonly non è null e il valore di thresholdOperatorReadonly?.value è ThresholdOperator.EQUAL\_TO.}{Esito atteso: thresholdOperatorReadonly non è null; il valore di thresholdOperatorReadonly?.value è ThresholdOperator.EQUAL\_TO; thresholdOperatorSelect è null.}{S}
\addTestI{Contesto \textit{AlarmConfigFormComponent}; caso 'in edit mode limita gli operatori al datapoint risolto e resetta operatori non supportati': verifica che component.thresholdOperatorOptions() coincide con [ThresholdOperator.EQUAL\_TO] e component.form.controls.thresholdOperator.value è null.}{Esito atteso: component.thresholdOperatorOptions() coincide con [ThresholdOperator.EQUAL\_TO]; component.form.controls.thresholdOperator.value è null.}{S}
\addTestI{Contesto \textit{AlarmConfigFormComponent}; caso 'in edit mode valida thresholdValue in base al datapoint risolto': verifica che il valore di component.form.controls.thresholdValue.errors?.['invalidEnumThreshold'] è true.}{Esito atteso: il valore di component.form.controls.thresholdValue.errors?.['invalidEnumThreshold'] è true.}{S}
\addTestI{Contesto \textit{AlarmConfigFormComponent}; caso '\textit{onSubmit} emette submittedForm in create mode con form valido': verifica che emitSpy viene invocato con validFormValue e emitSpy viene invocato 1 volte.}{Esito atteso: emitSpy viene invocato con validFormValue; emitSpy viene invocato 1 volte.}{S}
\addTestI{Contesto \textit{AlarmConfigFormComponent}; caso '\textit{onSubmit} non invia se il form e invalido': verifica che emitSpy non viene invocato.}{Esito atteso: emitSpy non viene invocato.}{S}
\addTestI{Contesto \textit{AlarmConfigFormComponent}; caso '\textit{onCancel} emette evento cancelled': verifica che emitSpy viene invocato 1 volte.}{Esito atteso: emitSpy viene invocato 1 volte.}{S}
\addTestI{Contesto \textit{AlarmConfigFormComponent}; caso 'espone opzioni enum per priorita e operatore': verifica che component.priorityOptions coincide con [ AlarmPriority.WHITE, AlarmPriority.GREEN, AlarmPriority.ORANGE, AlarmPriority.RED, ] e component.thresholdOperatorOptions() coincide con [ ThresholdOperator.GREATER\_THAN, ThresholdOperator.GREATER\_THAN\_OR\_EQUAL, ThresholdOperator.L....}{Esito atteso: component.priorityOptions coincide con [ AlarmPriority.WHITE, AlarmPriority.GREEN, AlarmPriority.ORANGE, AlarmPriority.RED, ]; component.thresholdOperatorOptions() coincide con [ ThresholdOperator.GREATER\_THAN, ThresholdOperator.GREATER\_THAN\_OR\_EQUAL, ThresholdOperator.L....}{S}
\addTestI{Contesto \textit{AlarmConfigFormComponent}; caso 'limita operatori e valore soglia quando il datapoint ha enum': verifica che component.thresholdOperatorOptions() coincide con [ThresholdOperator.EQUAL\_TO] e il valore di component.form.controls.thresholdOperator.invalid è true.}{Esito atteso: component.thresholdOperatorOptions() coincide con [ThresholdOperator.EQUAL\_TO]; il valore di component.form.controls.thresholdOperator.invalid è true; il valore di component.form.controls.thresholdValue.invalid è true; il valore di component.form.controls.thresholdValue.errors?.['invalidEnumThreshold'] è true.}{S}
\addTestI{Contesto \textit{AlarmConfigFormComponent}; caso 'accetta valore soglia booleano quando operatore è uguale': verifica che il valore di component.form.controls.thresholdOperator.valid è true e il valore di component.form.controls.thresholdValue.valid è true.}{Esito atteso: il valore di component.form.controls.thresholdOperator.valid è true; il valore di component.form.controls.thresholdValue.valid è true.}{S}
\addTestI{Contesto \textit{AlarmConfigFormComponent}; caso 'carica impianti da /plant/all deduplicando e ordinando per nome': verifica che component.plants().map((plant) => plant.id) coincide con ['plant-1', 'plant-2', 'plant-3'] e component.plantsLoadError() è null.}{Esito atteso: component.plants().map((plant) => plant.id) coincide con ['plant-1', 'plant-2', 'plant-3']; component.plantsLoadError() è null.}{S}
\addTestI{Contesto \textit{AlarmConfigFormComponent}; caso 'se /plant/all fallisce usa fallback available + plants assegnati ai reparti': verifica che component.plants().map((plant) => plant.id) coincide con ['plant-7', 'plant-8', 'plant-9'] e component.plantsLoadError() è null.}{Esito atteso: component.plants().map((plant) => plant.id) coincide con ['plant-7', 'plant-8', 'plant-9']; component.plantsLoadError() è null; wardApiStub.getPlantsByWardId viene invocato 1 volte.}{S}
\addTestI{Contesto \textit{AlarmConfigFormComponent}; caso 'se il plant viene deselezionato resetta opzioni dispositivo e blocca il campo dispositivo': verifica che il valore di component.deviceOptions()[0]?.id è 'sensor-1' e il valore di component.form.controls.deviceId.enabled è true.}{Esito atteso: il valore di component.deviceOptions()[0]?.id è 'sensor-1'; il valore di component.form.controls.deviceId.enabled è true; component.deviceOptions() coincide con []; il valore di component.form.controls.deviceId.value è ''.}{S}
\addTestI{Contesto \textit{AlarmConfigFormComponent}; caso 'in edit mode non consente di modificare il dispositivo associato': verifica che il valore di component.form.controls.deviceId.disabled è true e component.deviceOptions() coincide con [\{ id: 'sensor-9', label: 'sensor-9', datapoints: [] \}].}{Esito atteso: il valore di component.form.controls.deviceId.disabled è true; component.deviceOptions() coincide con [\{ id: 'sensor-9', label: 'sensor-9', datapoints: [] \}].}{S}
\addTestI{Contesto \textit{AlarmConfigFormComponent}; caso 'imposta errore quando il caricamento dispositivi fallisce': verifica che il valore di component.devicesLoadError() è 'Errore durante il caricamento dei dispositivi.' e il valore di component.form.controls.deviceId.disabled è true.}{Esito atteso: il valore di component.devicesLoadError() è 'Errore durante il caricamento dei dispositivi.'; il valore di component.form.controls.deviceId.disabled è true.}{S}
\addTestI{Contesto \textit{AlarmConfigForm} feature integration; caso 'RF88-OBL create mode carica impianti e abilita selezione dispositivo': verifica che component.plants().map((plant) => plant.id) coincide con ['plant-1', 'plant-2'] e il valore di component.form.controls.deviceId.disabled è true.}{Esito atteso: component.plants().map((plant) => plant.id) coincide con ['plant-1', 'plant-2']; il valore di component.form.controls.deviceId.disabled è true; apartmentApiStub.loadApartmentViewForPlantId viene invocato con 'plant-1'; il valore di component.form.controls.deviceId.enabled è true.}{S}
\addTestI{Contesto \textit{AlarmConfigForm} feature integration; caso 'RF89-OBL selezione datapoint enum limita operatore e valida threshold': verifica che component.thresholdOperatorOptions() coincide con [ThresholdOperator.EQUAL\_TO] e il valore di component.form.controls.thresholdOperator.errors?.['unsupportedOperator'] è true.}{Esito atteso: component.thresholdOperatorOptions() coincide con [ThresholdOperator.EQUAL\_TO]; il valore di component.form.controls.thresholdOperator.errors?.['unsupportedOperator'] è true; il valore di component.form.controls.thresholdValue.errors?.['invalidEnumThreshold'] è true.}{S}
\addTestI{Contesto \textit{AlarmConfigForm} feature integration; caso 'RF91-OBL edit mode blocca campi vincolati e mantiene nome originale su submit': verifica che il valore di component.form.controls.name.disabled è true e il valore di component.form.controls.deviceId.disabled è true.}{Esito atteso: il valore di component.form.controls.name.disabled è true; il valore di component.form.controls.deviceId.disabled è true; il valore di component.form.controls.datapointId.disabled è true; emitSpy viene invocato 1 volte.}{S}
\addTestI{Contesto \textit{AlarmConfigForm} feature integration; caso 'RF92-OBL annulla form emettendo evento cancelled': verifica che cancelSpy viene invocato 1 volte.}{Esito atteso: cancelSpy viene invocato 1 volte.}{S}
\addTestI{Contesto \textit{AlarmConfigPageComponent}; caso 'crea il componente': verifica che component è valorizzato.}{Esito atteso: component è valorizzato.}{S}
\addTestI{Contesto \textit{AlarmConfigPageComponent}; caso '\textit{ngOnInit} carica gli allarmi e aggiorna le righe tabellari': verifica che stateServiceStub.\textit{loadAlarmRules} viene invocato 1 volte e il valore di component.rows().length è 1.}{Esito atteso: stateServiceStub.\textit{loadAlarmRules} viene invocato 1 volte; il valore di component.rows().length è 1; il valore di component.error() è 'Errore test'.}{S}
\addTestI{Contesto \textit{AlarmConfigPageComponent}; caso '\textit{onCreateNew} apre la modale in create mode': verifica che il valore di component.isModalOpen() è true e component.editingRule() è null.}{Esito atteso: il valore di component.isModalOpen() è true; component.editingRule() è null; il valore di component.modalTitle() è 'Creazione allarme'.}{S}
\addTestI{Contesto \textit{AlarmConfigPageComponent}; caso '\textit{onEdit} apre la modale in edit mode con regola selezionata': verifica che il valore di component.isModalOpen() è true e il valore di component.editingRule()?.id è 'alarm-1'.}{Esito atteso: il valore di component.isModalOpen() è true; il valore di component.editingRule()?.id è 'alarm-1'; il valore di component.modalTitle() è 'Modifica allarme'.}{S}
\addTestI{Contesto \textit{AlarmConfigPageComponent}; caso '\textit{onEdit} ignora id non trovato': verifica che il valore di component.isModalOpen() è false e component.editingRule() è null.}{Esito atteso: il valore di component.isModalOpen() è false; component.editingRule() è null.}{S}
\addTestI{Contesto \textit{AlarmConfigPageComponent}; caso '\textit{onFormSubmitted} in create mode invoca create e chiude modale': verifica che stateServiceStub.\textit{createAlarmRule} viene invocato con formValue e il valore di component.isModalOpen() è false.}{Esito atteso: stateServiceStub.\textit{createAlarmRule} viene invocato con formValue; il valore di component.isModalOpen() è false.}{S}
\addTestI{Contesto \textit{AlarmConfigPageComponent}; caso '\textit{onFormSubmitted} in create mode mantiene aperta la modale se non arriva alcun esito': verifica che stateServiceStub.\textit{createAlarmRule} viene invocato con formValue e il valore di component.isModalOpen() è true.}{Esito atteso: stateServiceStub.\textit{createAlarmRule} viene invocato con formValue; il valore di component.isModalOpen() è true; component.editingRule() è null.}{S}
\addTestI{Contesto \textit{AlarmConfigPageComponent}; caso '\textit{onFormSubmitted} in edit mode invoca update e chiude modale': verifica che stateServiceStub.\textit{updateAlarmRule} viene invocato con 'alarm-1' e il valore di component.isModalOpen() è false.}{Esito atteso: stateServiceStub.\textit{updateAlarmRule} viene invocato con 'alarm-1'; il valore di component.isModalOpen() è false.}{S}
\addTestI{Contesto \textit{AlarmConfigPageComponent}; caso '\textit{onFormSubmitted} in edit mode mantiene aperta la modale se non arriva alcun esito': verifica che il valore di component.isModalOpen() è true e il valore di component.editingRule()?.id è 'alarm-1'.}{Esito atteso: il valore di component.isModalOpen() è true; il valore di component.editingRule()?.id è 'alarm-1'.}{S}
\addTestI{Contesto \textit{AlarmConfigPageComponent}; caso '\textit{onToggleEnabled} invoca \textit{toggleEnabled} con lo stato richiesto': verifica che stateServiceStub.\textit{toggleEnabled} viene invocato con 'alarm-1'.}{Esito atteso: stateServiceStub.\textit{toggleEnabled} viene invocato con 'alarm-1'.}{S}
\addTestI{Contesto \textit{AlarmConfigPageComponent}; caso '\textit{onToggleEnabled} blocca i click ripetuti finché la richiesta è in corso': verifica che il valore di component.pendingToggleRuleId() è 'alarm-1' e stateServiceStub.\textit{toggleEnabled} viene invocato 1 volte.}{Esito atteso: il valore di component.pendingToggleRuleId() è 'alarm-1'; stateServiceStub.\textit{toggleEnabled} viene invocato 1 volte; component.pendingToggleRuleId() è null; stateServiceStub.\textit{toggleEnabled} viene invocato 2 volte.}{S}
\addTestI{Contesto \textit{AlarmConfigPageComponent}; caso '\textit{onDelete} apre la modale di conferma senza chiamare subito la delete': verifica che component.pendingDelete() coincide con \{ id: 'alarm-1', name: 'Temperatura alta' \} e stateServiceStub.\textit{deleteAlarmRule} non viene invocato.}{Esito atteso: component.pendingDelete() coincide con \{ id: 'alarm-1', name: 'Temperatura alta' \}; stateServiceStub.\textit{deleteAlarmRule} non viene invocato.}{S}
\addTestI{Contesto \textit{AlarmConfigPageComponent}; caso '\textit{onDeleteConfirmed} invoca \textit{deleteAlarmRule} con id e chiude la modale di conferma': verifica che stateServiceStub.\textit{deleteAlarmRule} viene invocato con 'alarm-1' e component.pendingDelete() è null.}{Esito atteso: stateServiceStub.\textit{deleteAlarmRule} viene invocato con 'alarm-1'; component.pendingDelete() è null.}{S}
\addTestI{Contesto \textit{AlarmConfigPageComponent}; caso '\textit{onDeleteCancelled} chiude la modale di conferma senza eliminare': verifica che component.pendingDelete() è null e stateServiceStub.\textit{deleteAlarmRule} non viene invocato.}{Esito atteso: component.pendingDelete() è null; stateServiceStub.\textit{deleteAlarmRule} non viene invocato.}{S}
\addTestI{Contesto \textit{AlarmConfigPageComponent}; caso 'renderizza la modale di conferma cancellazione quando richiesta': verifica che confirmDialog non è null e content contiene 'Confermi l\textbackslash\{\}'eliminazione della soglia "Temperatura alta"?'.}{Esito atteso: confirmDialog non è null; content contiene 'Confermi l\textbackslash\{\}'eliminazione della soglia "Temperatura alta"?'.}{S}
\addTestI{Contesto \textit{AlarmConfigPageComponent}; caso '\textit{onModalClosed} chiude modale e resetta editingRule': verifica che il valore di component.isModalOpen() è false e component.editingRule() è null.}{Esito atteso: il valore di component.isModalOpen() è false; component.editingRule() è null.}{S}
\addTestI{Contesto \textit{AlarmConfigPageComponent}; caso '\textit{onCreateNew} resetta l eventuale regola in edit e apre la modale in create': verifica che il valore di component.isModalOpen() è true e component.editingRule() è null.}{Esito atteso: il valore di component.isModalOpen() è true; component.editingRule() è null.}{S}
\addTestI{Contesto \textit{AlarmConfigPageComponent}; caso 'renderizza la tabella con le righe e i pulsanti azione': verifica che il valore di tableRows.length è 1 e content contiene 'Temperatura alta'.}{Esito atteso: il valore di tableRows.length è 1; content contiene 'Temperatura alta'; content contiene 'Posizione'; content contiene 'Appartamento 1 - Soggiorno - Sensore porta'.}{S}
\addTestI{Contesto \textit{AlarmConfigPageComponent}; caso 'renderizza messaggio vuoto quando non ci sono regole': verifica che content contiene 'Nessuna regola configurata.'.}{Esito atteso: content contiene 'Nessuna regola configurata.'.}{S}
\addTestI{Contesto \textit{AlarmConfigPageComponent}; caso 'renderizza il form in modale quando la modale e aperta': verifica che modalForm non è null.}{Esito atteso: modalForm non è null.}{S}
\addTestI{Contesto AlarmConfiguration feature integration; caso 'RF85-OBL inizializza pagina e carica regole allarme': verifica che component è valorizzato e stateServiceStub.\textit{loadAlarmRules} viene invocato 1 volte.}{Esito atteso: component è valorizzato; stateServiceStub.\textit{loadAlarmRules} viene invocato 1 volte.}{S}
\addTestI{Contesto AlarmConfiguration feature integration; caso 'RF86-OBL create mode invia form al servizio e chiude modale': verifica che stateServiceStub.\textit{createAlarmRule} viene invocato con formValue e il valore di component.isModalOpen() è false.}{Esito atteso: stateServiceStub.\textit{createAlarmRule} viene invocato con formValue; il valore di component.isModalOpen() è false.}{S}
\addTestI{Contesto AlarmConfiguration feature integration; caso 'RF87-OBL edit mode mantiene modale aperta se update non emette esito': verifica che il valore di component.isModalOpen() è true e il valore di component.editingRule()?.id è 'alarm-1'.}{Esito atteso: il valore di component.isModalOpen() è true; il valore di component.editingRule()?.id è 'alarm-1'.}{S}
\addTestI{Contesto \textit{AlarmPageManagementComponent}; caso 'should create': verifica che component è valorizzato.}{Esito atteso: component è valorizzato.}{S}
\addTestI{Contesto \textit{AlarmPageManagementComponent}; caso '\textit{onResolve} delega a facade.\textit{resolveAlarm}': verifica che alarmManagementStub.\textit{resolveAlarm} viene invocato con 'active-2' e alarmManagementStub.\textit{resolveAlarm} viene invocato 1 volte.}{Esito atteso: alarmManagementStub.\textit{resolveAlarm} viene invocato con 'active-2'; alarmManagementStub.\textit{resolveAlarm} viene invocato 1 volte.}{S}
\addTestI{Contesto \textit{AlarmPageManagementComponent}; caso 'renderizza stato lista vuota quando non ci sono allarmi': verifica che nativeElement.querySelector('.alarm-management\_\_empty')?.textContent contiene 'Nessun allarme attivo al momento.' e il valore di nativeElement.querySelectorAll('tbody tr').length è 0.}{Esito atteso: nativeElement.querySelector('.alarm-management\_\_empty')?.textContent contiene 'Nessun allarme attivo al momento.'; il valore di nativeElement.querySelectorAll('tbody tr').length è 0.}{S}
\addTestI{Contesto \textit{AlarmPageManagementComponent}; caso 'renderizza errore e stato resolving quando presenti nel vm': verifica che nativeElement.querySelector('.alarm-management\_\_error')?.textContent contiene 'Errore durante la risoluzione' e nativeElement.querySelector('.alarm-management\_\_status')?.textContent contiene 'Risoluzione allarme in corso...'.}{Esito atteso: nativeElement.querySelector('.alarm-management\_\_error')?.textContent contiene 'Errore durante la risoluzione'; nativeElement.querySelector('.alarm-management\_\_status')?.textContent contiene 'Risoluzione allarme in corso...'.}{S}
\addTestI{Contesto \textit{AlarmPageManagementComponent}; caso 'renderizza tabella con una riga per ogni allarme': verifica che il valore di rows.length è 2 e nativeElement.textContent contiene 'Priorità'.}{Esito atteso: il valore di rows.length è 2; nativeElement.textContent contiene 'Priorità'; nativeElement.textContent contiene 'Corridoio Nord'; nativeElement.textContent contiene 'mock-elapsed:2026-03-24T10:00:00.000Z'.}{S}
\addTestI{Contesto \textit{AlarmPageManagementComponent}; caso 'nasconde la colonna gestore per operatore sanitario': verifica che nativeElement.textContent non contiene 'Gestore' e nativeElement.textContent non contiene 'oss\_1'.}{Esito atteso: nativeElement.textContent non contiene 'Gestore'; nativeElement.textContent non contiene 'oss\_1'.}{S}
\addTestI{Contesto \textit{AlarmPageManagementComponent}; caso 'disabilita solo il bottone associato al resolvingId corrente': verifica che il valore di firstButton?.disabled è true e il valore di secondButton?.disabled è false.}{Esito atteso: il valore di firstButton?.disabled è true; il valore di secondButton?.disabled è false.}{S}
\addTestI{Contesto \textit{AlarmPageManagementComponent}; caso 'renderizza paginazione coerente con vm e delega i click ai metodi del service': verifica che nativeElement.querySelector('.care-pagination\_\_status')?.textContent contiene 'Pagina 3' e il valore di previousButton?.disabled è false.}{Esito atteso: nativeElement.querySelector('.care-pagination\_\_status')?.textContent contiene 'Pagina 3'; il valore di previousButton?.disabled è false; il valore di nextButton?.disabled è true; alarmManagementStub.previousPage viene invocato 1 volte.}{S}
\addTestI{Contesto \textit{AlarmPageManagementComponent}; caso 'click su GESTISCI propaga resolve verso facade': verifica che alarmManagementStub.\textit{resolveAlarm} viene invocato con 'active-1' e alarmManagementStub.\textit{resolveAlarm} viene invocato 1 volte.}{Esito atteso: alarmManagementStub.\textit{resolveAlarm} viene invocato con 'active-1'; alarmManagementStub.\textit{resolveAlarm} viene invocato 1 volte.}{S}
\addTestI{Contesto \textit{AlarmPageManagementComponent}; caso 'mantiene la riga visibile se l'allarme e già gestito e disabilita l azione': verifica che il valore di rows.length è 1 e nativeElement.textContent contiene 'Non da gestire'.}{Esito atteso: il valore di rows.length è 1; nativeElement.textContent contiene 'Non da gestire'; managedButton non è null; il valore di managedButton?.disabled è true.}{S}
\addTestI{Contesto \textit{AlarmPageManagementComponent}; caso '\textit{onNextPage} e \textit{onPreviousPage} delegano al service': verifica che alarmManagementStub.nextPage viene invocato 1 volte e alarmManagementStub.previousPage viene invocato 1 volte.}{Esito atteso: alarmManagementStub.nextPage viene invocato 1 volte; alarmManagementStub.previousPage viene invocato 1 volte.}{S}
\addTestI{Contesto \textit{AlarmPageManagementComponent}; caso 'quando arriva \textit{refreshRequested} reinizializza la pagina': verifica che alarmManagementStub.initialize viene invocato 1 volte e alarmManagementStub.initialize viene invocato 2 volte.}{Esito atteso: alarmManagementStub.initialize viene invocato 1 volte; alarmManagementStub.initialize viene invocato 2 volte.}{S}
\addTestI{Contesto \textit{AlarmPageManagementComponent}; caso 'dopo destroy non risponde più agli eventi di refresh': verifica che alarmManagementStub.initialize viene invocato 1 volte.}{Esito atteso: alarmManagementStub.initialize viene invocato 1 volte.}{S}
\addTestI{Contesto AlarmManagement feature integration; caso 'RF49-OBL gestisce il flusso completo: render lista, resolve event e aggiornamento vm': verifica che alarmManagementStub.initialize viene invocato 1 volte e il valore di \textit{resolveButtons}.length è 2.}{Esito atteso: alarmManagementStub.initialize viene invocato 1 volte; il valore di \textit{resolveButtons}.length è 2; nativeElement.textContent contiene 'Allarmi attivi'; nativeElement.textContent contiene 'mock-elapsed:2026-03-24T10:00:00.000Z'.}{S}
\addTestI{Contesto AlarmManagement feature integration; caso 'RF60-OBL riesegue initialize quando arriva una richiesta di refresh esterna': verifica che alarmManagementStub.initialize viene invocato 1 volte e alarmManagementStub.initialize viene invocato 2 volte.}{Esito atteso: alarmManagementStub.initialize viene invocato 1 volte; alarmManagementStub.initialize viene invocato 2 volte.}{S}
\addTestI{Contesto \textit{AlarmFrequencyChartComponent}; caso 'espone metadati chart consistenti': verifica che il valore di component.chartType è 'bar' e component.description.length è maggiore di 0.}{Esito atteso: il valore di component.chartType è 'bar'; component.description.length è maggiore di 0.}{S}
\addTestI{Contesto \textit{AlarmFrequencyChartComponent} - con suggerimenti attivi; caso 'dovrebbe restringere il grafico al 70\% su schermi grandi': verifica che il contenitore del grafico applica la classe responsive lg:w-[70\%] quando i suggerimenti sono attivi.}{Esito atteso: l'elemento [data-test="chart-container"] include la classe lg:w-[70\%].}{S}
\addTestI{Contesto \textit{AlarmFrequencyChartComponent} - con suggerimenti attivi; caso 'dovrebbe mostrare la colonna dei suggerimenti': verifica che la sidebar dei suggerimenti viene renderizzata nel DOM quando la modalita suggerimenti è attiva.}{Esito atteso: l'elemento [data-test="suggestions-sidebar"] è presente e non risulta null.}{S}
\addTestI{Contesto \textit{AlarmFrequencyChartComponent} - con suggerimenti attivi; caso 'dovrebbe renderizzare il numero corretto di messaggi': verifica che la lista dei suggerimenti mostra esattamente il numero di messaggi previsto dai dati di input.}{Esito atteso: gli elementi [data-test="suggestion-item"] sono esattamente 2, in linea con i messaggi forniti.}{S}
\addTestI{Contesto \textit{AlarmFrequencyChartComponent} - senza suggerimenti; caso 'dovrebbe espandere il grafico al 100\% della larghezza': verifica che senza suggerimenti il contenitore del grafico occupa tutta la larghezza disponibile.}{Esito atteso: è presente un contenitore grafico con classe .w-full, quindi la vista occupa il 100\% della larghezza.}{S}
\addTestI{Contesto \textit{AlarmFrequencyChartComponent} - senza suggerimenti; caso 'NON dovrebbe esistere la colonna dei suggerimenti nel DOM': verifica che senza suggerimenti la colonna laterale non viene montata nel DOM.}{Esito atteso: la colonna suggerimenti (.lg:w-[30\%]) è assente e la query restituisce null.}{S}
\addTestI{Contesto \textit{AlarmsSentResolvedChartComponent}; caso 'espone metadati chart consistenti': verifica che il valore di component.chartType è 'bar' e component.description.length è maggiore di 0.}{Esito atteso: il valore di component.chartType è 'bar'; component.description.length è maggiore di 0.}{S}
\addTestI{Contesto \textit{AlarmsSentResolvedChartComponent} - con suggerimenti attivi; caso 'dovrebbe restringere il grafico al 70\% su schermi grandi': verifica che il contenitore del grafico applica la classe responsive lg:w-[70\%] quando i suggerimenti sono attivi.}{Esito atteso: l'elemento [data-test="chart-container"] include la classe lg:w-[70\%].}{S}
\addTestI{Contesto \textit{AlarmsSentResolvedChartComponent} - con suggerimenti attivi; caso 'dovrebbe mostrare la colonna dei suggerimenti': verifica che la sidebar dei suggerimenti viene renderizzata nel DOM quando la modalita suggerimenti è attiva.}{Esito atteso: l'elemento [data-test="suggestions-sidebar"] è presente e non risulta null.}{S}
\addTestI{Contesto \textit{AlarmsSentResolvedChartComponent} - con suggerimenti attivi; caso 'dovrebbe renderizzare il numero corretto di messaggi': verifica che la lista dei suggerimenti mostra esattamente il numero di messaggi previsto dai dati di input.}{Esito atteso: gli elementi [data-test="suggestion-item"] sono esattamente 2, in linea con i messaggi forniti.}{S}
\addTestI{Contesto \textit{AlarmsSentResolvedChartComponent} - senza suggerimenti; caso 'dovrebbe espandere il grafico al 100\% della larghezza': verifica che senza suggerimenti il contenitore del grafico occupa tutta la larghezza disponibile.}{Esito atteso: è presente un contenitore grafico con classe .w-full, quindi la vista occupa il 100\% della larghezza.}{S}
\addTestI{Contesto \textit{AlarmsSentResolvedChartComponent} - senza suggerimenti; caso 'NON dovrebbe esistere la colonna dei suggerimenti nel DOM': verifica che senza suggerimenti la colonna laterale non viene montata nel DOM.}{Esito atteso: la colonna suggerimenti (.lg:w-[30\%]) è assente e la query restituisce null.}{S}
\addTestI{Contesto \textit{EnergyConsumptionChartComponent}; caso 'espone metadati chart consistenti': verifica che il valore di component.chartType è 'line' e component.description.length è maggiore di 0.}{Esito atteso: il valore di component.chartType è 'line'; component.description.length è maggiore di 0.}{S}
\addTestI{Contesto \textit{EnergyConsumptionChartComponent} - con suggerimenti attivi; caso 'dovrebbe restringere il grafico al 70\% su schermi grandi': verifica che il contenitore del grafico applica la classe responsive lg:w-[70\%] quando i suggerimenti sono attivi.}{Esito atteso: l'elemento [data-test="chart-container"] include la classe lg:w-[70\%].}{S}
\addTestI{Contesto \textit{EnergyConsumptionChartComponent} - con suggerimenti attivi; caso 'dovrebbe mostrare la colonna dei suggerimenti': verifica che la sidebar dei suggerimenti viene renderizzata nel DOM quando la modalita suggerimenti è attiva.}{Esito atteso: l'elemento [data-test="suggestions-sidebar"] è presente e non risulta null.}{S}
\addTestI{Contesto \textit{EnergyConsumptionChartComponent} - con suggerimenti attivi; caso 'dovrebbe renderizzare il numero corretto di messaggi': verifica che la lista dei suggerimenti mostra esattamente il numero di messaggi previsto dai dati di input.}{Esito atteso: gli elementi [data-test="suggestion-item"] sono esattamente 2, in linea con i messaggi forniti.}{S}
\addTestI{Contesto \textit{EnergyConsumptionChartComponent} - senza suggerimenti; caso 'dovrebbe espandere il grafico al 100\% della larghezza': verifica che senza suggerimenti il contenitore del grafico occupa tutta la larghezza disponibile.}{Esito atteso: è presente un contenitore grafico con classe .w-full, quindi la vista occupa il 100\% della larghezza.}{S}
\addTestI{Contesto \textit{EnergyConsumptionChartComponent} - senza suggerimenti; caso 'NON dovrebbe esistere la colonna dei suggerimenti nel DOM': verifica che senza suggerimenti la colonna laterale non viene montata nel DOM.}{Esito atteso: la colonna suggerimenti (.lg:w-[30\%]) è assente e la query restituisce null.}{S}
\addTestI{Contesto \textit{FallFrequencyChartComponent}; caso 'espone metadati chart consistenti': verifica che il valore di component.chartType è 'bar' e component.description.length è maggiore di 0.}{Esito atteso: il valore di component.chartType è 'bar'; component.description.length è maggiore di 0.}{S}
\addTestI{Contesto \textit{FallFrequencyChartComponent} - con suggerimenti attivi; caso 'dovrebbe restringere il grafico al 70\% su schermi grandi': verifica che il contenitore del grafico applica la classe responsive lg:w-[70\%] quando i suggerimenti sono attivi.}{Esito atteso: l'elemento [data-test="chart-container"] include la classe lg:w-[70\%].}{S}
\addTestI{Contesto \textit{FallFrequencyChartComponent} - con suggerimenti attivi; caso 'dovrebbe mostrare la colonna dei suggerimenti': verifica che la sidebar dei suggerimenti viene renderizzata nel DOM quando la modalita suggerimenti è attiva.}{Esito atteso: l'elemento [data-test="suggestions-sidebar"] è presente e non risulta null.}{S}
\addTestI{Contesto \textit{FallFrequencyChartComponent} - con suggerimenti attivi; caso 'dovrebbe renderizzare il numero corretto di messaggi': verifica che la lista dei suggerimenti mostra esattamente il numero di messaggi previsto dai dati di input.}{Esito atteso: gli elementi [data-test="suggestion-item"] sono esattamente 2, in linea con i messaggi forniti.}{S}
\addTestI{Contesto \textit{FallFrequencyChartComponent} - senza suggerimenti; caso 'dovrebbe espandere il grafico al 100\% della larghezza': verifica che senza suggerimenti il contenitore del grafico occupa tutta la larghezza disponibile.}{Esito atteso: è presente un contenitore grafico con classe .w-full, quindi la vista occupa il 100\% della larghezza.}{S}
\addTestI{Contesto \textit{FallFrequencyChartComponent} - senza suggerimenti; caso 'NON dovrebbe esistere la colonna dei suggerimenti nel DOM': verifica che senza suggerimenti la colonna laterale non viene montata nel DOM.}{Esito atteso: la colonna suggerimenti (.lg:w-[30\%]) è assente e la query restituisce null.}{S}
\addTestI{Contesto \textit{PlantAnomaliesChartComponent}; caso 'espone metadati chart consistenti': verifica che il valore di component.chartType è 'bar' e component.description.length è maggiore di 0.}{Esito atteso: il valore di component.chartType è 'bar'; component.description.length è maggiore di 0.}{S}
\addTestI{Contesto \textit{PlantAnomaliesChartComponent} - con suggerimenti attivi; caso 'dovrebbe restringere il grafico al 70\% su schermi grandi': verifica che il contenitore del grafico applica la classe responsive lg:w-[70\%] quando i suggerimenti sono attivi.}{Esito atteso: l'elemento [data-test="chart-container"] include la classe lg:w-[70\%].}{S}
\addTestI{Contesto \textit{PlantAnomaliesChartComponent} - con suggerimenti attivi; caso 'dovrebbe mostrare la colonna dei suggerimenti': verifica che la sidebar dei suggerimenti viene renderizzata nel DOM quando la modalita suggerimenti è attiva.}{Esito atteso: l'elemento [data-test="suggestions-sidebar"] è presente e non risulta null.}{S}
\addTestI{Contesto \textit{PlantAnomaliesChartComponent} - con suggerimenti attivi; caso 'dovrebbe renderizzare il numero corretto di messaggi': verifica che la lista dei suggerimenti mostra esattamente il numero di messaggi previsto dai dati di input.}{Esito atteso: gli elementi [data-test="suggestion-item"] sono esattamente 2, in linea con i messaggi forniti.}{S}
\addTestI{Contesto \textit{PlantAnomaliesChartComponent} - senza suggerimenti; caso 'dovrebbe espandere il grafico al 100\% della larghezza': verifica che senza suggerimenti il contenitore del grafico occupa tutta la larghezza disponibile.}{Esito atteso: è presente un contenitore grafico con classe .w-full, quindi la vista occupa il 100\% della larghezza.}{S}
\addTestI{Contesto \textit{PlantAnomaliesChartComponent} - senza suggerimenti; caso 'NON dovrebbe esistere la colonna dei suggerimenti nel DOM': verifica che senza suggerimenti la colonna laterale non viene montata nel DOM.}{Esito atteso: la colonna suggerimenti (.lg:w-[30\%]) è assente e la query restituisce null.}{S}
\addTestI{Contesto \textit{PresenceDetectionChartComponent}; caso 'espone metadati chart consistenti': verifica che il valore di component.chartType è 'bar' e component.description.length è maggiore di 0.}{Esito atteso: il valore di component.chartType è 'bar'; component.description.length è maggiore di 0.}{S}
\addTestI{Contesto \textit{PresenceDetectionChartComponent} - con suggerimenti attivi; caso 'dovrebbe restringere il grafico al 70\% su schermi grandi': verifica che il contenitore del grafico applica la classe responsive lg:w-[70\%] quando i suggerimenti sono attivi.}{Esito atteso: l'elemento [data-test="chart-container"] include la classe lg:w-[70\%].}{S}
\addTestI{Contesto \textit{PresenceDetectionChartComponent} - con suggerimenti attivi; caso 'dovrebbe mostrare la colonna dei suggerimenti': verifica che la sidebar dei suggerimenti viene renderizzata nel DOM quando la modalita suggerimenti è attiva.}{Esito atteso: l'elemento [data-test="suggestions-sidebar"] è presente e non risulta null.}{S}
\addTestI{Contesto \textit{PresenceDetectionChartComponent} - con suggerimenti attivi; caso 'dovrebbe renderizzare il numero corretto di messaggi': verifica che la lista dei suggerimenti mostra esattamente il numero di messaggi previsto dai dati di input.}{Esito atteso: gli elementi [data-test="suggestion-item"] sono esattamente 2, in linea con i messaggi forniti.}{S}
\addTestI{Contesto \textit{PresenceDetectionChartComponent} - senza suggerimenti; caso 'dovrebbe espandere il grafico al 100\% della larghezza': verifica che senza suggerimenti il contenitore del grafico occupa tutta la larghezza disponibile.}{Esito atteso: è presente un contenitore grafico con classe .w-full, quindi la vista occupa il 100\% della larghezza.}{S}
\addTestI{Contesto \textit{PresenceDetectionChartComponent} - senza suggerimenti; caso 'NON dovrebbe esistere la colonna dei suggerimenti nel DOM': verifica che senza suggerimenti la colonna laterale non viene montata nel DOM.}{Esito atteso: la colonna suggerimenti (.lg:w-[30\%]) è assente e la query restituisce null.}{S}
\addTestI{Contesto \textit{ProlongedPresenceChartComponent}; caso 'espone metadati chart consistenti': verifica che il valore di component.chartType è 'bar' e component.description.length è maggiore di 0.}{Esito atteso: il valore di component.chartType è 'bar'; component.description.length è maggiore di 0.}{S}
\addTestI{Contesto \textit{ProlongedPresenceChartComponent} - con suggerimenti attivi; caso 'dovrebbe restringere il grafico al 70\% su schermi grandi': verifica che il contenitore del grafico applica la classe responsive lg:w-[70\%] quando i suggerimenti sono attivi.}{Esito atteso: l'elemento [data-test="chart-container"] include la classe lg:w-[70\%].}{S}
\addTestI{Contesto \textit{ProlongedPresenceChartComponent} - con suggerimenti attivi; caso 'dovrebbe mostrare la colonna dei suggerimenti': verifica che la sidebar dei suggerimenti viene renderizzata nel DOM quando la modalita suggerimenti è attiva.}{Esito atteso: l'elemento [data-test="suggestions-sidebar"] è presente e non risulta null.}{S}
\addTestI{Contesto \textit{ProlongedPresenceChartComponent} - con suggerimenti attivi; caso 'dovrebbe renderizzare il numero corretto di messaggi': verifica che la lista dei suggerimenti mostra esattamente il numero di messaggi previsto dai dati di input.}{Esito atteso: gli elementi [data-test="suggestion-item"] sono esattamente 2, in linea con i messaggi forniti.}{S}
\addTestI{Contesto \textit{ProlongedPresenceChartComponent} - senza suggerimenti; caso 'dovrebbe espandere il grafico al 100\% della larghezza': verifica che senza suggerimenti il contenitore del grafico occupa tutta la larghezza disponibile.}{Esito atteso: è presente un contenitore grafico con classe .w-full, quindi la vista occupa il 100\% della larghezza.}{S}
\addTestI{Contesto \textit{ProlongedPresenceChartComponent} - senza suggerimenti; caso 'NON dovrebbe esistere la colonna dei suggerimenti nel DOM': verifica che senza suggerimenti la colonna laterale non viene montata nel DOM.}{Esito atteso: la colonna suggerimenti (.lg:w-[30\%]) è assente e la query restituisce null.}{S}
\addTestI{Contesto \textit{TemperatureVariationsChartComponent}; caso 'espone metadati chart consistenti': verifica che il valore di component.chartType è 'line' e component.description.length è maggiore di 0.}{Esito atteso: il valore di component.chartType è 'line'; component.description.length è maggiore di 0.}{S}
\addTestI{Contesto \textit{TemperatureVariationsChartComponent} - con suggerimenti attivi; caso 'dovrebbe restringere il grafico al 70\% su schermi grandi': verifica che il contenitore del grafico applica la classe responsive lg:w-[70\%] quando i suggerimenti sono attivi.}{Esito atteso: l'elemento [data-test="chart-container"] include la classe lg:w-[70\%].}{S}
\addTestI{Contesto \textit{TemperatureVariationsChartComponent} - con suggerimenti attivi; caso 'dovrebbe mostrare la colonna dei suggerimenti': verifica che la sidebar dei suggerimenti viene renderizzata nel DOM quando la modalita suggerimenti è attiva.}{Esito atteso: l'elemento [data-test="suggestions-sidebar"] è presente e non risulta null.}{S}
\addTestI{Contesto \textit{TemperatureVariationsChartComponent} - con suggerimenti attivi; caso 'dovrebbe renderizzare il numero corretto di messaggi': verifica che la lista dei suggerimenti mostra esattamente il numero di messaggi previsto dai dati di input.}{Esito atteso: gli elementi [data-test="suggestion-item"] sono esattamente 2, in linea con i messaggi forniti.}{S}
\addTestI{Contesto \textit{TemperatureVariationsChartComponent} - senza suggerimenti; caso 'dovrebbe espandere il grafico al 100\% della larghezza': verifica che senza suggerimenti il contenitore del grafico occupa tutta la larghezza disponibile.}{Esito atteso: è presente un contenitore grafico con classe .w-full, quindi la vista occupa il 100\% della larghezza.}{S}
\addTestI{Contesto \textit{TemperatureVariationsChartComponent} - senza suggerimenti; caso 'NON dovrebbe esistere la colonna dei suggerimenti nel DOM': verifica che senza suggerimenti la colonna laterale non viene montata nel DOM.}{Esito atteso: la colonna suggerimenti (.lg:w-[30\%]) è assente e la query restituisce null.}{S}
\addTestI{Contesto \textit{ApartmentMonitorComponent}; caso 'crea il componente e sincronizza activeApartmentId dal feed apartment\$': verifica che component è valorizzato e il valore di component.activeApartmentId è 'ap-1'.}{Esito atteso: component è valorizzato; il valore di component.activeApartmentId è 'ap-1'; apartmentApiStub.\textit{getCurrentApartment} viene invocato.}{S}
\addTestI{Contesto \textit{ApartmentMonitorComponent}; caso '\textit{onApartmentSelected} ignora id vuoto o uguale': verifica che apartmentApiStub.setActivePlantId non viene invocato.}{Esito atteso: apartmentApiStub.setActivePlantId non viene invocato.}{S}
\addTestI{Contesto \textit{ApartmentMonitorComponent}; caso '\textit{onApartmentSelected} cambia appartamento, resetta errore e triggera refresh': verifica che apartmentApiStub.setActivePlantId viene invocato con 'ap-2' e il valore di component.error è ''.}{Esito atteso: apartmentApiStub.setActivePlantId viene invocato con 'ap-2'; il valore di component.error è ''.}{S}
\addTestI{Contesto \textit{ApartmentMonitorComponent}; caso 'refresh realtime azzera errore e ricarica apartment\$': verifica che il valore di component.error è '' e apartmentApiStub.\textit{getCurrentApartment}.mock.calls.length è maggiore di callsBefore.}{Esito atteso: il valore di component.error è ''; apartmentApiStub.\textit{getCurrentApartment}.mock.calls.length è maggiore di callsBefore.}{S}
\addTestI{Contesto ApartmentMonitor feature integration; caso 'RF72-OBL inizializza appartamento attivo dal feed e prepara la vista': verifica che component è valorizzato e apartmentApiStub.\textit{getCurrentApartment} viene invocato.}{Esito atteso: component è valorizzato; apartmentApiStub.\textit{getCurrentApartment} viene invocato; il valore di component.activeApartmentId è 'ap-1'.}{S}
\addTestI{Contesto ApartmentMonitor feature integration; caso 'RF73-OBL cambio appartamento aggiorna plant attivo e resetta errore': verifica che apartmentApiStub.setActivePlantId viene invocato con 'ap-2' e il valore di component.error è ''.}{Esito atteso: apartmentApiStub.setActivePlantId viene invocato con 'ap-2'; il valore di component.error è ''.}{S}
\addTestI{Contesto ApartmentMonitor feature integration; caso 'RF74-OBL evento realtime forza refresh del caricamento appartamento': verifica che apartmentApiStub.\textit{getCurrentApartment}.mock.calls.length è maggiore di beforeCalls.}{Esito atteso: apartmentApiStub.\textit{getCurrentApartment}.mock.calls.length è maggiore di beforeCalls.}{S}
\addTestI{Contesto \textit{AlarmWidgetComponent}; caso 'dovrebbe essere creato': verifica che component è valorizzato.}{Esito atteso: component è valorizzato.}{S}
\addTestI{Contesto \textit{AlarmWidgetComponent}; caso 'dovrebbe inizializzare il servizio al caricamento (\textit{ngOnInit})': verifica che mockAlarmService.initialize viene invocato.}{Esito atteso: mockAlarmService.initialize viene invocato.}{S}
\addTestI{Contesto \textit{AlarmWidgetComponent}; caso 'dovrebbe mostrare lo stato di caricamento se vm è null': verifica che compiled.textContent contiene 'Caricamento allarmi in corso'.}{Esito atteso: compiled.textContent contiene 'Caricamento allarmi in corso'.}{S}
\addTestI{Contesto \textit{AlarmWidgetComponent}; caso 'dovrebbe mostrare il messaggio "Nessun allarme" se la lista è vuota': verifica che compiled.textContent contiene 'Nessun allarme attivo al momento'.}{Esito atteso: compiled.textContent contiene 'Nessun allarme attivo al momento'.}{S}
\addTestI{Contesto \textit{AlarmWidgetComponent}; caso 'dovrebbe chiamare \textit{resolveAlarm} quando \textit{onResolve} viene invocato': verifica che mockAlarmService.\textit{resolveAlarm} viene invocato con alarmId.}{Esito atteso: mockAlarmService.\textit{resolveAlarm} viene invocato con alarmId.}{S}
\addTestI{Contesto \textit{AlarmWidgetComponent}; caso 'dovrebbe ricaricare i dati quando il refresh service emette': verifica che mockAlarmService.initialize viene invocato.}{Esito atteso: mockAlarmService.initialize viene invocato.}{S}
\addTestI{Contesto \textit{AlarmWidgetComponent}; caso 'dovrebbe calcolare le righe correttamente tramite il presenter': verifica che mockPresenter.\textit{toRows} viene invocato con alarms.}{Esito atteso: mockPresenter.\textit{toRows} viene invocato con alarms.}{S}
\addTestI{Contesto \textit{AlarmWidget} feature integration; caso 'RF44-OBL init widget invoca initialize service': verifica che component è valorizzato e alarmManagementStub.initialize viene invocato.}{Esito atteso: component è valorizzato; alarmManagementStub.initialize viene invocato.}{S}
\addTestI{Contesto \textit{AlarmWidget} feature integration; caso 'RF44-OBL click azione allarme delega \textit{resolveAlarm}': verifica che alarmManagementStub.\textit{resolveAlarm} viene invocato con 'a-1'.}{Esito atteso: alarmManagementStub.\textit{resolveAlarm} viene invocato con 'a-1'.}{S}
\addTestI{Contesto \textit{AlarmWidget} feature integration; caso 'RF44-OBL refresh esterno reinizializza la tabella': verifica che alarmManagementStub.initialize.mock.calls.length è maggiore di callsBefore.}{Esito atteso: alarmManagementStub.initialize.mock.calls.length è maggiore di callsBefore.}{S}
\addTestI{Contesto \textit{AnalyticsComponent}; caso 'viene creato': verifica che component è valorizzato.}{Esito atteso: component è valorizzato.}{S}
\addTestI{Contesto \textit{AnalyticsComponent}; caso 'carica il primo appartamento di default': verifica che mockAnalyticsService.\textit{getAllApartments} viene invocato e mockAnalyticsService.\textit{getAnalytics} viene invocato con '1'.}{Esito atteso: mockAnalyticsService.\textit{getAllApartments} viene invocato; mockAnalyticsService.\textit{getAnalytics} viene invocato con '1'.}{S}
\addTestI{Contesto \textit{AnalyticsComponent}; caso 'aggiorna grafici al cambio appartamento': verifica che mockAnalyticsService.\textit{getAnalytics} viene invocato con '2'.}{Esito atteso: mockAnalyticsService.\textit{getAnalytics} viene invocato con '2'.}{S}
\addTestI{Contesto \textit{AnalyticsComponent}; caso 'should return undefined if metric is not found': verifica che result è undefined.}{Esito atteso: result è undefined.}{S}
\addTestI{Contesto \textit{AnalyticsComponent}; caso 'should find the correct metric in \textit{getChartByMetric}': verifica che result è definito e il valore di result?.metric è 'plant-anomalies'.}{Esito atteso: result è definito; il valore di result?.metric è 'plant-anomalies'.}{S}
\addTestI{Contesto \textit{AnalyticsWidget} feature integration; caso 'RF45-OBL inizializza widget e carica analytics del primo appartamento': verifica che analyticsApiStub.\textit{getAllApartments} viene invocato 1 volte e analyticsApiStub.\textit{getAnalytics} viene invocato con 'plant-1'.}{Esito atteso: analyticsApiStub.\textit{getAllApartments} viene invocato 1 volte; analyticsApiStub.\textit{getAnalytics} viene invocato con 'plant-1'.}{S}
\addTestI{Contesto \textit{AnalyticsWidget} feature integration; caso 'RF45-OBL cambio appartamento aggiorna richiesta analytics': verifica che analyticsApiStub.\textit{getAnalytics} viene invocato con 'plant-2'.}{Esito atteso: analyticsApiStub.\textit{getAnalytics} viene invocato con 'plant-2'.}{S}
\addTestI{Contesto \textit{AnalyticsWidget} feature integration; caso 'RF45-OBL renderizza overview appartamento selezionato': verifica che overview non è null e fixture.nativeElement.textContent contiene 'Appartamento Monitorato'.}{Esito atteso: overview non è null; fixture.nativeElement.textContent contiene 'Appartamento Monitorato'.}{S}
\addTestI{Contesto Dashboard feature integration; caso 'RF43-OBL renderizza la dashboard': verifica che component è valorizzato.}{Esito atteso: component è valorizzato.}{S}
\addTestI{Contesto Dashboard feature integration; caso 'RF44-OBL visualizza widget allarmi': verifica che fixture.debugElement.query(By.css('app-alarms-widget')) non è null.}{Esito atteso: fixture.debugElement.query(By.css('app-alarms-widget')) non è null.}{S}
\addTestI{Contesto Dashboard feature integration; caso 'RF45-OBL visualizza widget analytics': verifica che fixture.debugElement.query(By.css('app-analytics-widget')) non è null.}{Esito atteso: fixture.debugElement.query(By.css('app-analytics-widget')) non è null.}{S}
\addTestI{Contesto \textit{PlantOverviewComponent}; caso '1. Dovrebbe mostrare il nome dell impianto': verifica che html.textContent contiene 'Casa Test'.}{Esito atteso: html.textContent contiene 'Casa Test'.}{S}
\addTestI{Contesto \textit{PlantOverviewComponent}; caso '2. Dovrebbe mostrare i nomi delle stanze': verifica che html.textContent contiene 'Soggiorno' e html.textContent contiene 'Camera'.}{Esito atteso: html.textContent contiene 'Soggiorno'; html.textContent contiene 'Camera'.}{S}
\addTestI{Contesto \textit{PlantOverviewComponent}; caso '3. Dovrebbe gestire il caso senza dispositivi': verifica che html.textContent contiene 'Nessun dispositivo configurato'.}{Esito atteso: html.textContent contiene 'Nessun dispositivo configurato'.}{S}
\addTestI{Contesto \textit{PlantOverviewComponent}; caso '4. Non dovrebbe mostrare nulla se apartment è null': verifica che html.querySelector('header') è null.}{Esito atteso: html.querySelector('header') è null.}{S}
\addTestI{Contesto \textit{EndpointTableComponent} deep integration; caso '\textit{groupRows} creates one group for same room id and \textit{trackBy} returns stable keys': verifica che il valore di groups.length è 1 e il valore di component.\textit{trackByGroup}(0, \{ roomId: 'room-1', roomName: '', rows: [] \}) è 'room-1'.}{Esito atteso: il valore di groups.length è 1; il valore di component.\textit{trackByGroup}(0, \{ roomId: 'room-1', roomName: '', rows: [] \}) è 'room-1'; component.\textit{trackByRow}(0, mappedRow) contiene 'room-1::device-1::dp-cmd-1'.}{S}
\addTestI{Contesto \textit{EndpointTableComponent} deep integration; caso '\textit{getCurrentValue} resolves exact datapoint match': verifica che il valore di exact è 'Off'.}{Esito atteso: il valore di exact è 'Off'.}{S}
\addTestI{Contesto \textit{EndpointTableComponent} deep integration; caso '\textit{getCurrentValue} uses first available value when exact and semantic matches are missing': verifica che il valore di component.\textit{getCurrentValue}(fallbackRow) è 'FallbackValue'.}{Esito atteso: il valore di component.\textit{getCurrentValue}(fallbackRow) è 'FallbackValue'.}{S}
\addTestI{Contesto \textit{EndpointTableComponent} deep integration; caso 'selected values default to first enum and can be changed': verifica che il valore di component.\textit{getSelectedValue}(mappedRow) è 'Cool' e il valore di component.\textit{getSelectedValue}(mappedRow) è 'Heat'.}{Esito atteso: il valore di component.\textit{getSelectedValue}(mappedRow) è 'Cool'; il valore di component.\textit{getSelectedValue}(mappedRow) è 'Heat'.}{S}
\addTestI{Contesto \textit{EndpointTableComponent} deep integration; caso '\textit{onWriteRow} success toggles executing state and exposes success feedback': verifica che il valore di component.isRowExecuting(mappedRow) è true e deviceApiMock.writeDatapointValue viene invocato con \{ datapointId: 'dp-cmd-1', value: 'Heat', \}.}{Esito atteso: il valore di component.isRowExecuting(mappedRow) è true; deviceApiMock.writeDatapointValue viene invocato con \{ datapointId: 'dp-cmd-1', value: 'Heat', \}; il valore di component.isRowExecuting(mappedRow) è false; il valore di component.\textit{getRowFeedback}(mappedRow)?.type è 'success'.}{S}
\addTestI{Contesto \textit{EndpointTableComponent} deep integration; caso '\textit{onWriteRow} ignores duplicate invocation while row is already executing': verifica che deviceApiMock.writeDatapointValue viene invocato 1 volte.}{Esito atteso: deviceApiMock.writeDatapointValue viene invocato 1 volte.}{S}
\addTestI{Contesto \textit{EndpointTableComponent} deep integration; caso '\textit{onWriteRow} with no selectable enum value does not call API': verifica che deviceApiMock.writeDatapointValue non viene invocato.}{Esito atteso: deviceApiMock.writeDatapointValue non viene invocato.}{S}
\addTestI{Contesto \textit{EndpointTableComponent} deep integration; caso '\textit{onWriteRow} error path sets error feedback and clears executing state': verifica che il valore di component.isRowExecuting(mappedRow) è false e il valore di component.\textit{getRowFeedback}(mappedRow)?.type è 'error'.}{Esito atteso: il valore di component.isRowExecuting(mappedRow) è false; il valore di component.\textit{getRowFeedback}(mappedRow)?.type è 'error'.}{S}
\addTestI{Contesto \textit{EndpointTableComponent} deep integration; caso 'feedback is auto-cleared after timeout': verifica che il valore di component.\textit{getRowFeedback}(mappedRow)?.type è 'success' e component.\textit{getRowFeedback}(mappedRow) è null.}{Esito atteso: il valore di component.\textit{getRowFeedback}(mappedRow)?.type è 'success'; component.\textit{getRowFeedback}(mappedRow) è null.}{S}
\addTestI{Contesto \textit{EndpointTableComponent} deep integration; caso 'active-plant-changed event resets transient state and triggers a refresh': verifica che il valore di component.\textit{getSelectedValue}(mappedRow) è 'Cool' e il valore di component.\textit{loadError} è ''.}{Esito atteso: il valore di component.\textit{getSelectedValue}(mappedRow) è 'Cool'; il valore di component.\textit{loadError} è ''; deviceApiMock.\textit{getWritableEndpointRows}.mock.calls.length è maggiore di callsBefore.}{S}
\addTestI{Contesto \textit{EndpointTableComponent} deep integration; caso 'polling periodically reloads current values': verifica che deviceApiMock.getCurrentValuePointsByDeviceIds.mock.calls.length è maggiore di callsBefore.}{Esito atteso: deviceApiMock.getCurrentValuePointsByDeviceIds.mock.calls.length è maggiore di callsBefore.}{S}
\addTestI{Contesto \textit{EndpointTableComponent} deep integration; caso 'load current values error branch keeps component usable and returns fallback dash': verifica che il valore di component.\textit{getCurrentValue}(mappedRow) è '-'.}{Esito atteso: il valore di component.\textit{getCurrentValue}(mappedRow) è '-'.}{S}
\addTestI{Contesto \textit{EndpointTableComponent}; caso 'creates component and loads writable endpoints': verifica che component è valorizzato e deviceApiMock.\textit{getWritableEndpointRows} viene invocato 1 volte.}{Esito atteso: component è valorizzato; deviceApiMock.\textit{getWritableEndpointRows} viene invocato 1 volte; deviceApiMock.getCurrentValuePointsByDeviceIds viene invocato 1 volte.}{S}
\addTestI{Contesto \textit{EndpointTableComponent}; caso 'renders mapped endpoint label for known sfeType': verifica che il valore di component.\textit{getEndpointLabel}(mappedRow) è 'Comando cambio modalita HVAC'.}{Esito atteso: il valore di component.\textit{getEndpointLabel}(mappedRow) è 'Comando cambio modalita HVAC'.}{S}
\addTestI{Contesto \textit{EndpointTableComponent}; caso 'falls back to a runtime humanized label when sfeType is unknown': verifica che il valore di component.\textit{getEndpointLabel}(unknownRow) è 'Comando future feature'.}{Esito atteso: il valore di component.\textit{getEndpointLabel}(unknownRow) è 'Comando future feature'.}{S}
\addTestI{Contesto \textit{EndpointTableComponent}; caso 'returns semantic current value by matching Cmd endpoint to State datapoint': verifica che il valore di component.\textit{getCurrentValue}(mappedRow) è 'Heat'.}{Esito atteso: il valore di component.\textit{getCurrentValue}(mappedRow) è 'Heat'.}{S}
\addTestI{Contesto \textit{EndpointTableComponent}; caso 'returns fallback when current datapoint value is unavailable': verifica che il valore di component.\textit{getCurrentValue}(unknownRow) è '-'.}{Esito atteso: il valore di component.\textit{getCurrentValue}(unknownRow) è '-'.}{S}
\addTestI{Contesto DeviceInteraction feature integration; caso 'RF79-OBL carica endpoint scrivibili e valori correnti alla init': verifica che component è valorizzato e deviceApiStub.\textit{getWritableEndpointRows} viene invocato 1 volte.}{Esito atteso: component è valorizzato; deviceApiStub.\textit{getWritableEndpointRows} viene invocato 1 volte; deviceApiStub.getCurrentValuePointsByDeviceIds viene invocato 1 volte.}{S}
\addTestI{Contesto DeviceInteraction feature integration; caso 'RF80-OBL espone label endpoint leggibili per sfeType noti e ignoti': verifica che il valore di component.\textit{getEndpointLabel}(mappedRow) è 'Comando cambio modalita HVAC' e il valore di component.\textit{getEndpointLabel}(unknownRow) è 'Comando future feature'.}{Esito atteso: il valore di component.\textit{getEndpointLabel}(mappedRow) è 'Comando cambio modalita HVAC'; il valore di component.\textit{getEndpointLabel}(unknownRow) è 'Comando future feature'.}{S}
\addTestI{Contesto DeviceInteraction feature integration; caso 'RF81-OBL risolve valore corrente dal datapoint state corrispondente': verifica che il valore di component.\textit{getCurrentValue}(mappedRow) è 'Heat' e il valore di component.\textit{getCurrentValue}(unknownRow) è '-'.}{Esito atteso: il valore di component.\textit{getCurrentValue}(mappedRow) è 'Heat'; il valore di component.\textit{getCurrentValue}(unknownRow) è '-'.}{S}
\addTestI{Contesto \textit{MyVimarPageComponent}; caso 'should create': verifica che component è valorizzato.}{Esito atteso: component è valorizzato.}{S}
\addTestI{Contesto \textit{MyVimarPageComponent}; caso 'carica lo stato account in \textit{ngOnInit}': verifica che serviceStub.\textit{getLinkedAccount} viene invocato 1 volte.}{Esito atteso: serviceStub.\textit{getLinkedAccount} viene invocato 1 volte.}{S}
\addTestI{Contesto \textit{MyVimarPageComponent}; caso 'delegates \textit{onLinkAccount} al service': verifica che serviceStub.initiateOAuth viene invocato 1 volte.}{Esito atteso: serviceStub.initiateOAuth viene invocato 1 volte.}{S}
\addTestI{Contesto \textit{MyVimarPageComponent}; caso '\textit{onUnlinkAccount} richiama il service e triggera refresh': verifica che serviceStub.unlinkAccount viene invocato 1 volte e serviceStub.\textit{getLinkedAccount} viene invocato 2 volte.}{Esito atteso: serviceStub.unlinkAccount viene invocato 1 volte; serviceStub.\textit{getLinkedAccount} viene invocato 2 volte.}{S}
\addTestI{Contesto \textit{MyVimarPageComponent}; caso 'imposta errore quando unlinkAccount fallisce': verifica che component.error contiene 'Errore durante la rimozione' e il valore di component.isLoading è false.}{Esito atteso: component.error contiene 'Errore durante la rimozione'; il valore di component.isLoading è false.}{S}
\addTestI{Contesto MyVimarIntegration feature integration; caso 'RF11-OBL carica stato account in init': verifica che component è valorizzato e serviceStub.\textit{getLinkedAccount} viene invocato 1 volte.}{Esito atteso: component è valorizzato; serviceStub.\textit{getLinkedAccount} viene invocato 1 volte.}{S}
\addTestI{Contesto MyVimarIntegration feature integration; caso 'RF12-OBL unlink account esegue refresh stato': verifica che serviceStub.unlinkAccount viene invocato 1 volte e serviceStub.\textit{getLinkedAccount} viene invocato 2 volte.}{Esito atteso: serviceStub.unlinkAccount viene invocato 1 volte; serviceStub.\textit{getLinkedAccount} viene invocato 2 volte.}{S}
\addTestI{Contesto MyVimarIntegration feature integration; caso 'RF13-OBL fallback errore quando unlink fallisce': verifica che consoleErrorSpy viene invocato e component.error contiene 'Errore durante la rimozione'.}{Esito atteso: consoleErrorSpy viene invocato; component.error contiene 'Errore durante la rimozione'; il valore di component.isLoading è false.}{S}
\addTestI{Contesto \textit{NotificationPageComponent}; caso 'should create': verifica che component è valorizzato.}{Esito atteso: component è valorizzato.}{S}
\addTestI{Contesto \textit{NotificationPageComponent}; caso 'renderizza uno stato vuoto accessibile quando non ci sono notifiche': verifica che emptyState è valorizzato e emptyState.textContent contiene 'Nessuna notifica disponibile'.}{Esito atteso: emptyState è valorizzato; emptyState.textContent contiene 'Nessuna notifica disponibile'; items ha lunghezza 0.}{S}
\addTestI{Contesto \textit{NotificationPageComponent}; caso 'renderizza la lista notifiche e passa i dati ai figli': verifica che summary.textContent contiene '2 notifiche non lette' e items ha lunghezza 2.}{Esito atteso: summary.textContent contiene '2 notifiche non lette'; items ha lunghezza 2; items[0].componentInstance.notification() coincide con notificationA; items[1].componentInstance.notification() coincide con notificationB.}{S}
\addTestI{Contesto \textit{NotificationPageComponent}; caso 'usa il singolare nel riepilogo quando unreadCount e 1': verifica che summary.textContent contiene '1 notifica non letta'.}{Esito atteso: summary.textContent contiene '1 notifica non letta'.}{S}
\addTestI{Contesto \textit{NotificationPageComponent}; caso 'aggiorna il DOM quando vm\$ emette un nuovo snapshot coerente': verifica che items ha lunghezza 1 e items ha lunghezza 2.}{Esito atteso: items ha lunghezza 1; items ha lunghezza 2.}{S}
\addTestI{Contesto \textit{NotificationPageComponent}; caso 'invoca \textit{clearAllNotifications} quando si clicca Cancella tutte': verifica che notificationServiceStub.\textit{clearAllNotifications} viene invocato con [ notificationA, notificationB, ].}{Esito atteso: notificationServiceStub.\textit{clearAllNotifications} viene invocato con [ notificationA, notificationB, ].}{S}
\addTestI{Contesto \textit{NotificationPageComponent}; caso 'evidenzia leggermente la notifica richiesta da query param focus': verifica che il valore di items[0].componentInstance.isHighlighted() è false e il valore di items[1].componentInstance.isHighlighted() è true.}{Esito atteso: il valore di items[0].componentInstance.isHighlighted() è false; il valore di items[1].componentInstance.isHighlighted() è true.}{S}
\addTestI{Contesto Notification feature integration; caso 'RF102-OPL visualizza stato vuoto quando non ci sono notifiche': verifica che component è valorizzato e emptyState è valorizzato.}{Esito atteso: component è valorizzato; emptyState è valorizzato; emptyState.textContent contiene 'Nessuna notifica disponibile'.}{S}
\addTestI{Contesto Notification feature integration; caso 'RF103-OPL renderizza lista notifiche e clear-all sul servizio': verifica che items ha lunghezza 2 e notificationServiceStub.\textit{clearAllNotifications} viene invocato con [notificationA, notificationB].}{Esito atteso: items ha lunghezza 2; notificationServiceStub.\textit{clearAllNotifications} viene invocato con [notificationA, notificationB].}{S}
\addTestI{Contesto Notification feature integration; caso 'RF104-OPL focus query param evidenzia la notifica corretta': verifica che il valore di items[0].componentInstance.isHighlighted() è false e il valore di items[1].componentInstance.isHighlighted() è true.}{Esito atteso: il valore di items[0].componentInstance.isHighlighted() è false; il valore di items[1].componentInstance.isHighlighted() è true.}{S}
\addTestI{Contesto \textit{NotificationTopbarPanel} feature integration; caso 'RF105-OPL stato vuoto mostra messaggio e nasconde clear-all': verifica che empty?.textContent contiene 'Nessuna notifica disponibile' e \textit{clearAll} è null.}{Esito atteso: empty?.textContent contiene 'Nessuna notifica disponibile'; \textit{clearAll} è null.}{S}
\addTestI{Contesto \textit{NotificationTopbarPanel} feature integration; caso 'RF103-OPL mostra max 6 notifiche e clear-all emette evento': verifica che il valore di items.length è 6 e \textit{clearSpy} viene invocato 1 volte.}{Esito atteso: il valore di items.length è 6; \textit{clearSpy} viene invocato 1 volte.}{S}
\addTestI{Contesto \textit{NotificationTopbarPanel} feature integration; caso 'RF104-OPL click su notifica e remove emettono gli id corretti': verifica che selectedSpy viene invocato con 'n-2' e \textit{removeSpy} viene invocato con 'n-2'.}{Esito atteso: selectedSpy viene invocato con 'n-2'; \textit{removeSpy} viene invocato con 'n-2'.}{S}
\addTestI{Contesto \textit{NotificationTopbarPanel} feature integration; caso 'RF102-OPL applica parsing priorita e titolo visualizzato': verifica che component.\textit{getPriority}(notifications[0]) non è null e il valore di component.\textit{getDisplayTitle}(notifications[0]) è 'Allarme antincendio'.}{Esito atteso: component.\textit{getPriority}(notifications[0]) non è null; il valore di component.\textit{getDisplayTitle}(notifications[0]) è 'Allarme antincendio'; il valore di component.\textit{getPrioritySymbol}(component.\textit{getPriority}(notifications[0])!) è '!'.}{S}
\addTestI{Contesto \textit{AuthBaseComponent}; caso '\textit{handleSuccess} naviga a first-access quando isFirstAccess e true': verifica che il valore di component.isLoading è false e component.errorType è null.}{Esito atteso: il valore di component.isLoading è false; component.errorType è null; routerMock.navigate viene invocato con ['/auth/first-access'].}{S}
\addTestI{Contesto \textit{AuthBaseComponent}; caso '\textit{handleSuccess} naviga a dashboard quando isFirstAccess e false': verifica che routerMock.navigate viene invocato con ['/dashboard'].}{Esito atteso: routerMock.navigate viene invocato con ['/dashboard'].}{S}
\addTestI{Contesto \textit{AuthBaseComponent}; caso '\textit{handleSuccess} naviga su returnUrl quando presente': verifica che routerMock.navigateByUrl viene invocato con '/vimar-link' e routerMock.navigate non viene invocato con ['/dashboard'].}{Esito atteso: routerMock.navigateByUrl viene invocato con '/vimar-link'; routerMock.navigate non viene invocato con ['/dashboard'].}{S}
\addTestI{Contesto \textit{AuthBaseComponent}; caso '\textit{handleError} imposta USERNAME\_OR\_PASSWORD\_WRONG su 400 e 401': verifica che il valore di component.errorType è AuthErrorType.USERNAME\_OR\_PASSWORD\_WRONG.}{Esito atteso: il valore di component.errorType è AuthErrorType.USERNAME\_OR\_PASSWORD\_WRONG.}{S}
\addTestI{Contesto \textit{AuthBaseComponent}; caso '\textit{handleError} imposta NEW\_PASSWORD\_EQUALS\_TEMP su 409': verifica che il valore di component.errorType è AuthErrorType.NEW\_PASSWORD\_EQUALS\_TEMP.}{Esito atteso: il valore di component.errorType è AuthErrorType.NEW\_PASSWORD\_EQUALS\_TEMP.}{S}
\addTestI{Contesto \textit{AuthBaseComponent}; caso '\textit{handleError} usa fallback USERNAME\_OR\_PASSWORD\_WRONG per errori non Http': verifica che il valore di component.errorType è AuthErrorType.USERNAME\_OR\_PASSWORD\_WRONG e il valore di component.isLoading è false.}{Esito atteso: il valore di component.errorType è AuthErrorType.USERNAME\_OR\_PASSWORD\_WRONG; il valore di component.isLoading è false.}{S}
\addTestI{Contesto \textit{AuthBaseComponent}; caso 'espone il messaggio umano per credenziali errate': verifica che il valore di component.loginErrorMessage è 'Utente non trovato: username o password errati.'.}{Esito atteso: il valore di component.loginErrorMessage è 'Utente non trovato: username o password errati.'.}{S}
\addTestI{Contesto \textit{LoginComponent}; caso 'crea il componente': verifica che component è valorizzato.}{Esito atteso: component è valorizzato.}{S}
\addTestI{Contesto \textit{LoginComponent}; caso 'non invoca login se il form non rispetta i validator': verifica che markAllAsTouchedSpy viene invocato e authServiceMock.login non viene invocato.}{Esito atteso: markAllAsTouchedSpy viene invocato; authServiceMock.login non viene invocato.}{S}
\addTestI{Contesto \textit{LoginComponent}; caso 'invoca login e naviga su dashboard se autenticazione ok': verifica che authServiceMock.login viene invocato con 'mrossi' e routerMock.navigate viene invocato con ['/dashboard'].}{Esito atteso: authServiceMock.login viene invocato con 'mrossi'; routerMock.navigate viene invocato con ['/dashboard']; component.errorType è null; il valore di component.isLoading è false.}{S}
\addTestI{Contesto \textit{LoginComponent}; caso 'invoca login anche con valori presenti solo nel DOM (autofill browser)': verifica che authServiceMock.login viene invocato con 'mrossi'.}{Esito atteso: authServiceMock.login viene invocato con 'mrossi'.}{S}
\addTestI{Contesto \textit{LoginComponent}; caso 'invoca login e naviga su returnUrl se presente': verifica che authServiceMock.login viene invocato con 'mrossi' e routerMock.navigateByUrl viene invocato con '/vimar-link'.}{Esito atteso: authServiceMock.login viene invocato con 'mrossi'; routerMock.navigateByUrl viene invocato con '/vimar-link'.}{S}
\addTestI{Contesto \textit{LoginComponent}; caso 'imposta errore USERNAME\_OR\_PASSWORD\_WRONG su 401': verifica che il valore di component.errorType è AuthErrorType.USERNAME\_OR\_PASSWORD\_WRONG e il valore di component.isLoading è false.}{Esito atteso: il valore di component.errorType è AuthErrorType.USERNAME\_OR\_PASSWORD\_WRONG; il valore di component.isLoading è false.}{S}
\addTestI{Contesto \textit{LoginComponent}; caso 'espone il messaggio umano già al primo tentativo fallito': verifica che il valore di component.isLoading è false e il valore di component.loginErrorMessage è 'Utente non trovato: username o password errati.'.}{Esito atteso: il valore di component.isLoading è false; il valore di component.loginErrorMessage è 'Utente non trovato: username o password errati.'.}{S}
\addTestI{Contesto \textit{LoginComponent}; caso 'renderizza il testo errore quando errorType è valorizzato': verifica che content contiene 'Utente non trovato: username o password errati.'.}{Esito atteso: content contiene 'Utente non trovato: username o password errati.'.}{S}
\addTestI{Contesto UserAuth feature integration; caso 'RF1-OBL login valido reindirizza in dashboard': verifica che authServiceStub.login viene invocato con 'mrossi' e routerStub.navigate viene invocato con ['/dashboard'].}{Esito atteso: authServiceStub.login viene invocato con 'mrossi'; routerStub.navigate viene invocato con ['/dashboard']; component.errorType è null.}{S}
\addTestI{Contesto UserAuth feature integration; caso 'RF2-OBL login invalido espone errore credenziali': verifica che il valore di component.errorType è AuthErrorType.USERNAME\_OR\_PASSWORD\_WRONG e component.loginErrorMessage contiene 'username o password errati'.}{Esito atteso: il valore di component.errorType è AuthErrorType.USERNAME\_OR\_PASSWORD\_WRONG; component.loginErrorMessage contiene 'username o password errati'.}{S}
\addTestI{Contesto UserAuth feature integration; caso 'RF3-OBL form invalido non invoca autenticazione': verifica che authServiceStub.login non viene invocato.}{Esito atteso: authServiceStub.login non viene invocato.}{S}
\addTestI{Contesto \textit{UserCreatedDialogComponent}; caso 'should create': verifica che component è valorizzato.}{Esito atteso: component è valorizzato.}{S}
\addTestI{Contesto \textit{UserCreatedDialogComponent}; caso 'mostra nome utente e password temporanea': verifica che textContent contiene 'Mario' e textContent contiene 'mrossi'.}{Esito atteso: textContent contiene 'Mario'; textContent contiene 'mrossi'; textContent contiene 'TempPass123'.}{S}
\addTestI{Contesto \textit{UserCreatedDialogComponent}; caso 'closeDialog emette evento closed': verifica che closedSpy viene invocato 1 volte.}{Esito atteso: closedSpy viene invocato 1 volte.}{S}
\addTestI{Contesto \textit{UserCreatedDialogComponent}; caso 'click sul bottone chiudi invoca closeDialog': verifica che closeSpy viene invocato 1 volte.}{Esito atteso: closeSpy viene invocato 1 volte.}{S}
\addTestI{Contesto \textit{UserListComponent}; caso 'should create': verifica che component è valorizzato.}{Esito atteso: component è valorizzato.}{S}
\addTestI{Contesto \textit{UserListComponent}; caso '\textit{onDelete} apre il dialog di conferma': verifica che component.pendingDeleteUser() coincide con user.}{Esito atteso: component.pendingDeleteUser() coincide con user.}{S}
\addTestI{Contesto \textit{UserListComponent}; caso '\textit{onDeleteConfirm} emette \textit{deleteUser} con id numerico': verifica che emitSpy viene invocato con 1 e emitSpy viene invocato 1 volte.}{Esito atteso: emitSpy viene invocato con 1; emitSpy viene invocato 1 volte; component.pendingDeleteUser() è null.}{S}
\addTestI{Contesto \textit{UserListComponent}; caso '\textit{onDeleteCancel} non emette \textit{deleteUser} e chiude il dialog': verifica che emitSpy non viene invocato e component.pendingDeleteUser() è null.}{Esito atteso: emitSpy non viene invocato; component.pendingDeleteUser() è null.}{S}
\addTestI{Contesto UserManagement feature integration; caso 'RF16-OBL inizializza pagina e carica elenco utenti': verifica che component è valorizzato e \textit{getUsersStub} viene invocato 1 volte.}{Esito atteso: component è valorizzato; \textit{getUsersStub} viene invocato 1 volte; il valore di component.isCreateFormOpen() è false.}{S}
\addTestI{Contesto UserManagement feature integration; caso 'RF17-OBL submit creazione utente aggiorna stato e aggiorna la lista': verifica che \textit{createUserStub} viene invocato con \{ name: 'Mario', surname: 'Rossi', username: 'mrossi', \} e il valore di component.createdResponse()?.tempPassword è 'TempPass123'.}{Esito atteso: \textit{createUserStub} viene invocato con \{ name: 'Mario', surname: 'Rossi', username: 'mrossi', \}; il valore di component.createdResponse()?.tempPassword è 'TempPass123'; \textit{getUsersStub} viene invocato 2 volte.}{S}
\addTestI{Contesto UserManagement feature integration; caso 'RF18-OBL errore 409 in creazione espone errore username in uso': verifica che il valore di component.formError() è UserManagementErrorType.USERNAME\_ALREADY\_IN\_USE.}{Esito atteso: il valore di component.formError() è UserManagementErrorType.USERNAME\_ALREADY\_IN\_USE.}{S}
\addTestI{Contesto \textit{AssignWardDialogComponent}; caso 'should create': verifica che component è valorizzato.}{Esito atteso: component è valorizzato.}{S}
\addTestI{Contesto \textit{AssignWardDialogComponent}; caso 'dovrebbe resettare plantId a null in \textit{ngOnInit}': verifica che component.form.controls.plantId.value è null.}{Esito atteso: component.form.controls.plantId.value è null.}{S}
\addTestI{Contesto \textit{AssignWardDialogComponent}; caso 'non dovrebbe emettere submitted se il form e invalido': verifica che il valore di component.form.invalid è true e il valore di component.form.controls.plantId.touched è true.}{Esito atteso: il valore di component.form.invalid è true; il valore di component.form.controls.plantId.touched è true; submittedSpy non viene invocato.}{S}
\addTestI{Contesto \textit{AssignWardDialogComponent}; caso 'dovrebbe emettere submitted con plantId quando il form e valido': verifica che submittedSpy viene invocato con \{ plantId: '101' \} e submittedSpy viene invocato 1 volte.}{Esito atteso: submittedSpy viene invocato con \{ plantId: '101' \}; submittedSpy viene invocato 1 volte.}{S}
\addTestI{Contesto \textit{AssignWardDialogComponent}; caso 'dovrebbe emettere cancelled quando \textit{onCancel} viene chiamato': verifica che cancelledSpy viene invocato 1 volte.}{Esito atteso: cancelledSpy viene invocato 1 volte.}{S}
\addTestI{Contesto \textit{WardManagementPageComponent}; caso 'should create': verifica che component è valorizzato.}{Esito atteso: component è valorizzato.}{S}
\addTestI{Contesto \textit{WardManagementPageComponent}; caso '\textit{ngOnInit} dovrebbe chiamare \textit{loadWards} e sincronizzare snapshots/error': verifica che storeStub.\textit{loadWards} viene invocato 1 volte e component.wardsSnapshot() coincide con [].}{Esito atteso: storeStub.\textit{loadWards} viene invocato 1 volte; component.wardsSnapshot() coincide con []; component.snackbarMessage() è null; component.wardsSnapshot() ha lunghezza 2.}{S}
\addTestI{Contesto \textit{WardManagementPageComponent}; caso 'quando wards diventa vuoto resetta selezioni e step mobile': verifica che component.selectedWardId() è null e component.selectedApartmentId() è null.}{Esito atteso: component.selectedWardId() è null; component.selectedApartmentId() è null; il valore di component.mobileStep() è 'wards'.}{S}
\addTestI{Contesto \textit{WardManagementPageComponent}; caso 'selectedApartment si riallinea al ward selezionato': verifica che component.selectedApartment() coincide con \{ id: '101', name: 'App. 101' \} e component.selectedApartment() è null.}{Esito atteso: component.selectedApartment() coincide con \{ id: '101', name: 'App. 101' \}; component.selectedApartment() è null.}{S}
\addTestI{Contesto \textit{WardManagementPageComponent}; caso '\textit{ngOnDestroy} dovrebbe interrompere aggiornamenti da stream': verifica che il valore di component.snackbarMessage() è 'Errore iniziale'.}{Esito atteso: il valore di component.snackbarMessage() è 'Errore iniziale'.}{S}
\addTestI{Contesto \textit{WardManagementPageComponent}; caso '\textit{dismissSnackbar} dovrebbe azzerare il messaggio': verifica che component.snackbarMessage() è null.}{Esito atteso: component.snackbarMessage() è null.}{S}
\addTestI{Contesto \textit{WardManagementPageComponent}; caso '\textit{onCreateWard} dovrebbe aprire dialog in create mode': verifica che il valore di component.selectedWardId() è wardA.id e il valore di component.wardDialogMode() è 'create'.}{Esito atteso: il valore di component.selectedWardId() è wardA.id; il valore di component.wardDialogMode() è 'create'.}{S}
\addTestI{Contesto \textit{WardManagementPageComponent}; caso '\textit{onCreateWardSubmit} dovrebbe delegare a store e chiudere dialog': verifica che storeStub.\textit{createWard} viene invocato con \{ name: 'Oncologia' \} e storeStub.\textit{createWard} viene invocato 1 volte.}{Esito atteso: storeStub.\textit{createWard} viene invocato con \{ name: 'Oncologia' \}; storeStub.\textit{createWard} viene invocato 1 volte; il valore di component.wardDialogMode() è 'closed'.}{S}
\addTestI{Contesto \textit{WardManagementPageComponent}; caso '\textit{onEditWard} e \textit{onEditWardSubmit} dovrebbero aggiornare il reparto selezionato': verifica che il valore di component.selectedWardId() è wardA.id e il valore di component.wardDialogMode() è 'edit'.}{Esito atteso: il valore di component.selectedWardId() è wardA.id; il valore di component.wardDialogMode() è 'edit'; storeStub.\textit{updateWard} viene invocato con 1; il valore di component.wardDialogMode() è 'closed'.}{S}
\addTestI{Contesto \textit{WardManagementPageComponent}; caso '\textit{onEditWardSubmit} non dovrebbe chiamare store se selectedWard è null': verifica che storeStub.\textit{updateWard} non viene invocato.}{Esito atteso: storeStub.\textit{updateWard} non viene invocato.}{S}
\addTestI{Contesto \textit{WardManagementPageComponent}; caso '\textit{onCloseWardDialog} dovrebbe chiudere dialog e pulire selezione': verifica che il valore di component.selectedWardId() è wardA.id e il valore di component.wardDialogMode() è 'closed'.}{Esito atteso: il valore di component.selectedWardId() è wardA.id; il valore di component.wardDialogMode() è 'closed'.}{S}
\addTestI{Contesto \textit{WardManagementPageComponent}; caso '\textit{onCloseAssignOperatorDialog} e \textit{onCloseAssignPlantDialog} ripuliscono stato di caricamento e lookup': verifica che component.operatorWardId() è null e component.availableOperatorsFromFetch() è null.}{Esito atteso: component.operatorWardId() è null; component.availableOperatorsFromFetch() è null; il valore di component.isLoadingAvailableOperators() è false; component.plantWardId() è null.}{S}
\addTestI{Contesto \textit{WardManagementPageComponent}; caso '\textit{selectWard} dovrebbe aggiornare ward attivo e step mobile': verifica che il valore di component.selectedWardId() è 2 e il valore di component.mobileStep() è 'apartments'.}{Esito atteso: il valore di component.selectedWardId() è 2; il valore di component.mobileStep() è 'apartments'.}{S}
\addTestI{Contesto \textit{WardManagementPageComponent}; caso '\textit{selectWard} sul ward già selezionato mantiene la selezione dell appartamento': verifica che il valore di component.selectedWardId() è wardA.id e il valore di component.selectedApartmentId() è '101'.}{Esito atteso: il valore di component.selectedWardId() è wardA.id; il valore di component.selectedApartmentId() è '101'; il valore di component.mobileStep() è 'apartments'.}{S}
\addTestI{Contesto \textit{WardManagementPageComponent}; caso '\textit{selectApartment} e step helpers dovrebbero aggiornare lo stato ui': verifica che il valore di component.selectedApartmentId() è '101' e il valore di component.mobileStep() è 'wards'.}{Esito atteso: il valore di component.selectedApartmentId() è '101'; il valore di component.mobileStep() è 'wards'; il valore di component.mobileStep() è 'apartments'; il valore di component.mobileStep() è 'operators'.}{S}
\addTestI{Contesto \textit{WardManagementPageComponent}; caso '\textit{onEditSelectedWard} non dovrebbe fare nulla senza ward selezionato': verifica che component.selectedWardId() è null e il valore di component.wardDialogMode() è 'closed'.}{Esito atteso: component.selectedWardId() è null; il valore di component.wardDialogMode() è 'closed'.}{S}
\addTestI{Contesto \textit{WardManagementPageComponent}; caso '\textit{onAssignOperatorSubmit} non dovrebbe chiamare store senza wardId': verifica che storeStub.\textit{assignOperator} non viene invocato.}{Esito atteso: storeStub.\textit{assignOperator} non viene invocato.}{S}
\addTestI{Contesto \textit{WardManagementPageComponent}; caso 'flow plant dovrebbe impostare wardId, submit e chiudere': verifica che il valore di component.plantWardId() è 2 e storeStub.\textit{getAvailablePlantsForWard} viene invocato con 2.}{Esito atteso: il valore di component.plantWardId() è 2; storeStub.\textit{getAvailablePlantsForWard} viene invocato con 2; component.availablePlants() coincide con [\{ id: '103', name: 'App. 103' \}]; storeStub.\textit{assignPlant} viene invocato con 2.}{S}
\addTestI{Contesto \textit{WardManagementPageComponent}; caso '\textit{onAssignPlantSubmit} non dovrebbe chiamare store senza wardId': verifica che storeStub.\textit{assignPlant} non viene invocato.}{Esito atteso: storeStub.\textit{assignPlant} non viene invocato.}{S}
\addTestI{Contesto \textit{WardManagementPageComponent}; caso 'confirmState + \textit{onConfirmDialogConfirmed} dovrebbero gestire i 3 rami': verifica che component.confirmState() coincide con \{ kind: 'delete-ward', wardId: 1 \} e storeStub.\textit{deleteWard} viene invocato con 1.}{Esito atteso: component.confirmState() coincide con \{ kind: 'delete-ward', wardId: 1 \}; storeStub.\textit{deleteWard} viene invocato con 1; component.confirmState() è null; component.confirmState() coincide con \{ kind: 'remove-operator', wardId: 1, userId: 2, \}.}{S}
\addTestI{Contesto \textit{WardManagementPageComponent}; caso '\textit{onConfirmDialogConfirmed} non dovrebbe fare nulla con stato nullo': verifica che storeStub.\textit{deleteWard} non viene invocato e storeStub.\textit{removeOperator} non viene invocato.}{Esito atteso: storeStub.\textit{deleteWard} non viene invocato; storeStub.\textit{removeOperator} non viene invocato; storeStub.\textit{removePlant} non viene invocato.}{S}
\addTestI{Contesto \textit{WardManagementPageComponent}; caso '\textit{onConfirmDialogCancelled} dovrebbe pulire confirmState': verifica che component.confirmState() è null.}{Esito atteso: component.confirmState() è null.}{S}
\addTestI{Contesto \textit{WardManagementPageComponent}; caso '\textit{getConfirmMessage} dovrebbe restituire il testo corretto per ogni stato': verifica che il valore di component.\textit{getConfirmMessage}() è 'Confermi questa operazione?' e il valore di component.\textit{getConfirmLabel}() è 'Conferma'.}{Esito atteso: il valore di component.\textit{getConfirmMessage}() è 'Confermi questa operazione?'; il valore di component.\textit{getConfirmLabel}() è 'Conferma'; il valore di component.\textit{getConfirmMessage}() è 'Confermi l\textbackslash\{\}'eliminazione del reparto?'; il valore di component.\textit{getConfirmMessage}() è 'Confermi la rimozione dell\textbackslash\{\}'operatore dal reparto?'.}{S}
\addTestI{Contesto \textit{WardManagementPageComponent}; caso '\textit{getOperatorDisplayName} dovrebbe comporre e trimmare nome e cognome': verifica che il valore di component.\textit{getOperatorDisplayName}('Mario', 'Rossi') è 'Mario Rossi' e il valore di component.\textit{getOperatorDisplayName}('Mario', '') è 'Mario'.}{Esito atteso: il valore di component.\textit{getOperatorDisplayName}('Mario', 'Rossi') è 'Mario Rossi'; il valore di component.\textit{getOperatorDisplayName}('Mario', '') è 'Mario'.}{S}
\addTestI{Contesto \textit{WardManagementPageComponent}; caso 'availablePlants dovrebbe deduplicare e escludere quelli già assegnati al ward selezionato': verifica che component.availablePlants() coincide con [].}{Esito atteso: component.availablePlants() coincide con [].}{S}
\addTestI{Contesto \textit{WardManagementPageComponent}; caso '\textit{onAssignPlant} non usa fallback locale se fetch ad-hoc fallisce': verifica che component.availablePlants() coincide con [] e il valore di component.isLoadingAvailablePlants() è false.}{Esito atteso: component.availablePlants() coincide con []; il valore di component.isLoadingAvailablePlants() è false.}{S}
\addTestI{Contesto WardManagement feature integration; caso 'RF27-OBL mantiene coerente la selezione del ward e aggiorna il nome in entrambe le sezioni': verifica che wardApiStub.\textit{getWards} viene invocato 1 volte e nativeElement.querySelector('section[aria-label="Dettagli reparto e appartamenti"] h2')?.textC... contiene ward1.name.}{Esito atteso: wardApiStub.\textit{getWards} viene invocato 1 volte; nativeElement.querySelector('section[aria-label="Dettagli reparto e appartamenti"] h2')?.textC... contiene ward1.name; il valore di wardButtons.length è 2; nativeElement.querySelector('section[aria-label="Dettagli reparto e appartamenti"] h2')?.textC... contiene ward2.name.}{S}
\addTestI{Contesto WardManagement feature integration; caso 'RF28-OBL aggiorna il nome del ward e sincronizza lista e pannello centrale': verifica che nativeElement.querySelector('section[aria-label="Dettagli reparto e appartamenti"] h2')?.textC... contiene ward1.name e \textit{getDialog}()?.textContent contiene 'Modifica reparto'.}{Esito atteso: nativeElement.querySelector('section[aria-label="Dettagli reparto e appartamenti"] h2')?.textC... contiene ward1.name; \textit{getDialog}()?.textContent contiene 'Modifica reparto'; wardApiStub.\textit{updateWard} viene invocato con 1; nativeElement.querySelector('section[aria-label="Dettagli reparto e appartamenti"] h2')?.textC... contiene 'Cardiologia A'.}{S}
\addTestI{Contesto WardManagement feature integration; caso 'RF29-OBL dopo la creazione mantiene cliccabile il ward appena creato': verifica che \textit{getDialog}()?.textContent contiene 'Crea reparto' e wardApiStub.\textit{createWard} viene invocato con \{ name: 'Oncologia' \}.}{Esito atteso: \textit{getDialog}()?.textContent contiene 'Crea reparto'; wardApiStub.\textit{createWard} viene invocato con \{ name: 'Oncologia' \}; il valore di wardButtons.length è 3; wardButtons[2].textContent contiene 'Oncologia'.}{S}
\addTestI{Contesto WardManagement feature integration; caso 'RF30-OBL assegna un operatore dal dialog e invoca il flusso di assegnazione': verifica che \textit{getDialog}()?.textContent contiene 'Assegna operatore sanitario' e wardApiStub.\textit{getAvailableOperators} viene invocato 1 volte.}{Esito atteso: \textit{getDialog}()?.textContent contiene 'Assegna operatore sanitario'; wardApiStub.\textit{getAvailableOperators} viene invocato 1 volte; \textit{assignOperatorDialog} è valorizzato; wardApiStub.\textit{assignOperatorToWard} viene invocato con 1.}{S}
\addTestI{Contesto WardManagement feature integration; caso 'RF31-OBL non mostra appartamenti assegnabili se fetch disponibile fallisce': verifica che \textit{getDialog}()?.textContent contiene 'Assegna appartamento' e wardApiStub.\textit{getAvailablePlants} viene invocato 1 volte.}{Esito atteso: \textit{getDialog}()?.textContent contiene 'Assegna appartamento'; wardApiStub.\textit{getAvailablePlants} viene invocato 1 volte; apartmentOptions non contiene 'App. 102'; apartmentOptions non contiene 'App. 101'.}{S}
\addTestI{Contesto WardManagement feature integration; caso 'RF32-OBL rimuove operatore e appartamento passando dal confirm dialog': verifica che \textit{getDialog}()?.textContent contiene 'rimozione dell\textbackslash\{\}'operatore' e wardApiStub.\textit{removeOperatorFromWard} viene invocato con 1.}{Esito atteso: \textit{getDialog}()?.textContent contiene 'rimozione dell\textbackslash\{\}'operatore'; wardApiStub.\textit{removeOperatorFromWard} viene invocato con 1; \textit{getDialog}()?.textContent contiene 'rimozione dell\textbackslash\{\}'appartamento'; wardApiStub.\textit{removePlantFromWard} viene invocato con 1.}{S}
\addTestI{Contesto WardManagement feature integration; caso 'RF33-OBL elimina un ward dopo conferma e aggiorna la selezione': verifica che \textit{getDialog}()?.textContent contiene 'eliminazione del reparto' e wardApiStub.\textit{deleteWard} viene invocato con 1.}{Esito atteso: \textit{getDialog}()?.textContent contiene 'eliminazione del reparto'; wardApiStub.\textit{deleteWard} viene invocato con 1; il valore di \textit{getWardButtons}().length è 1; nativeElement.querySelector('section[aria-label="Dettagli reparto e appartamenti"] h2')?.textC... contiene ward2.name.}{S}

\addTestI{Contesto \textit{Auth Integration Test}; caso 'should login and set refresh cookie with real auth flow': verifica che response.body contiene la proprietà accessToken e il set-cookie header è definito.}{Esito atteso: response.body contiene la proprietà accessToken; il set-cookie header è definito; il valore di HTTP status è 201.}{S}
\addTestI{Contesto \textit{Auth Integration Test}; caso 'should execute first-login end to end and persist credentials change': verifica che la password dell'utente viene aggiornata nel database e il valore di first\_access diventa false.}{Esito atteso: la password dell'utente viene aggiornata nel database; il valore di first\_access diventa false; il valore di HTTP status è 201.}{S}
\addTestI{Contesto \textit{Auth Integration Test}; caso 'should refresh token from cookie generated by /auth/login': verifica che il refresh endpoint ritorna un nuovo accessToken e utilizza il refresh cookie.}{Esito atteso: il refresh endpoint ritorna un nuovo accessToken; il nuovo token è valido; il valore di HTTP status è 201.}{S}
\addTestI{Contesto \textit{Auth Integration Test}; caso 'should return 401 on /auth/refresh without cookie': verifica che il refresh endpoint restituisce un errore quando il cookie non è presente.}{Esito atteso: il valore di HTTP status è 401; la risposta non contiene accessToken.}{S}

\addTestI{Contesto \textit{Alarms Integration Test}; caso 'should create alarm rule through real controller/service/repository': verifica che il POST /alarm-rules crea una nuova regola di allarme nel database.}{Esito atteso: response.body contiene la proprietà name con valore 'Too Hot'; il numero di regole nel database aumenta di 1; il valore di HTTP status è 201.}{S}
\addTestI{Contesto \textit{Alarms Integration Test}; caso 'should list alarm rules': verifica che il GET /alarm-rules ritorna una lista di regole di allarme.}{Esito atteso: response.body è un array non vuoto; le regole contengono la proprietà id; il valore di HTTP status è 200.}{S}
\addTestI{Contesto \textit{Alarms Integration Test}; caso 'should get alarm rule by id': verifica che il GET /alarm-rules/:id ritorna una specifica regola di allarme.}{Esito atteso: response.body contiene la proprietà id con valore cercato; la regola contiene tutte le proprietà richieste; il valore di HTTP status è 200.}{S}
\addTestI{Contesto \textit{Alarms Integration Test}; caso 'should list managed alarm events for a user': verifica che il GET /alarm-events/managed/:userId/:Ward/:offset ritorna gli allarmi risolti.}{Esito atteso: response.body è un array; todos gli elementi hanno resolution\_time valorizzato; il valore di HTTP status è 200.}{S}
\addTestI{Contesto \textit{Alarms Integration Test}; caso 'should list unmanaged alarm events for a user': verifica che il GET /alarm-events/unmanaged/:userId/:Ward/:offset ritorna gli allarmi attivi.}{Esito atteso: response.body è un array; gli allarmi hanno resolution\_time null; il valore di HTTP status è 200.}{S}
\addTestI{Contesto \textit{Alarms Integration Test}; caso 'should resolve alarm event and emit resolution flow': verifica che il PATCH /alarm-events/resolve risolve un allarme e aggiorna il database.}{Esito atteso: l'allarme nel database ha resolution\_time valorizzato; l'user\_id e user\_username vengono salvati; il valore di HTTP status è 200.}{S}
\addTestI{Contesto \textit{Alarms Integration Test}; caso 'should update and delete alarm rules through real service': verifica che il PUT e DELETE aggiornano/cancellano le regole nel database.}{Esito atteso: l'aggiornamento modifica le proprietà della regola; la cancellazione rimuove la regola dal database; il valore di HTTP status è appropriato.}{S}
\addTestI{Contesto \textit{Alarms Integration Test}; caso 'should execute checkAlarmRule flow and create alarm event': verifica che il flusso di controllo regola allarme crea un evento di allarme quando la soglia viene superata.}{Esito attesto: il numero di eventi allarme aumenta di 1; il nuovo evento ha i parametri corretti; il valore del alarm\_rule\_id corrisponde alla regola.}{S}

\addTestI{Contesto \textit{Users Integration Test}; caso 'should list users using real controller/service/adapter/repository chain': verifica che il GET /users ritorna l'elenco di tutti gli utenti.}{Esito atteso: response.body è un array con 2 elementi; il primo utente contiene le proprietà id, username, surname, name; il valore di HTTP status è 200.}{S}
\addTestI{Contesto \textit{Users Integration Test}; caso 'should create user and return generated tempPassword': verifica che il POST /users crea un nuovo utente e genera una password temporanea.}{Esito atteso: response.body contiene le proprietà id, username, tempPassword; il nuovo utente viene salvato nel database; il valore di HTTP status è 201.}{S}
\addTestI{Contesto \textit{Users Integration Test}; caso 'should validate payload on create user': verifica che il POST /users con payload non valido ritorna un errore di validazione.}{Esito atteso: il valore di HTTP status è 400; la risposta contiene i dettagli dell'errore di validazione.}{S}

\addTestI{Contesto \textit{Wards Integration Test}; caso 'should create and list wards with real service and adapters': verifica che il POST /wards crea un nuovo reparto e il GET /wards ritorna la lista aggiornata.}{Esito atteso: il numero di reparti aumenta di 1; il nuovo reparto è presente nella lista dei reparti; il valore di HTTP status è 201 per POST e 200 per GET.}{S}
\addTestI{Contesto \textit{Wards Integration Test}; caso 'should add and remove user from ward': verifica che il POST /wards-users-relationships aggiunge un utente al reparto e il DELETE lo rimuove.}{Esito atteso: dopo POST, l'utente è associato al reparto; dopo DELETE, l'utente non è più associato al reparto; il valore di HTTP status è appropriato.}{S}
\addTestI{Contesto \textit{Wards Integration Test}; caso 'should add and list plants for a ward': verifica che il POST /wards-plants-relationships aggiunge un impianto al reparto e il GET ritorna la lista.}{Esito atteso: il numero di impianti associati al reparto aumenta di 1; il nuovo impianto è presente nella lista; il valore di HTTP status è appropriato.}{S}

\addTestI{Contesto \textit{Plant Integration Test}; caso 'should retrieve plant by id with real repository': verifica che il GET /plant con plantid query param ritorna l'impianto corretto.}{Esito atteso: response.body contiene le proprietà id, name; il valore di HTTP status è 200; l'impianto ritornato corrisponde a quello cercato.}{S}
\addTestI{Contesto \textit{Plant Integration Test}; caso 'should return only available plants': verifica che il GET /plant/available ritorna solo i pianti disponibili.}{Esito atteso: response.body è un array che contiene solo pianti disponibili; il valore di HTTP status è 200.}{S}
\addTestI{Contesto \textit{Plant Integration Test}; caso 'should return all plants': verifica che il GET /plant/all ritorna tutti i pianti nel sistema.}{Esito atteso: response.body è un array con 2 elementi; il valore di HTTP status è 200.}{S}
\addTestI{Contesto \textit{Plant Integration Test}; caso 'should return 404 when plantid query param is missing': verifica che il GET /plant senza parametro ritorna un errore 404.}{Esito atteso: il valore di HTTP status è 404; la risposta contiene un messaggio di errore.}{S}
\addTestI{Contesto \textit{Plant Integration Test}; caso 'should return fallback message when no available plants are found': verifica che il GET /plant/available ritorna un messaggio fallback quando non ci sono pianti disponibili.}{Esito atteso: response.body contiene il messaggio 'No available plants found'; il valore di status code nel messaggio è 202.}{S}

\addTestI{Contesto \textit{Analytics Integration Test}; caso 'should retrieve all analytics given the plantId': verifica che il GET /analytics/:plantId ritorna una lista di grafici di analisi.}{Esito atteso: response.body è un array non vuoto; i grafici contengono le proprietà metric, unit, labels, series; il valore di HTTP status è 200.}{S}
\addTestI{Contesto \textit{Analytics Integration Test}; caso 'should return one PlotDto per strategy': verifica che ogni strategia di analisi ritorna un grafico.}{Esito atteso: il numero di grafici corrisponde al numero di strategie; ogni grafico ha proprietà distinct; il valore di HTTP status è 200.}{S}
\addTestI{Contesto \textit{Analytics Integration Test}; caso 'should return 500 when all strategies fail': verifica che quando tutte le strategie falliscono il sistema ritorna un errore 500.}{Esito atteso: il valore di HTTP status è 500; la risposta contiene un messaggio di errore.}{S}
\addTestI{Contesto \textit{Analytics Integration Test}; caso 'should return 404 when plantId query param is missing': verifica che il GET /analytics senza plantId ritorna un errore 404.}{Esito atteso: il valore di HTTP status è 404.}{S}
\addTestI{Contesto \textit{Analytics Integration Test}; caso 'should throw when suggestion service receives no suggestion from LLM port': verifica che quando l'LLM non ritorna suggerimenti il servizio genera un errore.}{Esito atteso: il servizio genera un errore con messaggio appropriato.}{S}

\addTestI{Contesto \textit{Cache Integration Test}; caso 'should process webhook and write plant cache using real cache flow': verifica che il POST /cache/update processa il webhook e scrive la cache dell'impianto.}{Esito atteso: response.body contiene success: true; il valore di HTTP status è 202; la cache dell'impianto viene scritta nel sistema.}{S}
\addTestI{Contesto \textit{Cache Integration Test}; caso 'should refresh all plants on fetched.tokens event': verifica che l'evento fetched.tokens aggiorna la cache di tutti gli impianti.}{Esito atteso: le cache di tutti i pianti vengono aggiornate; il flusso completa correttamente.}{S}
\addTestI{Contesto \textit{Cache Integration Test}; caso 'should reject updateCache use case when plantId is missing': verifica che il use case rifiuta la richiesta quando plantId manca.}{Esito atteso: il sistema genera un errore con messaggio 'PlantId is null'.}{S}
\addTestI{Contesto \textit{Cache Integration Test}; caso 'should handle webhook update failure with Error without crashing': verifica che gli errori nei webhook vengono gestiti senza bloccare il sistema.}{Esito atteso: response.body contiene success: true; il sistema continua a funzionare; il valore di HTTP status è 202.}{S}
\addTestI{Contesto \textit{Cache Integration Test}; caso 'should skip duplicate cache sync events while one sync is already running': verifica che gli eventi di sincronizzazione duplicati vengono ignorati mentre une sincronizzazione è in corso.}{Esito atteso: solo une sincronizzazione è eseguita; gli eventi duplicati non creano nuove sincronizzazioni.}{S}

\addTestI{Contesto \textit{Notifications Integration Test}; caso 'should handle alarm.activated event and persist notification': verifica che l'evento alarm.activated crea e persiste una notifica nel database.}{Esito atteso: la notifica viene scritta nel database; il valore di HTTP status è appropriato; il mock pool è invocato correttamente.}{S}
\addTestI{Contesto \textit{Notifications Integration Test}; caso 'should handle alarm.resolved event and persist notification': verifica che l'evento alarm.resolved crea una notifica di risoluzione.}{Esito atteso: la notifica viene scritta nel database; l'evento viene elaborato correttamente.}{S}
\addTestI{Contesto \textit{Notifications Integration Test}; caso 'should ignore invalid alarm payload without throwing': verifica che i payload non validi vengono ignorati senza errori.}{Esito atteso: il sistema non genera eccezioni; nessuna query viene eseguita per payload null.}{S}
\addTestI{Contesto \textit{Notifications Integration Test}; caso 'should reject ward notification when alarmEventId is missing': verifica che la notifica viene rifiutata quando alarmEventId manca.}{Esito atteso: il sistema genera un messaggio di errore appropriato.}{S}
\addTestI{Contesto \textit{Notifications Integration Test}; caso 'should reject resolution notification when wardId is missing': verifica che la notifica di risoluzione viene rifiutata quando wardId manca.}{Esito atteso: il sistema genera un messaggio di errore appropriato.}{S}

\addTestI{Contesto \textit{ApiAuthVimar Integration Test}; caso 'should return account status from real token service and cache repository': verifica che il GET /my-vimar/account ritorna lo stato dell'account MyVimar collegato.}{Esito atteso: response.body contiene le proprietà isLinked e email; il valore di HTTP status è 200.}{S}
\addTestI{Contesto \textit{ApiAuthVimar Integration Test}; caso 'should prepare and authorize OAuth ticket end-to-end': verifica che il flusso OAuth completo funziona: preparazione -> autorizzazione.}{Esito atteso: il POST /api/auth/prepare-oauth ritorna un ticket; il GET /api/auth/authorize reindirizza al servizio OAuth; il valore di HTTP status è 302 per il redirect.}{S}
\addTestI{Contesto \textit{ApiAuthVimar Integration Test}; caso 'should handle callback and persist tokens with real token flow': verifica che il callback OAuth persiste i token nel sistema.}{Esito atteso: il GET /my-vimar/callback elabora il codice OAuth e persiste i token; il valore di HTTP status è 302 per il redirect.}{S}
\addTestI{Contesto \textit{ApiAuthVimar Integration Test}; caso 'should disconnect account by clearing token cache': verifica che il DELETE /my-vimar/account disconnette l'account e pulisce la cache.}{Esito atteso: il valore di HTTP status è 200; il token viene rimosso dalla cache.}{S}
\addTestI{Contesto \textit{ApiAuthVimar Integration Test}; caso 'should return 401 when prepare-oauth user identity is invalid': verifica che la preparazione OAuth fallisce con identità utente non valida.}{Esito atteso: il valore di HTTP status è 401; la risposta non contiene un ticket.}{S}
\addTestI{Contesto \textit{ApiAuthVimar Integration Test}; caso 'should return 400 when authorize redirect\_url is missing': verifica che l'autorizzazione fallisce quando redirect\_url manca.}{Esito atteso: il valore di HTTP status è 400.}{S}
\addTestI{Contesto \textit{ApiAuthVimar Integration Test}; caso 'should return 401 when authorize ticket is invalid': verifica che l'autorizzazione fallisce con ticket non valido.}{Esito atteso: il valore di HTTP status è 401.}{S}

\addTestI{Contesto \textit{Subscription Integration Test}; caso 'should refresh subscriptions for all plants on cache.all.updated': verifica che l'evento cache.all.updated aggiorna le sottoscrizioni di tutti i pianti.}{Esito atteso: le sottoscrizioni di entrambi i pianti vengono aggiornate; il flusso completa correttamente.}{S}
\addTestI{Contesto \textit{Subscription Integration Test}; caso 'should refresh subscriptions for one plant on cache.updated': verifica che l'evento cache.updated aggiorna la sottoscrizione di un singolo impianto.}{Esito atteso: solo la sottoscrizione dell'impianto specificato viene aggiornata.}{S}
\addTestI{Contesto \textit{Subscription Integration Test}; caso 'should execute node cron flow with real service': verifica che il flusso cron dei nodi viene eseguito correttamente.}{Esito atteso: le sottoscrizioni dei nodi vengono aggiornate; il servizio completa senza errori.}{S}
\addTestI{Contesto \textit{Subscription Integration Test}; caso 'should skip duplicate full-cache refresh event while another refresh is running': verifica che gli eventi di refresh duplicati vengono ignorati quando un refresh è in corso.}{Esito atteso: solo un refresh è eseguito; gli eventi duplicati non creano nuovi refresh.}{S}
\addTestI{Contesto \textit{Subscription Integration Test}; caso 'should return false when datapoint refresh fails in refreshAllSubscription': verifica che il fallimento del refresh datapoint causa il ritorno di false.}{Esito atteso: il flusso continua; il valore di ritorno è false.}{S}

\addTestI{Contesto \textit{Device Integration Test}; caso 'should retrieve a device by ID and return correct response structure': verifica che il GET /device/:id ritorna il dispositivo con la struttura corretta.}{Esito atteso: response.body contiene le proprietà id, plantId, name, type, datapoints; è un array di datapoint; il valore di HTTP status è 200.}{S}
\addTestI{Contesto \textit{Device Integration Test}; caso 'should retrieve all devices for a plant': verifica che il GET /device/plant/:plantId ritorna i dispositivi dell'impianto.}{Esito atteso: response.body è un array di dispositivi; tutti i dispositivi appartengono al plantId specificato; il valore di HTTP status è 200.}{S}
\addTestI{Contesto \textit{Device Integration Test}; caso 'should accept write datapoint request and return 202': verifica che il POST /device accetta la richiesta di scrittura e ritorna 202.}{Esito atteso: il valore di HTTP status è 202; la richiesta viene accodata per l'elaborazione asincrona.}{S}
\addTestI{Contesto \textit{Device Integration Test}; caso 'should reject request without datapointId': verifica che il POST /device rifiuta le richieste senza datapointId.}{Esito atteso: il valore di HTTP status è 400; la risposta contiene un errore di validazione.}{S}
\addTestI{Contesto \textit{Device Integration Test}; caso 'should accept webhook payload with datapoint updates': verifica che il POST /device/update processa i dati di aggiornamento dal webhook.}{Esito atteso: il valore di HTTP status è 202; i datapoint vengono aggiornati; gli allarmi vengono controllati.}{S}
\addTestI{Contesto \textit{Device Integration Test}; caso 'should filter non-datapoint items from webhook payload': verifica che il webhook filtra gli elementi non datapoint.}{Esito atteso: solo i datapoint validi vengono elaborati; gli elementi non validi vengono ignorati.}{S}
\addTestI{Contesto \textit{Device Integration Test}; caso 'should retrieve current values for a device': verifica che il GET /device/:deviceId/value ritorna i valori attuali del dispositivo.}{Esito atteso: response.body è un array di valori datapoint; ogni valore contiene le proprietà necessarie; il valore di HTTP status è 200.}{S}
\addTestI{Contesto \textit{Device Integration Test}; caso 'should handle full device lifecycle: find -> write -> read values': verifica che il ciclo completo di gestione del dispositivo funziona correttamente.}{Esito atteso: il dispositivo viene trovato; il valore viene scritto; i valori sono letti correttamente; il valore di HTTP status è appropriato.}{S}
\addTestI{Contesto \textit{Device Integration Test}; caso 'should handle multiple devices in same plant': verifica che il sistema gestisce correttamente più dispositivi dello stesso impianto.}{Esito atteso: i dispositivi mantengono l'associazione corretta con l'impianto; le operazioni non interferiscono tra loro.}{S}
\addTestI{Contesto \textit{Device Integration Test}; caso 'should process webhook and trigger alarms concurrently': verifica che il webhook processa i dati e triggerizza allarmi in modo concorrente.}{Esito atteso: più dispositivi vengono elaborati contemporaneamente; gli allarmi vengono controllati per ogni datapoint; il flusso completa correttamente.}{S}

\makeTITable

\subsection{Test di unità}
\addTestU{Contesto appConfig initialization; caso 'runs \textit{EventSubscriptionService} initialization at bootstrap': verifica che \textit{initializeSpy} viene invocato con []. Esito atteso: \textit{initializeSpy} viene invocato con [].}{S}
\addTestU{Contesto App Routes; caso 'auth route configurata con lazy loading': verifica che route è definito e route?.\textit{loadChildren} è definito. Esito atteso: route è definito; route?.\textit{loadChildren} è definito; il valore di typeof route?.\textit{loadChildren} è 'function'.}{S}
\addTestU{Contesto App Routes; caso 'vimar-link route protetta da authGuard e roleGuard': verifica che route?.canActivate contiene authGuard e route?.canActivate contiene roleGuard. Esito atteso: route?.canActivate contiene authGuard; route?.canActivate contiene roleGuard.}{S}
\addTestU{Contesto App Routes; caso 'vimar-link route richiede ruolo AMMINISTRATORE': verifica che il valore di route?.data?.['requiredRole'] è UserRole.AMMINISTRATORE. Esito atteso: il valore di route?.data?.['requiredRole'] è UserRole.AMMINISTRATORE.}{S}
\addTestU{Contesto App Routes; caso 'route "" carica il main layout': verifica che route è definito e route?.canActivate contiene authGuard. Esito atteso: route è definito; route?.canActivate contiene authGuard; route?.\textit{loadChildren} è definito.}{S}
\addTestU{Contesto App Routes; caso 'route "**" fa redirect alla root': verifica che wildcardRoute è definito e il valore di wildcardRoute?.redirectTo è ''. Esito atteso: wildcardRoute è definito; il valore di wildcardRoute?.redirectTo è ''.}{S}
\addTestU{Contesto App; caso 'should create the app': verifica che app è valorizzato. Esito atteso: app è valorizzato.}{S}
\addTestU{Contesto App; caso 'should render router outlet': verifica che compiled.querySelector('router-outlet') è valorizzato. Esito atteso: compiled.querySelector('router-outlet') è valorizzato.}{S}
\addTestU{Contesto environment; caso 'espone la configurazione base attesa': verifica che il valore di typeof environment.production è 'boolean' e il valore di environment.apiBaseUrl.endsWith('/api') è true. Esito atteso: il valore di typeof environment.production è 'boolean'; il valore di environment.apiBaseUrl.endsWith('/api') è true.}{S}
\addTestU{Contesto Feature Routes Quick Coverage; caso 'copre alarm-configuration route e lazy component': verifica che il valore di route.path è '' e route.canActivate contiene authGuard. Esito atteso: il valore di route.path è ''; route.canActivate contiene authGuard; route.providers contiene \textit{AlarmConfigStateService}; il valore di typeof child?.\textit{loadComponent} è 'function'.}{S}
\addTestU{Contesto Feature Routes Quick Coverage; caso 'copre alarm-history route e lazy component': verifica che il valore di route.path è '' e route.canActivate contiene authGuard. Esito atteso: il valore di route.path è ''; route.canActivate contiene authGuard; il valore di typeof route.\textit{loadComponent} è 'function'; loaded è valorizzato.}{S}
\addTestU{Contesto Feature Routes Quick Coverage; caso 'copre alarm-management route e lazy component': verifica che il valore di route.path è '' e route.canActivate contiene authGuard. Esito atteso: il valore di route.path è ''; route.canActivate contiene authGuard; il valore di typeof route.\textit{loadComponent} è 'function'; loaded è valorizzato.}{S}
\addTestU{Contesto Feature Routes Quick Coverage; caso 'copre analytics route e lazy component': verifica che il valore di route.path è '' e il valore di typeof route.\textit{loadComponent} è 'function'. Esito atteso: il valore di route.path è ''; il valore di typeof route.\textit{loadComponent} è 'function'; loaded è valorizzato.}{S}
\addTestU{Contesto Feature Routes Quick Coverage; caso 'copre apartment-monitor route e lazy component': verifica che il valore di route.path è '' e route.canActivate contiene authGuard. Esito atteso: il valore di route.path è ''; route.canActivate contiene authGuard; il valore di typeof route.\textit{loadComponent} è 'function'; loaded è valorizzato.}{S}
\addTestU{Contesto Feature Routes Quick Coverage; caso 'copre dashboard route e lazy component': verifica che il valore di route.path è '' e il valore di typeof route.\textit{loadComponent} è 'function'. Esito atteso: il valore di route.path è ''; il valore di typeof route.\textit{loadComponent} è 'function'; loaded è valorizzato.}{S}
\addTestU{Contesto Feature Routes Quick Coverage; caso 'copre device-interaction routes e lazy component': verifica che DEVICE\_INTERACTION\_ROUTES ha lunghezza 2 e route.canActivate contiene authGuard. Esito atteso: DEVICE\_INTERACTION\_ROUTES ha lunghezza 2; route.canActivate contiene authGuard; il valore di typeof route.\textit{loadComponent} è 'function'; loaded è valorizzato.}{S}
\addTestU{Contesto Feature Routes Quick Coverage; caso 'copre my-vimar-integration route, provider binding e lazy component': verifica che il valore di route.path è '' e providerList contiene \textit{MyVimarCloudApiFeatureService}. Esito atteso: il valore di route.path è ''; providerList contiene \textit{MyVimarCloudApiFeatureService}; il valore di existingProvider?.useExisting è \textit{MyVimarCloudApiFeatureService}; child?.canActivate contiene authGuard.}{S}
\addTestU{Contesto Feature Routes Quick Coverage; caso 'copre notification route e lazy component': verifica che il valore di route.path è '' e route.canActivate contiene authGuard. Esito atteso: il valore di route.path è ''; route.canActivate contiene authGuard; il valore di typeof route.\textit{loadComponent} è 'function'; loaded è valorizzato.}{S}
\addTestU{Contesto Feature Routes Quick Coverage; caso 'copre user-management route e lazy component': verifica che il valore di route.path è '' e route.canActivate contiene authGuard. Esito atteso: il valore di route.path è ''; route.canActivate contiene authGuard; route.canActivate contiene adminGuard; il valore di typeof route.\textit{loadComponent} è 'function'.}{S}
\addTestU{Contesto Feature Routes Quick Coverage; caso 'copre ward-management route, providers e lazy component': verifica che il valore di route.path è '' e route.providers contiene \textit{WardStore}. Esito atteso: il valore di route.path è ''; route.providers contiene \textit{WardStore}; route.providers contiene \textit{WardOperationsService}; route.providers contiene \textit{AssignmentOperationsService}.}{S}
\addTestU{Contesto adminGuard; caso 'consente sempre l accesso': verifica che il valore di result è true. Esito atteso: il valore di result è true.}{S}
\addTestU{Contesto authGuard; caso 'restituisce true quando l'utente e autenticato': verifica che il valore di result è true e authServiceMock.restoreSessionFromRefresh non viene invocato. Esito atteso: il valore di result è true; authServiceMock.restoreSessionFromRefresh non viene invocato; routerMock.\textit{createUrlTree} non viene invocato.}{S}
\addTestU{Contesto authGuard; caso 'consente accesso quando restore sessione riesce': verifica che il valore di await firstValueFrom(result as ReturnType<typeof of>) è true e routerMock.\textit{createUrlTree} non viene invocato. Esito atteso: il valore di await firstValueFrom(result as ReturnType<typeof of>) è true; routerMock.\textit{createUrlTree} non viene invocato.}{S}
\addTestU{Contesto authGuard; caso 'redirect su /auth/login quando restore sessione fallisce': verifica che await firstValueFrom(result as ReturnType<typeof of>) coincide con \{ redirected: true \} e routerMock.\textit{createUrlTree} viene invocato con ['/auth/login']. Esito atteso: await firstValueFrom(result as ReturnType<typeof of>) coincide con \{ redirected: true \}; routerMock.\textit{createUrlTree} viene invocato con ['/auth/login'].}{S}
\addTestU{Contesto roleGuard; caso 'restituisce true quando la rotta non richiede ruolo': verifica che il valore di result è true e authServiceMock.hasRole non viene invocato. Esito atteso: il valore di result è true; authServiceMock.hasRole non viene invocato.}{S}
\addTestU{Contesto roleGuard; caso 'restituisce true quando il ruolo richiesto è presente': verifica che authServiceMock.hasRole viene invocato con UserRole.AMMINISTRATORE e il valore di result è true. Esito atteso: authServiceMock.hasRole viene invocato con UserRole.AMMINISTRATORE; il valore di result è true.}{S}
\addTestU{Contesto roleGuard; caso 'redirect su apartment-monitor quando il ruolo richiesto manca': verifica che routerMock.\textit{createUrlTree} viene invocato con ['/apartment-monitor'] e result coincide con \{ redirected: true \}. Esito atteso: routerMock.\textit{createUrlTree} viene invocato con ['/apartment-monitor']; result coincide con \{ redirected: true \}.}{S}
\addTestU{Contesto authInterceptor; caso 'inoltra la richiesta invariata quando il token non esiste': verifica che il valore di capturedRequest è request e il valore di capturedRequest?.headers.has('Authorization') è false. Esito atteso: il valore di capturedRequest è request; il valore di capturedRequest?.headers.has('Authorization') è false.}{S}
\addTestU{Contesto authInterceptor; caso 'aggiunge header Authorization Bearer quando il token esiste': verifica che il valore di capturedRequest non è request e il valore di capturedRequest?.headers.get('Authorization') è 'Bearer jwt-123'. Esito atteso: il valore di capturedRequest non è request; il valore di capturedRequest?.headers.get('Authorization') è 'Bearer jwt-123'.}{S}
\addTestU{Contesto authInterceptor; caso 'aggiunge header Authorization anche su /auth/first-login': verifica che il valore di capturedRequest?.headers.get('Authorization') è 'Bearer jwt-123'. Esito atteso: il valore di capturedRequest?.headers.get('Authorization') è 'Bearer jwt-123'.}{S}
\addTestU{Contesto authInterceptor; caso 'su 401 esegue refresh e ritenta con il nuovo access token': verifica che authServiceMock.\textit{refreshAccessToken} viene invocato 1 volte e forwardedRequests ha lunghezza 2. Esito atteso: authServiceMock.\textit{refreshAccessToken} viene invocato 1 volte; forwardedRequests ha lunghezza 2; il valore di forwardedRequests[0].headers.get('Authorization') è 'Bearer jwt-old'; il valore di forwardedRequests[1].headers.get('Authorization') è 'Bearer jwt-new'.}{S}
\addTestU{Contesto authInterceptor; caso 'non tenta refresh quando la 401 arriva da un endpoint auth': verifica che authServiceMock.\textit{refreshAccessToken} non viene invocato. Esito atteso: authServiceMock.\textit{refreshAccessToken} non viene invocato.}{S}
\addTestU{Contesto authInterceptor; caso 'non tenta refresh quando la 401 arriva da /auth/first-login': verifica che authServiceMock.\textit{refreshAccessToken} non viene invocato. Esito atteso: authServiceMock.\textit{refreshAccessToken} non viene invocato.}{S}
\addTestU{Contesto \textit{ApiErrorDisplayService}; caso 'usa fallback per errori non HTTP quando nonHttpStrategy non è message': verifica che il valore di message è 'Messaggio fallback'. Esito atteso: il valore di message è 'Messaggio fallback'.}{S}
\addTestU{Contesto \textit{ApiErrorDisplayService}; caso 'usa il messaggio degli errori non HTTP quando nonHttpStrategy e message': verifica che il valore di message è 'runtime failure'. Esito atteso: il valore di message è 'runtime failure'.}{S}
\addTestU{Contesto \textit{ApiErrorDisplayService}; caso 'mappa status 0 con actionLabel': verifica che il valore di message è 'Impossibile contattare il server durante caricare gli allarmi.'. Esito atteso: il valore di message è 'Impossibile contattare il server durante caricare gli allarmi.'.}{S}
\addTestU{Contesto \textit{ApiErrorDisplayService}; caso 'mappa status 401 con actionLabel': verifica che il valore di message è 'Sessione scaduta o non valida. Effettua nuovamente il login per aggiornare la vista.'. Esito atteso: il valore di message è 'Sessione scaduta o non valida. Effettua nuovamente il login per aggiornare la vista.'.}{S}
\addTestU{Contesto \textit{ApiErrorDisplayService}; caso 'mappa status 403 con forbiddenMessage custom se presente': verifica che il valore di message è 'Accesso negato dalla policy applicativa.'. Esito atteso: il valore di message è 'Accesso negato dalla policy applicativa.'.}{S}
\addTestU{Contesto \textit{ApiErrorDisplayService}; caso 'mappa status 403 usando actionLabel quando forbiddenMessage non è presente': verifica che il valore di message è 'Non hai i permessi necessari per eseguire l operazione.'. Esito atteso: il valore di message è 'Non hai i permessi necessari per eseguire l operazione.'.}{S}
\addTestU{Contesto \textit{ApiErrorDisplayService}; caso 'mappa status 500 con actionLabel': verifica che il valore di message è 'Errore interno del server durante completare la richiesta. Riprova tra poco.'. Esito atteso: il valore di message è 'Errore interno del server durante completare la richiesta. Riprova tra poco.'.}{S}
\addTestU{Contesto \textit{ApiErrorDisplayService}; caso 'usa error.error string quando disponibile per status non speciali': verifica che il valore di message è 'Payload non valido'. Esito atteso: il valore di message è 'Payload non valido'.}{S}
\addTestU{Contesto \textit{ApiErrorDisplayService}; caso 'usa error.error.message array quando disponibile per status non speciali': verifica che il valore di message è 'campo obbligatorio; formato non valido'. Esito atteso: il valore di message è 'campo obbligatorio; formato non valido'.}{S}
\addTestU{Contesto \textit{ApiErrorDisplayService}; caso 'usa fallback quando payload HTTP non contiene messaggi utili': verifica che il valore di message è 'Messaggio fallback'. Esito atteso: il valore di message è 'Messaggio fallback'.}{S}
\addTestU{Contesto \textit{BreadcrumbService}; caso 'should be created': verifica che service è valorizzato. Esito atteso: service è valorizzato.}{S}
\addTestU{Contesto \textit{BreadcrumbService}; caso 'should expose a breadcrumbs\$ observable': verifica che service.breadcrumbs\$ è definito. Esito atteso: service.breadcrumbs\$ è definito.}{S}
\addTestU{Contesto \textit{BreadcrumbService}; caso 'should emit an empty array immediately when there are no children': verifica che crumbs coincide con []. Esito atteso: crumbs coincide con [].}{S}
\addTestU{Contesto \textit{BreadcrumbService}; caso 'should emit breadcrumbs immediately on subscription without waiting for navigation': verifica che il valore di crumbs.length è 1 e crumbs[0] coincide con \{ label: 'Home', url: '/home' \}. Esito atteso: il valore di crumbs.length è 1; crumbs[0] coincide con \{ label: 'Home', url: '/home' \}.}{S}
\addTestU{Contesto \textit{BreadcrumbService}; caso 'should re-emit breadcrumbs on \textit{NavigationEnd}': verifica che il valore di emissions.length è 2. Esito atteso: il valore di emissions.length è 2.}{S}
\addTestU{Contesto \textit{BreadcrumbService}; caso 'should ignore router events that are not \textit{NavigationEnd}': verifica che crumbs coincide con [] e crumbsAfter coincide con []. Esito atteso: crumbs coincide con []; crumbsAfter coincide con [].}{S}
\addTestU{Contesto \textit{BreadcrumbService}; caso 'should build a single breadcrumb for a one-level route': verifica che crumbs coincide con [\{ label: 'About', url: '/about' \}]. Esito atteso: crumbs coincide con [\{ label: 'About', url: '/about' \}].}{S}
\addTestU{Contesto \textit{BreadcrumbService}; caso 'should skip routes without a breadcrumb data property': verifica che crumbs coincide con []. Esito atteso: crumbs coincide con [].}{S}
\addTestU{Contesto \textit{BreadcrumbService}; caso 'should skip routes with an empty url segment even if breadcrumb label is set': verifica che crumbs coincide con []. Esito atteso: crumbs coincide con [].}{S}
\addTestU{Contesto \textit{BreadcrumbService}; caso 'should handle multiple sibling top-level routes': verifica che crumbs coincide con [ \{ label: 'Products', url: '/products' \}, \{ label: 'Contact', url: '/contact' \}, ]. Esito atteso: crumbs coincide con [ \{ label: 'Products', url: '/products' \}, \{ label: 'Contact', url: '/contact' \}, ].}{S}
\addTestU{Contesto \textit{BreadcrumbService}; caso 'should build breadcrumbs for two-level nested routes': verifica che crumbs coincide con [ \{ label: 'Products', url: '/products' \}, \{ label: 'Detail', url: '/products/detail' \}, ]. Esito atteso: crumbs coincide con [ \{ label: 'Products', url: '/products' \}, \{ label: 'Detail', url: '/products/detail' \}, ].}{S}
\addTestU{Contesto \textit{BreadcrumbService}; caso 'should skip intermediate segments that have no breadcrumb label': verifica che crumbs coincide con [\{ label: 'Detail', url: '/products/detail' \}]. Esito atteso: crumbs coincide con [\{ label: 'Detail', url: '/products/detail' \}].}{S}
\addTestU{Contesto \textit{BreadcrumbService}; caso 'should concatenate multi-segment url paths correctly': verifica che crumbs coincide con [\{ label: 'Locale', url: '/en/us' \}]. Esito atteso: crumbs coincide con [\{ label: 'Locale', url: '/en/us' \}].}{S}
\addTestU{Contesto \textit{BreadcrumbService}; caso 'should not double-slash when a parent segment is empty': verifica che il valore di crumbs[0].url è '/settings' e crumbs[0].url non contiene '//'. Esito atteso: il valore di crumbs[0].url è '/settings'; crumbs[0].url non contiene '//'.}{S}
\addTestU{Contesto \textit{InternalAuthService}; caso 'esegue login, salva sessione e rende autenticato l'utente': verifica che il valore di request.request.method è 'POST' e il valore di request.request.withCredentials è true. Esito atteso: il valore di request.request.method è 'POST'; il valore di request.request.withCredentials è true; request.request.body coincide con \{ username: 'mrossi', password: loginCredential \}; il valore di session.username è 'mrossi'.}{S}
\addTestU{Contesto \textit{InternalAuthService}; caso 'chiama endpoint first-login con payload corretto': verifica che il valore di request.request.method è 'POST' e il valore di request.request.withCredentials è true. Esito atteso: il valore di request.request.method è 'POST'; il valore di request.request.withCredentials è true; request.request.body coincide con \{ username: 'mrossi', tempPassword: temporaryCredential, password: nextCredential, \}; il valore di session.username è 'mrossi'.}{S}
\addTestU{Contesto \textit{InternalAuthService}; caso 'chiama refresh, aggiorna il token e restituisce il nuovo accessToken': verifica che il valore di request.request.method è 'POST' e il valore di request.request.withCredentials è true. Esito atteso: il valore di request.request.method è 'POST'; il valore di request.request.withCredentials è true; request.request.body coincide con \{\}; il valore di refreshedToken è accessToken.}{S}
\addTestU{Contesto \textit{InternalAuthService}; caso 'ripristina sessione da refresh cookie quando token in memoria manca': verifica che il valore di request.request.method è 'POST' e il valore di request.request.withCredentials è true. Esito atteso: il valore di request.request.method è 'POST'; il valore di request.request.withCredentials è true; il valore di await restorePromise è true; il valore di service.isAuthenticated() è true.}{S}
\addTestU{Contesto \textit{InternalAuthService}; caso 'restituisce false quando restore sessione fallisce': verifica che il valore di await restorePromise è false e il valore di service.isAuthenticated() è false. Esito atteso: il valore di await restorePromise è false; il valore di service.isAuthenticated() è false.}{S}
\addTestU{Contesto \textit{InternalAuthService}; caso 'esegue una sola richiesta refresh con restore concorrenti': verifica che il valore di await restorePromiseOne è true e il valore di await restorePromiseTwo è true. Esito atteso: il valore di await restorePromiseOne è true; il valore di await restorePromiseTwo è true.}{S}
\addTestU{Contesto \textit{InternalAuthService}; caso 'resetta token e utente al logout': verifica che service.\textit{getToken}() è null e service.\textit{getRole}() è null. Esito atteso: service.\textit{getToken}() è null; service.\textit{getRole}() è null; il valore di service.isAuthenticated() è false.}{S}
\addTestU{Contesto \textit{InternalAuthService}; caso 'chiama /auth/logout con credenziali e resetta sessione locale': verifica che il valore di request.request.method è 'POST' e il valore di request.request.withCredentials è true. Esito atteso: il valore di request.request.method è 'POST'; il valore di request.request.withCredentials è true; request.request.body coincide con \{\}; service.\textit{getToken}() è null.}{S}
\addTestU{Contesto \textit{InternalAuthService}; caso 'espone stream utente corrente con valore iniziale null': verifica che currentUser è null. Esito atteso: currentUser è null.}{S}
\addTestU{Contesto \textit{AnalyticsComponent}; caso 'viene creato': verifica che component è valorizzato. Esito atteso: component è valorizzato.}{S}
\addTestU{Contesto \textit{AnalyticsComponent}; caso 'carica il primo appartamento di default': verifica che mockAnalyticsService.\textit{getAllApartments} viene invocato e mockAnalyticsService.\textit{getAnalytics} viene invocato con '1'. Esito atteso: mockAnalyticsService.\textit{getAllApartments} viene invocato; mockAnalyticsService.\textit{getAnalytics} viene invocato con '1'.}{F}
\addTestU{Contesto \textit{AnalyticsComponent}; caso 'aggiorna grafici al cambio appartamento': verifica che mockAnalyticsService.\textit{getAnalytics} viene invocato con '2'. Esito atteso: mockAnalyticsService.\textit{getAnalytics} viene invocato con '2'.}{F}
\addTestU{Contesto \textit{AnalyticsComponent}; caso 'should return undefined if metric is not found': verifica che result è undefined. Esito atteso: result è undefined.}{S}
\addTestU{Contesto \textit{AnalyticsComponent}; caso 'should find the correct metric in \textit{getChartByMetric}': verifica che result è definito e il valore di result?.metric è 'plant-anomalies'. Esito atteso: result è definito; il valore di result?.metric è 'plant-anomalies'.}{S}
\addTestU{Contesto \textit{MainLayoutComponent}; caso 'carica i \textit{NavItem} con il ruolo corretto': verifica che mockNavService.\textit{getNavItems} viene invocato con UserRole.OPERATORE\_SANITARIO e component.navItems.length è maggiore di 0. Esito atteso: mockNavService.\textit{getNavItems} viene invocato con UserRole.OPERATORE\_SANITARIO; component.navItems.length è maggiore di 0.}{S}
\addTestU{Contesto \textit{MainLayoutComponent}; caso 'reindirizza al login quando non esiste sessione utente': verifica che component.navItems coincide con [] e navigateSpy viene invocato con ['/auth/login']. Esito atteso: component.navItems coincide con []; navigateSpy viene invocato con ['/auth/login'].}{S}
\addTestU{Contesto \textit{MainLayoutComponent}; caso 'mostra avviso topbar per admin senza account MyVimar collegato': verifica che mockMyVimarService.\textit{getLinkedAccount} viene invocato 1 volte e il valore di component.showVimarAssociationWarning è true. Esito atteso: mockMyVimarService.\textit{getLinkedAccount} viene invocato 1 volte; il valore di component.showVimarAssociationWarning è true.}{S}
\addTestU{Contesto \textit{MainLayoutComponent}; caso 'non mostra avviso topbar quando account MyVimar risulta collegato': verifica che il valore di component.showVimarAssociationWarning è false. Esito atteso: il valore di component.showVimarAssociationWarning è false.}{S}
\addTestU{Contesto \textit{MainLayoutComponent}; caso 'inverte correttamente isCollapsed con il segnale ricevuto': verifica che il valore di component.isCollapsed è true e il valore di component.isCollapsed è false. Esito atteso: il valore di component.isCollapsed è true; il valore di component.isCollapsed è false.}{S}
\addTestU{Contesto \textit{MainLayoutComponent}; caso 'invoca correttamente logout': verifica che mockAuthService.logoutFromBackend viene invocato e navigateSpy viene invocato con ['/auth/login']. Esito atteso: mockAuthService.logoutFromBackend viene invocato; navigateSpy viene invocato con ['/auth/login'].}{S}
\addTestU{Contesto \textit{MainLayoutComponent}; caso 'apre il pannello profilo e carica stato MyVimar': verifica che il valore di component.isProfilePanelOpen è true e mockMyVimarService.\textit{getLinkedAccount} viene invocato 2 volte. Esito atteso: il valore di component.isProfilePanelOpen è true; mockMyVimarService.\textit{getLinkedAccount} viene invocato 2 volte; component.vimarAccount coincide con \{ email: '', isLinked: false \}.}{S}
\addTestU{Contesto \textit{MainLayoutComponent}; caso 'apre il pannello profilo anche per utenti non admin ma senza caricare stato MyVimar': verifica che il valore di component.isProfilePanelOpen è true e mockMyVimarService.\textit{getLinkedAccount} non viene invocato. Esito atteso: il valore di component.isProfilePanelOpen è true; mockMyVimarService.\textit{getLinkedAccount} non viene invocato.}{S}
\addTestU{Contesto \textit{MainLayoutComponent}; caso 'naviga a vimar-link dalla sezione profilo': verifica che navigateSpy viene invocato con ['/vimar-link']. Esito atteso: navigateSpy viene invocato con ['/vimar-link'].}{S}
\addTestU{Contesto \textit{MainLayoutComponent}; caso 'richiede refresh allarmi attivi quando si clicca la stessa voce menu della pagina corrente': verifica che mockAlarmManagementRefreshService.\textit{requestRefresh} viene invocato 1 volte. Esito atteso: mockAlarmManagementRefreshService.\textit{requestRefresh} viene invocato 1 volte.}{S}
\addTestU{Contesto \textit{MainLayoutComponent}; caso 'dovrebbe collegarsi correttamente al conteggio delle notifiche non lette': verifica che il valore di result è mockCount. Esito atteso: il valore di result è mockCount.}{S}
\addTestU{Contesto \textit{MainLayoutComponent}; caso 'apre automaticamente il pannello notifiche quando arriva una nuova notifica realtime': verifica che il valore di component.isNotificationPanelOpen è false e il valore di component.isNotificationPanelOpen è true. Esito atteso: il valore di component.isNotificationPanelOpen è false; il valore di component.isNotificationPanelOpen è true.}{S}
\addTestU{Contesto Main Layout Routes; caso 'carica \textit{MainLayoutComponent} in modalità lazy': verifica che rootRoute?.\textit{loadComponent} è definito e il valore di typeof rootRoute?.\textit{loadComponent} è 'function'. Esito atteso: rootRoute?.\textit{loadComponent} è definito; il valore di typeof rootRoute?.\textit{loadComponent} è 'function'.}{S}
\addTestU{Contesto Main Layout Routes; caso 'dashboard è figlio': verifica che dashboard è definito e dashboard?.canActivate contiene authGuard. Esito atteso: dashboard è definito; dashboard?.canActivate contiene authGuard.}{S}
\addTestU{Contesto Main Layout Routes - Gestione Allarmi (Nested Routes); caso 'alarms route esiste e ha dei figli': verifica che alarmsRoute è definito e alarmsRoute?.children è definito. Esito atteso: alarmsRoute è definito; alarmsRoute?.children è definito.}{S}
\addTestU{Contesto Main Layout Routes - Gestione Allarmi (Nested Routes); caso 'configurazione allarmi riservata ad AMMINISTRATORE': verifica che config?.canActivate contiene roleGuard e il valore di config?.data?.['requiredRole'] è UserRole.AMMINISTRATORE. Esito atteso: config?.canActivate contiene roleGuard; il valore di config?.data?.['requiredRole'] è UserRole.AMMINISTRATORE.}{S}
\addTestU{Contesto Main Layout Routes - Gestione Allarmi (Nested Routes); caso 'gestione e storico allarmi richiedono authGuard': verifica che mgmt?.canActivate contiene authGuard e history?.canActivate contiene authGuard. Esito atteso: mgmt?.canActivate contiene authGuard; history?.canActivate contiene authGuard.}{S}
\addTestU{Contesto Main Layout Routes - Sicurezza e Ruoli; caso 'routes riservate con roleGuard e ruolo AMMINISTRATORE': verifica che route?.canActivate contiene roleGuard e il valore di route?.data?.['requiredRole'] è UserRole.AMMINISTRATORE. Esito atteso: route?.canActivate contiene roleGuard; il valore di route?.data?.['requiredRole'] è UserRole.AMMINISTRATORE.}{S}
\addTestU{Contesto Main Layout Routes; caso 'wildcard reindirizza a dashboard': verifica che il valore di wildcard?.redirectTo è 'dashboard'. Esito atteso: il valore di wildcard?.redirectTo è 'dashboard'.}{S}
\addTestU{Contesto user-auth.routes; caso 'espone i path login e first-access con lazy load dei componenti': verifica che il valore di typeof loginRoute?.\textit{loadComponent} è 'function' e il valore di typeof firstAccessRoute?.\textit{loadComponent} è 'function'. Esito atteso: il valore di typeof loginRoute?.\textit{loadComponent} è 'function'; il valore di typeof firstAccessRoute?.\textit{loadComponent} è 'function'; loginComponent è definito; firstAccessComponent è definito.}{S}
\addTestU{Contesto user-auth.routes; caso 'contiene redirect di default verso login': verifica che il valore di defaultRedirect?.pathMatch è 'full' e il valore di defaultRedirect?.redirectTo è 'login'. Esito atteso: il valore di defaultRedirect?.pathMatch è 'full'; il valore di defaultRedirect?.redirectTo è 'login'.}{S}
\addTestU{Contesto device-taxonomy; caso 'resolves thermostat when type token is explicit': verifica che il valore di resolved è DeviceType.THERMOSTAT. Esito atteso: il valore di resolved è DeviceType.THERMOSTAT.}{S}
\addTestU{Contesto device-taxonomy; caso 'infers thermostat from endpoint sfeType when raw type is ambiguous': verifica che il valore di resolved è DeviceType.THERMOSTAT. Esito atteso: il valore di resolved è DeviceType.THERMOSTAT.}{S}
\addTestU{Contesto device-taxonomy; caso 'infers thermostat from HVAC command endpoint when raw type is ambiguous': verifica che il valore di resolved è DeviceType.THERMOSTAT. Esito atteso: il valore di resolved è DeviceType.THERMOSTAT.}{S}
\addTestU{Contesto device-taxonomy; caso 'infers entrance door from italian device name when raw type is ambiguous': verifica che il valore di resolved è DeviceType.ENTRANCE\_DOOR. Esito atteso: il valore di resolved è DeviceType.ENTRANCE\_DOOR.}{S}
\addTestU{Contesto device-taxonomy; caso 'infers scenario devices from name instead of falling back to light': verifica che il valore di resolved è DeviceType.ALARM\_BUTTON. Esito atteso: il valore di resolved è DeviceType.ALARM\_BUTTON.}{S}
\addTestU{Contesto device-taxonomy; caso 'falls back to unknown for unknown device signatures': verifica che il valore di resolved è DeviceType.UNKNOWN. Esito atteso: il valore di resolved è DeviceType.UNKNOWN.}{S}
\addTestU{Contesto device-taxonomy; caso 'returns readable label for known endpoint code': verifica che il valore di \textit{getEndpointLabel}('SFE\_Cmd\_ChangeOverMode') è 'Comando cambio modalita HVAC'. Esito atteso: il valore di \textit{getEndpointLabel}('SFE\_Cmd\_ChangeOverMode') è 'Comando cambio modalita HVAC'.}{S}
\addTestU{Contesto device-taxonomy; caso 'returns readable label for fall state endpoint code': verifica che il valore di \textit{getEndpointLabel}('SFE\_State\_Fall') è 'Stato caduta rilevata'. Esito atteso: il valore di \textit{getEndpointLabel}('SFE\_State\_Fall') è 'Stato caduta rilevata'.}{S}
\addTestU{Contesto device-taxonomy; caso 'returns readable labels for newly observed runtime thermostat commands': verifica che il valore di \textit{getEndpointLabel}('SFE\_Cmd\_HVACMode') è 'Comando modalita HVAC' e il valore di \textit{getEndpointLabel}('SFE\_Cmd\_OffBehaviour') è 'Comando comportamento in assenza'. Esito atteso: il valore di \textit{getEndpointLabel}('SFE\_Cmd\_HVACMode') è 'Comando modalita HVAC'; il valore di \textit{getEndpointLabel}('SFE\_Cmd\_OffBehaviour') è 'Comando comportamento in assenza'; il valore di \textit{getEndpointLabel}('SFE\_Cmd\_OnBehaviour') è 'Comando comportamento in presenza'.}{S}
\addTestU{Contesto device-taxonomy; caso 'returns readable label for active scene command': verifica che il valore di \textit{getEndpointLabel}('SFE\_Cmd\_DownKey\_ActiveScene') è 'Comando attivazione scenario'. Esito atteso: il valore di \textit{getEndpointLabel}('SFE\_Cmd\_DownKey\_ActiveScene') è 'Comando attivazione scenario'.}{S}
\addTestU{Contesto device-taxonomy; caso 'returns readable label for timed dynamic light command': verifica che il valore di \textit{getEndpointLabel}('SFE\_Cmd\_TimedDynamicMode') è 'Comando modalita dinamica temporizzata'. Esito atteso: il valore di \textit{getEndpointLabel}('SFE\_Cmd\_TimedDynamicMode') è 'Comando modalita dinamica temporizzata'.}{S}
\addTestU{Contesto device-taxonomy; caso 'returns raw endpoint code for unknown values': verifica che il valore di \textit{getEndpointLabel}('SFE\_Cmd\_FutureFeature') è 'Comando future feature'. Esito atteso: il valore di \textit{getEndpointLabel}('SFE\_Cmd\_FutureFeature') è 'Comando future feature'.}{S}
\addTestU{Contesto device-taxonomy; caso 'tracks unknown endpoint codes discovered at runtime': verifica che \textit{getDiscoveredUnmappedSfeCodes}() contiene 'SFE\_Cmd\_FutureFeature' e \textit{getDiscoveredUnmappedSfeCodes}() contiene 'SFE\_Cmd\_CompletelyNewOne'. Esito atteso: \textit{getDiscoveredUnmappedSfeCodes}() contiene 'SFE\_Cmd\_FutureFeature'; \textit{getDiscoveredUnmappedSfeCodes}() contiene 'SFE\_Cmd\_CompletelyNewOne'.}{S}
\addTestU{Contesto device-taxonomy; caso 'returns expected device type label': verifica che il valore di \textit{getDeviceTypeLabel}(DeviceType.PRESENCE\_SENSOR) è 'Sensore presenza'. Esito atteso: il valore di \textit{getDeviceTypeLabel}(DeviceType.PRESENCE\_SENSOR) è 'Sensore presenza'.}{S}
\addTestU{Contesto device-taxonomy; caso 'returns expected fallback label for unknown device type': verifica che il valore di \textit{getDeviceTypeLabel}(DeviceType.UNKNOWN) è 'Sconosciuto'. Esito atteso: il valore di \textit{getDeviceTypeLabel}(DeviceType.UNKNOWN) è 'Sconosciuto'.}{S}
\addTestU{Contesto \textit{ElapsedTimePipe}; caso 'ritorna il valore originale se la data non è valida': verifica che il valore di pipe.transform('not-a-date') è 'not-a-date'. Esito atteso: il valore di pipe.transform('not-a-date') è 'not-a-date'.}{S}
\addTestU{Contesto \textit{ElapsedTimePipe}; caso 'mostra tra poco per timestamp futuri': verifica che il valore di pipe.transform('2026-03-24T12:01:00.000Z') è 'tra poco'. Esito atteso: il valore di pipe.transform('2026-03-24T12:01:00.000Z') è 'tra poco'.}{S}
\addTestU{Contesto \textit{ElapsedTimePipe}; caso 'mostra secondi compatti per intervalli inferiori a 60 secondi': verifica che il valore di pipe.transform('2026-03-24T11:59:30.000Z') è '30s fa' e il valore di pipe.transform('2026-03-24T11:59:01.000Z') è '59s fa'. Esito atteso: il valore di pipe.transform('2026-03-24T11:59:30.000Z') è '30s fa'; il valore di pipe.transform('2026-03-24T11:59:01.000Z') è '59s fa'.}{S}
\addTestU{Contesto \textit{ElapsedTimePipe}; caso 'mostra minuti compatti da 60 secondi a meno di 1 ora': verifica che il valore di pipe.transform('2026-03-24T11:59:00.000Z') è '1m fa' e il valore di pipe.transform('2026-03-24T11:45:00.000Z') è '15m fa'. Esito atteso: il valore di pipe.transform('2026-03-24T11:59:00.000Z') è '1m fa'; il valore di pipe.transform('2026-03-24T11:45:00.000Z') è '15m fa'.}{S}
\addTestU{Contesto \textit{ElapsedTimePipe}; caso 'mostra ore e minuti compatti per intervalli inferiori a 1 giorno': verifica che il valore di pipe.transform('2026-03-24T11:00:00.000Z') è '1h fa' e il valore di pipe.transform('2026-03-24T10:48:00.000Z') è '1h 12m fa'. Esito atteso: il valore di pipe.transform('2026-03-24T11:00:00.000Z') è '1h fa'; il valore di pipe.transform('2026-03-24T10:48:00.000Z') è '1h 12m fa'; il valore di pipe.transform('2026-03-24T09:00:00.000Z') è '3h fa'.}{S}
\addTestU{Contesto \textit{ElapsedTimePipe}; caso 'mostra giorni compatti per intervalli pari o superiori a 1 giorno': verifica che il valore di pipe.transform('2026-03-23T12:00:00.000Z') è '1g fa' e il valore di pipe.transform('2026-03-19T12:00:00.000Z') è '5g fa'. Esito atteso: il valore di pipe.transform('2026-03-23T12:00:00.000Z') è '1g fa'; il valore di pipe.transform('2026-03-19T12:00:00.000Z') è '5g fa'.}{S}
\addTestU{Contesto \textit{AlarmApiService}; caso '\textit{getAlarmRules} chiama GET /alarm-rules e restituisce la lista': verifica che result coincide con [alarm] e result ha lunghezza 1. Esito atteso: result coincide con [alarm]; result ha lunghezza 1; il valore di request.request.method è 'GET'.}{S}
\addTestU{Contesto \textit{AlarmApiService}; caso '\textit{getAlarmRule} chiama GET /alarm-rules/:id e restituisce il dettaglio': verifica che result coincide con alarm e il valore di result.id è 'alarm-1'. Esito atteso: result coincide con alarm; il valore di result.id è 'alarm-1'; il valore di request.request.method è 'GET'.}{S}
\addTestU{Contesto \textit{AlarmApiService}; caso '\textit{getAlarmRule} codifica id con caratteri riservati': verifica che result coincide con alarm e il valore di request.request.method è 'GET'. Esito atteso: result coincide con alarm; il valore di request.request.method è 'GET'.}{S}
\addTestU{Contesto \textit{AlarmApiService}; caso '\textit{createAlarmRule} chiama POST /alarm-rules con payload corretto': verifica che result coincide con alarm e il valore di request.request.method è 'POST'. Esito atteso: result coincide con alarm; il valore di request.request.method è 'POST'; request.request.body coincide con payload.}{S}
\addTestU{Contesto \textit{AlarmApiService}; caso '\textit{updateAlarmRule} chiama PUT /alarm-rules/:id con payload corretto': verifica che result coincide con alarm e il valore di request.request.method è 'PUT'. Esito atteso: result coincide con alarm; il valore di request.request.method è 'PUT'; request.request.body coincide con payload.}{S}
\addTestU{Contesto \textit{AlarmApiService}; caso '\textit{deleteAlarmRule} chiama DELETE /alarm-rules/:id': verifica che result è null e il valore di request.request.method è 'DELETE'. Esito atteso: result è null; il valore di request.request.method è 'DELETE'.}{S}
\addTestU{Contesto \textit{AlarmApiService}; caso '\textit{getActiveAlarms} chiama GET /alarm-events/unmanaged/:userId/:limit/:offset e restituisce la lista': verifica che result coincide con [activeAlarm] e result ha lunghezza 1. Esito atteso: result coincide con [activeAlarm]; result ha lunghezza 1; il valore di request.request.method è 'GET'.}{S}
\addTestU{Contesto \textit{AlarmApiService}; caso '\textit{getResolvedAlarms} chiama GET /alarm-events/managed/:userId/:limit/:offset e restituisce la lista': verifica che result coincide con [activeAlarm] e il valore di request.request.method è 'GET'. Esito atteso: result coincide con [activeAlarm]; il valore di request.request.method è 'GET'.}{S}
\addTestU{Contesto \textit{AlarmApiService}; caso 'getAlarmEventById chiama GET /alarm-events/:id e restituisce il dettaglio': verifica che result coincide con activeAlarm e il valore di request.request.method è 'GET'. Esito atteso: result coincide con activeAlarm; il valore di request.request.method è 'GET'.}{S}
\addTestU{Contesto \textit{AlarmApiService}; caso '\textit{getActiveAlarms} valida userId, limit e offset interi': verifica che () => service.\textit{getActiveAlarms}(1.5) genera l'errore atteso e () => service.\textit{getActiveAlarms}(7, 2.5, 0) genera l'errore atteso. Esito atteso: () => service.\textit{getActiveAlarms}(1.5) genera l'errore atteso; () => service.\textit{getActiveAlarms}(7, 2.5, 0) genera l'errore atteso; () => service.\textit{getActiveAlarms}(7, 2, 1.5) genera l'errore atteso.}{S}
\addTestU{Contesto \textit{AlarmApiService}; caso '\textit{getResolvedAlarms} valida userId, limit e offset interi': verifica che () => service.\textit{getResolvedAlarms}(1.5) genera l'errore atteso e () => service.\textit{getResolvedAlarms}(7, 2.5, 0) genera l'errore atteso. Esito atteso: () => service.\textit{getResolvedAlarms}(1.5) genera l'errore atteso; () => service.\textit{getResolvedAlarms}(7, 2.5, 0) genera l'errore atteso; () => service.\textit{getResolvedAlarms}(7, 2, 1.5) genera l'errore atteso.}{S}
\addTestU{Contesto \textit{AlarmApiService}; caso '\textit{resolveAlarm} chiama PATCH /alarm-events/resolve': verifica che result è null e il valore di request.request.method è 'PATCH'. Esito atteso: result è null; il valore di request.request.method è 'PATCH'; request.request.body coincide con \{ alarmId: 'active-1', userId: 7 \}.}{S}
\addTestU{Contesto \textit{AlarmApiService}; caso '\textit{resolveAlarm} invia alarmId nel body anche con caratteri riservati': verifica che result è null e il valore di request.request.method è 'PATCH'. Esito atteso: result è null; il valore di request.request.method è 'PATCH'; request.request.body coincide con \{ alarmId: 'active/1', userId: 9 \}.}{S}
\addTestU{Contesto \textit{AlarmApiService}; caso '\textit{getAlarmRules} restituisce lista vuota senza errori di parsing': verifica che result coincide con [] e il valore di request.request.method è 'GET'. Esito atteso: result coincide con []; il valore di request.request.method è 'GET'.}{S}
\addTestU{Contesto \textit{AlarmApiService}; caso 'propaga l'errore HTTP al chiamante': verifica che il valore di receivedStatus è 404. Esito atteso: il valore di receivedStatus è 404.}{S}
\addTestU{Contesto \textit{AlarmManagementRefreshService}; caso 'espone un observable hot che emette alla richiesta di refresh': verifica che dopo risoluzione, \textit{refreshPromise} è undefined. Esito atteso: dopo risoluzione, \textit{refreshPromise} è undefined.}{S}
\addTestU{Contesto \textit{AlarmManagementRefreshService}; caso 'non bufferizza refresh precedenti per nuovi subscriber': verifica che nextEmission non viene invocato e nextEmission viene invocato 1 volte. Esito atteso: nextEmission non viene invocato; nextEmission viene invocato 1 volte.}{S}
\addTestU{Contesto \textit{AlarmStateService}; caso 'inizializza lo stato con allarmi e notifiche vuoti': verifica che il valore di alarmsHistory.at(-1) è 0 e il valore di notificationsHistory.at(-1) è 0. Esito atteso: il valore di alarmsHistory.at(-1) è 0; il valore di notificationsHistory.at(-1) è 0.}{S}
\addTestU{Contesto \textit{AlarmStateService}; caso 'aggiunge un nuovo allarme attivo quando viene chiamato \textit{onAlarmTriggered}': verifica che il valore di latestAlarms.length è 1 e il valore di latestAlarms[0].id è alarmEventA.id. Esito atteso: il valore di latestAlarms.length è 1; il valore di latestAlarms[0].id è alarmEventA.id; il valore di latestAlarms[0].alarmRuleId è alarmEventA.alarmRuleId; il valore di latestAlarms[0].alarmName è alarmEventA.alarmName.}{S}
\addTestU{Contesto \textit{AlarmStateService}; caso 'aggiorna un allarme esistente senza duplicarlo per id istanza': verifica che il valore di latestAlarms.length è 1 e il valore di latestAlarms[0].id è alarmEventA.id. Esito atteso: il valore di latestAlarms.length è 1; il valore di latestAlarms[0].id è alarmEventA.id; il valore di latestAlarms[0].alarmName è 'Nome allarme aggiornato'; il valore di latestAlarms[0].priority è AlarmPriority.WHITE.}{S}
\addTestU{Contesto \textit{AlarmStateService}; caso 'rimuove solo l'allarme target quando viene chiamato \textit{onAlarmResolved}': verifica che latestActiveAlarmIds coincide con [alarmEventB.id]. Esito atteso: latestActiveAlarmIds coincide con [alarmEventB.id].}{S}
\addTestU{Contesto \textit{AlarmStateService}; caso 'antepone le notifiche in modo che la più recente sia la prima': verifica che latestNotificationIds coincide con [notificationB.notificationId, notificationA.notificationId]. Esito atteso: latestNotificationIds coincide con [notificationB.notificationId, notificationA.notificationId].}{S}
\addTestU{Contesto \textit{AlarmStateService}; caso '\textit{removeNotification} elimina una sola notifica e aggiorna il contatore': verifica che latestNotificationIds coincide con [notificationB.notificationId] e il valore di unreadNotificationsCount è 1. Esito atteso: latestNotificationIds coincide con [notificationB.notificationId]; il valore di unreadNotificationsCount è 1.}{S}
\addTestU{Contesto \textit{AlarmStateService}; caso '\textit{clearNotifications} svuota la lista notifiche e porta il contatore a zero': verifica che latestNotificationIds coincide con [] e il valore di unreadNotificationsCount è 0. Esito atteso: latestNotificationIds coincide con []; il valore di unreadNotificationsCount è 0.}{S}
\addTestU{Contesto \textit{AlarmStateService}; caso 'una notifica dismissata non viene reintrodotta quando arriva di nuovo': verifica che latestNotificationIds coincide con []. Esito atteso: latestNotificationIds coincide con [].}{S}
\addTestU{Contesto \textit{AlarmStateService}; caso '\textit{setActiveAlarms} imposta snapshot iniziale degli allarmi attivi': verifica che latestActiveAlarmIds coincide con ['snapshot-1']. Esito atteso: latestActiveAlarmIds coincide con ['snapshot-1'].}{S}
\addTestU{Contesto \textit{AlarmStateService}; caso '\textit{setActiveAlarms} preserva eventuali allarmi locali non inclusi nello snapshot': verifica che latestActiveAlarmIds coincide con ['snapshot-1', alarmEventA.id]. Esito atteso: latestActiveAlarmIds coincide con ['snapshot-1', alarmEventA.id].}{S}
\addTestU{Contesto \textit{AlarmStateService}; caso '\textit{setActiveAlarms} non reintroduce allarmi già risolti localmente': verifica che latestActiveAlarmIds coincide con []. Esito atteso: latestActiveAlarmIds coincide con [].}{S}
\addTestU{Contesto \textit{AlarmStateService}; caso '\textit{setActiveAlarms} in replace sostituisce la pagina corrente senza mantenere elementi precedenti': verifica che latestActiveAlarmIds coincide con ['page-2-a']. Esito atteso: latestActiveAlarmIds coincide con ['page-2-a'].}{S}
\addTestU{Contesto \textit{AlarmStateService}; caso 'emette conteggi derivati corretti per allarmi e notifiche non lette': verifica che il valore di activeAlarmsCount è 2 e il valore di unreadNotificationsCount è 1. Esito atteso: il valore di activeAlarmsCount è 2; il valore di unreadNotificationsCount è 1; il valore di activeAlarmsCount è 1; il valore di unreadNotificationsCount è 2.}{S}
\addTestU{Contesto \textit{EventSubscriptionService}; caso 'crea la connessione socket in initialize usando API\_BASE\_URL normalizzato': verifica che socketIoFactoryMock viene invocato con 'http://api.example.local/ws'. Esito atteso: socketIoFactoryMock viene invocato con 'http://api.example.local/ws'.}{S}
\addTestU{Contesto \textit{EventSubscriptionService}; caso 'risolve il namespace websocket ws quando API\_BASE\_URL termina con api': verifica che socketIoFactoryMock viene invocato con 'http://localhost/ws'. Esito atteso: socketIoFactoryMock viene invocato con 'http://localhost/ws'.}{S}
\addTestU{Contesto \textit{EventSubscriptionService}; caso 'aggiorna lo stato connessione su connect, reconnect attempt e disconnect': verifica che statuses coincide con [ ConnectionStatus.DISCONNECTED, ConnectionStatus.CONNECTED, ConnectionStatus.RECONNECTING, Co.... Esito atteso: statuses coincide con [ ConnectionStatus.DISCONNECTED, ConnectionStatus.CONNECTED, ConnectionStatus.RECONNECTING, Co....}{S}
\addTestU{Contesto \textit{EventSubscriptionService}; caso 'esegue join e leave delle room con deduplicazione': verifica che fakeSocket.emit viene invocato con 'join-ward' e fakeSocket.emit viene invocato con 'leave-ward'. Esito atteso: fakeSocket.emit viene invocato con 'join-ward'; fakeSocket.emit viene invocato con 'leave-ward'; il valore di joinWardCalls.length è 2; il valore di leaveWardCalls.length è 1.}{S}
\addTestU{Contesto \textit{EventSubscriptionService}; caso 'riesegue il join delle room già sottoscritte al connect': verifica che il valore di joinCallsBeforeReconnect è 1 e il valore di joinCallsAfterReconnect.length è 2. Esito atteso: il valore di joinCallsBeforeReconnect è 1; il valore di joinCallsAfterReconnect.length è 2; joinCallsAfterReconnect.at(-1) coincide con ['join-ward', 'ward-a'].}{S}
\addTestU{Contesto \textit{EventSubscriptionService}; caso 'bootstrap delle room da plant all al login quando cache ward e vuota': verifica che httpClientSpy.get viene invocato con 'http://api.example.local/plant/all' e joinWardCalls coincide con [ ['join-ward', '21'], ['join-ward', '22'], ]. Esito atteso: httpClientSpy.get viene invocato con 'http://api.example.local/plant/all'; joinWardCalls coincide con [ ['join-ward', '21'], ['join-ward', '22'], ].}{S}
\addTestU{Contesto \textit{EventSubscriptionService}; caso 'non interrompe il realtime quando bootstrap plant all fallisce': verifica che httpClientSpy.get viene invocato 1 volte e alarmStateSpy.\textit{onAlarmTriggered} viene invocato con expect.objectContaining(\{ id: 'active-alarm-bootstrap-fallback', \}). Esito atteso: httpClientSpy.get viene invocato 1 volte; alarmStateSpy.\textit{onAlarmTriggered} viene invocato con expect.objectContaining(\{ id: 'active-alarm-bootstrap-fallback', \}); alarmStateSpy.\textit{onNotificationReceived} viene invocato con expect.objectContaining(\{ notificationId: 'alarm-triggered-active-alarm-bootstrap-fallback', \}).}{S}
\addTestU{Contesto \textit{EventSubscriptionService}; caso 'evita bootstrap ripetuto per lo stesso utente nella stessa sessione': verifica che httpClientSpy.get viene invocato 1 volte. Esito atteso: httpClientSpy.get viene invocato 1 volte.}{S}
\addTestU{Contesto \textit{EventSubscriptionService}; caso '\textit{refreshWardRoomSubscription} riallinea le room joinate senza reload pagina': verifica che httpClientSpy.get viene invocato 2 volte e fakeSocket.emit viene invocato con 'join-ward'. Esito atteso: httpClientSpy.get viene invocato 2 volte; fakeSocket.emit viene invocato con 'join-ward'; fakeSocket.emit viene invocato con 'leave-ward'.}{S}
\addTestU{Contesto \textit{EventSubscriptionService}; caso 'ignora eventi raw malformati e payload backend non validi': verifica che alarmStateSpy.\textit{onAlarmTriggered} non viene invocato e alarmStateSpy.\textit{onAlarmResolved} non viene invocato. Esito atteso: alarmStateSpy.\textit{onAlarmTriggered} non viene invocato; alarmStateSpy.\textit{onAlarmResolved} non viene invocato; alarmStateSpy.\textit{onNotificationReceived} non viene invocato.}{S}
\addTestU{Contesto \textit{EventSubscriptionService}; caso 'disconnette il socket e rimuove i listener alla distruzione': verifica che fakeSocket.\textit{removeAllListeners} viene invocato e fakeSocket.disconnect viene invocato. Esito atteso: fakeSocket.\textit{removeAllListeners} viene invocato; fakeSocket.disconnect viene invocato.}{S}
\addTestU{Contesto \textit{RealtimeAlarmEventNormalizerService}; caso '\textit{parseBackendTriggeredPayload} accetta il contratto backend nativo': verifica che parsed coincide con \{ alarmEventId: 'alarm-2', wardId: 3, alarmRuleId: 'rule-2', \}. Esito atteso: parsed coincide con \{ alarmEventId: 'alarm-2', wardId: 3, alarmRuleId: 'rule-2', \}.}{S}
\addTestU{Contesto \textit{RealtimeAlarmEventNormalizerService}; caso '\textit{parseBackendTriggeredPayload} accetta wardId numerico espresso come stringa': verifica che parsed coincide con \{ alarmEventId: 'alarm-3', wardId: 12, alarmRuleId: 'rule-3', \}. Esito atteso: parsed coincide con \{ alarmEventId: 'alarm-3', wardId: 12, alarmRuleId: 'rule-3', \}.}{S}
\addTestU{Contesto \textit{RealtimeAlarmEventNormalizerService}; caso '\textit{parseBackendResolvedPayload} estrae alarmEventId e wardId opzionale': verifica che withWardId coincide con \{ alarmEventId: 'alarm-4', wardId: 8 \} e withoutWardId coincide con \{ alarmEventId: 'alarm-5', wardId: undefined \}. Esito atteso: withWardId coincide con \{ alarmEventId: 'alarm-4', wardId: 8 \}; withoutWardId coincide con \{ alarmEventId: 'alarm-5', wardId: undefined \}.}{S}
\addTestU{Contesto \textit{WardRealtimeCacheService}; caso 'mergeWardIds deduplica e normalizza valori stringa/numero': verifica che service.getWardIds('user-1') coincide con ['1', '2', '3']. Esito atteso: service.getWardIds('user-1') coincide con ['1', '2', '3'].}{S}
\addTestU{Contesto \textit{WardRealtimeCacheService}; caso 'mantiene cache separata per utente': verifica che service.getWardIds('user-a') coincide con ['10'] e service.getWardIds('user-b') coincide con ['11', '12']. Esito atteso: service.getWardIds('user-a') coincide con ['10']; service.getWardIds('user-b') coincide con ['11', '12'].}{S}
\addTestU{Contesto \textit{WardRealtimeCacheService}; caso 'clearWardIds rimuove solo l'utente target': verifica che service.getWardIds('user-a') coincide con [] e service.getWardIds('user-b') coincide con ['11']. Esito atteso: service.getWardIds('user-a') coincide con []; service.getWardIds('user-b') coincide con ['11'].}{S}
\addTestU{Contesto \textit{WardRealtimeCacheService}; caso 'recupera cache vuota in caso di JSON non valido': verifica che service.getWardIds('user-z') coincide con []. Esito atteso: service.getWardIds('user-z') coincide con [].}{S}
\addTestU{Contesto \textit{WardSocketRoomCoordinatorService}; caso '\textit{activateUser} carica wardIds da cache e propone join': verifica che actions.roomsToLeave coincide con [] e actions.roomsToJoin coincide con ['10', '11']. Esito atteso: actions.roomsToLeave coincide con []; actions.roomsToJoin coincide con ['10', '11']; service.\textit{getJoinedRooms}() coincide con ['10', '11'].}{S}
\addTestU{Contesto \textit{WardSocketRoomCoordinatorService}; caso '\textit{requestJoinRoom} deduplica e persiste su cache': verifica che il valore di service.\textit{requestJoinRoom}(' 20 ') è '20' e service.\textit{requestJoinRoom}('20') è null. Esito atteso: il valore di service.\textit{requestJoinRoom}(' 20 ') è '20'; service.\textit{requestJoinRoom}('20') è null; cacheSpy.setWardIds viene invocato.}{S}
\addTestU{Contesto \textit{WardSocketRoomCoordinatorService}; caso '\textit{deactivateUser} rimuove room runtime e pulisce cache utente corrente': verifica che roomsToLeave coincide con ['30'] e service.\textit{getJoinedRooms}() coincide con []. Esito atteso: roomsToLeave coincide con ['30']; service.\textit{getJoinedRooms}() coincide con []; cacheSpy.clearWardIds viene invocato con 'user-1'.}{S}
\addTestU{Contesto \textit{WardSocketRoomCoordinatorService}; caso 'switch utente restituisce room da lasciare e room cache del nuovo utente': verifica che actions.roomsToLeave coincide con ['10', '30'] e actions.roomsToJoin coincide con ['40', '50']. Esito atteso: actions.roomsToLeave coincide con ['10', '30']; actions.roomsToJoin coincide con ['40', '50']; service.\textit{getJoinedRooms}() coincide con ['40', '50'].}{S}
\addTestU{Contesto \textit{NotificationApiService}; caso '\textit{getNotificationsHistory} chiama managed/unmanaged e mappa verso \textit{\textit{NotificationEvent}}': verifica che result coincide con expected e il valore di managedRequest.request.method è 'GET'. Esito atteso: result coincide con expected; il valore di managedRequest.request.method è 'GET'; il valore di unmanagedRequest.request.method è 'GET'.}{S}
\addTestU{Contesto \textit{NotificationApiService}; caso '\textit{getNotificationsHistory} emette una volta e completa': verifica che il valore di emissions è 1 e il valore di completed è true. Esito atteso: il valore di emissions è 1; il valore di completed è true.}{S}
\addTestU{Contesto \textit{NotificationApiService}; caso '\textit{getNotificationsHistory} e cold: due sottoscrizioni producono due coppie di request HTTP': verifica che requests ha lunghezza 4 e il valore di request.request.method è 'GET'. Esito atteso: requests ha lunghezza 4; il valore di request.request.method è 'GET'.}{S}
\addTestU{Contesto \textit{NotificationApiService}; caso 'se non c'è sessione attiva restituisce [] senza chiamate HTTP': verifica che result coincide con [] e requests ha lunghezza 0. Esito atteso: result coincide con []; requests ha lunghezza 0.}{S}
\addTestU{Contesto \textit{AlarmRuleFormMapper}; caso 'mappa operatore >= come GREATER\_THAN\_OR\_EQUAL': verifica che il valore di mapper.\textit{toFormValue}(rule).thresholdOperator è ThresholdOperator.GREATER\_THAN\_OR\_EQUAL. Esito atteso: il valore di mapper.\textit{toFormValue}(rule).thresholdOperator è ThresholdOperator.GREATER\_THAN\_OR\_EQUAL.}{S}
\addTestU{Contesto \textit{AlarmRuleFormMapper}; caso 'mappa operatore <= come LESS\_THAN\_OR\_EQUAL': verifica che il valore di mapper.\textit{toFormValue}(rule).thresholdOperator è ThresholdOperator.LESS\_THAN\_OR\_EQUAL. Esito atteso: il valore di mapper.\textit{toFormValue}(rule).thresholdOperator è ThresholdOperator.LESS\_THAN\_OR\_EQUAL.}{S}
\addTestU{Contesto \textit{AlarmRuleFormMapper}; caso 'fallisce su operatori sconosciuti': verifica che () => mapper.\textit{toFormValue}(rule) genera l'errore atteso. Esito atteso: () => mapper.\textit{toFormValue}(rule) genera l'errore atteso.}{S}
\addTestU{Contesto \textit{AlarmRuleRequestMapper}; caso '\textit{toUpdateRequest} serializza operatore >= senza degradarlo': verifica che il valore di result.thresholdOperator è '>='. Esito atteso: il valore di result.thresholdOperator è '>='.}{S}
\addTestU{Contesto \textit{AlarmRuleRequestMapper}; caso '\textit{toCreateRequest} serializza operatore <= senza degradarlo': verifica che il valore di result.thresholdOperator è '<='. Esito atteso: il valore di result.thresholdOperator è '<='.}{S}
\addTestU{Contesto \textit{AlarmRuleRequestMapper}; caso '\textit{toCreateRequest} normalizza soglia booleana in on/off': verifica che l'asserzione mapper.\textit{toCreateRequest}(\textit{onFormValue}) \textit{toMatchObject} \{ plantId: 'plant-1', datapointId: 'dp-4', thresholdOperator: '=', thresholdValue: 'on', \} risulta soddisfatta e l'asserzione mapper.\textit{toCreateRequest}(offFormValue) \textit{toMatchObject} \{ plantId: 'plant-1', datapointId: 'dp-4', thresholdOperator: '=', thresholdValue: 'off', \} risulta soddisfatta. Esito atteso: l'asserzione mapper.\textit{toCreateRequest}(\textit{onFormValue}) \textit{toMatchObject} \{ plantId: 'plant-1', datapointId: 'dp-4', thresholdOperator: '=', thresholdValue: 'on', \} risulta soddisfatta; l'asserzione mapper.\textit{toCreateRequest}(offFormValue) \textit{toMatchObject} \{ plantId: 'plant-1', datapointId: 'dp-4', thresholdOperator: '=', thresholdValue: 'off', \} risulta soddisfatta.}{S}
\addTestU{Contesto \textit{AlarmRuleRequestMapper}; caso '\textit{toCreateRequest} fallisce se thresholdValue e vuoto dopo trim': verifica che () => mapper.\textit{toCreateRequest}(invalidFormValue) genera l'errore atteso. Esito atteso: () => mapper.\textit{toCreateRequest}(invalidFormValue) genera l'errore atteso.}{S}
\addTestU{Contesto \textit{AlarmRuleRequestMapper}; caso '\textit{toCreateRequest} fallisce se thresholdOperator è null': verifica che () => mapper.\textit{toCreateRequest}(invalidFormValue) genera l'errore atteso. Esito atteso: () => mapper.\textit{toCreateRequest}(invalidFormValue) genera l'errore atteso.}{S}
\addTestU{Contesto \textit{AlarmRuleRequestMapper}; caso '\textit{toToggleRequest} fallisce se name e vuoto': verifica che () => mapper.\textit{toToggleRequest}(invalidRule, true) genera l'errore atteso. Esito atteso: () => mapper.\textit{toToggleRequest}(invalidRule, true) genera l'errore atteso.}{S}
\addTestU{Contesto \textit{AlarmRuleRequestMapper}; caso '\textit{toToggleRequest} normalizza operatori con padding dal backend': verifica che il valore di result.thresholdOperator è '>'. Esito atteso: il valore di result.thresholdOperator è '>'.}{S}
\addTestU{Contesto \textit{AlarmConfigFormFieldOrchestratorHelper}; caso 'applyModeState in edit mode disabilita campi non modificabili': verifica che il valore di formLike.controls.name.disabled è true e il valore di formLike.controls.plantId.disabled è true. Esito atteso: il valore di formLike.controls.name.disabled è true; il valore di formLike.controls.plantId.disabled è true; il valore di formLike.controls.deviceId.disabled è true; il valore di formLike.controls.datapointId.disabled è true.}{S}
\addTestU{Contesto \textit{AlarmConfigFormFieldOrchestratorHelper}; caso '\textit{resolveDeviceSelection} ritorna solo datapoint leggibili e auto-selezione se singolo': verifica che il valore di selection.normalizedDeviceId è 'device-1' e selection.readableDatapoints coincide con [readableDatapoint]. Esito atteso: il valore di selection.normalizedDeviceId è 'device-1'; selection.readableDatapoints coincide con [readableDatapoint]; il valore di selection.autoSelectedDatapointId è 'dp-readable'.}{S}
\addTestU{Contesto \textit{AlarmConfigFormFieldOrchestratorHelper}; caso '\textit{resolveDatapointSelection} normalizza id e restituisce il datapoint atteso': verifica che il valore di selection.normalizedDatapointId è 'dp-1' e selection.selectedDatapoint coincide con datapoint. Esito atteso: il valore di selection.normalizedDatapointId è 'dp-1'; selection.selectedDatapoint coincide con datapoint.}{S}
\addTestU{Contesto \textit{AlarmConfigFormValidationHelper}; caso 'in create mode con enum forza operatore uguale e valida threshold su enum': verifica che operatorControl.value è null e il valore di valueControl.errors?.['invalidEnumThreshold'] è true. Esito atteso: operatorControl.value è null; il valore di valueControl.errors?.['invalidEnumThreshold'] è true.}{S}
\addTestU{Contesto \textit{AlarmConfigFormValidationHelper}; caso 'in create mode con datapoint numerico applica validazione numerica': verifica che il valore di operatorControl.value è ThresholdOperator.GREATER\_THAN e valueControl.errors?.['pattern'] è valorizzato. Esito atteso: il valore di operatorControl.value è ThresholdOperator.GREATER\_THAN; valueControl.errors?.['pattern'] è valorizzato.}{S}
\addTestU{Contesto \textit{AlarmConfigFormValidationHelper}; caso 'in edit mode con datapoint enum applica lo stesso vincolo operatore della create': verifica che operatorControl.value è null e il valore di valueControl.errors?.['invalidEnumThreshold'] è true. Esito atteso: operatorControl.value è null; il valore di valueControl.errors?.['invalidEnumThreshold'] è true.}{S}
\addTestU{Contesto \textit{AlarmConfigFormValidationHelper}; caso 'in edit mode con datapoint numerico applica validazione numerica': verifica che valueControl.errors?.['pattern'] è valorizzato. Esito atteso: valueControl.errors?.['pattern'] è valorizzato.}{S}
\addTestU{Contesto \textit{AlarmConfigFormValidationHelper}; caso 'in edit mode senza metadati datapoint mantiene come ammesso solo l operatore corrente': verifica che il valore di operatorControl.errors?.['unsupportedOperator'] è true. Esito atteso: il valore di operatorControl.errors?.['unsupportedOperator'] è true.}{S}
\addTestU{Contesto \textit{AlarmConfigFormStateService}; caso '\textit{ensurePlantsLoaded} in create mode carica e deduplica gli impianti': verifica che service.plants().map((plant) => plant.id) coincide con ['plant-1', 'plant-2', 'plant-3'] e service.plantsLoadError() è null. Esito atteso: service.plants().map((plant) => plant.id) coincide con ['plant-1', 'plant-2', 'plant-3']; service.plantsLoadError() è null.}{S}
\addTestU{Contesto \textit{AlarmConfigFormStateService}; caso '\textit{ensurePlantsLoaded} in edit mode non esegue chiamate': verifica che wardApiStub.\textit{getAvailablePlants} non viene invocato e wardApiStub.\textit{getWards} non viene invocato. Esito atteso: wardApiStub.\textit{getAvailablePlants} non viene invocato; wardApiStub.\textit{getWards} non viene invocato.}{S}
\addTestU{Contesto \textit{AlarmConfigFormStateService}; caso '\textit{loadDeviceOptionsByPlant} ritorna opzioni dispositivi e pulisce loading state': verifica che apartmentApiStub.loadApartmentViewForPlantId viene invocato con 'plant-1' e result coincide con options. Esito atteso: apartmentApiStub.loadApartmentViewForPlantId viene invocato con 'plant-1'; result coincide con options; il valore di service.isDevicesLoading() è false; service.devicesLoadError() è null.}{S}
\addTestU{Contesto \textit{AlarmConfigFormStateService}; caso '\textit{loadDeviceOptionsByPlant} in errore restituisce lista vuota e imposta messaggio': verifica che result coincide con [] e il valore di service.isDevicesLoading() è false. Esito atteso: result coincide con []; il valore di service.isDevicesLoading() è false; il valore di service.devicesLoadError() è 'Errore durante il caricamento dei dispositivi.'.}{S}
\addTestU{Contesto \textit{AlarmConfigFormStateService}; caso '\textit{resolveDatapointForEdit} usa il catalogo piante e restituisce il datapoint trovato': verifica che apartmentApiStub.\textit{getAllPlants} viene invocato 1 volte e datapointExtractionStub.findDatapointByDeviceAndDatapointId viene invocato con [\{ id: 'plant-1', name: 'Apt', rooms: [] \}]. Esito atteso: apartmentApiStub.\textit{getAllPlants} viene invocato 1 volte; datapointExtractionStub.findDatapointByDeviceAndDatapointId viene invocato con [\{ id: 'plant-1', name: 'Apt', rooms: [] \}]; result coincide con expectedDatapoint.}{S}
\addTestU{Contesto \textit{AlarmConfigFormStateService}; caso '\textit{resolveDatapointForEdit} ritorna null se il catalogo piante fallisce': verifica che result è null. Esito atteso: result è null.}{S}
\addTestU{Contesto \textit{AlarmConfigStateService}; caso 'inizializza stato con alarms vuoti e nessun errore': verifica che await firstValueFrom(service.alarms\$) coincide con [] e await firstValueFrom(service.error\$) è null. Esito atteso: await firstValueFrom(service.alarms\$) coincide con []; await firstValueFrom(service.error\$) è null.}{S}
\addTestU{Contesto \textit{AlarmConfigStateService}; caso '\textit{loadAlarmRules} aggiorna stato locale quando API risponde con successo': verifica che apiStub.\textit{getAlarmRules} viene invocato 1 volte e alarms coincide con [alarmA, alarmB]. Esito atteso: apiStub.\textit{getAlarmRules} viene invocato 1 volte; alarms coincide con [alarmA, alarmB]; alarms ha lunghezza 2; await firstValueFrom(service.error\$) è null.}{S}
\addTestU{Contesto \textit{AlarmConfigStateService}; caso '\textit{loadAlarmRules} espone messaggio errore e pulisce error precedente al retry': verifica che il valore di await firstValueFrom(service.error\$) è 'Errore durante il caricamento degli allarmi.' e await firstValueFrom(service.error\$) è null. Esito atteso: il valore di await firstValueFrom(service.error\$) è 'Errore durante il caricamento degli allarmi.'; await firstValueFrom(service.error\$) è null; await firstValueFrom(service.alarms\$) coincide con [alarmA].}{S}
\addTestU{Contesto \textit{AlarmConfigStateService}; caso 'getAlarmRuleById inoltra il valore di API in caso di successo': verifica che apiStub.\textit{getAlarmRule} viene invocato con 'alarm-1' e result coincide con alarmA. Esito atteso: apiStub.\textit{getAlarmRule} viene invocato con 'alarm-1'; result coincide con alarmA.}{S}
\addTestU{Contesto \textit{AlarmConfigStateService}; caso 'getAlarmRuleById in errore imposta messaggio e completa senza emissioni': verifica che il valore di emitted è false e il valore di await firstValueFrom(service.error\$) è 'Errore durante il recupero dell\textbackslash\{\}'allarme.'. Esito atteso: il valore di emitted è false; il valore di await firstValueFrom(service.error\$) è 'Errore durante il recupero dell\textbackslash\{\}'allarme.'.}{S}
\addTestU{Contesto \textit{AlarmConfigStateService}; caso '\textit{createAlarmRule} con form invalido non chiama API e imposta errore di validazione': verifica che il valore di emitted è false e apiStub.\textit{createAlarmRule} non viene invocato. Esito atteso: il valore di emitted è false; apiStub.\textit{createAlarmRule} non viene invocato; il valore di await firstValueFrom(service.error\$) è 'Dati del form non validi per la creazione dell\textbackslash\{\}'allarme.'.}{S}
\addTestU{Contesto \textit{AlarmConfigStateService}; caso '\textit{createAlarmRule} in errore API imposta messaggio e non altera lo stato locale': verifica che il valore di emitted è false e il valore di await firstValueFrom(service.error\$) è 'Errore durante la creazione dell\textbackslash\{\}'allarme.'. Esito atteso: il valore di emitted è false; il valore di await firstValueFrom(service.error\$) è 'Errore durante la creazione dell\textbackslash\{\}'allarme.'; await firstValueFrom(service.alarms\$) coincide con [alarmA]; \textit{refreshServiceStub}.\textit{requestRefresh} non viene invocato.}{S}
\addTestU{Contesto \textit{AlarmConfigStateService}; caso '\textit{createAlarmRule} pulisce l'errore precedente quando una creazione successiva va a buon fine': verifica che il valore di await firstValueFrom(service.error\$) è 'Dati del form non validi per la creazione dell\textbackslash\{\}'allarme.' e await firstValueFrom(service.error\$) è null. Esito atteso: il valore di await firstValueFrom(service.error\$) è 'Dati del form non validi per la creazione dell\textbackslash\{\}'allarme.'; await firstValueFrom(service.error\$) è null; await firstValueFrom(service.alarms\$) coincide con [alarmB].}{S}
\addTestU{Contesto \textit{AlarmConfigStateService}; caso '\textit{updateAlarmRule} mantiene il nome originale nel payload update e sostituisce l'allarme nello stato': verifica che apiStub.\textit{updateAlarmRule} viene invocato con 'alarm-1' e await firstValueFrom(service.alarms\$) coincide con [updatedAlarm, alarmB]. Esito atteso: apiStub.\textit{updateAlarmRule} viene invocato con 'alarm-1'; await firstValueFrom(service.alarms\$) coincide con [updatedAlarm, alarmB]; \textit{refreshServiceStub}.\textit{requestRefresh} viene invocato 1 volte.}{S}
\addTestU{Contesto \textit{AlarmConfigStateService}; caso '\textit{updateAlarmRule} con form invalido non chiama API e imposta errore': verifica che il valore di emitted è false e apiStub.\textit{updateAlarmRule} non viene invocato. Esito atteso: il valore di emitted è false; apiStub.\textit{updateAlarmRule} non viene invocato; il valore di await firstValueFrom(service.error\$) è 'Dati del form non validi per l\textbackslash\{\}'aggiornamento dell\textbackslash\{\}'allarme.'.}{S}
\addTestU{Contesto \textit{AlarmConfigStateService}; caso '\textit{updateAlarmRule} in errore API imposta messaggio e mantiene lo stato invariato': verifica che il valore di emitted è false e il valore di await firstValueFrom(service.error\$) è 'Errore durante l\textbackslash\{\}'aggiornamento dell\textbackslash\{\}'allarme.'. Esito atteso: il valore di emitted è false; il valore di await firstValueFrom(service.error\$) è 'Errore durante l\textbackslash\{\}'aggiornamento dell\textbackslash\{\}'allarme.'; await firstValueFrom(service.alarms\$) coincide con [alarmA, alarmB]; \textit{refreshServiceStub}.\textit{requestRefresh} non viene invocato.}{S}
\addTestU{Contesto \textit{AlarmConfigStateService}; caso '\textit{updateAlarmRule} senza allarme locale usa il nome del form nel payload': verifica che apiStub.\textit{updateAlarmRule} viene invocato con 'missing-id'. Esito atteso: apiStub.\textit{updateAlarmRule} viene invocato con 'missing-id'.}{S}
\addTestU{Contesto \textit{AlarmConfigStateService}; caso '\textit{updateAlarmRule} sostituisce la riga locale anche se il backend ritorna un nuovo id': verifica che alarms coincide con [updatedWithNewId, alarmB] e il valore di alarms.some((alarm) => alarm.id === 'alarm-1') è false. Esito atteso: alarms coincide con [updatedWithNewId, alarmB]; il valore di alarms.some((alarm) => alarm.id === 'alarm-1') è false; il valore di alarms.some((alarm) => alarm.id === 'alarm-1-v2') è true; \textit{refreshServiceStub}.\textit{requestRefresh} viene invocato 1 volte.}{S}
\addTestU{Contesto \textit{AlarmConfigStateService}; caso '\textit{toggleEnabled} aggiorna enabled costruendo payload dai dati locali': verifica che apiStub.\textit{updateAlarmRule} viene invocato con 'alarm-1' e await firstValueFrom(service.alarms\$) coincide con [toggledAlarm]. Esito atteso: apiStub.\textit{updateAlarmRule} viene invocato con 'alarm-1'; await firstValueFrom(service.alarms\$) coincide con [toggledAlarm]; \textit{refreshServiceStub}.\textit{requestRefresh} viene invocato 1 volte.}{S}
\addTestU{Contesto \textit{AlarmConfigStateService}; caso '\textit{toggleEnabled} su allarme assente non chiama API e imposta errore': verifica che il valore di emitted è false e apiStub.\textit{updateAlarmRule} non viene invocato. Esito atteso: il valore di emitted è false; apiStub.\textit{updateAlarmRule} non viene invocato; il valore di await firstValueFrom(service.error\$) è 'Allarme non trovato nello stato locale.'; \textit{refreshServiceStub}.\textit{requestRefresh} non viene invocato.}{S}
\addTestU{Contesto \textit{AlarmConfigStateService}; caso '\textit{toggleEnabled} in errore API imposta messaggio e mantiene lo stato invariato': verifica che il valore di emitted è false e il valore di await firstValueFrom(service.error\$) è 'Errore durante la modifica dello stato dell\textbackslash\{\}'allarme.'. Esito atteso: il valore di emitted è false; il valore di await firstValueFrom(service.error\$) è 'Errore durante la modifica dello stato dell\textbackslash\{\}'allarme.'; await firstValueFrom(service.alarms\$) coincide con [alarmA]; \textit{refreshServiceStub}.\textit{requestRefresh} non viene invocato.}{S}
\addTestU{Contesto \textit{AlarmConfigStateService}; caso '\textit{toggleEnabled} aggiorna subito lo stato locale anche se il backend ritorna un nuovo id': verifica che alarms coincide con [toggledWithNewId, alarmB] e il valore di alarms.some((alarm) => alarm.id === 'alarm-1') è false. Esito atteso: alarms coincide con [toggledWithNewId, alarmB]; il valore di alarms.some((alarm) => alarm.id === 'alarm-1') è false; il valore di alarms.some((alarm) => alarm.id === 'alarm-1-v2') è true.}{S}
\addTestU{Contesto \textit{AlarmConfigStateService}; caso '\textit{deleteAlarmRule} rimuove l'allarme locale quando API conferma': verifica che il valore di emitted è true e apiStub.\textit{deleteAlarmRule} viene invocato con 'alarm-1'. Esito atteso: il valore di emitted è true; apiStub.\textit{deleteAlarmRule} viene invocato con 'alarm-1'; await firstValueFrom(service.alarms\$) coincide con [alarmB]; \textit{refreshServiceStub}.\textit{requestRefresh} viene invocato 1 volte.}{S}
\addTestU{Contesto \textit{AlarmConfigStateService}; caso '\textit{deleteAlarmRule} in errore API imposta messaggio e mantiene lo stato invariato': verifica che il valore di emitted è false e il valore di await firstValueFrom(service.error\$) è 'Errore durante l\textbackslash\{\}'eliminazione dell\textbackslash\{\}'allarme.'. Esito atteso: il valore di emitted è false; il valore di await firstValueFrom(service.error\$) è 'Errore durante l\textbackslash\{\}'eliminazione dell\textbackslash\{\}'allarme.'; await firstValueFrom(service.alarms\$) coincide con [alarmA, alarmB]; \textit{refreshServiceStub}.\textit{requestRefresh} non viene invocato.}{S}
\addTestU{Contesto \textit{AlarmConfigStateService}; caso '\textit{deleteAlarmRule} pulisce errore precedente quando un delete successivo riesce': verifica che il valore di await firstValueFrom(service.error\$) è 'Errore durante l\textbackslash\{\}'eliminazione dell\textbackslash\{\}'allarme.' e await firstValueFrom(service.error\$) è null. Esito atteso: il valore di await firstValueFrom(service.error\$) è 'Errore durante l\textbackslash\{\}'eliminazione dell\textbackslash\{\}'allarme.'; await firstValueFrom(service.error\$) è null; await firstValueFrom(service.alarms\$) coincide con [alarmB]; \textit{refreshServiceStub}.\textit{requestRefresh} viene invocato 1 volte.}{S}
\addTestU{Contesto \textit{AlarmConfigTablePresenterService}; caso 'normalizza gli spazi intorno ai separatori della posizione': verifica che il valore di row.position è 'Appartamento 2 - Cucina - Sensore fumo'. Esito atteso: il valore di row.position è 'Appartamento 2 - Cucina - Sensore fumo'.}{S}
\addTestU{Contesto \textit{AlarmConfigTablePresenterService}; caso 'applica fallback su posizione quando e vuota': verifica che il valore di row.position è '-'. Esito atteso: il valore di row.position è '-'.}{S}
\addTestU{Contesto \textit{AlarmConfigTablePresenterService}; caso 'mantiene fallback orario quando il valore non è ISO/HH:mm:ss': verifica che il valore di row.armingTime è 'xx:yy' e il valore di row.dearmingTime è 'foo'. Esito atteso: il valore di row.armingTime è 'xx:yy'; il valore di row.dearmingTime è 'foo'.}{S}
\addTestU{Contesto \textit{DeviceDatapointExtractionService}; caso '\textit{extractDeviceOptions} crea label stanza-dispositivo e normalizza enum duplicati': verifica che options ha lunghezza 1 e il valore di options[0].id è 'dev-1'. Esito atteso: options ha lunghezza 1; il valore di options[0].id è 'dev-1'; il valore di options[0].label è 'Soggiorno - Sensore UWB'; options[0].datapoints[0].enum coincide con ['Off', 'On'].}{S}
\addTestU{Contesto \textit{DeviceDatapointExtractionService}; caso '\textit{findReadableDatapoints} ritorna solo datapoint leggibili': verifica che datapoints.map((datapoint) => datapoint.id) coincide con ['dp-1', 'dp-2'] e il valore di datapoints.every((datapoint) => datapoint.readable) è true. Esito atteso: datapoints.map((datapoint) => datapoint.id) coincide con ['dp-1', 'dp-2']; il valore di datapoints.every((datapoint) => datapoint.readable) è true.}{S}
\addTestU{Contesto \textit{DeviceDatapointExtractionService}; caso '\textit{findReadableDatapoints} ritorna lista vuota se il device non esiste': verifica che datapoints coincide con []. Esito atteso: datapoints coincide con [].}{S}
\addTestU{Contesto \textit{DeviceDatapointExtractionService}; caso '\textit{getAllowedOperators} limita a uguale quando il datapoint ha enum': verifica che service.\textit{getAllowedOperators}(enumDatapoint) coincide con [ThresholdOperator.EQUAL\_TO] e service.\textit{getAllowedOperators}(numericDatapoint) coincide con [ ThresholdOperator.GREATER\_THAN, ThresholdOperator.GREATER\_THAN\_OR\_EQUAL, ThresholdOperator.L.... Esito atteso: service.\textit{getAllowedOperators}(enumDatapoint) coincide con [ThresholdOperator.EQUAL\_TO]; service.\textit{getAllowedOperators}(numericDatapoint) coincide con [ ThresholdOperator.GREATER\_THAN, ThresholdOperator.GREATER\_THAN\_OR\_EQUAL, ThresholdOperator.L....}{S}
\addTestU{Contesto \textit{DeviceDatapointExtractionService}; caso 'isThresholdValueValid valida enum case-insensitive e soglie numeriche': verifica che il valore di service.isThresholdValueValid(enumDatapoint, 'on') è true e il valore di service.isThresholdValueValid(enumDatapoint, 'Unknown') è false. Esito atteso: il valore di service.isThresholdValueValid(enumDatapoint, 'on') è true; il valore di service.isThresholdValueValid(enumDatapoint, 'Unknown') è false; il valore di service.isThresholdValueValid(numericDatapoint, '10.5') è true; il valore di service.isThresholdValueValid(numericDatapoint, 'abc') è false.}{S}
\addTestU{Contesto \textit{DeviceDatapointExtractionService}; caso 'findDatapointByDeviceAndDatapointId trova e normalizza il datapoint dal catalogo plant': verifica che datapoint non è null e il valore di datapoint?.id è 'dp-lookup'. Esito atteso: datapoint non è null; il valore di datapoint?.id è 'dp-lookup'; datapoint?.enum coincide con ['Off', 'On'].}{S}
\addTestU{Contesto \textit{DeviceDatapointExtractionService}; caso 'findDatapointByDeviceAndDatapointId ritorna null se non trova corrispondenze': verifica che datapoint è null. Esito atteso: datapoint è null.}{S}
\addTestU{Contesto \textit{AlarmHistoryPageComponent}; caso 'crea il componente e inizializza il service': verifica che component è valorizzato e alarmHistoryServiceStub.initialize viene invocato 1 volte. Esito atteso: component è valorizzato; alarmHistoryServiceStub.initialize viene invocato 1 volte.}{S}
\addTestU{Contesto \textit{AlarmHistoryPageComponent}; caso 'mappa e ordina le righe per closedAt desc': verifica che rows ha lunghezza 2 e il valore di rows[0].id è 'a-2'. Esito atteso: rows ha lunghezza 2; il valore di rows[0].id è 'a-2'; il valore di rows[1].id è 'a-1'; il valore di rows[0].status è 'Gestito'.}{S}
\addTestU{Contesto \textit{AlarmHistoryPageComponent}; caso '\textit{toRow} applica fallback e stato per allarme non risolto': verifica che il valore di row.status è 'Da gestire' e il valore di row.closedAt è 'sconosciuto'. Esito atteso: il valore di row.status è 'Da gestire'; il valore di row.closedAt è 'sconosciuto'; il valore di row.handlerUsername è 'sconosciuto'.}{S}
\addTestU{Contesto \textit{AlarmHistoryPageComponent}; caso 'delega next e previous page al service': verifica che alarmHistoryServiceStub.nextPage viene invocato 1 volte e alarmHistoryServiceStub.previousPage viene invocato 1 volte. Esito atteso: alarmHistoryServiceStub.nextPage viene invocato 1 volte; alarmHistoryServiceStub.previousPage viene invocato 1 volte.}{S}
\addTestU{Contesto \textit{AlarmManagementTablePresenterService}; caso 'mappa un allarme aperto con azione gestibile': verifica che il valore di row.id è 'active-1' e il valore di row.status è 'Da gestire'. Esito atteso: il valore di row.id è 'active-1'; il valore di row.status è 'Da gestire'; il valore di row.isManaged è false; il valore di row.isResolving è false.}{S}
\addTestU{Contesto \textit{AlarmManagementTablePresenterService}; caso 'mappa un allarme già gestito come non azionabile': verifica che il valore di row.status è 'Non da gestire' e il valore di row.isManaged è true. Esito atteso: il valore di row.status è 'Non da gestire'; il valore di row.isManaged è true; il valore di row.isActionDisabled è true; il valore di row.actionLabel è 'GESTITO'.}{S}
\addTestU{Contesto \textit{AlarmManagementTablePresenterService}; caso 'quando l'allarme e in risoluzione mostra stato e label specifici': verifica che il valore di row.isResolving è true e il valore di row.isActionDisabled è true. Esito atteso: il valore di row.isResolving è true; il valore di row.isActionDisabled è true; il valore di row.actionLabel è 'GESTIONE...'; il valore di row.status è 'Da gestire'.}{S}
\addTestU{Contesto \textit{AlarmManagementTablePresenterService}; caso 'applica fallback su posizione, manager e orari invalidi': verifica che il valore di row.location è '-' e il valore di row.device è '-'. Esito atteso: il valore di row.location è '-'; il valore di row.device è '-'; il valore di row.openedAt è 'abcde12345'.}{S}
\addTestU{Contesto \textit{AlarmManagementTablePresenterService}; caso 'normalizza nome allarme vuoto per testo riga e aria label': verifica che il valore di row.name è 'senza nome' e il valore di row.actionAriaLabel è 'Gestisci allarme senza nome'. Esito atteso: il valore di row.name è 'senza nome'; il valore di row.actionAriaLabel è 'Gestisci allarme senza nome'.}{S}
\addTestU{Contesto \textit{AlarmManagementTablePresenterService}; caso 'gestisce posizione null senza eccezioni runtime': verifica che il valore di row.location è '-'. Esito atteso: il valore di row.location è '-'.}{S}
\addTestU{Contesto \textit{AlarmManagementService}; caso 'non carica lo snapshot iniziale automaticamente alla creazione': verifica che alarmApiStub.\textit{getActiveAlarms} non viene invocato e alarmStateStub.\textit{setActiveAlarms} non viene invocato. Esito atteso: alarmApiStub.\textit{getActiveAlarms} non viene invocato; alarmStateStub.\textit{setActiveAlarms} non viene invocato.}{S}
\addTestU{Contesto \textit{AlarmManagementService}; caso 'initialize per amministratore mostra tutti gli allarmi senza filtri frontend': verifica che alarmApiStub.\textit{getActiveAlarms} viene invocato 2 volte e alarmApiStub.\textit{getActiveAlarms} viene invocato con 99. Esito atteso: alarmApiStub.\textit{getActiveAlarms} viene invocato 2 volte; alarmApiStub.\textit{getActiveAlarms} viene invocato con 99; l'asserzione alarmApiStub.\textit{getActiveAlarms} \textit{toHaveBeenNthCalledWith} 2 risulta soddisfatta; alarmStateStub.\textit{setActiveAlarms} viene invocato con initialSnapshot.}{S}
\addTestU{Contesto \textit{AlarmManagementService}; caso 'initialize per OSS usa endpoint unmanaged con userId numerico': verifica che alarmApiStub.\textit{getActiveAlarms} viene invocato 2 volte e alarmApiStub.\textit{getActiveAlarms} viene invocato con 7. Esito atteso: alarmApiStub.\textit{getActiveAlarms} viene invocato 2 volte; alarmApiStub.\textit{getActiveAlarms} viene invocato con 7; l'asserzione alarmApiStub.\textit{getActiveAlarms} \textit{toHaveBeenNthCalledWith} 2 risulta soddisfatta; alarmStateStub.\textit{setActiveAlarms} viene invocato con initialSnapshot.}{S}
\addTestU{Contesto \textit{AlarmManagementService}; caso 'se initialize fallisce non propaga eccezioni e valorizza l'errore nel vm': verifica che () => service.initialize() non genera l'errore atteso e alarmStateStub.\textit{setActiveAlarms} non viene invocato. Esito atteso: () => service.initialize() non genera l'errore atteso; alarmStateStub.\textit{setActiveAlarms} non viene invocato; il valore di vm.\textit{resolveError} è 'bootstrap error'.}{S}
\addTestU{Contesto \textit{AlarmManagementService}; caso '\textit{resolveAlarm} segue il flusso successo: resolving true, API call, update state, resolving false': verifica che il valore di vm.isResolving è true e il valore di vm.resolvingId è alarmA.id. Esito atteso: il valore di vm.isResolving è true; il valore di vm.resolvingId è alarmA.id; vm.\textit{resolveError} è null; il valore di vm.isResolving è false.}{S}
\addTestU{Contesto \textit{AlarmManagementService}; caso 'prima della resolve il vm riflette solo \textit{AlarmStateService}': verifica che latestVm.alarms coincide con [alarmB] e vmHistory.at(-1) coincide con [alarmB]. Esito atteso: latestVm.alarms coincide con [alarmB]; vmHistory.at(-1) coincide con [alarmB].}{S}
\addTestU{Contesto \textit{AlarmManagementService}; caso 'dopo resolve non mantiene righe locali e riallinea lo stato allo snapshot backend': verifica che vm.alarms coincide con [] e alarmStateStub.\textit{setActiveAlarms} viene invocato con []. Esito atteso: vm.alarms coincide con []; alarmStateStub.\textit{setActiveAlarms} viene invocato con [].}{S}
\addTestU{Contesto \textit{AlarmManagementService}; caso 'dopo resolve ricarica la pagina corrente dal backend': verifica che il valore di vm.alarms.length è 1 e il valore di vm.alarms[0].id è alarmA.id. Esito atteso: il valore di vm.alarms.length è 1; il valore di vm.alarms[0].id è alarmA.id; l'asserzione alarmApiStub.\textit{getActiveAlarms} \textit{toHaveBeenNthCalledWith} 1 risulta soddisfatta; l'asserzione alarmApiStub.\textit{getActiveAlarms} \textit{toHaveBeenNthCalledWith} 2 risulta soddisfatta.}{S}
\addTestU{Contesto \textit{AlarmManagementService}; caso 'dopo resolve su pagina non iniziale vuota torna automaticamente alla pagina precedente': verifica che l'asserzione alarmApiStub.\textit{getActiveAlarms} \textit{toHaveBeenNthCalledWith} 1 risulta soddisfatta e l'asserzione alarmApiStub.\textit{getActiveAlarms} \textit{toHaveBeenNthCalledWith} 2 risulta soddisfatta. Esito atteso: l'asserzione alarmApiStub.\textit{getActiveAlarms} \textit{toHaveBeenNthCalledWith} 1 risulta soddisfatta; l'asserzione alarmApiStub.\textit{getActiveAlarms} \textit{toHaveBeenNthCalledWith} 2 risulta soddisfatta; l'asserzione alarmApiStub.\textit{getActiveAlarms} \textit{toHaveBeenNthCalledWith} 3 risulta soddisfatta; l'asserzione alarmApiStub.\textit{getActiveAlarms} \textit{toHaveBeenNthCalledWith} 4 risulta soddisfatta.}{S}
\addTestU{Contesto \textit{AlarmManagementService}; caso 'nextPage usa limit 6 e incrementa offset quando la pagina corrente e piena': verifica che l'asserzione alarmApiStub.\textit{getActiveAlarms} \textit{toHaveBeenNthCalledWith} 1 risulta soddisfatta e l'asserzione alarmApiStub.\textit{getActiveAlarms} \textit{toHaveBeenNthCalledWith} 2 risulta soddisfatta. Esito atteso: l'asserzione alarmApiStub.\textit{getActiveAlarms} \textit{toHaveBeenNthCalledWith} 1 risulta soddisfatta; l'asserzione alarmApiStub.\textit{getActiveAlarms} \textit{toHaveBeenNthCalledWith} 2 risulta soddisfatta; l'asserzione alarmApiStub.\textit{getActiveAlarms} \textit{toHaveBeenNthCalledWith} 3 risulta soddisfatta; l'asserzione alarmApiStub.\textit{getActiveAlarms} \textit{toHaveBeenNthCalledWith} 4 risulta soddisfatta.}{S}
\addTestU{Contesto \textit{AlarmManagementService}; caso 'nextPage non fa nulla quando canGoNext e false': verifica che alarmApiStub.\textit{getActiveAlarms} viene invocato 2 volte e l'asserzione alarmApiStub.\textit{getActiveAlarms} \textit{toHaveBeenNthCalledWith} 1 risulta soddisfatta. Esito atteso: alarmApiStub.\textit{getActiveAlarms} viene invocato 2 volte; l'asserzione alarmApiStub.\textit{getActiveAlarms} \textit{toHaveBeenNthCalledWith} 1 risulta soddisfatta; l'asserzione alarmApiStub.\textit{getActiveAlarms} \textit{toHaveBeenNthCalledWith} 2 risulta soddisfatta; il valore di vm.currentPage è 1.}{S}
\addTestU{Contesto \textit{AlarmManagementService}; caso 'la pagina successiva con meno di 6 elementi imposta canGoNext a false': verifica che il valore di vm.currentPage è 2 e il valore di vm.canGoNext è false. Esito atteso: il valore di vm.currentPage è 2; il valore di vm.canGoNext è false; vm.alarms.map((alarm) => alarm.id) coincide con ['active-7', 'active-8'].}{S}
\addTestU{Contesto \textit{AlarmManagementService}; caso 'previousPage decrementa l'offset fino a 0': verifica che l'asserzione alarmApiStub.\textit{getActiveAlarms} \textit{toHaveBeenNthCalledWith} 1 risulta soddisfatta e l'asserzione alarmApiStub.\textit{getActiveAlarms} \textit{toHaveBeenNthCalledWith} 2 risulta soddisfatta. Esito atteso: l'asserzione alarmApiStub.\textit{getActiveAlarms} \textit{toHaveBeenNthCalledWith} 1 risulta soddisfatta; l'asserzione alarmApiStub.\textit{getActiveAlarms} \textit{toHaveBeenNthCalledWith} 2 risulta soddisfatta; l'asserzione alarmApiStub.\textit{getActiveAlarms} \textit{toHaveBeenNthCalledWith} 3 risulta soddisfatta; l'asserzione alarmApiStub.\textit{getActiveAlarms} \textit{toHaveBeenNthCalledWith} 4 risulta soddisfatta.}{S}
\addTestU{Contesto \textit{AlarmManagementService}; caso 'previousPage non fa nulla se l'offset e già a zero': verifica che alarmApiStub.\textit{getActiveAlarms} viene invocato 2 volte e l'asserzione alarmApiStub.\textit{getActiveAlarms} \textit{toHaveBeenNthCalledWith} 1 risulta soddisfatta. Esito atteso: alarmApiStub.\textit{getActiveAlarms} viene invocato 2 volte; l'asserzione alarmApiStub.\textit{getActiveAlarms} \textit{toHaveBeenNthCalledWith} 1 risulta soddisfatta; l'asserzione alarmApiStub.\textit{getActiveAlarms} \textit{toHaveBeenNthCalledWith} 2 risulta soddisfatta.}{S}
\addTestU{Contesto \textit{AlarmManagementService}; caso 'in errore resetta resolving e valorizza \textit{resolveError} nel vm': verifica che il valore di vm.isResolving è false e vm.resolvingId è null. Esito atteso: il valore di vm.isResolving è false; vm.resolvingId è null; il valore di vm.\textit{resolveError} è 'timeout rete'; alarmStateStub.\textit{onAlarmResolved} non viene invocato.}{S}
\addTestU{Contesto \textit{AlarmManagementService}; caso '\textit{resolveAlarm} con userId non numerico non chiama API e valorizza errore esplicito': verifica che il valore di vm.isResolving è false e vm.resolvingId è null. Esito atteso: il valore di vm.isResolving è false; vm.resolvingId è null; il valore di vm.\textit{resolveError} è "Utente corrente non valido per risolvere l'allarme."; alarmApiStub.\textit{resolveAlarm} non viene invocato.}{S}
\addTestU{Contesto \textit{AlarmManagementService}; caso 'initialize con errore non strutturato usa fallback di caricamento': verifica che il valore di vm.\textit{resolveError} è 'Errore durante il caricamento degli allarmi attivi.'. Esito atteso: il valore di vm.\textit{resolveError} è 'Errore durante il caricamento degli allarmi attivi.'.}{S}
\addTestU{Contesto \textit{AlarmManagementService}; caso 'initialize con userId non numerico non chiama API e valorizza errore esplicito': verifica che alarmApiStub.\textit{getActiveAlarms} non viene invocato e alarmStateStub.\textit{setActiveAlarms} non viene invocato. Esito atteso: alarmApiStub.\textit{getActiveAlarms} non viene invocato; alarmStateStub.\textit{setActiveAlarms} non viene invocato; il valore di vm.\textit{resolveError} è 'Utente corrente non valido per caricare gli allarmi attivi.'.}{S}
\addTestU{Contesto \textit{AlarmManagementService}; caso 'consente richieste concorrenti e termina resolving solo a completamento di tutte': verifica che il valore di vm.isResolving è true e il valore di vm.resolvingId è 'active-2'. Esito atteso: il valore di vm.isResolving è true; il valore di vm.resolvingId è 'active-2'; il valore di vm.isResolving è false; vm.resolvingId è null.}{S}
\addTestU{Contesto \textit{AlarmManagementService}; caso 'resetta eventuale errore precedente quando parte una nuova \textit{resolveAlarm}': verifica che il valore di vm.\textit{resolveError} è 'errore precedente' e vm.\textit{resolveError} è null. Esito atteso: il valore di vm.\textit{resolveError} è 'errore precedente'; vm.\textit{resolveError} è null.}{S}
\addTestU{Contesto \textit{AlarmHistoryService}; caso 'initialize carica la prima pagina con offset 0': verifica che alarmApiStub.\textit{getResolvedAlarms} viene invocato con 99. Esito atteso: alarmApiStub.\textit{getResolvedAlarms} viene invocato con 99.}{S}
\addTestU{Contesto \textit{AlarmHistoryService}; caso 'initialize con lista vuota non imposta canGoNext': verifica che il valore di vm.alarms.length è 0 e vm.canGoNext è falsy. Esito atteso: il valore di vm.alarms.length è 0; vm.canGoNext è falsy.}{S}
\addTestU{Contesto \textit{AlarmHistoryService}; caso 'initialize con pageLimit+1 elementi imposta canGoNext e fa lo slice': verifica che vm.canGoNext è valorizzato e il valore di vm.alarms.length è 6. Esito atteso: vm.canGoNext è valorizzato; il valore di vm.alarms.length è 6.}{S}
\addTestU{Contesto \textit{AlarmHistoryService}; caso 'initialize con userId non numerico non chiama API': verifica che alarmApiStub.\textit{getResolvedAlarms} non viene invocato. Esito atteso: alarmApiStub.\textit{getResolvedAlarms} non viene invocato.}{S}
\addTestU{Contesto \textit{AlarmHistoryService}; caso 'initialize non propaga eccezioni al chiamante': verifica che () => service.initialize() non genera l'errore atteso. Esito atteso: () => service.initialize() non genera l'errore atteso.}{S}
\addTestU{Contesto \textit{AlarmHistoryService}; caso 'nextPage avanza alla pagina successiva': verifica che il valore di vm.currentPage è 2 e il valore di vm.pageOffset è 6. Esito atteso: il valore di vm.currentPage è 2; il valore di vm.pageOffset è 6; vm.canGoPrevious è valorizzato; vm.alarms.map((a) => a.id) coincide con ['resolved-10'].}{S}
\addTestU{Contesto \textit{AlarmHistoryService}; caso 'nextPage non fa nulla se canGoNext è false': verifica che alarmApiStub.\textit{getResolvedAlarms} viene invocato 1 volte. Esito atteso: alarmApiStub.\textit{getResolvedAlarms} viene invocato 1 volte.}{S}
\addTestU{Contesto \textit{AlarmHistoryService}; caso 'pagina successiva con meno di pageLimit elementi imposta canGoNext a false': verifica che vm.canGoNext è falsy e vm.alarms.map((a) => a.id) coincide con ['resolved-10', 'resolved-11']. Esito atteso: vm.canGoNext è falsy; vm.alarms.map((a) => a.id) coincide con ['resolved-10', 'resolved-11'].}{S}
\addTestU{Contesto \textit{AlarmHistoryService}; caso 'previousPage torna alla pagina precedente': verifica che il valore di vm.currentPage è 1 e il valore di vm.pageOffset è 0. Esito atteso: il valore di vm.currentPage è 1; il valore di vm.pageOffset è 0; vm.canGoPrevious è falsy; l'asserzione alarmApiStub.\textit{getResolvedAlarms} \textit{toHaveBeenNthCalledWith} 1 risulta soddisfatta.}{S}
\addTestU{Contesto \textit{AlarmHistoryService}; caso 'previousPage non fa nulla se l'offset è già a zero': verifica che alarmApiStub.\textit{getResolvedAlarms} viene invocato 1 volte. Esito atteso: alarmApiStub.\textit{getResolvedAlarms} viene invocato 1 volte.}{S}
\addTestU{Contesto \textit{AnalyticsApiService}; caso 'should be created': verifica che service è valorizzato. Esito atteso: service è valorizzato.}{S}
\addTestU{Contesto \textit{AnalyticsApiService} - \textit{getAllApartments}; caso 'should fetch all apartments and map them correctly': verifica che il valore di apartments.length è 2 e apartments[0] coincide con \{ id: '1', name: 'Apartment 1' \}. Esito atteso: il valore di apartments.length è 2; apartments[0] coincide con \{ id: '1', name: 'Apartment 1' \}; l'asserzione apartments[0] \textit{toHaveProperty} 'extraField' risulta non soddisfatta; il valore di req.request.method è 'GET'.}{S}
\addTestU{Contesto \textit{AnalyticsApiService} - \textit{getAnalytics}; caso 'should fetch analytics and map the complex DTO structure': verifica che il valore di result.apartmentId è apartmentId e il valore di result.analyticsInfo.length è 1. Esito atteso: il valore di result.apartmentId è apartmentId; il valore di result.analyticsInfo.length è 1; il valore di info.title è 'Test Chart'; info.suggestions.messages coincide con ['Suggestion 1'].}{S}
\addTestU{Contesto \textit{ApartmentApiService}; caso 'RF83-OBL abilita un appartamento nel Sistema': verifica che il valore di request.request.method è 'PATCH' e request.request.body coincide con \{\}. Esito atteso: il valore di request.request.method è 'PATCH'; request.request.body coincide con \{\}.}{S}
\addTestU{Contesto \textit{ApartmentApiService}; caso 'RF84-OBL disabilita un appartamento nel Sistema': verifica che il valore di request.request.method è 'PATCH' e request.request.body coincide con \{\}. Esito atteso: il valore di request.request.method è 'PATCH'; request.request.body coincide con \{\}.}{S}
\addTestU{Contesto Action Card Components; caso '\textit{AlarmButtonCardComponent} emette actionRequested con payload corretto': verifica che emitSpy viene invocato con \{ roomId: 'room-1', deviceId: 'device-1', action: 'toggle', \} e emitSpy viene invocato 1 volte. Esito atteso: emitSpy viene invocato con \{ roomId: 'room-1', deviceId: 'device-1', action: 'toggle', \}; emitSpy viene invocato 1 volte.}{S}
\addTestU{Contesto Action Card Components; caso '\textit{BlindCardComponent} emette actionRequested con payload corretto': verifica che emitSpy viene invocato con \{ roomId: 'room-1', deviceId: 'device-1', action: 'toggle', \} e emitSpy viene invocato 1 volte. Esito atteso: emitSpy viene invocato con \{ roomId: 'room-1', deviceId: 'device-1', action: 'toggle', \}; emitSpy viene invocato 1 volte.}{S}
\addTestU{Contesto Action Card Components; caso '\textit{EntranceDoorCardComponent} emette actionRequested con payload corretto': verifica che emitSpy viene invocato con \{ roomId: 'room-1', deviceId: 'device-1', action: 'toggle', \} e emitSpy viene invocato 1 volte. Esito atteso: emitSpy viene invocato con \{ roomId: 'room-1', deviceId: 'device-1', action: 'toggle', \}; emitSpy viene invocato 1 volte.}{S}
\addTestU{Contesto Action Card Components; caso '\textit{FallSensorCardComponent} emette actionRequested con payload corretto': verifica che emitSpy viene invocato con \{ roomId: 'room-1', deviceId: 'device-1', action: 'toggle', \} e emitSpy viene invocato 1 volte. Esito atteso: emitSpy viene invocato con \{ roomId: 'room-1', deviceId: 'device-1', action: 'toggle', \}; emitSpy viene invocato 1 volte.}{S}
\addTestU{Contesto Action Card Components; caso '\textit{LightCardComponent} emette actionRequested con payload corretto': verifica che emitSpy viene invocato con \{ roomId: 'room-1', deviceId: 'device-1', action: 'toggle', \} e emitSpy viene invocato 1 volte. Esito atteso: emitSpy viene invocato con \{ roomId: 'room-1', deviceId: 'device-1', action: 'toggle', \}; emitSpy viene invocato 1 volte.}{S}
\addTestU{Contesto Action Card Components; caso '\textit{PresenceSensorCardComponent} emette actionRequested con payload corretto': verifica che emitSpy viene invocato con \{ roomId: 'room-1', deviceId: 'device-1', action: 'toggle', \} e emitSpy viene invocato 1 volte. Esito atteso: emitSpy viene invocato con \{ roomId: 'room-1', deviceId: 'device-1', action: 'toggle', \}; emitSpy viene invocato 1 volte.}{S}
\addTestU{Contesto Action Card Components; caso '\textit{ThermostatCardComponent} emette actionRequested con payload corretto': verifica che emitSpy viene invocato con \{ roomId: 'room-1', deviceId: 'device-1', action: 'toggle', \} e emitSpy viene invocato 1 volte. Esito atteso: emitSpy viene invocato con \{ roomId: 'room-1', deviceId: 'device-1', action: 'toggle', \}; emitSpy viene invocato 1 volte.}{S}
\addTestU{Contesto Action Card Components; caso '\textit{ThermostatCardComponent} espone una label valida per il tipo': verifica che l'asserzione component.\textit{getTypeLabel}(DeviceType.THERMOSTAT) \textit{toBeTypeOf} 'string' risulta soddisfatta e component.\textit{getTypeLabel}(DeviceType.THERMOSTAT).length è maggiore di 0. Esito atteso: l'asserzione component.\textit{getTypeLabel}(DeviceType.THERMOSTAT) \textit{toBeTypeOf} 'string' risulta soddisfatta; component.\textit{getTypeLabel}(DeviceType.THERMOSTAT).length è maggiore di 0.}{S}
\addTestU{Contesto \textit{DeviceApiService}; caso 'mappa i datapoint dal payload plant nella stanza richiesta': verifica che il valore di request.request.method è 'GET' e il valore di room.id è 'room-1'. Esito atteso: il valore di request.request.method è 'GET'; il valore di room.id è 'room-1'; room.devices ha lunghezza 1; room.devices[0].datapoints ha lunghezza 1.}{S}
\addTestU{Contesto \textit{DeviceApiService}; caso 'invia la scrittura datapoint su POST /api/device con payload corretto': verifica che il valore di request.request.method è 'POST' e request.request.body coincide con \{ datapointId: 'dp-1', value: 'On' \}. Esito atteso: il valore di request.request.method è 'POST'; request.request.body coincide con \{ datapointId: 'dp-1', value: 'On' \}.}{S}
\addTestU{Contesto \textit{DeviceApiService}; caso 'propaga errore quando la scrittura datapoint fallisce': verifica che dopo rigetto, writePromise è valorizzato. Esito atteso: dopo rigetto, writePromise è valorizzato.}{S}
\addTestU{Contesto \textit{DeviceApiService}; caso 'quando activePlantId manca, risolve il plant dall endpoint /plant/all': verifica che il valore di allPlantsRequest.request.method è 'GET' e il valore di plantByIdRequest.request.method è 'GET'. Esito atteso: il valore di allPlantsRequest.request.method è 'GET'; il valore di plantByIdRequest.request.method è 'GET'; rows coincide con []; il valore di storage.\textit{getItem}('activePlantId') è 'plant-2'.}{S}
\addTestU{Contesto \textit{DeviceApiService}; caso 'esclude dispositivi energy measure dalla tabella endpoint': verifica che il valore di request.request.method è 'GET' e rows ha lunghezza 1. Esito atteso: il valore di request.request.method è 'GET'; rows ha lunghezza 1; il valore di rows[0].deviceId è 'allowed-device'.}{S}
\addTestU{Contesto \textit{DeviceApiService}; caso 'deduplica dispositivi con stesso nome mantenendo quello a priorita maggiore': verifica che il valore di request.request.method è 'GET' e rows ha lunghezza 1. Esito atteso: il valore di request.request.method è 'GET'; rows ha lunghezza 1; il valore di rows[0].deviceId è 'priority-device'.}{S}
\addTestU{Contesto \textit{DeviceApiService}; caso 'recupera i valori correnti dai device unici e li mappa per datapoint': verifica che il valore di req1.request.method è 'GET' e il valore di req2.request.method è 'GET'. Esito atteso: il valore di req1.request.method è 'GET'; il valore di req2.request.method è 'GET'; il valore di values.get('dp-1') è 'On'; il valore di values.get('dp-2') è 21.}{S}
\addTestU{Contesto \textit{DeviceApiService}; caso 'continua il merge valori anche se un device value endpoint fallisce': verifica che il valore di values.get('dp-ok') è 'Off'. Esito atteso: il valore di values.get('dp-ok') è 'Off'.}{S}
\addTestU{Contesto \textit{NavService}; caso 'viene creato': verifica che service è valorizzato. Esito atteso: service è valorizzato.}{S}
\addTestU{Contesto \textit{NavService}; caso 'restituisce tutti i NavItems per AMMINISTRATORE': verifica che items.length è maggiore di 0 e il valore di items.some(i => i.label === "Dashboard") è true. Esito atteso: items.length è maggiore di 0; il valore di items.some(i => i.label === "Dashboard") è true; il valore di items.some(i => i.route.includes('dashboard')) è true; il valore di items.some(i => i.label === "Configurazione Allarmi") è true.}{S}
\addTestU{Contesto \textit{NavService}; caso 'restituisce solo i NavItems filtrati per OPERATORE SANITARIO': verifica che OSSItems.length è maggiore di 0 e il valore di roleReqItems.length è 0. Esito atteso: OSSItems.length è maggiore di 0; il valore di roleReqItems.length è 0; il valore di items.some(i => i.label === "Dashboard") è true; il valore di items.some(i => i.route.includes('dashboard')) è true.}{S}
\addTestU{Contesto \textit{MyVimarCloudApiFeatureService}; caso 'chiama GET /my-vimar/account in \textit{getLinkedAccount}': verifica che account coincide con \{ email: 'admin@example.com', isLinked: true \} e il valore di request.request.method è 'GET'. Esito atteso: account coincide con \{ email: 'admin@example.com', isLinked: true \}; il valore di request.request.method è 'GET'.}{S}
\addTestU{Contesto \textit{MyVimarCloudApiFeatureService}; caso 'fallback su /api/vimar-account quando /my-vimar/account risponde 404': verifica che account coincide con \{ email: 'admin@example.com', isLinked: true \} e il valore di firstRequest.request.method è 'GET'. Esito atteso: account coincide con \{ email: 'admin@example.com', isLinked: true \}; il valore di firstRequest.request.method è 'GET'; il valore di fallbackRequest.request.method è 'GET'.}{S}
\addTestU{Contesto \textit{MyVimarCloudApiFeatureService}; caso 'ritorna account scollegato quando entrambi gli endpoint non esistono': verifica che account coincide con \{ email: '', isLinked: false \}. Esito atteso: account coincide con \{ email: '', isLinked: false \}.}{S}
\addTestU{Contesto \textit{MyVimarCloudApiFeatureService}; caso 'chiama DELETE /my-vimar/account in unlinkAccount': verifica che result è null e il valore di request.request.method è 'DELETE'. Esito atteso: result è null; il valore di request.request.method è 'DELETE'.}{S}
\addTestU{Contesto \textit{MyVimarCloudApiFeatureService}; caso 'fallback su DELETE /api/vimar-account quando /my-vimar/account non esiste': verifica che result è null e il valore di firstRequest.request.method è 'DELETE'. Esito atteso: result è null; il valore di firstRequest.request.method è 'DELETE'; il valore di fallbackRequest.request.method è 'DELETE'.}{S}
\addTestU{Contesto \textit{MyVimarCloudApiFeatureService}; caso 'richiede ticket OAuth e poi naviga su /api/auth/authorize': verifica che il valore di request.request.method è 'POST' e il valore di mockLocation.href è `\$\{baseUrl\}/api/auth/authorize?ticket=ticket-123\&redirect\_url=http\%3A\%2F\%2Flocalhost\%3A4200\%2F.... Esito atteso: il valore di request.request.method è 'POST'; il valore di mockLocation.href è `\$\{baseUrl\}/api/auth/authorize?ticket=ticket-123\&redirect\_url=http\%3A\%2F\%2Flocalhost\%3A4200\%2F....}{S}
\addTestU{Contesto \textit{MyVimarCloudApiFeatureService}; caso 'mostra alert quando il backend restituisce ticket vuoto': verifica che mockAlert viene invocato 1 volte e mockAlert viene invocato con 'Ticket OAuth non valido. Riprova tra qualche istante.'. Esito atteso: mockAlert viene invocato 1 volte; mockAlert viene invocato con 'Ticket OAuth non valido. Riprova tra qualche istante.'; il valore di mockLocation.href è 'http://localhost:4200/vimar-link'.}{S}
\addTestU{Contesto \textit{MyVimarCloudApiFeatureService}; caso 'mostra alert con messaggio mappato quando prepare-oauth fallisce': verifica che apiErrorDisplayServiceStub.\textit{toMessage} viene invocato 1 volte e mockAlert viene invocato 1 volte. Esito atteso: apiErrorDisplayServiceStub.\textit{toMessage} viene invocato 1 volte; mockAlert viene invocato 1 volte; mockAlert viene invocato con 'Errore avvio OAuth'; il valore di mockLocation.href è 'http://localhost:4200/vimar-link'.}{S}
\addTestU{Contesto \textit{NotificationService}; caso 'espone vm iniziale vuoto quando storico e stream in-session sono vuoti': verifica che vm coincide con \{ notifications: [], unreadCount: 0, \}. Esito atteso: vm coincide con \{ notifications: [], unreadCount: 0, \}.}{S}
\addTestU{Contesto \textit{NotificationService}; caso 'combina storico e in-session, deduplica per notificationId e ordina per sentAt desc': verifica che il valore di vm.unreadCount è 2 e vm.notifications coincide con [sharedInSession, inSessionOnly, historyOlder]. Esito atteso: il valore di vm.unreadCount è 2; vm.notifications coincide con [sharedInSession, inSessionOnly, historyOlder]; vm.notifications ha lunghezza 3.}{S}
\addTestU{Contesto \textit{NotificationService}; caso 'in errore sullo storico applica fallback a [] e mantiene i dati in-session': verifica che vm.notifications coincide con [inSessionOnly] e il valore di vm.unreadCount è 1. Esito atteso: vm.notifications coincide con [inSessionOnly]; il valore di vm.unreadCount è 1.}{S}
\addTestU{Contesto \textit{NotificationService}; caso 'propaga nuovi unreadCount senza alterare la lista notifiche': verifica che vm.notifications coincide con [historyOlder] e il valore di vm.unreadCount è 5. Esito atteso: vm.notifications coincide con [historyOlder]; il valore di vm.unreadCount è 5.}{S}
\addTestU{Contesto \textit{NotificationService}; caso 'aggiorna il vm quando arrivano nuove notifiche in-session dopo la prima emissione': verifica che vm.notifications coincide con [inSessionOnly, historyOlder] e il valore di vm.unreadCount è 0. Esito atteso: vm.notifications coincide con [inSessionOnly, historyOlder]; il valore di vm.unreadCount è 0.}{S}
\addTestU{Contesto \textit{NotificationService}; caso 'deduplica tenendo la versione in-session quando notificationId è presente anche nello storico': verifica che vm.notifications coincide con [sharedInSession] e vm.notifications ha lunghezza 1. Esito atteso: vm.notifications coincide con [sharedInSession]; vm.notifications ha lunghezza 1.}{S}
\addTestU{Contesto \textit{NotificationService}; caso 'gestisce timestamp non validi senza lanciare errori e li ordina in fondo': verifica che vm.notifications coincide con [historyOlder, invalidTimestampNotification]. Esito atteso: vm.notifications coincide con [historyOlder, invalidTimestampNotification].}{S}
\addTestU{Contesto \textit{NotificationService}; caso 'riusa la stessa pipeline vm con più subscriber senza richiamare di nuovo lo storico': verifica che vmA coincide con vmB e notificationApiStub.\textit{getNotificationsHistory} viene invocato 1 volte. Esito atteso: vmA coincide con vmB; notificationApiStub.\textit{getNotificationsHistory} viene invocato 1 volte.}{S}
\addTestU{Contesto \textit{NotificationService}; caso 'mantiene uno snapshot coerente per subscriber tardivi senza ricaricare lo storico': verifica che vmLateSubscriber.notifications coincide con [inSessionOnly, historyOlder] e notificationApiStub.\textit{getNotificationsHistory} viene invocato 1 volte. Esito atteso: vmLateSubscriber.notifications coincide con [inSessionOnly, historyOlder]; notificationApiStub.\textit{getNotificationsHistory} viene invocato 1 volte.}{S}
\addTestU{Contesto \textit{NotificationService}; caso 'non muta gli array sorgente ricevuti da API e \textit{AlarmStateService}': verifica che historicList coincide con [sharedHistoric, historyOlder] e inSessionList coincide con [sharedInSession, inSessionOnly]. Esito atteso: historicList coincide con [sharedHistoric, historyOlder]; inSessionList coincide con [sharedInSession, inSessionOnly].}{S}
\addTestU{Contesto \textit{NotificationService}; caso 'deduplica anche duplicati interni alla stessa sorgente mantenendo l ultimo elemento': verifica che vm.notifications coincide con [duplicatedSessionB, duplicatedHistoryB] e vm.notifications ha lunghezza 2. Esito atteso: vm.notifications coincide con [duplicatedSessionB, duplicatedHistoryB]; vm.notifications ha lunghezza 2.}{S}
\addTestU{Contesto \textit{NotificationService}; caso 'mantiene coerenza: unreadCount segue lo stream dedicato anche quando la deduplica riduce la lista': verifica che vm.notifications ha lunghezza 1 e il valore di vm.unreadCount è 5. Esito atteso: vm.notifications ha lunghezza 1; il valore di vm.unreadCount è 5.}{S}
\addTestU{Contesto \textit{NotificationService}; caso 'filtra le notifiche dismissate localmente dall archivio unificato': verifica che vm.notifications.map((notification) => notification.notificationId) coincide con [ 'session-1', 'shared-id', ]. Esito atteso: vm.notifications.map((notification) => notification.notificationId) coincide con [ 'session-1', 'shared-id', ].}{S}
\addTestU{Contesto \textit{NotificationService}; caso '\textit{removeNotification} delega dismiss + remove allo stato realtime': verifica che alarmStateStub.\textit{dismissNotification} viene invocato con 'n-100' e alarmStateStub.\textit{removeNotification} viene invocato con 'n-100'. Esito atteso: alarmStateStub.\textit{dismissNotification} viene invocato con 'n-100'; alarmStateStub.\textit{removeNotification} viene invocato con 'n-100'.}{S}
\addTestU{Contesto \textit{NotificationService}; caso '\textit{clearAllNotifications} delega dismiss multiplo + clear realtime': verifica che alarmStateStub.\textit{dismissNotifications} viene invocato con [ historyOlder.notificationId, inSessionOnly.notificationId, ] e alarmStateStub.\textit{clearNotifications} viene invocato 1 volte. Esito atteso: alarmStateStub.\textit{dismissNotifications} viene invocato con [ historyOlder.notificationId, inSessionOnly.notificationId, ]; alarmStateStub.\textit{clearNotifications} viene invocato 1 volte.}{S}
\addTestU{Contesto \textit{UserApiService}; caso '\textit{deleteUser} chiama DELETE /users/:id con id numerico': verifica che result è null e il valore di request.request.method è 'DELETE'. Esito atteso: result è null; il valore di request.request.method è 'DELETE'.}{S}
\addTestU{Contesto \textit{UserApiService}; caso '\textit{createUser} decodifica tempPassword base64 prima di esporla al chiamante': verifica che il valore di result.tempPassword è 'TempPass123' e il valore di request.request.method è 'POST'. Esito atteso: il valore di result.tempPassword è 'TempPass123'; il valore di request.request.method è 'POST'; request.request.body coincide con dto.}{S}
\addTestU{Contesto \textit{UserApiService}; caso '\textit{createUser} mantiene tempPassword invariata se non è base64 valida': verifica che il valore di result.tempPassword è 'TempPass123'. Esito atteso: il valore di result.tempPassword è 'TempPass123'.}{S}
\addTestU{Contesto \textit{AssignmentOperationsService}; caso '\textit{assignOperator} esegue operazione e poi ricarica i wards': verifica che result è undefined e apiStub.\textit{assignOperatorToWard} viene invocato con 1. Esito atteso: result è undefined; apiStub.\textit{assignOperatorToWard} viene invocato con 1; l'asserzione apiStub.\textit{getWards} \textit{toHaveBeenCalledOnce} risulta soddisfatta; apiStub.getPlantsByWardId viene invocato con 1.}{S}
\addTestU{Contesto \textit{AssignmentOperationsService}; caso '\textit{removeOperator} in errore espone errore testo da \textit{HttpErrorResponse}.error': verifica che il valore di emitted è false e storeStub.\textit{setError} viene invocato con 'Operatore non assegnato'. Esito atteso: il valore di emitted è false; storeStub.\textit{setError} viene invocato con 'Operatore non assegnato'; storeStub.\textit{setError} viene invocato 1 volte; apiStub.\textit{getWards} non viene invocato.}{S}
\addTestU{Contesto \textit{AssignmentOperationsService}; caso '\textit{getAvailablePlantsForWard} filtra i plant già assegnati in qualunque ward': verifica che result coincide con [\{ id: '203', name: 'App. 203' \}] e l'asserzione apiStub.\textit{getAvailablePlants} \textit{toHaveBeenCalledOnce} risulta soddisfatta. Esito atteso: result coincide con [\{ id: '203', name: 'App. 203' \}]; l'asserzione apiStub.\textit{getAvailablePlants} \textit{toHaveBeenCalledOnce} risulta soddisfatta.}{S}
\addTestU{Contesto \textit{AssignmentOperationsService}; caso 'dopo rimozione un appartamento torna disponibile': verifica che result coincide con [\{ id: '202', name: 'App. 202' \}] e result coincide con [ \{ id: '201', name: 'App. 201' \}, \{ id: '202', name: 'App. 202' \}, ]. Esito atteso: result coincide con [\{ id: '202', name: 'App. 202' \}]; result coincide con [ \{ id: '201', name: 'App. 201' \}, \{ id: '202', name: 'App. 202' \}, ].}{S}
\addTestU{Contesto \textit{AssignmentOperationsService}; caso 'dopo eliminazione reparto, gli appartamenti del reparto eliminato tornano disponibili': verifica che result coincide con [\{ id: '203', name: 'App. 203' \}] e result coincide con [ \{ id: '201', name: 'App. 201' \}, \{ id: '203', name: 'App. 203' \}, ]. Esito atteso: result coincide con [\{ id: '203', name: 'App. 203' \}]; result coincide con [ \{ id: '201', name: 'App. 201' \}, \{ id: '203', name: 'App. 203' \}, ].}{S}
\addTestU{Contesto \textit{AssignmentOperationsService}; caso '\textit{assignPlant} in errore usa \textit{HttpErrorResponse}.message quando non c'è error string': verifica che messageArg.length è maggiore di 0 e apiStub.\textit{getWards} non viene invocato. Esito atteso: messageArg.length è maggiore di 0; apiStub.\textit{getWards} non viene invocato; storeStub.\textit{setWards} non viene invocato.}{S}
\addTestU{Contesto \textit{AssignmentOperationsService}; caso '\textit{removePlant} in errore non-http usa fallback message': verifica che storeStub.\textit{setError} viene invocato con 'Operazione di assegnazione non riuscita.' e storeStub.\textit{setError} viene invocato 1 volte. Esito atteso: storeStub.\textit{setError} viene invocato con 'Operazione di assegnazione non riuscita.'; storeStub.\textit{setError} viene invocato 1 volte; apiStub.\textit{getWards} non viene invocato.}{S}
\addTestU{Contesto \textit{WardApiService}; caso 'chiama GET /api/wards in \textit{getWards}': verifica che result coincide con [ \{ id: 1, name: 'Cardiologia', \}, ] e result ha lunghezza 1. Esito atteso: result coincide con [ \{ id: 1, name: 'Cardiologia', \}, ]; result ha lunghezza 1; il valore di request.request.method è 'GET'.}{S}
\addTestU{Contesto \textit{WardApiService}; caso 'chiama GET /api/wards-users-relationships/:wardId in getOperatorsByWardId': verifica che result coincide con [\{ id: 1, username: 'mrossi' \}] e il valore di request.request.method è 'GET'. Esito atteso: result coincide con [\{ id: 1, username: 'mrossi' \}]; il valore di request.request.method è 'GET'.}{S}
\addTestU{Contesto \textit{WardApiService}; caso 'chiama GET /users in \textit{getAvailableOperators}': verifica che result coincide con [\{ id: 2, username: 'lverdi' \}] e il valore di request.request.method è 'GET'. Esito atteso: result coincide con [\{ id: 2, username: 'lverdi' \}]; il valore di request.request.method è 'GET'.}{S}
\addTestU{Contesto \textit{WardApiService}; caso 'chiama GET /api/wards-plants-relationships/:wardId in getPlantsByWardId': verifica che result coincide con [\{ id: '101', name: 'App. 101' \}] e il valore di request.request.method è 'GET'. Esito atteso: result coincide con [\{ id: '101', name: 'App. 101' \}]; il valore di request.request.method è 'GET'.}{S}
\addTestU{Contesto \textit{WardApiService}; caso 'chiama GET /api/plant/all in \textit{getAvailablePlants}': verifica che result coincide con [\{ id: '103', name: 'App. 103' \}] e il valore di request.request.method è 'GET'. Esito atteso: result coincide con [\{ id: '103', name: 'App. 103' \}]; il valore di request.request.method è 'GET'.}{S}
\addTestU{Contesto \textit{WardApiService}; caso 'chiama POST /api/wards in \textit{createWard}': verifica che result coincide con ward e il valore di result.id è 1. Esito atteso: result coincide con ward; il valore di result.id è 1; il valore di request.request.method è 'POST'; request.request.body coincide con dto.}{S}
\addTestU{Contesto \textit{WardApiService}; caso 'chiama PUT /api/wards/:id in \textit{updateWard}': verifica che result coincide con ward e il valore di request.request.method è 'PUT'. Esito atteso: result coincide con ward; il valore di request.request.method è 'PUT'; request.request.body coincide con \{ name: 'Cardio 2' \}.}{S}
\addTestU{Contesto \textit{WardApiService}; caso 'chiama DELETE /api/wards-users-relationships/:wardId/:userId': verifica che result è null e il valore di request.request.method è 'DELETE'. Esito atteso: result è null; il valore di request.request.method è 'DELETE'.}{S}
\addTestU{Contesto \textit{WardApiService}; caso 'chiama POST /api/wards-plants-relationships con wardId e plantId': verifica che result è null e il valore di request.request.method è 'POST'. Esito atteso: result è null; il valore di request.request.method è 'POST'; request.request.body coincide con \{ wardId: 1, plantId: '102' \}.}{S}
\addTestU{Contesto \textit{WardApiService}; caso 'chiama DELETE /api/wards/:id in \textit{deleteWard}': verifica che result è null e il valore di request.request.method è 'DELETE'. Esito atteso: result è null; il valore di request.request.method è 'DELETE'.}{S}
\addTestU{Contesto \textit{WardApiService}; caso 'chiama POST /api/wards-users-relationships con wardId e userId': verifica che result è null e il valore di request.request.method è 'POST'. Esito atteso: result è null; il valore di request.request.method è 'POST'; request.request.body coincide con \{ wardId: 1, userId: 9 \}.}{S}
\addTestU{Contesto \textit{WardApiService}; caso 'chiama DELETE /api/wards-plants-relationships/:plantId': verifica che result è null e il valore di request.request.method è 'DELETE'. Esito atteso: result è null; il valore di request.request.method è 'DELETE'.}{S}
\addTestU{Contesto \textit{WardApiService}; caso 'mantiene compatibilita con \textit{assignPlantToWard} mappando plantId -> plantId': verifica che result è null e il valore di request.request.method è 'POST'. Esito atteso: result è null; il valore di request.request.method è 'POST'; request.request.body coincide con \{ wardId: 1, plantId: '103' \}.}{S}
\addTestU{Contesto \textit{WardManagementStore}; caso '\textit{loadWards} imposta loading e delega a \textit{WardOperationsService}': verifica che wardStoreStub.\textit{setLoading} viene invocato con true e wardStoreStub.\textit{setLoading} viene invocato 1 volte. Esito atteso: wardStoreStub.\textit{setLoading} viene invocato con true; wardStoreStub.\textit{setLoading} viene invocato 1 volte; l'asserzione wardOperationsStub.\textit{loadWards} \textit{toHaveBeenCalledOnce} risulta soddisfatta.}{S}
\addTestU{Contesto \textit{WardManagementStore}; caso 'create/update/delete delegano al service corretto': verifica che wardOperationsStub.\textit{createWard} viene invocato con \{ name: 'Cardio' \} e wardOperationsStub.\textit{updateWard} viene invocato con 1. Esito atteso: wardOperationsStub.\textit{createWard} viene invocato con \{ name: 'Cardio' \}; wardOperationsStub.\textit{updateWard} viene invocato con 1; wardOperationsStub.\textit{deleteWard} viene invocato con 1; wardStoreStub.\textit{setLoading} viene invocato 3 volte.}{S}
\addTestU{Contesto \textit{WardManagementStore}; caso 'assign/remove operator e plant delegano ad \textit{AssignmentOperationsService}': verifica che assignmentOperationsStub.\textit{assignOperator} viene invocato con 1 e assignmentOperationsStub.\textit{removeOperator} viene invocato con 1. Esito atteso: assignmentOperationsStub.\textit{assignOperator} viene invocato con 1; assignmentOperationsStub.\textit{removeOperator} viene invocato con 1; assignmentOperationsStub.\textit{assignPlant} viene invocato con 1; assignmentOperationsStub.\textit{removePlant} viene invocato con 1.}{S}
\addTestU{Contesto \textit{WardManagementStore}; caso '\textit{getAvailablePlantsForWard} delega ad \textit{AssignmentOperationsService}': verifica che assignmentOperationsStub.\textit{getAvailablePlantsForWard} viene invocato con 10 e result coincide con [\{ id: '200', name: 'App. 200' \}]. Esito atteso: assignmentOperationsStub.\textit{getAvailablePlantsForWard} viene invocato con 10; result coincide con [\{ id: '200', name: 'App. 200' \}]; wardStoreStub.\textit{setError} non viene invocato.}{S}
\addTestU{Contesto \textit{WardManagementStore}; caso '\textit{getAvailablePlantsForWard} in errore ritorna null e setta errore': verifica che result è null e wardStoreStub.\textit{setError} viene invocato con 'fetch failed'. Esito atteso: result è null; wardStoreStub.\textit{setError} viene invocato con 'fetch failed'.}{S}
\addTestU{Contesto \textit{WardManagementStore}; caso '\textit{getAvailablePlantsForWard} con errore non strutturato usa il messaggio di fallback': verifica che result è null e wardStoreStub.\textit{setError} viene invocato con 'Operazione plant non riuscita.'. Esito atteso: result è null; wardStoreStub.\textit{setError} viene invocato con 'Operazione plant non riuscita.'.}{S}
\addTestU{Contesto \textit{WardManagementStore}; caso '\textit{getAvailableUsersForWard} in errore ritorna null e setta errore': verifica che result è null e wardStoreStub.\textit{setError} viene invocato con 'users fetch failed'. Esito atteso: result è null; wardStoreStub.\textit{setError} viene invocato con 'users fetch failed'.}{S}
\addTestU{Contesto \textit{WardOperationsService}; caso '\textit{loadWards} aggiorna store e termina senza valore': verifica che result è undefined e l'asserzione apiStub.\textit{getWards} \textit{toHaveBeenCalledOnce} risulta soddisfatta. Esito atteso: result è undefined; l'asserzione apiStub.\textit{getWards} \textit{toHaveBeenCalledOnce} risulta soddisfatta; apiStub.getPlantsByWardId viene invocato con 1; apiStub.getOperatorsByWardId viene invocato con 1.}{S}
\addTestU{Contesto \textit{WardOperationsService}; caso '\textit{loadWards} con lista vuota non richiama le relationship API': verifica che l'asserzione apiStub.\textit{getWards} \textit{toHaveBeenCalledOnce} risulta soddisfatta e apiStub.getPlantsByWardId non viene invocato. Esito atteso: l'asserzione apiStub.\textit{getWards} \textit{toHaveBeenCalledOnce} risulta soddisfatta; apiStub.getPlantsByWardId non viene invocato; apiStub.getOperatorsByWardId non viene invocato; storeStub.\textit{setWards} viene invocato con [].}{S}
\addTestU{Contesto \textit{WardOperationsService}; caso '\textit{createWard} con 409 imposta messaggio funzionale': verifica che il valore di emitted è false e storeStub.addWard non viene invocato. Esito atteso: il valore di emitted è false; storeStub.addWard non viene invocato; storeStub.\textit{setError} viene invocato con 'Esiste già un reparto con questo nome.'; storeStub.\textit{setError} viene invocato 1 volte.}{S}
\addTestU{Contesto \textit{WardOperationsService}; caso '\textit{createWard} normalizza payload minimale mantenendo il modello Ward coerente': verifica che storeStub.addWard viene invocato con \{ id: 11, name: 'Oncologia', apartments: [], operators: [], \} e storeStub.\textit{setLoading} viene invocato con false. Esito atteso: storeStub.addWard viene invocato con \{ id: 11, name: 'Oncologia', apartments: [], operators: [], \}; storeStub.\textit{setLoading} viene invocato con false.}{S}
\addTestU{Contesto \textit{WardOperationsService}; caso '\textit{createWard} preserva apartments/operators dallo snapshot se il backend non li restituisce': verifica che storeStub.addWard viene invocato con \{ id: 1, name: 'Cardiologia nuova', apartments: ward.apartments, operators: ward.operators, \}. Esito atteso: storeStub.addWard viene invocato con \{ id: 1, name: 'Cardiologia nuova', apartments: ward.apartments, operators: ward.operators, \}.}{S}
\addTestU{Contesto \textit{WardOperationsService}; caso '\textit{updateWard} sostituisce il ward nel store': verifica che apiStub.\textit{updateWard} viene invocato con 1 e storeStub.\textit{replaceWard} viene invocato con updatedWard. Esito atteso: apiStub.\textit{updateWard} viene invocato con 1; storeStub.\textit{replaceWard} viene invocato con updatedWard; storeStub.\textit{replaceWard} viene invocato 1 volte; storeStub.\textit{setLoading} viene invocato con false.}{S}
\addTestU{Contesto \textit{WardOperationsService}; caso '\textit{updateWard} preserva apartments/operators quando backend restituisce solo id e name': verifica che storeStub.\textit{replaceWard} viene invocato con \{ id: 1, name: 'Cardiologia X', apartments: ward.apartments, operators: ward.operators, \} e storeStub.\textit{setLoading} viene invocato con false. Esito atteso: storeStub.\textit{replaceWard} viene invocato con \{ id: 1, name: 'Cardiologia X', apartments: ward.apartments, operators: ward.operators, \}; storeStub.\textit{setLoading} viene invocato con false.}{S}
\addTestU{Contesto \textit{WardOperationsService}; caso '\textit{updateWard} gestisce errore HTTP con messaggio testuale del server': verifica che storeStub.\textit{setError} viene invocato con 'update failed' e storeStub.\textit{replaceWard} non viene invocato. Esito atteso: storeStub.\textit{setError} viene invocato con 'update failed'; storeStub.\textit{replaceWard} non viene invocato.}{S}
\addTestU{Contesto \textit{WardOperationsService}; caso '\textit{deleteWard} rimuove il ward dal store': verifica che apiStub.\textit{deleteWard} viene invocato con 1 e storeStub.\textit{removeWard} viene invocato con 1. Esito atteso: apiStub.\textit{deleteWard} viene invocato con 1; storeStub.\textit{removeWard} viene invocato con 1; storeStub.\textit{removeWard} viene invocato 1 volte; storeStub.\textit{setLoading} viene invocato con false.}{S}
\addTestU{Contesto \textit{WardOperationsService}; caso '\textit{deleteWard} usa il fallback generico quando l'errore non è strutturato': verifica che storeStub.\textit{setError} viene invocato con 'Operazione sui reparti non riuscita.' e storeStub.\textit{removeWard} non viene invocato. Esito atteso: storeStub.\textit{setError} viene invocato con 'Operazione sui reparti non riuscita.'; storeStub.\textit{removeWard} non viene invocato.}{S}
\addTestU{Contesto \textit{WardOperationsService}; caso '\textit{createWard} usa il fallback generico quando l'errore non è strutturato': verifica che storeStub.\textit{setError} viene invocato con 'Operazione sui reparti non riuscita.'. Esito atteso: storeStub.\textit{setError} viene invocato con 'Operazione sui reparti non riuscita.'.}{S}
\addTestU{Contesto \textit{WardOperationsService}; caso 'in caso di errore generico usa il fallback message': verifica che storeStub.\textit{setError} viene invocato con 'Operazione sui reparti non riuscita.' e storeStub.\textit{setError} viene invocato 1 volte. Esito atteso: storeStub.\textit{setError} viene invocato con 'Operazione sui reparti non riuscita.'; storeStub.\textit{setError} viene invocato 1 volte; storeStub.\textit{setWards} non viene invocato.}{S}
\addTestU{Contesto \textit{WardOperationsService}; caso '\textit{loadWards} normalizza apartments dal dto': verifica che storeStub.\textit{setWards} viene invocato con [ \{ ...hydratedWard, apartments: [\{ id: '101', name: 'App. 101' \}], \}, ]. Esito atteso: storeStub.\textit{setWards} viene invocato con [ \{ ...hydratedWard, apartments: [\{ id: '101', name: 'App. 101' \}], \}, ].}{S}
\addTestU{Contesto \textit{WardStore}; caso 'inizializza stato vuoto': verifica che await firstValueFrom(store.wards\$) coincide con [] e il valore di await firstValueFrom(store.isLoading\$) è false. Esito atteso: await firstValueFrom(store.wards\$) coincide con []; il valore di await firstValueFrom(store.isLoading\$) è false; await firstValueFrom(store.error\$) è null.}{S}
\addTestU{Contesto \textit{WardStore}; caso '\textit{setWards} sovrascrive i wards': verifica che wards coincide con [wardA, wardB] e wards ha lunghezza 2. Esito atteso: wards coincide con [wardA, wardB]; wards ha lunghezza 2.}{S}
\addTestU{Contesto \textit{WardStore}; caso 'addWard aggiunge il reparto e resetta error/loading': verifica che wards coincide con [wardA] e wards ha lunghezza 1. Esito atteso: wards coincide con [wardA]; wards ha lunghezza 1; il valore di await firstValueFrom(store.isLoading\$) è false; await firstValueFrom(store.error\$) è null.}{S}
\addTestU{Contesto \textit{WardStore}; caso '\textit{replaceWard} aggiorna il reparto target': verifica che wards coincide con [wardA, renamed] e wards ha lunghezza 2. Esito atteso: wards coincide con [wardA, renamed]; wards ha lunghezza 2.}{S}
\addTestU{Contesto \textit{WardStore}; caso '\textit{removeWard} elimina il reparto target': verifica che wards coincide con [wardB] e wards ha lunghezza 1. Esito atteso: wards coincide con [wardB]; wards ha lunghezza 1.}{S}
\addTestU{Contesto \textit{WardStore}; caso '\textit{patchPlant} aggiorna solo l appartamento richiesto': verifica che wards ha lunghezza 2 e wards[0].apartments[0] coincide con wardA.apartments[0]. Esito atteso: wards ha lunghezza 2; wards[0].apartments[0] coincide con wardA.apartments[0]; wards[1].apartments[0] coincide con \{ id: '102', name: 'App. 102A' \}.}{S}
\addTestU{Contesto \textit{AlarmActionButtonComponent}; caso 'renderizza label e aria-label di fallback': verifica che button.textContent contiene 'GESTISCI' e il valore di button.\textit{getAttribute}('aria-label') è 'GESTISCI'. Esito atteso: button.textContent contiene 'GESTISCI'; il valore di button.\textit{getAttribute}('aria-label') è 'GESTISCI'; il valore di button.disabled è false.}{S}
\addTestU{Contesto \textit{AlarmActionButtonComponent}; caso 'usa aria-label esplicita quando presente': verifica che il valore di button.\textit{getAttribute}('aria-label') è 'Gestisci allarme Antipanico'. Esito atteso: il valore di button.\textit{getAttribute}('aria-label') è 'Gestisci allarme Antipanico'.}{S}
\addTestU{Contesto \textit{AlarmActionButtonComponent}; caso 'applica la classe corretta per il tone primary': verifica che component.buttonClass() contiene 'care-btn' e component.buttonClass() contiene 'care-btn--primary'. Esito atteso: component.buttonClass() contiene 'care-btn'; component.buttonClass() contiene 'care-btn--primary'.}{S}
\addTestU{Contesto \textit{AlarmActionButtonComponent}; caso 'applica la classe corretta per il tone danger': verifica che component.buttonClass() contiene 'care-btn' e component.buttonClass() contiene 'care-btn--danger'. Esito atteso: component.buttonClass() contiene 'care-btn'; component.buttonClass() contiene 'care-btn--danger'.}{S}
\addTestU{Contesto \textit{AlarmActionButtonComponent}; caso 'non emette click quando disabilitato': verifica che emitSpy non viene invocato e il valore di (fixture.nativeElement.querySelector('button') as HTMLButtonElement).disabled è true. Esito atteso: emitSpy non viene invocato; il valore di (fixture.nativeElement.querySelector('button') as HTMLButtonElement).disabled è true.}{S}
\addTestU{Contesto \textit{AlarmActionButtonComponent}; caso 'emette pressed al click quando abilitato': verifica che emitSpy viene invocato 1 volte. Esito atteso: emitSpy viene invocato 1 volte.}{S}
\addTestU{Contesto \textit{AlarmPriorityIndicatorComponent}; caso 'mappa la priorita rossa': verifica che il valore di component.ui().label è 'Alta' e il valore di nativeElement.querySelector('[aria-label]')?.\textit{getAttribute}('aria-label') è 'Priorita Alta'. Esito atteso: il valore di component.ui().label è 'Alta'; il valore di nativeElement.querySelector('[aria-label]')?.\textit{getAttribute}('aria-label') è 'Priorita Alta'; nativeElement.textContent contiene 'triangolo'.}{S}
\addTestU{Contesto \textit{AlarmPriorityIndicatorComponent}; caso 'mappa la priorita arancione': verifica che il valore di component.ui().label è 'Media' e fixture.nativeElement.textContent contiene '!'. Esito atteso: il valore di component.ui().label è 'Media'; fixture.nativeElement.textContent contiene '!'.}{S}
\addTestU{Contesto \textit{AlarmPriorityIndicatorComponent}; caso 'mappa la priorita verde': verifica che il valore di component.ui().label è 'Bassa' e fixture.nativeElement.textContent contiene 'pallino'. Esito atteso: il valore di component.ui().label è 'Bassa'; fixture.nativeElement.textContent contiene 'pallino'.}{S}
\addTestU{Contesto \textit{AlarmPriorityIndicatorComponent}; caso 'mappa la priorita bianca': verifica che il valore di component.ui().label è 'Informativa' e fixture.nativeElement.textContent contiene 'i'. Esito atteso: il valore di component.ui().label è 'Informativa'; fixture.nativeElement.textContent contiene 'i'.}{S}
\addTestU{Contesto \textit{AlarmTableShellComponent}; caso 'renderizza intestazione, caption e contenuto proiettato': verifica che nativeElement.textContent contiene 'Allarmi attivi' e il valore di nativeElement.querySelector('table')?.\textit{getAttribute}('aria-label') è 'Tabella test'. Esito atteso: nativeElement.textContent contiene 'Allarmi attivi'; il valore di nativeElement.querySelector('table')?.\textit{getAttribute}('aria-label') è 'Tabella test'; nativeElement.textContent contiene 'Contenuto proiettato'.}{S}
\addTestU{Contesto \textit{AlarmTableShellComponent}; caso 'mostra il pulsante create e emette \textit{createRequested} al click': verifica che \textit{createButton} non è null e fixture.componentInstance.\textit{onCreateRequested} viene invocato 1 volte. Esito atteso: \textit{createButton} non è null; fixture.componentInstance.\textit{onCreateRequested} viene invocato 1 volte.}{S}
\addTestU{Contesto \textit{AlarmToggleSwitchComponent}; caso 'mostra lo stato iniziale non attivo': verifica che il valore di button.\textit{getAttribute}('aria-pressed') è 'false' e component.\textit{trackClass}() contiene 'border-slate-500'. Esito atteso: il valore di button.\textit{getAttribute}('aria-pressed') è 'false'; component.\textit{trackClass}() contiene 'border-slate-500'; component.thumbClass() contiene 'translate-x-0.5'.}{S}
\addTestU{Contesto \textit{AlarmToggleSwitchComponent}; caso 'mostra lo stato attivo e la classe corretta': verifica che il valore di button.\textit{getAttribute}('aria-pressed') è 'true' e component.\textit{trackClass}() contiene 'border-yellow-700'. Esito atteso: il valore di button.\textit{getAttribute}('aria-pressed') è 'true'; component.\textit{trackClass}() contiene 'border-yellow-700'; component.thumbClass() contiene 'translate-x-4'.}{S}
\addTestU{Contesto \textit{AlarmToggleSwitchComponent}; caso 'emette il nuovo valore quando viene togglato': verifica che emitSpy viene invocato con true e emitSpy viene invocato 1 volte. Esito atteso: emitSpy viene invocato con true; emitSpy viene invocato 1 volte.}{S}
\addTestU{Contesto \textit{AlarmToggleSwitchComponent}; caso 'non emette quando e disabilitato': verifica che emitSpy non viene invocato e il valore di (fixture.nativeElement.querySelector('button') as HTMLButtonElement).disabled è true. Esito atteso: emitSpy non viene invocato; il valore di (fixture.nativeElement.querySelector('button') as HTMLButtonElement).disabled è true.}{S}
\addTestU{Contesto \textit{AlarmToggleSwitchComponent}; caso 'usa aria-label custom quando fornita': verifica che il valore di (fixture.nativeElement.querySelector('button') as HTMLButtonElement).\textit{getAttribute}('aria-label') è 'Attiva allarme'. Esito atteso: il valore di (fixture.nativeElement.querySelector('button') as HTMLButtonElement).\textit{getAttribute}('aria-label') è 'Attiva allarme'.}{S}
\addTestU{Contesto \textit{ConfirmDialogComponent}; caso 'renderizza i valori di default': verifica che nativeElement.textContent contiene 'Confermi questa operazione?' e nativeElement.textContent contiene 'Conferma'. Esito atteso: nativeElement.textContent contiene 'Confermi questa operazione?'; nativeElement.textContent contiene 'Conferma'; nativeElement.textContent contiene 'Annulla'.}{S}
\addTestU{Contesto \textit{ConfirmDialogComponent}; caso 'emette confirmed e cancelled ai click corretti': verifica che cancelledSpy viene invocato 1 volte e confirmedSpy viene invocato 1 volte. Esito atteso: cancelledSpy viene invocato 1 volte; confirmedSpy viene invocato 1 volte.}{S}
\addTestU{Contesto \textit{ConfirmDialogComponent}; caso 'supporta label personalizzate': verifica che nativeElement.textContent contiene 'Procedere con la cancellazione?' e nativeElement.textContent contiene 'Sì, elimina'. Esito atteso: nativeElement.textContent contiene 'Procedere con la cancellazione?'; nativeElement.textContent contiene 'Sì, elimina'; nativeElement.textContent contiene 'No, torna indietro'.}{S}
\addTestU{Contesto \textit{ModalShellComponent}; caso 'non renderizza il dialog quando open e false': verifica che fixture.nativeElement.querySelector('dialog') è null. Esito atteso: fixture.nativeElement.querySelector('dialog') è null.}{S}
\addTestU{Contesto \textit{ModalShellComponent}; caso 'renderizza il dialog e la descrizione quando open e true': verifica che dialog non è null e il valore di dialog?.\textit{getAttribute}('aria-describedby') è 'app-modal-shell-description'. Esito atteso: dialog non è null; il valore di dialog?.\textit{getAttribute}('aria-describedby') è 'app-modal-shell-description'; fixture.nativeElement.textContent contiene 'Finestra test'; fixture.nativeElement.textContent contiene 'Descrizione di test'.}{S}
\addTestU{Contesto \textit{ModalShellComponent}; caso 'non espone aria-describedby quando la descrizione e nulla': verifica che (fixture.nativeElement.querySelector('dialog') as HTMLDialogElement | null)?.\textit{getAttribute}('ari... è null. Esito atteso: (fixture.nativeElement.querySelector('dialog') as HTMLDialogElement | null)?.\textit{getAttribute}('ari... è null.}{S}
\addTestU{Contesto \textit{ModalShellComponent}; caso 'emette closed al click sul pulsante di chiusura': verifica che emitSpy viene invocato 1 volte. Esito atteso: emitSpy viene invocato 1 volte.}{S}
\addTestU{Contesto \textit{AlarmConfigFormComponent}; caso 'crea il componente': verifica che component è valorizzato. Esito atteso: component è valorizzato.}{S}
\addTestU{Contesto \textit{AlarmConfigFormComponent}; caso '\textit{buildForm} applica i validatori required ai campi richiesti': verifica che il valore di component.form.controls.plantId.invalid è true e il valore di component.form.controls.deviceId.invalid è true. Esito atteso: il valore di component.form.controls.plantId.invalid è true; il valore di component.form.controls.deviceId.invalid è true; il valore di component.form.controls.datapointId.invalid è true; il valore di component.form.controls.priority.invalid è true.}{S}
\addTestU{Contesto \textit{AlarmConfigFormComponent}; caso 'in create mode mantiene il dispositivo bloccato finché non viene selezionato un plant': verifica che il valore di component.form.controls.deviceId.disabled è true e apartmentApiStub.loadApartmentViewForPlantId viene invocato con 'plant-1'. Esito atteso: il valore di component.form.controls.deviceId.disabled è true; apartmentApiStub.loadApartmentViewForPlantId viene invocato con 'plant-1'; il valore di component.form.controls.deviceId.enabled è true; il valore di component.deviceOptions()[0]?.id è 'sensor-1'.}{S}
\addTestU{Contesto \textit{AlarmConfigFormComponent}; caso 'mostra solo datapoint leggibili dopo la selezione dispositivo': verifica che component.datapointOptions().map((datapoint) => datapoint.id) coincide con [ 'dp-readable-1', 'dp-readable-2', ] e il valore di component.datapointOptions().every((datapoint) => datapoint.readable) è true. Esito atteso: component.datapointOptions().map((datapoint) => datapoint.id) coincide con [ 'dp-readable-1', 'dp-readable-2', ]; il valore di component.datapointOptions().every((datapoint) => datapoint.readable) è true.}{S}
\addTestU{Contesto \textit{AlarmConfigFormComponent}; caso 'in create mode non mostra il campo posizione read-only': verifica che positionInput è null. Esito atteso: positionInput è null.}{S}
\addTestU{Contesto \textit{AlarmConfigFormComponent}; caso 'in edit mode mostra il campo posizione read-only valorizzato': verifica che positionInput non è null e il valore di positionInput?.value è 'Appartamento 2 - Ingresso - Sensore porta'. Esito atteso: positionInput non è null; il valore di positionInput?.value è 'Appartamento 2 - Ingresso - Sensore porta'.}{S}
\addTestU{Contesto \textit{AlarmConfigFormComponent}; caso 'in edit mode non mostra il campo datapoint nel form': verifica che datapointInput è null. Esito atteso: datapointInput è null.}{S}
\addTestU{Contesto \textit{AlarmConfigFormComponent}; caso 'in edit mode mostra thresholdOperator come campo readonly': verifica che thresholdOperatorReadonly non è null e il valore di thresholdOperatorReadonly?.value è ThresholdOperator.EQUAL\_TO. Esito atteso: thresholdOperatorReadonly non è null; il valore di thresholdOperatorReadonly?.value è ThresholdOperator.EQUAL\_TO; thresholdOperatorSelect è null.}{S}
\addTestU{Contesto \textit{AlarmConfigFormComponent}; caso 'in edit mode limita gli operatori al datapoint risolto e resetta operatori non supportati': verifica che component.thresholdOperatorOptions() coincide con [ThresholdOperator.EQUAL\_TO] e component.form.controls.thresholdOperator.value è null. Esito atteso: component.thresholdOperatorOptions() coincide con [ThresholdOperator.EQUAL\_TO]; component.form.controls.thresholdOperator.value è null.}{S}
\addTestU{Contesto \textit{AlarmConfigFormComponent}; caso 'in edit mode valida thresholdValue in base al datapoint risolto': verifica che il valore di component.form.controls.thresholdValue.errors?.['invalidEnumThreshold'] è true. Esito atteso: il valore di component.form.controls.thresholdValue.errors?.['invalidEnumThreshold'] è true.}{S}
\addTestU{Contesto \textit{AlarmConfigFormComponent}; caso '\textit{onSubmit} emette submittedForm in create mode con form valido': verifica che emitSpy viene invocato con validFormValue e emitSpy viene invocato 1 volte. Esito atteso: emitSpy viene invocato con validFormValue; emitSpy viene invocato 1 volte.}{S}
\addTestU{Contesto \textit{AlarmConfigFormComponent}; caso '\textit{onSubmit} non invia se il form e invalido': verifica che emitSpy non viene invocato. Esito atteso: emitSpy non viene invocato.}{S}
\addTestU{Contesto \textit{AlarmConfigFormComponent}; caso '\textit{onCancel} emette evento cancelled': verifica che emitSpy viene invocato 1 volte. Esito atteso: emitSpy viene invocato 1 volte.}{S}
\addTestU{Contesto \textit{AlarmConfigFormComponent}; caso 'espone opzioni enum per priorita e operatore': verifica che component.priorityOptions coincide con [ AlarmPriority.WHITE, AlarmPriority.GREEN, AlarmPriority.ORANGE, AlarmPriority.RED, ] e component.thresholdOperatorOptions() coincide con [ ThresholdOperator.GREATER\_THAN, ThresholdOperator.GREATER\_THAN\_OR\_EQUAL, ThresholdOperator.L.... Esito atteso: component.priorityOptions coincide con [ AlarmPriority.WHITE, AlarmPriority.GREEN, AlarmPriority.ORANGE, AlarmPriority.RED, ]; component.thresholdOperatorOptions() coincide con [ ThresholdOperator.GREATER\_THAN, ThresholdOperator.GREATER\_THAN\_OR\_EQUAL, ThresholdOperator.L....}{S}
\addTestU{Contesto \textit{AlarmConfigFormComponent}; caso 'limita operatori e valore soglia quando il datapoint ha enum': verifica che component.thresholdOperatorOptions() coincide con [ThresholdOperator.EQUAL\_TO] e il valore di component.form.controls.thresholdOperator.invalid è true. Esito atteso: component.thresholdOperatorOptions() coincide con [ThresholdOperator.EQUAL\_TO]; il valore di component.form.controls.thresholdOperator.invalid è true; il valore di component.form.controls.thresholdValue.invalid è true; il valore di component.form.controls.thresholdValue.errors?.['invalidEnumThreshold'] è true.}{S}
\addTestU{Contesto \textit{AlarmConfigFormComponent}; caso 'accetta valore soglia booleano quando operatore è uguale': verifica che il valore di component.form.controls.thresholdOperator.valid è true e il valore di component.form.controls.thresholdValue.valid è true. Esito atteso: il valore di component.form.controls.thresholdOperator.valid è true; il valore di component.form.controls.thresholdValue.valid è true.}{S}
\addTestU{Contesto \textit{AlarmConfigFormComponent}; caso 'carica impianti da /plant/all deduplicando e ordinando per nome': verifica che component.plants().map((plant) => plant.id) coincide con ['plant-1', 'plant-2', 'plant-3'] e component.plantsLoadError() è null. Esito atteso: component.plants().map((plant) => plant.id) coincide con ['plant-1', 'plant-2', 'plant-3']; component.plantsLoadError() è null.}{S}
\addTestU{Contesto \textit{AlarmConfigFormComponent}; caso 'se /plant/all fallisce usa fallback available + plants assegnati ai reparti': verifica che component.plants().map((plant) => plant.id) coincide con ['plant-7', 'plant-8', 'plant-9'] e component.plantsLoadError() è null. Esito atteso: component.plants().map((plant) => plant.id) coincide con ['plant-7', 'plant-8', 'plant-9']; component.plantsLoadError() è null; wardApiStub.getPlantsByWardId viene invocato 1 volte.}{S}
\addTestU{Contesto \textit{AlarmConfigFormComponent}; caso 'se il plant viene deselezionato resetta opzioni dispositivo e blocca il campo dispositivo': verifica che il valore di component.deviceOptions()[0]?.id è 'sensor-1' e il valore di component.form.controls.deviceId.enabled è true. Esito atteso: il valore di component.deviceOptions()[0]?.id è 'sensor-1'; il valore di component.form.controls.deviceId.enabled è true; component.deviceOptions() coincide con []; il valore di component.form.controls.deviceId.value è ''.}{S}
\addTestU{Contesto \textit{AlarmConfigFormComponent}; caso 'in edit mode non consente di modificare il dispositivo associato': verifica che il valore di component.form.controls.deviceId.disabled è true e component.deviceOptions() coincide con [\{ id: 'sensor-9', label: 'sensor-9', datapoints: [] \}]. Esito atteso: il valore di component.form.controls.deviceId.disabled è true; component.deviceOptions() coincide con [\{ id: 'sensor-9', label: 'sensor-9', datapoints: [] \}].}{S}
\addTestU{Contesto \textit{AlarmConfigFormComponent}; caso 'imposta errore quando il caricamento dispositivi fallisce': verifica che il valore di component.devicesLoadError() è 'Errore durante il caricamento dei dispositivi.' e il valore di component.form.controls.deviceId.disabled è true. Esito atteso: il valore di component.devicesLoadError() è 'Errore durante il caricamento dei dispositivi.'; il valore di component.form.controls.deviceId.disabled è true.}{S}
\addTestU{Contesto \textit{AlarmConfigPageComponent}; caso 'crea il componente': verifica che component è valorizzato. Esito atteso: component è valorizzato.}{S}
\addTestU{Contesto \textit{AlarmConfigPageComponent}; caso '\textit{ngOnInit} carica gli allarmi e aggiorna le righe tabellari': verifica che stateServiceStub.\textit{loadAlarmRules} viene invocato 1 volte e il valore di component.rows().length è 1. Esito atteso: stateServiceStub.\textit{loadAlarmRules} viene invocato 1 volte; il valore di component.rows().length è 1; il valore di component.error() è 'Errore test'.}{S}
\addTestU{Contesto \textit{AlarmConfigPageComponent}; caso '\textit{onCreateNew} apre la modale in create mode': verifica che il valore di component.isModalOpen() è true e component.editingRule() è null. Esito atteso: il valore di component.isModalOpen() è true; component.editingRule() è null; il valore di component.modalTitle() è 'Creazione allarme'.}{S}
\addTestU{Contesto \textit{AlarmConfigPageComponent}; caso '\textit{onEdit} apre la modale in edit mode con regola selezionata': verifica che il valore di component.isModalOpen() è true e il valore di component.editingRule()?.id è 'alarm-1'. Esito atteso: il valore di component.isModalOpen() è true; il valore di component.editingRule()?.id è 'alarm-1'; il valore di component.modalTitle() è 'Modifica allarme'.}{S}
\addTestU{Contesto \textit{AlarmConfigPageComponent}; caso '\textit{onEdit} ignora id non trovato': verifica che il valore di component.isModalOpen() è false e component.editingRule() è null. Esito atteso: il valore di component.isModalOpen() è false; component.editingRule() è null.}{S}
\addTestU{Contesto \textit{AlarmConfigPageComponent}; caso '\textit{onFormSubmitted} in create mode invoca create e chiude modale': verifica che stateServiceStub.\textit{createAlarmRule} viene invocato con formValue e il valore di component.isModalOpen() è false. Esito atteso: stateServiceStub.\textit{createAlarmRule} viene invocato con formValue; il valore di component.isModalOpen() è false.}{S}
\addTestU{Contesto \textit{AlarmConfigPageComponent}; caso '\textit{onFormSubmitted} in create mode mantiene aperta la modale se non arriva alcun esito': verifica che stateServiceStub.\textit{createAlarmRule} viene invocato con formValue e il valore di component.isModalOpen() è true. Esito atteso: stateServiceStub.\textit{createAlarmRule} viene invocato con formValue; il valore di component.isModalOpen() è true; component.editingRule() è null.}{S}
\addTestU{Contesto \textit{AlarmConfigPageComponent}; caso '\textit{onFormSubmitted} in edit mode invoca update e chiude modale': verifica che stateServiceStub.\textit{updateAlarmRule} viene invocato con 'alarm-1' e il valore di component.isModalOpen() è false. Esito atteso: stateServiceStub.\textit{updateAlarmRule} viene invocato con 'alarm-1'; il valore di component.isModalOpen() è false.}{S}
\addTestU{Contesto \textit{AlarmConfigPageComponent}; caso '\textit{onFormSubmitted} in edit mode mantiene aperta la modale se non arriva alcun esito': verifica che il valore di component.isModalOpen() è true e il valore di component.editingRule()?.id è 'alarm-1'. Esito atteso: il valore di component.isModalOpen() è true; il valore di component.editingRule()?.id è 'alarm-1'.}{S}
\addTestU{Contesto \textit{AlarmConfigPageComponent}; caso '\textit{onToggleEnabled} invoca \textit{toggleEnabled} con lo stato richiesto': verifica che stateServiceStub.\textit{toggleEnabled} viene invocato con 'alarm-1'. Esito atteso: stateServiceStub.\textit{toggleEnabled} viene invocato con 'alarm-1'.}{S}
\addTestU{Contesto \textit{AlarmConfigPageComponent}; caso '\textit{onToggleEnabled} blocca i click ripetuti finché la richiesta è in corso': verifica che il valore di component.pendingToggleRuleId() è 'alarm-1' e stateServiceStub.\textit{toggleEnabled} viene invocato 1 volte. Esito atteso: il valore di component.pendingToggleRuleId() è 'alarm-1'; stateServiceStub.\textit{toggleEnabled} viene invocato 1 volte; component.pendingToggleRuleId() è null; stateServiceStub.\textit{toggleEnabled} viene invocato 2 volte.}{S}
\addTestU{Contesto \textit{AlarmConfigPageComponent}; caso '\textit{onDelete} apre la modale di conferma senza chiamare subito la delete': verifica che component.pendingDelete() coincide con \{ id: 'alarm-1', name: 'Temperatura alta' \} e stateServiceStub.\textit{deleteAlarmRule} non viene invocato. Esito atteso: component.pendingDelete() coincide con \{ id: 'alarm-1', name: 'Temperatura alta' \}; stateServiceStub.\textit{deleteAlarmRule} non viene invocato.}{S}
\addTestU{Contesto \textit{AlarmConfigPageComponent}; caso '\textit{onDeleteConfirmed} invoca \textit{deleteAlarmRule} con id e chiude la modale di conferma': verifica che stateServiceStub.\textit{deleteAlarmRule} viene invocato con 'alarm-1' e component.pendingDelete() è null. Esito atteso: stateServiceStub.\textit{deleteAlarmRule} viene invocato con 'alarm-1'; component.pendingDelete() è null.}{S}
\addTestU{Contesto \textit{AlarmConfigPageComponent}; caso '\textit{onDeleteCancelled} chiude la modale di conferma senza eliminare': verifica che component.pendingDelete() è null e stateServiceStub.\textit{deleteAlarmRule} non viene invocato. Esito atteso: component.pendingDelete() è null; stateServiceStub.\textit{deleteAlarmRule} non viene invocato.}{S}
\addTestU{Contesto \textit{AlarmConfigPageComponent}; caso 'renderizza la modale di conferma cancellazione quando richiesta': verifica che confirmDialog non è null e content contiene 'Confermi l\textbackslash\{\}'eliminazione della soglia "Temperatura alta"?'. Esito atteso: confirmDialog non è null; content contiene 'Confermi l\textbackslash\{\}'eliminazione della soglia "Temperatura alta"?'.}{S}
\addTestU{Contesto \textit{AlarmConfigPageComponent}; caso '\textit{onModalClosed} chiude modale e resetta editingRule': verifica che il valore di component.isModalOpen() è false e component.editingRule() è null. Esito atteso: il valore di component.isModalOpen() è false; component.editingRule() è null.}{S}
\addTestU{Contesto \textit{AlarmConfigPageComponent}; caso '\textit{onCreateNew} resetta l eventuale regola in edit e apre la modale in create': verifica che il valore di component.isModalOpen() è true e component.editingRule() è null. Esito atteso: il valore di component.isModalOpen() è true; component.editingRule() è null.}{S}
\addTestU{Contesto \textit{AlarmConfigPageComponent}; caso 'renderizza la tabella con le righe e i pulsanti azione': verifica che il valore di tableRows.length è 1 e content contiene 'Temperatura alta'. Esito atteso: il valore di tableRows.length è 1; content contiene 'Temperatura alta'; content contiene 'Posizione'; content contiene 'Appartamento 1 - Soggiorno - Sensore porta'.}{S}
\addTestU{Contesto \textit{AlarmConfigPageComponent}; caso 'renderizza messaggio vuoto quando non ci sono regole': verifica che content contiene 'Nessuna regola configurata.'. Esito atteso: content contiene 'Nessuna regola configurata.'.}{S}
\addTestU{Contesto \textit{AlarmConfigPageComponent}; caso 'renderizza il form in modale quando la modale e aperta': verifica che modalForm non è null. Esito atteso: modalForm non è null.}{S}
\addTestU{Contesto \textit{AlarmItemComponent}; caso 'should create': verifica che component è valorizzato. Esito atteso: component è valorizzato.}{S}
\addTestU{Contesto \textit{AlarmItemComponent}; caso 'priorityUi mappa correttamente label per tutti i livelli': verifica che il valore di component.vm().priorityLabel è 'Alta' e il valore di component.vm().priorityLabel è 'Media'. Esito atteso: il valore di component.vm().priorityLabel è 'Alta'; il valore di component.vm().priorityLabel è 'Media'; il valore di component.vm().priorityLabel è 'Bassa'; il valore di component.vm().priorityLabel è 'Informativa'.}{S}
\addTestU{Contesto \textit{AlarmItemComponent}; caso '\textit{onResolveClick} emette alarm.id': verifica che emitSpy viene invocato con 'active-1' e emitSpy viene invocato 1 volte. Esito atteso: emitSpy viene invocato con 'active-1'; emitSpy viene invocato 1 volte.}{S}
\addTestU{Contesto \textit{AlarmItemComponent}; caso 'renderizza nome allarme, priorita label e tempo trasformato dalla pipe mockata': verifica che nativeElement.querySelector('.alarm-item\_\_title')?.textContent contiene 'Allarme antipanico in "Camera 101"' e nativeElement.querySelector('.alarm-item\_\_priority')?.textContent contiene 'Alta'. Esito atteso: nativeElement.querySelector('.alarm-item\_\_title')?.textContent contiene 'Allarme antipanico in "Camera 101"'; nativeElement.querySelector('.alarm-item\_\_priority')?.textContent contiene 'Alta'; metaValues.item(1)?.textContent contiene 'mock-elapsed:2026-03-24T10:00:00.000Z'; \textit{resolveButton}?.\textit{getAttribute}('aria-label') contiene 'Risolvi allarme Allarme antipanico'.}{S}
\addTestU{Contesto \textit{AlarmItemComponent}; caso 'quando isResolving e true disabilita il bottone e mostra lo stato di avanzamento': verifica che button non è null e il valore di button?.hasAttribute('disabled') è true. Esito atteso: button non è null; il valore di button?.hasAttribute('disabled') è true; button?.textContent contiene 'Risoluzione...'; nativeElement.querySelector('.alarm-item\_\_status')?.textContent contiene 'Operazione in corso'.}{S}
\addTestU{Contesto \textit{AlarmItemComponent}; caso 'quando isResolving e false abilita il bottone e non mostra lo stato di avanzamento': verifica che button non è null e il valore di button?.hasAttribute('disabled') è false. Esito atteso: button non è null; il valore di button?.hasAttribute('disabled') è false; button?.textContent contiene 'Risolvi'; nativeElement.querySelector('.alarm-item\_\_status') è null.}{S}
\addTestU{Contesto \textit{AlarmItemComponent}; caso 'click sul bottone invia evento resolve con alarm.id': verifica che emitSpy viene invocato con 'active-1' e emitSpy viene invocato 1 volte. Esito atteso: emitSpy viene invocato con 'active-1'; emitSpy viene invocato 1 volte.}{S}
\addTestU{Contesto \textit{AlarmItemComponent}; caso 'usa fallback per nome e posizione quando i campi sono vuoti o solo spazi': verifica che il valore di component.vm().alarmTitle è 'senza nome in "posizione sconosciuta"' e il valore di component.vm().articleAriaLabel è 'Allarme senza nome'. Esito atteso: il valore di component.vm().alarmTitle è 'senza nome in "posizione sconosciuta"'; il valore di component.vm().articleAriaLabel è 'Allarme senza nome'; il valore di component.vm().\textit{resolveButtonAriaLabel} è 'Risolvi allarme senza nome'.}{S}
\addTestU{Contesto \textit{AlarmItemComponent}; caso 'quando e in resolving aggiorna aria-label del bottone in modo specifico': verifica che il valore di component.vm().\textit{resolveButtonText} è 'Risoluzione...' e button?.\textit{getAttribute}('aria-label') contiene 'Risoluzione in corso per Allarme antipanico'. Esito atteso: il valore di component.vm().\textit{resolveButtonText} è 'Risoluzione...'; button?.\textit{getAttribute}('aria-label') contiene 'Risoluzione in corso per Allarme antipanico'.}{S}
\addTestU{Contesto \textit{AlarmItemComponent}; caso 'renderizza nomi con caratteri speciali come testo senza errori runtime': verifica che title contiene '<script>alert("xss")</script> in "Sala monitoraggio"'. Esito atteso: title contiene '<script>alert("xss")</script> in "Sala monitoraggio"'.}{S}
\addTestU{Contesto \textit{AlarmPageManagementComponent}; caso 'should create': verifica che component è valorizzato. Esito atteso: component è valorizzato.}{S}
\addTestU{Contesto \textit{AlarmPageManagementComponent}; caso '\textit{onResolve} delega a facade.\textit{resolveAlarm}': verifica che alarmManagementStub.\textit{resolveAlarm} viene invocato con 'active-2' e alarmManagementStub.\textit{resolveAlarm} viene invocato 1 volte. Esito atteso: alarmManagementStub.\textit{resolveAlarm} viene invocato con 'active-2'; alarmManagementStub.\textit{resolveAlarm} viene invocato 1 volte.}{S}
\addTestU{Contesto \textit{AlarmPageManagementComponent}; caso 'renderizza stato lista vuota quando non ci sono allarmi': verifica che nativeElement.querySelector('.alarm-management\_\_empty')?.textContent contiene 'Nessun allarme attivo al momento.' e il valore di nativeElement.querySelectorAll('tbody tr').length è 0. Esito atteso: nativeElement.querySelector('.alarm-management\_\_empty')?.textContent contiene 'Nessun allarme attivo al momento.'; il valore di nativeElement.querySelectorAll('tbody tr').length è 0.}{S}
\addTestU{Contesto \textit{AlarmPageManagementComponent}; caso 'renderizza errore e stato resolving quando presenti nel vm': verifica che nativeElement.querySelector('.alarm-management\_\_error')?.textContent contiene 'Errore durante la risoluzione' e nativeElement.querySelector('.alarm-management\_\_status')?.textContent contiene 'Risoluzione allarme in corso...'. Esito atteso: nativeElement.querySelector('.alarm-management\_\_error')?.textContent contiene 'Errore durante la risoluzione'; nativeElement.querySelector('.alarm-management\_\_status')?.textContent contiene 'Risoluzione allarme in corso...'.}{S}
\addTestU{Contesto \textit{AlarmPageManagementComponent}; caso 'renderizza tabella con una riga per ogni allarme': verifica che il valore di rows.length è 2 e nativeElement.textContent contiene 'Priorità'. Esito atteso: il valore di rows.length è 2; nativeElement.textContent contiene 'Priorità'; nativeElement.textContent contiene 'Corridoio Nord'; nativeElement.textContent contiene 'mock-elapsed:2026-03-24T10:00:00.000Z'.}{S}
\addTestU{Contesto \textit{AlarmPageManagementComponent}; caso 'nasconde la colonna gestore per operatore sanitario': verifica che nativeElement.textContent non contiene 'Gestore' e nativeElement.textContent non contiene 'oss\_1'. Esito atteso: nativeElement.textContent non contiene 'Gestore'; nativeElement.textContent non contiene 'oss\_1'.}{S}
\addTestU{Contesto \textit{AlarmPageManagementComponent}; caso 'disabilita solo il bottone associato al resolvingId corrente': verifica che il valore di firstButton?.disabled è true e il valore di secondButton?.disabled è false. Esito atteso: il valore di firstButton?.disabled è true; il valore di secondButton?.disabled è false.}{S}
\addTestU{Contesto \textit{AlarmPageManagementComponent}; caso 'renderizza paginazione coerente con vm e delega i click ai metodi del service': verifica che nativeElement.querySelector('.care-pagination\_\_status')?.textContent contiene 'Pagina 3' e il valore di previousButton?.disabled è false. Esito atteso: nativeElement.querySelector('.care-pagination\_\_status')?.textContent contiene 'Pagina 3'; il valore di previousButton?.disabled è false; il valore di nextButton?.disabled è true; alarmManagementStub.previousPage viene invocato 1 volte.}{S}
\addTestU{Contesto \textit{AlarmPageManagementComponent}; caso 'click su GESTISCI propaga resolve verso facade': verifica che alarmManagementStub.\textit{resolveAlarm} viene invocato con 'active-1' e alarmManagementStub.\textit{resolveAlarm} viene invocato 1 volte. Esito atteso: alarmManagementStub.\textit{resolveAlarm} viene invocato con 'active-1'; alarmManagementStub.\textit{resolveAlarm} viene invocato 1 volte.}{S}
\addTestU{Contesto \textit{AlarmPageManagementComponent}; caso 'mantiene la riga visibile se l'allarme e già gestito e disabilita l azione': verifica che il valore di rows.length è 1 e nativeElement.textContent contiene 'Non da gestire'. Esito atteso: il valore di rows.length è 1; nativeElement.textContent contiene 'Non da gestire'; managedButton non è null; il valore di managedButton?.disabled è true.}{S}
\addTestU{Contesto \textit{AlarmPageManagementComponent}; caso '\textit{onNextPage} e \textit{onPreviousPage} delegano al service': verifica che alarmManagementStub.nextPage viene invocato 1 volte e alarmManagementStub.previousPage viene invocato 1 volte. Esito atteso: alarmManagementStub.nextPage viene invocato 1 volte; alarmManagementStub.previousPage viene invocato 1 volte.}{S}
\addTestU{Contesto \textit{AlarmPageManagementComponent}; caso 'quando arriva \textit{refreshRequested} reinizializza la pagina': verifica che alarmManagementStub.initialize viene invocato 1 volte e alarmManagementStub.initialize viene invocato 2 volte. Esito atteso: alarmManagementStub.initialize viene invocato 1 volte; alarmManagementStub.initialize viene invocato 2 volte.}{S}
\addTestU{Contesto \textit{AlarmPageManagementComponent}; caso 'dopo destroy non risponde più agli eventi di refresh': verifica che alarmManagementStub.initialize viene invocato 1 volte. Esito atteso: alarmManagementStub.initialize viene invocato 1 volte.}{S}
\addTestU{Contesto \textit{AlarmsSentResolvedChartComponent}; caso 'espone metadati chart consistenti': verifica che il valore di component.chartType è 'bar' e component.description.length è maggiore di 0. Esito atteso: il valore di component.chartType è 'bar'; component.description.length è maggiore di 0.}{S}
\addTestU{Contesto \textit{AlarmsSentResolvedChartComponent} - con suggerimenti attivi; caso 'dovrebbe restringere il grafico al 70\% su schermi grandi': verifica che il contenitore del grafico applica la classe responsive lg:w-[70\%] quando i suggerimenti sono attivi. Esito atteso: l'elemento [data-test="chart-container"] include la classe lg:w-[70\%].}{S}
\addTestU{Contesto \textit{AlarmsSentResolvedChartComponent} - con suggerimenti attivi; caso 'dovrebbe mostrare la colonna dei suggerimenti': verifica che la sidebar dei suggerimenti viene renderizzata nel DOM quando la modalita suggerimenti è attiva. Esito atteso: l'elemento [data-test="suggestions-sidebar"] è presente e non risulta null.}{S}
\addTestU{Contesto \textit{AlarmsSentResolvedChartComponent} - con suggerimenti attivi; caso 'dovrebbe renderizzare il numero corretto di messaggi': verifica che la lista dei suggerimenti mostra esattamente il numero di messaggi previsto dai dati di input. Esito atteso: gli elementi [data-test="suggestion-item"] sono esattamente 2, in linea con i messaggi forniti.}{S}
\addTestU{Contesto \textit{AlarmsSentResolvedChartComponent} - senza suggerimenti; caso 'dovrebbe espandere il grafico al 100\% della larghezza': verifica che senza suggerimenti il contenitore del grafico occupa tutta la larghezza disponibile. Esito atteso: è presente un contenitore grafico con classe .w-full, quindi la vista occupa il 100\% della larghezza.}{S}
\addTestU{Contesto \textit{AlarmsSentResolvedChartComponent} - senza suggerimenti; caso 'NON dovrebbe esistere la colonna dei suggerimenti nel DOM': verifica che senza suggerimenti la colonna laterale non viene montata nel DOM. Esito atteso: la colonna suggerimenti (.lg:w-[30\%]) è assente e la query restituisce null.}{S}
\addTestU{Contesto \textit{AlarmFrequencyChartComponent}; caso 'espone metadati chart consistenti': verifica che il valore di component.chartType è 'bar' e component.description.length è maggiore di 0. Esito atteso: il valore di component.chartType è 'bar'; component.description.length è maggiore di 0.}{S}
\addTestU{Contesto \textit{AlarmFrequencyChartComponent} - con suggerimenti attivi; caso 'dovrebbe restringere il grafico al 70\% su schermi grandi': verifica che il contenitore del grafico applica la classe responsive lg:w-[70\%] quando i suggerimenti sono attivi. Esito atteso: l'elemento [data-test="chart-container"] include la classe lg:w-[70\%].}{S}
\addTestU{Contesto \textit{AlarmFrequencyChartComponent} - con suggerimenti attivi; caso 'dovrebbe mostrare la colonna dei suggerimenti': verifica che la sidebar dei suggerimenti viene renderizzata nel DOM quando la modalita suggerimenti è attiva. Esito atteso: l'elemento [data-test="suggestions-sidebar"] è presente e non risulta null.}{S}
\addTestU{Contesto \textit{AlarmFrequencyChartComponent} - con suggerimenti attivi; caso 'dovrebbe renderizzare il numero corretto di messaggi': verifica che la lista dei suggerimenti mostra esattamente il numero di messaggi previsto dai dati di input. Esito atteso: gli elementi [data-test="suggestion-item"] sono esattamente 2, in linea con i messaggi forniti.}{S}
\addTestU{Contesto \textit{AlarmFrequencyChartComponent} - senza suggerimenti; caso 'dovrebbe espandere il grafico al 100\% della larghezza': verifica che senza suggerimenti il contenitore del grafico occupa tutta la larghezza disponibile. Esito atteso: è presente un contenitore grafico con classe .w-full, quindi la vista occupa il 100\% della larghezza.}{S}
\addTestU{Contesto \textit{AlarmFrequencyChartComponent} - senza suggerimenti; caso 'NON dovrebbe esistere la colonna dei suggerimenti nel DOM': verifica che senza suggerimenti la colonna laterale non viene montata nel DOM. Esito atteso: la colonna suggerimenti (.lg:w-[30\%]) è assente e la query restituisce null.}{S}
\addTestU{Contesto \textit{EnergyConsumptionChartComponent}; caso 'espone metadati chart consistenti': verifica che il valore di component.chartType è 'line' e component.description.length è maggiore di 0. Esito atteso: il valore di component.chartType è 'line'; component.description.length è maggiore di 0.}{S}
\addTestU{Contesto \textit{EnergyConsumptionChartComponent} - con suggerimenti attivi; caso 'dovrebbe restringere il grafico al 70\% su schermi grandi': verifica che il contenitore del grafico applica la classe responsive lg:w-[70\%] quando i suggerimenti sono attivi. Esito atteso: l'elemento [data-test="chart-container"] include la classe lg:w-[70\%].}{S}
\addTestU{Contesto \textit{EnergyConsumptionChartComponent} - con suggerimenti attivi; caso 'dovrebbe mostrare la colonna dei suggerimenti': verifica che la sidebar dei suggerimenti viene renderizzata nel DOM quando la modalita suggerimenti è attiva. Esito atteso: l'elemento [data-test="suggestions-sidebar"] è presente e non risulta null.}{S}
\addTestU{Contesto \textit{EnergyConsumptionChartComponent} - con suggerimenti attivi; caso 'dovrebbe renderizzare il numero corretto di messaggi': verifica che la lista dei suggerimenti mostra esattamente il numero di messaggi previsto dai dati di input. Esito atteso: gli elementi [data-test="suggestion-item"] sono esattamente 2, in linea con i messaggi forniti.}{S}
\addTestU{Contesto \textit{EnergyConsumptionChartComponent} - senza suggerimenti; caso 'dovrebbe espandere il grafico al 100\% della larghezza': verifica che senza suggerimenti il contenitore del grafico occupa tutta la larghezza disponibile. Esito atteso: è presente un contenitore grafico con classe .w-full, quindi la vista occupa il 100\% della larghezza.}{S}
\addTestU{Contesto \textit{EnergyConsumptionChartComponent} - senza suggerimenti; caso 'NON dovrebbe esistere la colonna dei suggerimenti nel DOM': verifica che senza suggerimenti la colonna laterale non viene montata nel DOM. Esito atteso: la colonna suggerimenti (.lg:w-[30\%]) è assente e la query restituisce null.}{S}
\addTestU{Contesto \textit{FallFrequencyChartComponent}; caso 'espone metadati chart consistenti': verifica che il valore di component.chartType è 'bar' e component.description.length è maggiore di 0. Esito atteso: il valore di component.chartType è 'bar'; component.description.length è maggiore di 0.}{S}
\addTestU{Contesto \textit{FallFrequencyChartComponent} - con suggerimenti attivi; caso 'dovrebbe restringere il grafico al 70\% su schermi grandi': verifica che il contenitore del grafico applica la classe responsive lg:w-[70\%] quando i suggerimenti sono attivi. Esito atteso: l'elemento [data-test="chart-container"] include la classe lg:w-[70\%].}{S}
\addTestU{Contesto \textit{FallFrequencyChartComponent} - con suggerimenti attivi; caso 'dovrebbe mostrare la colonna dei suggerimenti': verifica che la sidebar dei suggerimenti viene renderizzata nel DOM quando la modalita suggerimenti è attiva. Esito atteso: l'elemento [data-test="suggestions-sidebar"] è presente e non risulta null.}{S}
\addTestU{Contesto \textit{FallFrequencyChartComponent} - con suggerimenti attivi; caso 'dovrebbe renderizzare il numero corretto di messaggi': verifica che la lista dei suggerimenti mostra esattamente il numero di messaggi previsto dai dati di input. Esito atteso: gli elementi [data-test="suggestion-item"] sono esattamente 2, in linea con i messaggi forniti.}{S}
\addTestU{Contesto \textit{FallFrequencyChartComponent} - senza suggerimenti; caso 'dovrebbe espandere il grafico al 100\% della larghezza': verifica che senza suggerimenti il contenitore del grafico occupa tutta la larghezza disponibile. Esito atteso: è presente un contenitore grafico con classe .w-full, quindi la vista occupa il 100\% della larghezza.}{S}
\addTestU{Contesto \textit{FallFrequencyChartComponent} - senza suggerimenti; caso 'NON dovrebbe esistere la colonna dei suggerimenti nel DOM': verifica che senza suggerimenti la colonna laterale non viene montata nel DOM. Esito atteso: la colonna suggerimenti (.lg:w-[30\%]) è assente e la query restituisce null.}{S}
\addTestU{Contesto \textit{PlantAnomaliesChartComponent}; caso 'espone metadati chart consistenti': verifica che il valore di component.chartType è 'bar' e component.description.length è maggiore di 0. Esito atteso: il valore di component.chartType è 'bar'; component.description.length è maggiore di 0.}{S}
\addTestU{Contesto \textit{PlantAnomaliesChartComponent} - con suggerimenti attivi; caso 'dovrebbe restringere il grafico al 70\% su schermi grandi': verifica che il contenitore del grafico applica la classe responsive lg:w-[70\%] quando i suggerimenti sono attivi. Esito atteso: l'elemento [data-test="chart-container"] include la classe lg:w-[70\%].}{S}
\addTestU{Contesto \textit{PlantAnomaliesChartComponent} - con suggerimenti attivi; caso 'dovrebbe mostrare la colonna dei suggerimenti': verifica che la sidebar dei suggerimenti viene renderizzata nel DOM quando la modalita suggerimenti è attiva. Esito atteso: l'elemento [data-test="suggestions-sidebar"] è presente e non risulta null.}{S}
\addTestU{Contesto \textit{PlantAnomaliesChartComponent} - con suggerimenti attivi; caso 'dovrebbe renderizzare il numero corretto di messaggi': verifica che la lista dei suggerimenti mostra esattamente il numero di messaggi previsto dai dati di input. Esito atteso: gli elementi [data-test="suggestion-item"] sono esattamente 2, in linea con i messaggi forniti.}{S}
\addTestU{Contesto \textit{PlantAnomaliesChartComponent} - senza suggerimenti; caso 'dovrebbe espandere il grafico al 100\% della larghezza': verifica che senza suggerimenti il contenitore del grafico occupa tutta la larghezza disponibile. Esito atteso: è presente un contenitore grafico con classe .w-full, quindi la vista occupa il 100\% della larghezza.}{S}
\addTestU{Contesto \textit{PlantAnomaliesChartComponent} - senza suggerimenti; caso 'NON dovrebbe esistere la colonna dei suggerimenti nel DOM': verifica che senza suggerimenti la colonna laterale non viene montata nel DOM. Esito atteso: la colonna suggerimenti (.lg:w-[30\%]) è assente e la query restituisce null.}{S}
\addTestU{Contesto \textit{PresenceDetectionChartComponent}; caso 'espone metadati chart consistenti': verifica che il valore di component.chartType è 'bar' e component.description.length è maggiore di 0. Esito atteso: il valore di component.chartType è 'bar'; component.description.length è maggiore di 0.}{S}
\addTestU{Contesto \textit{PresenceDetectionChartComponent} - con suggerimenti attivi; caso 'dovrebbe restringere il grafico al 70\% su schermi grandi': verifica che il contenitore del grafico applica la classe responsive lg:w-[70\%] quando i suggerimenti sono attivi. Esito atteso: l'elemento [data-test="chart-container"] include la classe lg:w-[70\%].}{S}
\addTestU{Contesto \textit{PresenceDetectionChartComponent} - con suggerimenti attivi; caso 'dovrebbe mostrare la colonna dei suggerimenti': verifica che la sidebar dei suggerimenti viene renderizzata nel DOM quando la modalita suggerimenti è attiva. Esito atteso: l'elemento [data-test="suggestions-sidebar"] è presente e non risulta null.}{S}
\addTestU{Contesto \textit{PresenceDetectionChartComponent} - con suggerimenti attivi; caso 'dovrebbe renderizzare il numero corretto di messaggi': verifica che la lista dei suggerimenti mostra esattamente il numero di messaggi previsto dai dati di input. Esito atteso: gli elementi [data-test="suggestion-item"] sono esattamente 2, in linea con i messaggi forniti.}{S}
\addTestU{Contesto \textit{PresenceDetectionChartComponent} - senza suggerimenti; caso 'dovrebbe espandere il grafico al 100\% della larghezza': verifica che senza suggerimenti il contenitore del grafico occupa tutta la larghezza disponibile. Esito atteso: è presente un contenitore grafico con classe .w-full, quindi la vista occupa il 100\% della larghezza.}{S}
\addTestU{Contesto \textit{PresenceDetectionChartComponent} - senza suggerimenti; caso 'NON dovrebbe esistere la colonna dei suggerimenti nel DOM': verifica che senza suggerimenti la colonna laterale non viene montata nel DOM. Esito atteso: la colonna suggerimenti (.lg:w-[30\%]) è assente e la query restituisce null.}{S}
\addTestU{Contesto \textit{ProlongedPresenceChartComponent}; caso 'espone metadati chart consistenti': verifica che il valore di component.chartType è 'bar' e component.description.length è maggiore di 0. Esito atteso: il valore di component.chartType è 'bar'; component.description.length è maggiore di 0.}{S}
\addTestU{Contesto \textit{ProlongedPresenceChartComponent} - con suggerimenti attivi; caso 'dovrebbe restringere il grafico al 70\% su schermi grandi': verifica che il contenitore del grafico applica la classe responsive lg:w-[70\%] quando i suggerimenti sono attivi. Esito atteso: l'elemento [data-test="chart-container"] include la classe lg:w-[70\%].}{S}
\addTestU{Contesto \textit{ProlongedPresenceChartComponent} - con suggerimenti attivi; caso 'dovrebbe mostrare la colonna dei suggerimenti': verifica che la sidebar dei suggerimenti viene renderizzata nel DOM quando la modalita suggerimenti è attiva. Esito atteso: l'elemento [data-test="suggestions-sidebar"] è presente e non risulta null.}{S}
\addTestU{Contesto \textit{ProlongedPresenceChartComponent} - con suggerimenti attivi; caso 'dovrebbe renderizzare il numero corretto di messaggi': verifica che la lista dei suggerimenti mostra esattamente il numero di messaggi previsto dai dati di input. Esito atteso: gli elementi [data-test="suggestion-item"] sono esattamente 2, in linea con i messaggi forniti.}{S}
\addTestU{Contesto \textit{ProlongedPresenceChartComponent} - senza suggerimenti; caso 'dovrebbe espandere il grafico al 100\% della larghezza': verifica che senza suggerimenti il contenitore del grafico occupa tutta la larghezza disponibile. Esito atteso: è presente un contenitore grafico con classe .w-full, quindi la vista occupa il 100\% della larghezza.}{S}
\addTestU{Contesto \textit{ProlongedPresenceChartComponent} - senza suggerimenti; caso 'NON dovrebbe esistere la colonna dei suggerimenti nel DOM': verifica che senza suggerimenti la colonna laterale non viene montata nel DOM. Esito atteso: la colonna suggerimenti (.lg:w-[30\%]) è assente e la query restituisce null.}{S}
\addTestU{Contesto \textit{TemperatureVariationsChartComponent}; caso 'espone metadati chart consistenti': verifica che il valore di component.chartType è 'line' e component.description.length è maggiore di 0. Esito atteso: il valore di component.chartType è 'line'; component.description.length è maggiore di 0.}{S}
\addTestU{Contesto \textit{TemperatureVariationsChartComponent} - con suggerimenti attivi; caso 'dovrebbe restringere il grafico al 70\% su schermi grandi': verifica che il contenitore del grafico applica la classe responsive lg:w-[70\%] quando i suggerimenti sono attivi. Esito atteso: l'elemento [data-test="chart-container"] include la classe lg:w-[70\%].}{S}
\addTestU{Contesto \textit{TemperatureVariationsChartComponent} - con suggerimenti attivi; caso 'dovrebbe mostrare la colonna dei suggerimenti': verifica che la sidebar dei suggerimenti viene renderizzata nel DOM quando la modalita suggerimenti è attiva. Esito atteso: l'elemento [data-test="suggestions-sidebar"] è presente e non risulta null.}{S}
\addTestU{Contesto \textit{TemperatureVariationsChartComponent} - con suggerimenti attivi; caso 'dovrebbe renderizzare il numero corretto di messaggi': verifica che la lista dei suggerimenti mostra esattamente il numero di messaggi previsto dai dati di input. Esito atteso: gli elementi [data-test="suggestion-item"] sono esattamente 2, in linea con i messaggi forniti.}{S}
\addTestU{Contesto \textit{TemperatureVariationsChartComponent} - senza suggerimenti; caso 'dovrebbe espandere il grafico al 100\% della larghezza': verifica che senza suggerimenti il contenitore del grafico occupa tutta la larghezza disponibile. Esito atteso: è presente un contenitore grafico con classe .w-full, quindi la vista occupa il 100\% della larghezza.}{S}
\addTestU{Contesto \textit{TemperatureVariationsChartComponent} - senza suggerimenti; caso 'NON dovrebbe esistere la colonna dei suggerimenti nel DOM': verifica che senza suggerimenti la colonna laterale non viene montata nel DOM. Esito atteso: la colonna suggerimenti (.lg:w-[30\%]) è assente e la query restituisce null.}{S}
\addTestU{Contesto \textit{ApartmentMonitorComponent}; caso 'crea il componente e sincronizza activeApartmentId dal feed apartment\$': verifica che component è valorizzato e il valore di component.activeApartmentId è 'ap-1'. Esito atteso: component è valorizzato; il valore di component.activeApartmentId è 'ap-1'; apartmentApiStub.\textit{getCurrentApartment} viene invocato.}{S}
\addTestU{Contesto \textit{ApartmentMonitorComponent}; caso '\textit{onApartmentSelected} ignora id vuoto o uguale': verifica che apartmentApiStub.setActivePlantId non viene invocato. Esito atteso: apartmentApiStub.setActivePlantId non viene invocato.}{S}
\addTestU{Contesto \textit{ApartmentMonitorComponent}; caso '\textit{onApartmentSelected} cambia appartamento, resetta errore e triggera refresh': verifica che apartmentApiStub.setActivePlantId viene invocato con 'ap-2' e il valore di component.error è ''. Esito atteso: apartmentApiStub.setActivePlantId viene invocato con 'ap-2'; il valore di component.error è ''.}{S}
\addTestU{Contesto \textit{ApartmentMonitorComponent}; caso 'refresh realtime azzera errore e ricarica apartment\$': verifica che il valore di component.error è '' e apartmentApiStub.\textit{getCurrentApartment}.mock.calls.length è maggiore di callsBefore. Esito atteso: il valore di component.error è ''; apartmentApiStub.\textit{getCurrentApartment}.mock.calls.length è maggiore di callsBefore.}{S}
\addTestU{Contesto \textit{RoomListComponent}; caso 'crea il componente': verifica che component è valorizzato. Esito atteso: component è valorizzato.}{S}
\addTestU{Contesto \textit{RoomListComponent}; caso '\textit{onRoomSelect} emette l'id selezionato': verifica che emitSpy viene invocato con 'room-77' e emitSpy viene invocato 1 volte. Esito atteso: emitSpy viene invocato con 'room-77'; emitSpy viene invocato 1 volte.}{S}
\addTestU{Contesto \textit{AlarmWidgetComponent}; caso 'dovrebbe essere creato': verifica che component è valorizzato. Esito atteso: component è valorizzato.}{S}
\addTestU{Contesto \textit{AlarmWidgetComponent}; caso 'dovrebbe inizializzare il servizio al caricamento (\textit{ngOnInit})': verifica che mockAlarmService.initialize viene invocato. Esito atteso: mockAlarmService.initialize viene invocato.}{S}
\addTestU{Contesto \textit{AlarmWidgetComponent}; caso 'dovrebbe mostrare lo stato di caricamento se vm è null': verifica che compiled.textContent contiene 'Caricamento allarmi in corso'. Esito atteso: compiled.textContent contiene 'Caricamento allarmi in corso'.}{S}
\addTestU{Contesto \textit{AlarmWidgetComponent}; caso 'dovrebbe mostrare il messaggio "Nessun allarme" se la lista è vuota': verifica che compiled.textContent contiene 'Nessun allarme attivo al momento'. Esito atteso: compiled.textContent contiene 'Nessun allarme attivo al momento'.}{S}
\addTestU{Contesto \textit{AlarmWidgetComponent}; caso 'dovrebbe chiamare \textit{resolveAlarm} quando \textit{onResolve} viene invocato': verifica che mockAlarmService.\textit{resolveAlarm} viene invocato con alarmId. Esito atteso: mockAlarmService.\textit{resolveAlarm} viene invocato con alarmId.}{S}
\addTestU{Contesto \textit{AlarmWidgetComponent}; caso 'dovrebbe ricaricare i dati quando il refresh service emette': verifica che mockAlarmService.initialize viene invocato. Esito atteso: mockAlarmService.initialize viene invocato.}{S}
\addTestU{Contesto \textit{AlarmWidgetComponent}; caso 'dovrebbe calcolare le righe correttamente tramite il presenter': verifica che mockPresenter.\textit{toRows} viene invocato con alarms. Esito atteso: mockPresenter.\textit{toRows} viene invocato con alarms.}{S}
\addTestU{Contesto \textit{AnalyticsComponent}; caso 'viene creato': verifica che component è valorizzato. Esito atteso: component è valorizzato.}{S}
\addTestU{Contesto \textit{AnalyticsComponent}; caso 'carica il primo appartamento di default': verifica che mockAnalyticsService.\textit{getAllApartments} viene invocato e mockAnalyticsService.\textit{getAnalytics} viene invocato con '1'. Esito atteso: mockAnalyticsService.\textit{getAllApartments} viene invocato; mockAnalyticsService.\textit{getAnalytics} viene invocato con '1'.}{S}
\addTestU{Contesto \textit{AnalyticsComponent}; caso 'aggiorna grafici al cambio appartamento': verifica che mockAnalyticsService.\textit{getAnalytics} viene invocato con '2'. Esito atteso: mockAnalyticsService.\textit{getAnalytics} viene invocato con '2'.}{S}
\addTestU{Contesto \textit{AnalyticsComponent}; caso 'should return undefined if metric is not found': verifica che result è undefined. Esito atteso: result è undefined.}{S}
\addTestU{Contesto \textit{AnalyticsComponent}; caso 'should find the correct metric in \textit{getChartByMetric}': verifica che result è definito e il valore di result?.metric è 'plant-anomalies'. Esito atteso: result è definito; il valore di result?.metric è 'plant-anomalies'.}{S}
\addTestU{Contesto \textit{DashboardComponent}; caso 'dovrebbe creare il componente dashboard': verifica che component è valorizzato. Esito atteso: component è valorizzato.}{S}
\addTestU{Contesto \textit{DashboardComponent}; caso 'dovrebbe renderizzare il widget degli allarmi tramite il selettore': verifica che alarmWidget non è null. Esito atteso: alarmWidget non è null.}{S}
\addTestU{Contesto \textit{DashboardComponent}; caso 'dovrebbe renderizzare il widget degli analytics tramite il selettore': verifica che analyticsWidget non è null. Esito atteso: analyticsWidget non è null.}{S}
\addTestU{Contesto \textit{PlantOverviewComponent}; caso '1. Dovrebbe mostrare il nome dell impianto': verifica che html.textContent contiene 'Casa Test'. Esito atteso: html.textContent contiene 'Casa Test'.}{S}
\addTestU{Contesto \textit{PlantOverviewComponent}; caso '2. Dovrebbe mostrare i nomi delle stanze': verifica che html.textContent contiene 'Soggiorno' e html.textContent contiene 'Camera'. Esito atteso: html.textContent contiene 'Soggiorno'; html.textContent contiene 'Camera'.}{S}
\addTestU{Contesto \textit{PlantOverviewComponent}; caso '3. Dovrebbe gestire il caso senza dispositivi': verifica che html.textContent contiene 'Nessun dispositivo configurato'. Esito atteso: html.textContent contiene 'Nessun dispositivo configurato'.}{S}
\addTestU{Contesto \textit{PlantOverviewComponent}; caso '4. Non dovrebbe mostrare nulla se apartment è null': verifica che html.querySelector('header') è null. Esito atteso: html.querySelector('header') è null.}{S}
\addTestU{Contesto \textit{DeviceCardComponent}; caso 'filtra solo datapoint scrivibili con valori enum': verifica che writable ha lunghezza 1 e il valore di writable[0].id è 'dp-writable'. Esito atteso: writable ha lunghezza 1; il valore di writable[0].id è 'dp-writable'; il valore di component.hasWritableDatapoints() è true.}{S}
\addTestU{Contesto \textit{DeviceCardComponent}; caso '\textit{getSelectedValue} usa fallback iniziale e poi mantiene il valore selezionato': verifica che il valore di firstRead è 'OFF' e il valore di secondRead è 'ON'. Esito atteso: il valore di firstRead è 'OFF'; il valore di secondRead è 'ON'.}{S}
\addTestU{Contesto \textit{DeviceCardComponent}; caso '\textit{onWriteDatapoint} non emette se isExecuting e true': verifica che emitSpy non viene invocato. Esito atteso: emitSpy non viene invocato.}{S}
\addTestU{Contesto \textit{DeviceCardComponent}; caso '\textit{onWriteDatapoint} non emette se non c'è valore selezionabile': verifica che emitSpy non viene invocato. Esito atteso: emitSpy non viene invocato.}{S}
\addTestU{Contesto \textit{DeviceCardComponent}; caso '\textit{onWriteDatapoint} emette payload completo quando c'è un valore': verifica che emitSpy viene invocato con \{ roomId: 'room-1', deviceId: 'device-1', datapointId: 'dp-writable', value: 'ON', \} e emitSpy viene invocato 1 volte. Esito atteso: emitSpy viene invocato con \{ roomId: 'room-1', deviceId: 'device-1', datapointId: 'dp-writable', value: 'ON', \}; emitSpy viene invocato 1 volte.}{S}
\addTestU{Contesto \textit{DeviceCardComponent}; caso 'espone label non vuote per tipo e datapoint': verifica che component.\textit{getTypeLabel}(DeviceType.LIGHT).length è maggiore di 0 e component.\textit{getDatapointLabel}(writableDatapoint).length è maggiore di 0. Esito atteso: component.\textit{getTypeLabel}(DeviceType.LIGHT).length è maggiore di 0; component.\textit{getDatapointLabel}(writableDatapoint).length è maggiore di 0.}{S}
\addTestU{Contesto \textit{EndpointTableComponent}; caso 'creates component and loads writable endpoints': verifica che component è valorizzato e deviceApiMock.\textit{getWritableEndpointRows} viene invocato 1 volte. Esito atteso: component è valorizzato; deviceApiMock.\textit{getWritableEndpointRows} viene invocato 1 volte; deviceApiMock.getCurrentValuePointsByDeviceIds viene invocato 1 volte.}{S}
\addTestU{Contesto \textit{EndpointTableComponent}; caso 'renders mapped endpoint label for known sfeType': verifica che il valore di component.\textit{getEndpointLabel}(mappedRow) è 'Comando cambio modalita HVAC'. Esito atteso: il valore di component.\textit{getEndpointLabel}(mappedRow) è 'Comando cambio modalita HVAC'.}{S}
\addTestU{Contesto \textit{EndpointTableComponent}; caso 'falls back to a runtime humanized label when sfeType is unknown': verifica che il valore di component.\textit{getEndpointLabel}(unknownRow) è 'Comando future feature'. Esito atteso: il valore di component.\textit{getEndpointLabel}(unknownRow) è 'Comando future feature'.}{S}
\addTestU{Contesto \textit{EndpointTableComponent}; caso 'returns semantic current value by matching Cmd endpoint to State datapoint': verifica che il valore di component.\textit{getCurrentValue}(mappedRow) è 'Heat'. Esito atteso: il valore di component.\textit{getCurrentValue}(mappedRow) è 'Heat'.}{S}
\addTestU{Contesto \textit{EndpointTableComponent}; caso 'returns fallback when current datapoint value is unavailable': verifica che il valore di component.\textit{getCurrentValue}(unknownRow) è '-'. Esito atteso: il valore di component.\textit{getCurrentValue}(unknownRow) è '-'.}{S}
\addTestU{Contesto \textit{RoomDetailComponent}; caso 'should create': verifica che component è valorizzato. Esito atteso: component è valorizzato.}{S}
\addTestU{Contesto \textit{RoomDetailComponent}; caso 'mostra il titolo panoramica della pagina': verifica che heading contiene 'Panoramica allarmi e dispositivi'. Esito atteso: heading contiene 'Panoramica allarmi e dispositivi'.}{S}
\addTestU{Contesto \textit{RoomDetailComponent}; caso 'renderizza la tabella endpoint embedded': verifica che endpointTable è valorizzato. Esito atteso: endpointTable è valorizzato.}{S}
\addTestU{Contesto \textit{SidebarComponent}; caso 'mostra un link per ogni navItem con indirizzi corretti': verifica che il valore di links.length è 3 e links[0].textContent contiene 'test1'. Esito atteso: il valore di links.length è 3; links[0].textContent contiene 'test1'; il valore di links[0].\textit{getAttribute}('href') è '/path1'; links[1].textContent contiene 'test2'.}{S}
\addTestU{Contesto \textit{SidebarComponent}; caso 'mostra messaggio se navItems è vuoto': verifica che content contiene 'Menu di navigazione vuoto!'. Esito atteso: content contiene 'Menu di navigazione vuoto!'.}{S}
\addTestU{Contesto \textit{SidebarComponent}; caso 'non mostra controllo di chiusura sidebar': verifica che button è null. Esito atteso: button è null.}{S}
\addTestU{Contesto \textit{SidebarComponent}; caso 'emette navItemSelected con la route selezionata': verifica che spy viene invocato 1 volte e spy viene invocato con 'alarms/alarm-management'. Esito atteso: spy viene invocato 1 volte; spy viene invocato con 'alarms/alarm-management'.}{S}
\addTestU{Contesto \textit{TopbarComponent}; caso 'dovrebbe essere creato': verifica che component è valorizzato. Esito atteso: component è valorizzato.}{S}
\addTestU{Contesto \textit{TopbarComponent}; caso 'dovrebbe emettere hamburgerClicked quando viene premuto il menu': verifica che spy viene invocato. Esito atteso: spy viene invocato.}{S}
\addTestU{Contesto \textit{TopbarComponent}; caso 'dovrebbe emettere profileClicked quando viene premuto il nome utente': verifica che spy viene invocato. Esito atteso: spy viene invocato.}{S}
\addTestU{Contesto \textit{TopbarComponent}; caso 'emette profileClicked anche quando utente non admin': verifica che spy viene invocato. Esito atteso: spy viene invocato.}{S}
\addTestU{Contesto \textit{TopbarComponent}; caso 'evidenzia in giallo il profilo quando attivo': verifica che il valore di profileButton.classList.contains('bg-amber-300') è true. Esito atteso: il valore di profileButton.classList.contains('bg-amber-300') è true.}{S}
\addTestU{Contesto \textit{TopbarComponent}; caso 'mostra avviso MyVimar quando richiesto': verifica che warning è valorizzato e warning.textContent?.\textit{toLowerCase}() contiene 'account da associare a myvimar'. Esito atteso: warning è valorizzato; warning.textContent?.\textit{toLowerCase}() contiene 'account da associare a myvimar'.}{S}
\addTestU{Contesto \textit{MyVimarAccountStatusComponent}; caso 'should create': verifica che component è valorizzato. Esito atteso: component è valorizzato.}{S}
\addTestU{Contesto \textit{MyVimarAccountStatusComponent}; caso 'mostra pulsante di collegamento quando account non collegato': verifica che text contiene 'Nessun account MyVimar collegato.' e text contiene 'Collega account MyVimar'. Esito atteso: text contiene 'Nessun account MyVimar collegato.'; text contiene 'Collega account MyVimar'.}{S}
\addTestU{Contesto \textit{MyVimarAccountStatusComponent}; caso 'mostra email e pulsante di rimozione quando account collegato': verifica che text contiene 'admin@example.com' e text contiene 'Rimuovi account'. Esito atteso: text contiene 'admin@example.com'; text contiene 'Rimuovi account'.}{S}
\addTestU{Contesto \textit{MyVimarAccountStatusComponent}; caso 'emette linkClicked quando viene premuto il pulsante collega': verifica che emitSpy viene invocato 1 volte. Esito atteso: emitSpy viene invocato 1 volte.}{S}
\addTestU{Contesto \textit{MyVimarAccountStatusComponent}; caso 'emette unlinkClicked quando viene premuto il pulsante rimuovi': verifica che emitSpy viene invocato 1 volte. Esito atteso: emitSpy viene invocato 1 volte.}{S}
\addTestU{Contesto \textit{MyVimarPageComponent}; caso 'should create': verifica che component è valorizzato. Esito atteso: component è valorizzato.}{S}
\addTestU{Contesto \textit{MyVimarPageComponent}; caso 'carica lo stato account in \textit{ngOnInit}': verifica che serviceStub.\textit{getLinkedAccount} viene invocato 1 volte. Esito atteso: serviceStub.\textit{getLinkedAccount} viene invocato 1 volte.}{S}
\addTestU{Contesto \textit{MyVimarPageComponent}; caso 'delegates \textit{onLinkAccount} al service': verifica che serviceStub.initiateOAuth viene invocato 1 volte. Esito atteso: serviceStub.initiateOAuth viene invocato 1 volte.}{S}
\addTestU{Contesto \textit{MyVimarPageComponent}; caso '\textit{onUnlinkAccount} richiama il service e triggera refresh': verifica che serviceStub.unlinkAccount viene invocato 1 volte e serviceStub.\textit{getLinkedAccount} viene invocato 2 volte. Esito atteso: serviceStub.unlinkAccount viene invocato 1 volte; serviceStub.\textit{getLinkedAccount} viene invocato 2 volte.}{S}
\addTestU{Contesto \textit{MyVimarPageComponent}; caso 'imposta errore quando unlinkAccount fallisce': verifica che component.error contiene 'Errore durante la rimozione' e il valore di component.isLoading è false. Esito atteso: component.error contiene 'Errore durante la rimozione'; il valore di component.isLoading è false.}{S}
\addTestU{Contesto \textit{NotificationBadgeComponent}; caso 'should create': verifica che component è valorizzato. Esito atteso: component è valorizzato.}{S}
\addTestU{Contesto \textit{NotificationBadgeComponent}; caso 'nasconde il badge quando count e 0': verifica che badge è null. Esito atteso: badge è null.}{S}
\addTestU{Contesto \textit{NotificationBadgeComponent} - mostra il badge quando count; caso '0 con testo e label accessibile': verifica che con count positivo il badge viene mostrato con testo numerico e attributo aria-label coerente. Esito atteso: il badge è visibile, mostra testo '3' e aria-label '3 notifiche non lette'.}{S}
\addTestU{Contesto \textit{NotificationBadgeComponent}; caso 'normalizza valori negativi nascondendo il badge': verifica che badge è null. Esito atteso: badge è null.}{S}
\addTestU{Contesto \textit{NotificationBadgeComponent}; caso 'limita la visualizzazione a 99+ per count elevati': verifica che il valore di badge.textContent.trim() è '99+' e il valore di badge.\textit{getAttribute}('aria-label') è '120 notifiche non lette'. Esito atteso: il valore di badge.textContent.trim() è '99+'; il valore di badge.\textit{getAttribute}('aria-label') è '120 notifiche non lette'.}{S}
\addTestU{Contesto \textit{NotificationItemComponent}; caso 'should create': verifica che component è valorizzato. Esito atteso: component è valorizzato.}{S}
\addTestU{Contesto \textit{NotificationItemComponent}; caso 'renderizza titolo e tempo relativo': verifica che title.textContent contiene 'Rilevato movimento in soggiorno' e il valore di time.textContent.trim() è '1m fa'. Esito atteso: title.textContent contiene 'Rilevato movimento in soggiorno'; il valore di time.textContent.trim() è '1m fa'.}{S}
\addTestU{Contesto \textit{NotificationItemComponent}; caso 'espone il timestamp originale in datetime per accessibilita e machine-readability': verifica che il valore di time.\textit{getAttribute}('datetime') è notification.sentAt. Esito atteso: il valore di time.\textit{getAttribute}('datetime') è notification.sentAt.}{S}
\addTestU{Contesto \textit{NotificationItemComponent}; caso 'aggiorna il contenuto quando cambia input notification': verifica che title.textContent contiene 'Temperatura elevata in cucina' e il valore di time.textContent.trim() è '40s fa'. Esito atteso: title.textContent contiene 'Temperatura elevata in cucina'; il valore di time.textContent.trim() è '40s fa'.}{S}
\addTestU{Contesto \textit{NotificationItemComponent}; caso 'emette \textit{removeClicked} quando si clicca la X': verifica che emitSpy viene invocato con 'n-1' e emitSpy viene invocato 1 volte. Esito atteso: emitSpy viene invocato con 'n-1'; emitSpy viene invocato 1 volte.}{S}
\addTestU{Contesto \textit{NotificationPageComponent}; caso 'should create': verifica che component è valorizzato. Esito atteso: component è valorizzato.}{S}
\addTestU{Contesto \textit{NotificationPageComponent}; caso 'renderizza uno stato vuoto accessibile quando non ci sono notifiche': verifica che emptyState è valorizzato e emptyState.textContent contiene 'Nessuna notifica disponibile'. Esito atteso: emptyState è valorizzato; emptyState.textContent contiene 'Nessuna notifica disponibile'; items ha lunghezza 0.}{S}
\addTestU{Contesto \textit{NotificationPageComponent}; caso 'renderizza la lista notifiche e passa i dati ai figli': verifica che summary.textContent contiene '2 notifiche non lette' e items ha lunghezza 2. Esito atteso: summary.textContent contiene '2 notifiche non lette'; items ha lunghezza 2; items[0].componentInstance.notification() coincide con notificationA; items[1].componentInstance.notification() coincide con notificationB.}{S}
\addTestU{Contesto \textit{NotificationPageComponent}; caso 'usa il singolare nel riepilogo quando unreadCount e 1': verifica che summary.textContent contiene '1 notifica non letta'. Esito atteso: summary.textContent contiene '1 notifica non letta'.}{S}
\addTestU{Contesto \textit{NotificationPageComponent}; caso 'aggiorna il DOM quando vm\$ emette un nuovo snapshot coerente': verifica che items ha lunghezza 1 e items ha lunghezza 2. Esito atteso: items ha lunghezza 1; items ha lunghezza 2.}{S}
\addTestU{Contesto \textit{NotificationPageComponent}; caso 'invoca \textit{clearAllNotifications} quando si clicca Cancella tutte': verifica che notificationServiceStub.\textit{clearAllNotifications} viene invocato con [ notificationA, notificationB, ]. Esito atteso: notificationServiceStub.\textit{clearAllNotifications} viene invocato con [ notificationA, notificationB, ].}{S}
\addTestU{Contesto \textit{NotificationPageComponent}; caso 'evidenzia leggermente la notifica richiesta da query param focus': verifica che il valore di items[0].componentInstance.isHighlighted() è false e il valore di items[1].componentInstance.isHighlighted() è true. Esito atteso: il valore di items[0].componentInstance.isHighlighted() è false; il valore di items[1].componentInstance.isHighlighted() è true.}{S}
\addTestU{Contesto \textit{NotificationTopbarPanelComponent}; caso 'mostra stato vuoto quando non ci sono notifiche': verifica che empty è valorizzato e empty.textContent?.\textit{toLowerCase}() contiene 'nessuna notifica'. Esito atteso: empty è valorizzato; empty.textContent?.\textit{toLowerCase}() contiene 'nessuna notifica'.}{S}
\addTestU{Contesto \textit{NotificationTopbarPanelComponent}; caso 'emette viewAllClicked quando si clicca su Vedi tutte': verifica che emitSpy viene invocato 1 volte. Esito atteso: emitSpy viene invocato 1 volte.}{S}
\addTestU{Contesto \textit{NotificationTopbarPanelComponent}; caso 'emette \textit{removeClicked} quando si clicca sulla X di una notifica': verifica che \textit{removeSpy} viene invocato con 'n-1' e \textit{removeSpy} viene invocato 1 volte. Esito atteso: \textit{removeSpy} viene invocato con 'n-1'; \textit{removeSpy} viene invocato 1 volte.}{S}
\addTestU{Contesto \textit{NotificationTopbarPanelComponent}; caso 'emette notificationSelected quando si clicca su un item notifica in preview': verifica che selectedSpy viene invocato con 'n-1' e selectedSpy viene invocato 1 volte. Esito atteso: selectedSpy viene invocato con 'n-1'; selectedSpy viene invocato 1 volte.}{S}
\addTestU{Contesto \textit{NotificationTopbarPanelComponent}; caso 'emette \textit{clearAllClicked} quando si clicca Cancella tutte': verifica che \textit{clearAllSpy} viene invocato 1 volte. Esito atteso: \textit{clearAllSpy} viene invocato 1 volte.}{S}
\addTestU{Contesto \textit{AuthBaseComponent}; caso '\textit{handleSuccess} naviga a first-access quando isFirstAccess e true': verifica che il valore di component.isLoading è false e component.errorType è null. Esito atteso: il valore di component.isLoading è false; component.errorType è null; routerMock.navigate viene invocato con ['/auth/first-access'].}{S}
\addTestU{Contesto \textit{AuthBaseComponent}; caso '\textit{handleSuccess} naviga a dashboard quando isFirstAccess e false': verifica che routerMock.navigate viene invocato con ['/dashboard']. Esito atteso: routerMock.navigate viene invocato con ['/dashboard'].}{S}
\addTestU{Contesto \textit{AuthBaseComponent}; caso '\textit{handleSuccess} naviga su returnUrl quando presente': verifica che routerMock.navigateByUrl viene invocato con '/vimar-link' e routerMock.navigate non viene invocato con ['/dashboard']. Esito atteso: routerMock.navigateByUrl viene invocato con '/vimar-link'; routerMock.navigate non viene invocato con ['/dashboard'].}{S}
\addTestU{Contesto \textit{AuthBaseComponent}; caso '\textit{handleError} imposta USERNAME\_OR\_PASSWORD\_WRONG su 400 e 401': verifica che il valore di component.errorType è AuthErrorType.USERNAME\_OR\_PASSWORD\_WRONG. Esito atteso: il valore di component.errorType è AuthErrorType.USERNAME\_OR\_PASSWORD\_WRONG.}{S}
\addTestU{Contesto \textit{AuthBaseComponent}; caso '\textit{handleError} imposta NEW\_PASSWORD\_EQUALS\_TEMP su 409': verifica che il valore di component.errorType è AuthErrorType.NEW\_PASSWORD\_EQUALS\_TEMP. Esito atteso: il valore di component.errorType è AuthErrorType.NEW\_PASSWORD\_EQUALS\_TEMP.}{S}
\addTestU{Contesto \textit{AuthBaseComponent}; caso '\textit{handleError} usa fallback USERNAME\_OR\_PASSWORD\_WRONG per errori non Http': verifica che il valore di component.errorType è AuthErrorType.USERNAME\_OR\_PASSWORD\_WRONG e il valore di component.isLoading è false. Esito atteso: il valore di component.errorType è AuthErrorType.USERNAME\_OR\_PASSWORD\_WRONG; il valore di component.isLoading è false.}{S}
\addTestU{Contesto \textit{AuthBaseComponent}; caso 'espone il messaggio umano per credenziali errate': verifica che il valore di component.loginErrorMessage è 'Utente non trovato: username o password errati.'. Esito atteso: il valore di component.loginErrorMessage è 'Utente non trovato: username o password errati.'.}{S}
\addTestU{Contesto \textit{FirstAccessComponent}; caso 'crea il componente': verifica che component è valorizzato. Esito atteso: component è valorizzato.}{S}
\addTestU{Contesto \textit{FirstAccessComponent}; caso 'renderizza il contenuto informativo del primo accesso': verifica che content contiene 'Primo accesso' e content contiene 'Imposta la tua password personale'. Esito atteso: content contiene 'Primo accesso'; content contiene 'Imposta la tua password personale'; content contiene 'password temporanea'; content contiene 'Conferma e continua'.}{S}
\addTestU{Contesto \textit{FirstAccessComponent}; caso 'imposta errore di validazione quando la nuova password coincide con la temporanea': verifica che il valore di component.firstAccessForm.errors?.['sameAsTemporaryPassword'] è true e authServiceMock.\textit{setFirstAccessPassword} non viene invocato. Esito atteso: il valore di component.firstAccessForm.errors?.['sameAsTemporaryPassword'] è true; authServiceMock.\textit{setFirstAccessPassword} non viene invocato.}{S}
\addTestU{Contesto \textit{FirstAccessComponent}; caso 'invia \textit{setFirstAccessPassword} e naviga su dashboard in caso di successo': verifica che authServiceMock.\textit{setFirstAccessPassword} viene invocato con 'mrossi' e routerMock.navigate viene invocato con ['/dashboard']. Esito atteso: authServiceMock.\textit{setFirstAccessPassword} viene invocato con 'mrossi'; routerMock.navigate viene invocato con ['/dashboard']; il valore di component.isLoading è false.}{S}
\addTestU{Contesto \textit{FirstAccessComponent}; caso 'imposta errore USERNAME\_OR\_TEMP\_PASSWORD\_WRONG in caso di errore API': verifica che il valore di component.errorType è AuthErrorType.USERNAME\_OR\_TEMP\_PASSWORD\_WRONG e il valore di component.isLoading è false. Esito atteso: il valore di component.errorType è AuthErrorType.USERNAME\_OR\_TEMP\_PASSWORD\_WRONG; il valore di component.isLoading è false.}{S}
\addTestU{Contesto \textit{FirstAccessComponent}; caso 'renderizza un messaggio umano quando username o password temporanea sono errati': verifica che content contiene 'Username o password temporanea errati.'. Esito atteso: content contiene 'Username o password temporanea errati.'.}{S}
\addTestU{Contesto \textit{FirstAccessComponent}; caso 'marca il form touched quando submit su form invalido': verifica che markAllAsTouchedSpy viene invocato e authServiceMock.\textit{setFirstAccessPassword} non viene invocato. Esito atteso: markAllAsTouchedSpy viene invocato; authServiceMock.\textit{setFirstAccessPassword} non viene invocato.}{S}
\addTestU{Contesto \textit{FirstAccessComponent}; caso 'non invia la richiesta quando la nuova password e troppo corta': verifica che component.firstAccessForm.controls.newPassword.errors?.['minlength'] è valorizzato e authServiceMock.\textit{setFirstAccessPassword} non viene invocato. Esito atteso: component.firstAccessForm.controls.newPassword.errors?.['minlength'] è valorizzato; authServiceMock.\textit{setFirstAccessPassword} non viene invocato.}{S}
\addTestU{Contesto \textit{FirstAccessComponent}; caso 'renderizza il messaggio di validazione quando la nuova password coincide con la temporanea': verifica che content contiene 'La nuova password deve essere diversa da quella temporanea.'. Esito atteso: content contiene 'La nuova password deve essere diversa da quella temporanea.'.}{S}
\addTestU{Contesto \textit{LoginComponent}; caso 'crea il componente': verifica che component è valorizzato. Esito atteso: component è valorizzato.}{S}
\addTestU{Contesto \textit{LoginComponent}; caso 'non invoca login se il form non rispetta i validator': verifica che markAllAsTouchedSpy viene invocato e authServiceMock.login non viene invocato. Esito atteso: markAllAsTouchedSpy viene invocato; authServiceMock.login non viene invocato.}{S}
\addTestU{Contesto \textit{LoginComponent}; caso 'invoca login e naviga su dashboard se autenticazione ok': verifica che authServiceMock.login viene invocato con 'mrossi' e routerMock.navigate viene invocato con ['/dashboard']. Esito atteso: authServiceMock.login viene invocato con 'mrossi'; routerMock.navigate viene invocato con ['/dashboard']; component.errorType è null; il valore di component.isLoading è false.}{S}
\addTestU{Contesto \textit{LoginComponent}; caso 'invoca login anche con valori presenti solo nel DOM (autofill browser)': verifica che authServiceMock.login viene invocato con 'mrossi'. Esito atteso: authServiceMock.login viene invocato con 'mrossi'.}{S}
\addTestU{Contesto \textit{LoginComponent}; caso 'invoca login e naviga su returnUrl se presente': verifica che authServiceMock.login viene invocato con 'mrossi' e routerMock.navigateByUrl viene invocato con '/vimar-link'. Esito atteso: authServiceMock.login viene invocato con 'mrossi'; routerMock.navigateByUrl viene invocato con '/vimar-link'.}{S}
\addTestU{Contesto \textit{LoginComponent}; caso 'imposta errore USERNAME\_OR\_PASSWORD\_WRONG su 401': verifica che il valore di component.errorType è AuthErrorType.USERNAME\_OR\_PASSWORD\_WRONG e il valore di component.isLoading è false. Esito atteso: il valore di component.errorType è AuthErrorType.USERNAME\_OR\_PASSWORD\_WRONG; il valore di component.isLoading è false.}{S}
\addTestU{Contesto \textit{LoginComponent}; caso 'espone il messaggio umano già al primo tentativo fallito': verifica che il valore di component.isLoading è false e il valore di component.loginErrorMessage è 'Utente non trovato: username o password errati.'. Esito atteso: il valore di component.isLoading è false; il valore di component.loginErrorMessage è 'Utente non trovato: username o password errati.'.}{S}
\addTestU{Contesto \textit{LoginComponent}; caso 'renderizza il testo errore quando errorType è valorizzato': verifica che content contiene 'Utente non trovato: username o password errati.'. Esito atteso: content contiene 'Utente non trovato: username o password errati.'.}{S}
\addTestU{Contesto \textit{CreateUserForm}; caso 'should create': verifica che component è valorizzato. Esito atteso: component è valorizzato.}{S}
\addTestU{Contesto \textit{CreateUserForm}; caso 'submit emette formSubmit quando il form e valido': verifica che submitSpy viene invocato con \{ name: 'Mario', surname: 'Rossi', username: 'mrossi', \} e submitSpy viene invocato 1 volte. Esito atteso: submitSpy viene invocato con \{ name: 'Mario', surname: 'Rossi', username: 'mrossi', \}; submitSpy viene invocato 1 volte.}{S}
\addTestU{Contesto \textit{CreateUserForm}; caso 'submit non emette e marca touched quando il form e invalido': verifica che il valore di component.form.invalid è true e il valore di component.form.controls['name'].touched è true. Esito atteso: il valore di component.form.invalid è true; il valore di component.form.controls['name'].touched è true; il valore di component.form.controls['surname'].touched è true; il valore di component.form.controls['username'].touched è true.}{S}
\addTestU{Contesto \textit{CreateUserForm}; caso 'applica validazioni required e minLength': verifica che il valore di component.form.controls['name'].hasError('required') è true e il valore di component.form.controls['surname'].hasError('minlength') è true. Esito atteso: il valore di component.form.controls['name'].hasError('required') è true; il valore di component.form.controls['surname'].hasError('minlength') è true; il valore di component.form.controls['username'].hasError('minlength') è true.}{S}
\addTestU{Contesto \textit{CreateUserForm}; caso 'reset pulisce il form': verifica che component.form.value coincide con \{ name: '', surname: '', username: '', \}. Esito atteso: component.form.value coincide con \{ name: '', surname: '', username: '', \}.}{S}
\addTestU{Contesto \textit{CreateUserForm}; caso '\textit{requestClose} emette closeRequest': verifica che closeSpy viene invocato 1 volte. Esito atteso: closeSpy viene invocato 1 volte.}{S}
\addTestU{Contesto \textit{CreateUserForm}; caso '\textit{requestOpen} emette openRequest': verifica che openSpy viene invocato 1 volte. Esito atteso: openSpy viene invocato 1 volte.}{S}
\addTestU{Contesto \textit{UserCreatedDialogComponent}; caso 'should create': verifica che component è valorizzato. Esito atteso: component è valorizzato.}{S}
\addTestU{Contesto \textit{UserCreatedDialogComponent}; caso 'mostra nome utente e password temporanea': verifica che textContent contiene 'Mario' e textContent contiene 'mrossi'. Esito atteso: textContent contiene 'Mario'; textContent contiene 'mrossi'; textContent contiene 'TempPass123'.}{S}
\addTestU{Contesto \textit{UserCreatedDialogComponent}; caso 'closeDialog emette evento closed': verifica che closedSpy viene invocato 1 volte. Esito atteso: closedSpy viene invocato 1 volte.}{S}
\addTestU{Contesto \textit{UserCreatedDialogComponent}; caso 'click sul bottone chiudi invoca closeDialog': verifica che closeSpy viene invocato 1 volte. Esito atteso: closeSpy viene invocato 1 volte.}{S}
\addTestU{Contesto \textit{UserListComponent}; caso 'should create': verifica che component è valorizzato. Esito atteso: component è valorizzato.}{S}
\addTestU{Contesto \textit{UserListComponent}; caso '\textit{onDelete} apre il dialog di conferma': verifica che component.pendingDeleteUser() coincide con user. Esito atteso: component.pendingDeleteUser() coincide con user.}{S}
\addTestU{Contesto \textit{UserListComponent}; caso '\textit{onDeleteConfirm} emette \textit{deleteUser} con id numerico': verifica che emitSpy viene invocato con 1 e emitSpy viene invocato 1 volte. Esito atteso: emitSpy viene invocato con 1; emitSpy viene invocato 1 volte; component.pendingDeleteUser() è null.}{S}
\addTestU{Contesto \textit{UserListComponent}; caso '\textit{onDeleteCancel} non emette \textit{deleteUser} e chiude il dialog': verifica che emitSpy non viene invocato e component.pendingDeleteUser() è null. Esito atteso: emitSpy non viene invocato; component.pendingDeleteUser() è null.}{S}
\addTestU{Contesto UserManagementPage; caso 'should create': verifica che component è valorizzato e \textit{getUsersMock} viene invocato 1 volte. Esito atteso: component è valorizzato; \textit{getUsersMock} viene invocato 1 volte.}{S}
\addTestU{Contesto UserManagementPage; caso 'mostra di default solo la tabella e non il form OSS': verifica che il valore di component.isCreateFormOpen() è false e nativeElement.textContent contiene 'Elenco operatori sanitari'. Esito atteso: il valore di component.isCreateFormOpen() è false; nativeElement.textContent contiene 'Elenco operatori sanitari'; panel non è null; panel?.className contiene 'max-h-0'.}{S}
\addTestU{Contesto UserManagementPage; caso 'apre e chiude il form OSS con i pulsanti dedicati': verifica che il valore di component.isCreateFormOpen() è true e il valore di component.isCreateFormOpen() è false. Esito atteso: il valore di component.isCreateFormOpen() è true; il valore di component.isCreateFormOpen() è false.}{S}
\addTestU{Contesto UserManagementPage; caso '\textit{onFormSubmit} in successo imposta createdResponse e triggera refresh lista': verifica che \textit{createUserMock} viene invocato con \{ name: 'Mario', surname: 'Rossi', username: 'mrossi', \} e il valore di component.createdResponse()?.tempPassword è 'TempPass123'. Esito atteso: \textit{createUserMock} viene invocato con \{ name: 'Mario', surname: 'Rossi', username: 'mrossi', \}; il valore di component.createdResponse()?.tempPassword è 'TempPass123'; component.createdUser() coincide con \{ name: 'Mario', surname: 'Rossi', username: 'mrossi', \}; component.formError() è null.}{S}
\addTestU{Contesto UserManagementPage; caso '\textit{onFormSubmit} con 409 imposta USERNAME\_ALREADY\_IN\_USE': verifica che il valore di component.formError() è UserManagementErrorType.USERNAME\_ALREADY\_IN\_USE. Esito atteso: il valore di component.formError() è UserManagementErrorType.USERNAME\_ALREADY\_IN\_USE.}{S}
\addTestU{Contesto UserManagementPage; caso '\textit{onFormSubmit} con errore generico imposta OTHER\_ERROR': verifica che il valore di component.formError() è UserManagementErrorType.OTHER\_ERROR. Esito atteso: il valore di component.formError() è UserManagementErrorType.OTHER\_ERROR.}{S}
\addTestU{Contesto UserManagementPage; caso '\textit{onUserDeleted} delega al service e triggera refresh lista': verifica che \textit{deleteUserMock} viene invocato con 1 e \textit{getUsersMock} viene invocato 2 volte. Esito atteso: \textit{deleteUserMock} viene invocato con 1; \textit{getUsersMock} viene invocato 2 volte.}{S}
\addTestU{Contesto UserManagementPage; caso '\textit{onDialogClosed} azzera createdResponse': verifica che component.createdResponse() è null e component.createdUser() è null. Esito atteso: component.createdResponse() è null; component.createdUser() è null.}{S}
\addTestU{Contesto UserManagementPage; caso 'users\$ ritorna [] se il caricamento utenti fallisce': verifica che users coincide con []. Esito atteso: users coincide con [].}{S}
\addTestU{Contesto \textit{AssignOperatorDialogComponent}; caso 'should create': verifica che component è valorizzato. Esito atteso: component è valorizzato.}{S}
\addTestU{Contesto \textit{AssignOperatorDialogComponent}; caso 'dovrebbe resettare userId in \textit{ngOnInit}': verifica che component.form.controls.userId.value è null. Esito atteso: component.form.controls.userId.value è null.}{S}
\addTestU{Contesto \textit{AssignOperatorDialogComponent}; caso 'non dovrebbe emettere submitted se il form e invalido': verifica che il valore di component.form.invalid è true e il valore di component.form.controls.userId.touched è true. Esito atteso: il valore di component.form.invalid è true; il valore di component.form.controls.userId.touched è true; submittedSpy non viene invocato.}{S}
\addTestU{Contesto \textit{AssignOperatorDialogComponent}; caso 'dovrebbe emettere submitted con userId quando il form e valido': verifica che submittedSpy viene invocato con \{ userId: 1 \} e submittedSpy viene invocato 1 volte. Esito atteso: submittedSpy viene invocato con \{ userId: 1 \}; submittedSpy viene invocato 1 volte.}{S}
\addTestU{Contesto \textit{AssignOperatorDialogComponent}; caso 'dovrebbe emettere cancelled quando \textit{onCancel} viene chiamato': verifica che cancelledSpy viene invocato 1 volte. Esito atteso: cancelledSpy viene invocato 1 volte.}{S}
\addTestU{Contesto \textit{AssignWardDialogComponent}; caso 'should create': verifica che component è valorizzato. Esito atteso: component è valorizzato.}{S}
\addTestU{Contesto \textit{AssignWardDialogComponent}; caso 'dovrebbe resettare plantId a null in \textit{ngOnInit}': verifica che component.form.controls.plantId.value è null. Esito atteso: component.form.controls.plantId.value è null.}{S}
\addTestU{Contesto \textit{AssignWardDialogComponent}; caso 'non dovrebbe emettere submitted se il form e invalido': verifica che il valore di component.form.invalid è true e il valore di component.form.controls.plantId.touched è true. Esito atteso: il valore di component.form.invalid è true; il valore di component.form.controls.plantId.touched è true; submittedSpy non viene invocato.}{S}
\addTestU{Contesto \textit{AssignWardDialogComponent}; caso 'dovrebbe emettere submitted con plantId quando il form e valido': verifica che submittedSpy viene invocato con \{ plantId: '101' \} e submittedSpy viene invocato 1 volte. Esito atteso: submittedSpy viene invocato con \{ plantId: '101' \}; submittedSpy viene invocato 1 volte.}{S}
\addTestU{Contesto \textit{AssignWardDialogComponent}; caso 'dovrebbe emettere cancelled quando \textit{onCancel} viene chiamato': verifica che cancelledSpy viene invocato 1 volte. Esito atteso: cancelledSpy viene invocato 1 volte.}{S}
\addTestU{Contesto \textit{WardCardComponent}; caso 'should create': verifica che component è valorizzato. Esito atteso: component è valorizzato.}{S}
\addTestU{Contesto \textit{WardCardComponent}; caso '\textit{toggleExpanded} dovrebbe alternare lo stato espanso': verifica che il valore di component.isExpanded() è false e il valore di component.isExpanded() è true. Esito atteso: il valore di component.isExpanded() è false; il valore di component.isExpanded() è true.}{S}
\addTestU{Contesto \textit{WardCardComponent}; caso '\textit{onEditWardClick} dovrebbe emettere editWard con il ward corrente': verifica che editWardSpy viene invocato con ward e editWardSpy viene invocato 1 volte. Esito atteso: editWardSpy viene invocato con ward; editWardSpy viene invocato 1 volte.}{S}
\addTestU{Contesto \textit{WardCardComponent}; caso '\textit{onDeleteWardClick} dovrebbe emettere \textit{deleteWard} con ward.id': verifica che \textit{deleteWardSpy} viene invocato con 1 e \textit{deleteWardSpy} viene invocato 1 volte. Esito atteso: \textit{deleteWardSpy} viene invocato con 1; \textit{deleteWardSpy} viene invocato 1 volte.}{S}
\addTestU{Contesto \textit{WardCardComponent}; caso '\textit{onAssignOperatorClick} dovrebbe emettere \textit{assignOperator} con ward.id': verifica che \textit{assignOperatorSpy} viene invocato con 1 e \textit{assignOperatorSpy} viene invocato 1 volte. Esito atteso: \textit{assignOperatorSpy} viene invocato con 1; \textit{assignOperatorSpy} viene invocato 1 volte.}{S}
\addTestU{Contesto \textit{WardCardComponent}; caso '\textit{onRemoveOperatorClick} dovrebbe emettere payload completo': verifica che \textit{removeOperatorSpy} viene invocato con \{ wardId: 1, userId: 2, \} e \textit{removeOperatorSpy} viene invocato 1 volte. Esito atteso: \textit{removeOperatorSpy} viene invocato con \{ wardId: 1, userId: 2, \}; \textit{removeOperatorSpy} viene invocato 1 volte.}{S}
\addTestU{Contesto \textit{WardCardComponent}; caso '\textit{onAssignPlantClick} dovrebbe emettere \textit{assignPlant} con ward.id': verifica che \textit{assignPlantSpy} viene invocato con 1 e \textit{assignPlantSpy} viene invocato 1 volte. Esito atteso: \textit{assignPlantSpy} viene invocato con 1; \textit{assignPlantSpy} viene invocato 1 volte.}{S}
\addTestU{Contesto \textit{WardCardComponent}; caso '\textit{onRemovePlantClick} dovrebbe emettere payload completo': verifica che \textit{removePlantSpy} viene invocato con \{ wardId: 1, plantId: '102', \} e \textit{removePlantSpy} viene invocato 1 volte. Esito atteso: \textit{removePlantSpy} viene invocato con \{ wardId: 1, plantId: '102', \}; \textit{removePlantSpy} viene invocato 1 volte.}{S}
\addTestU{Contesto \textit{WardCardComponent}; caso 'non dovrebbe emettere eventi legati al ward quando ward è null': verifica che editWardSpy non viene invocato e \textit{deleteWardSpy} non viene invocato. Esito atteso: editWardSpy non viene invocato; \textit{deleteWardSpy} non viene invocato; \textit{assignOperatorSpy} non viene invocato; \textit{removeOperatorSpy} non viene invocato.}{S}
\addTestU{Contesto \textit{WardFormDialogComponent}; caso 'should create': verifica che component è valorizzato. Esito atteso: component è valorizzato.}{S}
\addTestU{Contesto \textit{WardFormDialogComponent}; caso 'dovrebbe precompilare il nome in edit mode': verifica che il valore di component.isEditMode() è true e il valore di component.form.controls.name.value è 'Cardiologia'. Esito atteso: il valore di component.isEditMode() è true; il valore di component.form.controls.name.value è 'Cardiologia'.}{S}
\addTestU{Contesto \textit{WardFormDialogComponent}; caso 'non dovrebbe emettere submitted se nome vuoto': verifica che il valore di component.form.invalid è true e il valore di component.form.controls.name.touched è true. Esito atteso: il valore di component.form.invalid è true; il valore di component.form.controls.name.touched è true; submittedSpy non viene invocato.}{S}
\addTestU{Contesto \textit{WardFormDialogComponent}; caso 'non dovrebbe emettere submitted se il nome supera 100 caratteri': verifica che il valore di component.form.invalid è true e component.form.controls.name.errors?.['maxlength'] è valorizzato. Esito atteso: il valore di component.form.invalid è true; component.form.controls.name.errors?.['maxlength'] è valorizzato; submittedSpy non viene invocato.}{S}
\addTestU{Contesto \textit{WardFormDialogComponent}; caso 'dovrebbe emettere submitted con nome trimmato se form valido': verifica che submittedSpy viene invocato con \{ name: 'Psicopazzia' \} e submittedSpy viene invocato 1 volte. Esito atteso: submittedSpy viene invocato con \{ name: 'Psicopazzia' \}; submittedSpy viene invocato 1 volte.}{S}
\addTestU{Contesto \textit{WardFormDialogComponent}; caso 'dovrebbe emettere cancelled quando \textit{onCancel} viene chiamato': verifica che cancelledSpy viene invocato 1 volte. Esito atteso: cancelledSpy viene invocato 1 volte.}{S}
\addTestU{Contesto \textit{WardManagementPageComponent}; caso 'should create': verifica che component è valorizzato. Esito atteso: component è valorizzato.}{S}
\addTestU{Contesto \textit{WardManagementPageComponent}; caso '\textit{ngOnInit} dovrebbe chiamare \textit{loadWards} e sincronizzare snapshots/error': verifica che storeStub.\textit{loadWards} viene invocato 1 volte e component.wardsSnapshot() coincide con []. Esito atteso: storeStub.\textit{loadWards} viene invocato 1 volte; component.wardsSnapshot() coincide con []; component.snackbarMessage() è null; component.wardsSnapshot() ha lunghezza 2.}{S}
\addTestU{Contesto \textit{WardManagementPageComponent}; caso 'quando wards diventa vuoto resetta selezioni e step mobile': verifica che component.selectedWardId() è null e component.selectedApartmentId() è null. Esito atteso: component.selectedWardId() è null; component.selectedApartmentId() è null; il valore di component.mobileStep() è 'wards'.}{S}
\addTestU{Contesto \textit{WardManagementPageComponent}; caso 'selectedApartment si riallinea al ward selezionato': verifica che component.selectedApartment() coincide con \{ id: '101', name: 'App. 101' \} e component.selectedApartment() è null. Esito atteso: component.selectedApartment() coincide con \{ id: '101', name: 'App. 101' \}; component.selectedApartment() è null.}{S}
\addTestU{Contesto \textit{WardManagementPageComponent}; caso '\textit{ngOnDestroy} dovrebbe interrompere aggiornamenti da stream': verifica che il valore di component.snackbarMessage() è 'Errore iniziale'. Esito atteso: il valore di component.snackbarMessage() è 'Errore iniziale'.}{S}
\addTestU{Contesto \textit{WardManagementPageComponent}; caso '\textit{dismissSnackbar} dovrebbe azzerare il messaggio': verifica che component.snackbarMessage() è null. Esito atteso: component.snackbarMessage() è null.}{S}
\addTestU{Contesto \textit{WardManagementPageComponent}; caso '\textit{onCreateWard} dovrebbe aprire dialog in create mode': verifica che il valore di component.selectedWardId() è wardA.id e il valore di component.wardDialogMode() è 'create'. Esito atteso: il valore di component.selectedWardId() è wardA.id; il valore di component.wardDialogMode() è 'create'.}{S}
\addTestU{Contesto \textit{WardManagementPageComponent}; caso '\textit{onCreateWardSubmit} dovrebbe delegare a store e chiudere dialog': verifica che storeStub.\textit{createWard} viene invocato con \{ name: 'Oncologia' \} e storeStub.\textit{createWard} viene invocato 1 volte. Esito atteso: storeStub.\textit{createWard} viene invocato con \{ name: 'Oncologia' \}; storeStub.\textit{createWard} viene invocato 1 volte; il valore di component.wardDialogMode() è 'closed'.}{S}
\addTestU{Contesto \textit{WardManagementPageComponent}; caso '\textit{onEditWard} e \textit{onEditWardSubmit} dovrebbero aggiornare il reparto selezionato': verifica che il valore di component.selectedWardId() è wardA.id e il valore di component.wardDialogMode() è 'edit'. Esito atteso: il valore di component.selectedWardId() è wardA.id; il valore di component.wardDialogMode() è 'edit'; storeStub.\textit{updateWard} viene invocato con 1; il valore di component.wardDialogMode() è 'closed'.}{S}
\addTestU{Contesto \textit{WardManagementPageComponent}; caso '\textit{onEditWardSubmit} non dovrebbe chiamare store se selectedWard è null': verifica che storeStub.\textit{updateWard} non viene invocato. Esito atteso: storeStub.\textit{updateWard} non viene invocato.}{S}
\addTestU{Contesto \textit{WardManagementPageComponent}; caso '\textit{onCloseWardDialog} dovrebbe chiudere dialog e pulire selezione': verifica che il valore di component.selectedWardId() è wardA.id e il valore di component.wardDialogMode() è 'closed'. Esito atteso: il valore di component.selectedWardId() è wardA.id; il valore di component.wardDialogMode() è 'closed'.}{S}
\addTestU{Contesto \textit{WardManagementPageComponent}; caso '\textit{onCloseAssignOperatorDialog} e \textit{onCloseAssignPlantDialog} ripuliscono stato di caricamento e lookup': verifica che component.operatorWardId() è null e component.availableOperatorsFromFetch() è null. Esito atteso: component.operatorWardId() è null; component.availableOperatorsFromFetch() è null; il valore di component.isLoadingAvailableOperators() è false; component.plantWardId() è null.}{S}
\addTestU{Contesto \textit{WardManagementPageComponent}; caso '\textit{selectWard} dovrebbe aggiornare ward attivo e step mobile': verifica che il valore di component.selectedWardId() è 2 e il valore di component.mobileStep() è 'apartments'. Esito atteso: il valore di component.selectedWardId() è 2; il valore di component.mobileStep() è 'apartments'.}{S}
\addTestU{Contesto \textit{WardManagementPageComponent}; caso '\textit{selectWard} sul ward già selezionato mantiene la selezione dell appartamento': verifica che il valore di component.selectedWardId() è wardA.id e il valore di component.selectedApartmentId() è '101'. Esito atteso: il valore di component.selectedWardId() è wardA.id; il valore di component.selectedApartmentId() è '101'; il valore di component.mobileStep() è 'apartments'.}{S}
\addTestU{Contesto \textit{WardManagementPageComponent}; caso '\textit{selectApartment} e step helpers dovrebbero aggiornare lo stato ui': verifica che il valore di component.selectedApartmentId() è '101' e il valore di component.mobileStep() è 'wards'. Esito atteso: il valore di component.selectedApartmentId() è '101'; il valore di component.mobileStep() è 'wards'; il valore di component.mobileStep() è 'apartments'; il valore di component.mobileStep() è 'operators'.}{S}
\addTestU{Contesto \textit{WardManagementPageComponent}; caso '\textit{onEditSelectedWard} non dovrebbe fare nulla senza ward selezionato': verifica che component.selectedWardId() è null e il valore di component.wardDialogMode() è 'closed'. Esito atteso: component.selectedWardId() è null; il valore di component.wardDialogMode() è 'closed'.}{S}
\addTestU{Contesto \textit{WardManagementPageComponent}; caso '\textit{onAssignOperatorSubmit} non dovrebbe chiamare store senza wardId': verifica che storeStub.\textit{assignOperator} non viene invocato. Esito atteso: storeStub.\textit{assignOperator} non viene invocato.}{S}
\addTestU{Contesto \textit{WardManagementPageComponent}; caso 'flow plant dovrebbe impostare wardId, submit e chiudere': verifica che il valore di component.plantWardId() è 2 e storeStub.\textit{getAvailablePlantsForWard} viene invocato con 2. Esito atteso: il valore di component.plantWardId() è 2; storeStub.\textit{getAvailablePlantsForWard} viene invocato con 2; component.availablePlants() coincide con [\{ id: '103', name: 'App. 103' \}]; storeStub.\textit{assignPlant} viene invocato con 2.}{S}
\addTestU{Contesto \textit{WardManagementPageComponent}; caso '\textit{onAssignPlantSubmit} non dovrebbe chiamare store senza wardId': verifica che storeStub.\textit{assignPlant} non viene invocato. Esito atteso: storeStub.\textit{assignPlant} non viene invocato.}{S}
\addTestU{Contesto \textit{WardManagementPageComponent}; caso 'confirmState + \textit{onConfirmDialogConfirmed} dovrebbero gestire i 3 rami': verifica che component.confirmState() coincide con \{ kind: 'delete-ward', wardId: 1 \} e storeStub.\textit{deleteWard} viene invocato con 1. Esito atteso: component.confirmState() coincide con \{ kind: 'delete-ward', wardId: 1 \}; storeStub.\textit{deleteWard} viene invocato con 1; component.confirmState() è null; component.confirmState() coincide con \{ kind: 'remove-operator', wardId: 1, userId: 2, \}.}{S}
\addTestU{Contesto \textit{WardManagementPageComponent}; caso '\textit{onConfirmDialogConfirmed} non dovrebbe fare nulla con stato nullo': verifica che storeStub.\textit{deleteWard} non viene invocato e storeStub.\textit{removeOperator} non viene invocato. Esito atteso: storeStub.\textit{deleteWard} non viene invocato; storeStub.\textit{removeOperator} non viene invocato; storeStub.\textit{removePlant} non viene invocato.}{S}
\addTestU{Contesto \textit{WardManagementPageComponent}; caso '\textit{onConfirmDialogCancelled} dovrebbe pulire confirmState': verifica che component.confirmState() è null. Esito atteso: component.confirmState() è null.}{S}
\addTestU{Contesto \textit{WardManagementPageComponent}; caso '\textit{getConfirmMessage} dovrebbe restituire il testo corretto per ogni stato': verifica che il valore di component.\textit{getConfirmMessage}() è 'Confermi questa operazione?' e il valore di component.\textit{getConfirmLabel}() è 'Conferma'. Esito atteso: il valore di component.\textit{getConfirmMessage}() è 'Confermi questa operazione?'; il valore di component.\textit{getConfirmLabel}() è 'Conferma'; il valore di component.\textit{getConfirmMessage}() è 'Confermi l\textbackslash\{\}'eliminazione del reparto?'; il valore di component.\textit{getConfirmMessage}() è 'Confermi la rimozione dell\textbackslash\{\}'operatore dal reparto?'.}{S}
\addTestU{Contesto \textit{WardManagementPageComponent}; caso '\textit{getOperatorDisplayName} dovrebbe comporre e trimmare nome e cognome': verifica che il valore di component.\textit{getOperatorDisplayName}('Mario', 'Rossi') è 'Mario Rossi' e il valore di component.\textit{getOperatorDisplayName}('Mario', '') è 'Mario'. Esito atteso: il valore di component.\textit{getOperatorDisplayName}('Mario', 'Rossi') è 'Mario Rossi'; il valore di component.\textit{getOperatorDisplayName}('Mario', '') è 'Mario'.}{S}
\addTestU{Contesto \textit{WardManagementPageComponent}; caso 'availablePlants dovrebbe deduplicare e escludere quelli già assegnati al ward selezionato': verifica che component.availablePlants() coincide con []. Esito atteso: component.availablePlants() coincide con [].}{S}
\addTestU{Contesto \textit{WardManagementPageComponent}; caso '\textit{onAssignPlant} non usa fallback locale se fetch ad-hoc fallisce': verifica che component.availablePlants() coincide con [] e il valore di component.isLoadingAvailablePlants() è false. Esito atteso: component.availablePlants() coincide con []; il valore di component.isLoadingAvailablePlants() è false.}{S}
\addTestU{Contesto \textit{WardRowComponent}; caso 'should create': verifica che component è valorizzato. Esito atteso: component è valorizzato.}{S}
\addTestU{Contesto \textit{WardRowComponent}; caso 'accetta plant nullo': verifica che component.plant() è null. Esito atteso: component.plant() è null.}{S}
\addTestU{Contesto \textit{WardRowComponent}; caso 'accetta plant valorizzato con id stringa': verifica che component.plant() coincide con \{ id: '101', name: 'App' \}. Esito atteso: component.plant() coincide con \{ id: '101', name: 'App' \}.}{S}

\subsubsection{Backend - View4Life}
\addTestU{Contesto \textit{Plot}; caso 'should be defined': verifica che un'istanza di Plot è valorizzata. Esito atteso: nuova istanza Plot è definita.}{S}
\addTestU{Contesto \textit{Plot}; caso 'should return correct title, metric, unit, labels e series': verifica che Plot restituisce le proprietà passate al costruttore. Esito atteso: getTitle(), getMetric(), getUnit(), getLabels(), getSeries() ritornano i valori corretti.}{S}
\addTestU{Contesto \textit{Plot}; caso 'should return undefined suggestion when not set': verifica che getSuggestion() è undefined prima di setSuggestion. Esito atteso: getSuggestion() ritorna undefined.}{S}
\addTestU{Contesto \textit{Plot}; caso 'should set and return suggestion': verifica che setSuggestion() aggiorna il valore e getSuggestion() lo ritorna. Esito atteso: dopo setSuggestion(), getSuggestion() ritorna il valore impostato.}{S}

\addTestU{Contesto \textit{ValidationPipe}; caso 'should be defined': verifica che il pipe è valorizzato. Esito atteso: ValidationPipe è definito.}{S}
\addTestU{Contesto \textit{ValidationPipe}; caso 'should return value if no metatype': verifica che il pipe passa il valore quando non c'è metatype. Esito atteso: transform ritorna il valore non modificato.}{S}
\addTestU{Contesto \textit{ValidationPipe}; caso 'should skip validation for primitive types': verifica che il pipe skip primitive types come String. Esito atteso: transform ritorna il valore per tipologie primitive.}{S}
\addTestU{Contesto \textit{ValidationPipe}; caso 'should pass validation for valid DTO': verifica che il pipe accetta DTO validi. Esito atteso: transform ritorna il DTO passato.}{S}
\addTestU{Contesto \textit{ValidationPipe}; caso 'should throw error for invalid DTO': verifica che il pipe rigetta DTO con proprietà non valide. Esito atteso: transform lancia BadRequestException.}{S}
\addTestU{Contesto \textit{ValidationPipe}; caso 'should throw error when required field is missing': verifica che il pipe rigetta DTO senza campi obbligatori. Esito atteso: transform lancia BadRequestException.}{S}

\addTestU{Contesto \textit{AlarmRulesService}; caso 'should return all alarm rules': verifica che il servizio ritorna la lista di tutte le regole. Esito atteso: getAllAlarmRules non è un array vuoto.}{S}
\addTestU{Contesto \textit{AlarmRulesService}; caso 'should return empty array when no rules exist': verifica che il servizio ritorna array vuoto quando non ci sono regole. Esito atteso: getAllAlarmRules ritorna un array vuoto.}{S}
\addTestU{Contesto \textit{AlarmRulesService}; caso 'should return alarm rule by id': verifica che il servizio ritorna una regola specifica. Esito atteso: getAlarmRuleById(id) ritorna la regola con id corretto.}{S}
\addTestU{Contesto \textit{AlarmRulesService}; caso 'should return null when rule not found': verifica che il servizio ritorna null quando regola non esiste. Esito atteso: getAlarmRuleById(invalid) ritorna null.}{S}
\addTestU{Contesto \textit{AlarmRulesService}; caso 'should create alarm rule successfully': verifica che il servizio crea una nuova regola. Esito atteso: createAlarmRule() aggiunge la regola al database.}{S}
\addTestU{Contesto \textit{AlarmRulesService}; caso 'should update alarm rule successfully': verifica che il servizio aggiorna una regola esistente. Esito atteso: updateAlarmRule() modifica le proprietà della regola.}{S}
\addTestU{Contesto \textit{AlarmRulesService}; caso 'should delete alarm rule successfully': verifica che il servizio cancella una regola. Esito atteso: deleteAlarmRule() rimuove la regola dal database.}{S}
\addTestU{Contesto \textit{AlarmRulesService}; caso 'should emit alarm.activated event when alarm matches': verifica che il servizio emette evento quando soglia è superata. Esito atteso: checkAlarmRule emette alarm.activated event.}{S}

\addTestU{Contesto \textit{AlarmEventsService}; caso 'should return alarm event by id': verifica che il servizio ritorna un evento specifico. Esito atteso: getAlarmEventById(id) ritorna l'evento con id corretto.}{S}
\addTestU{Contesto \textit{AlarmEventsService}; caso 'should return all alarm events': verifica che il servizio ritorna lista di tutti gli eventi. Esito atteso: getAllAlarmEvents non è un array vuoto.}{S}
\addTestU{Contesto \textit{AlarmEventsService}; caso 'should return managed alarm events by user': verifica che il servizio filtra gli allarmi risolti dell'utente. Esito atteso: getManagedAlarmEventsByUserId() ritorna solo allarmi con resolution\_time.}{S}
\addTestU{Contesto \textit{AlarmEventsService}; caso 'should return unmanaged alarm events by user': verifica che il servizio filtra gli allarmi attivi dell'utente. Esito atteso: getUnmanagedAlarmEventsByUserId() ritorna solo allarmi con resolution\_time null.}{S}
\addTestU{Contesto \textit{AlarmEventsService}; caso 'should resolve alarm event and emit event': verifica che il servizio risolve un allarme e emette evento. Esito atteso: resolveAlarmEvent() aggiorna resolution\_time e emette alarm.resolved.}{S}

\addTestU{Contesto \textit{WardDomain}; caso 'should be defined': verifica che Ward può essere istanziato. Esito atteso: new Ward(1, 'Neuro') è definito.}{S}

\addTestU{Contesto \textit{AnalyticsStrategy}; caso 'plant-anomalies returns plot with correct data': verifica che la strategia di anomalie assume correttamente i dati. Esito atteso: execute() ritorna Plot con serie e label allineate.}{S}
\addTestU{Contesto \textit{AnalyticsStrategy}; caso 'sensor-long-presence counts sessions over 30 minutes': verifica che la strategia conta sessioni prolungate > 30 min. Esito atteso: execute() conta solo una sessione per evento continuativo.}{S}
\addTestU{Contesto \textit{AnalyticsStrategy}; caso 'ward-resolved-alarm aggregates total and resolved': verifica che la strategia aggrega allarmi totali e risolti. Esito atteso: execute() ritorna due serie: una totale, una risolta.}{S}

\makeTUTable


\section{Cruscotto di qualità}
\subsection{MPC1 e MPC2 - Planned Value e Earned Value}
\begin{figure}[H]
    \centering
    \includegraphics[width=1\textwidth]{./grafici/MPC1_2.pdf}
    \caption{Andamento PV e EV in base al periodo}
    \label{fig:mpc1-mpc2}
\end{figure}

L'andamento del grafico denota un progressivo aumento del lavoro completato. 
Fino allo Sprint 3 rimane al di sotto della percentuale prevista.
Grazie ad una successiva ripianificazione la percentuale di lavoro svolto è di poco al di sopra del pianificato.
Durante i periodi degli Sprint 8 e 9, il valore di lavoro completato è ritornato leggermente al di sotto del preventivo, a causa di alcune sovrastime nella pianificazione.
Allo Sprint 11 entrambe le curve hanno raggiunto il Budget At Completion, confermando il completamento del progetto nei tempi e nei costi previsti.

\subsection{MPC3 e MPC7 - Actual Cost e Estimate To Complete}
\begin{figure}[H]
    \centering
    \includegraphics[width=1\textwidth]{./grafici/MPC3_7.pdf}
    \caption{Andamento AC e ETC in base al periodo}
    \label{fig:mpc3-mpc7}
\end{figure}

Il grafico mostra che il budget impiegato è coerente con la percentuale di lavoro completata, senza evidenziare particolari anomalie nell'andamento.
Allo Sprint 11 l'Estimate To Complete si è azzerato e il costo consuntivato finale si è attestato leggermente al di sotto del BAC, indicando una gestione complessivamente efficiente delle risorse economiche.

\subsection{MPC4 e MPC5 - Cost Performance Index e Schedule Performance Index}
\begin{figure}[H]
    \centering
    \includegraphics[width=1\textwidth]{./grafici/MPC4_5.pdf}
    \caption{Andamento CPI e SPI in base al periodo}
    \label{fig:mpc4-mpc5}
\end{figure}

Dal grafico si desume che, nella fase iniziale, il gruppo ha gestito correttamente il budget, sebbene le attività risultassero significativamente in ritardo. 
Successivamente, la situazione si è stabilizzata su valori superiori alla soglia minima accettabile fino a circa il quarto sprint, momento in cui il trend si è invertito: le attività sono state completate in anticipo rispetto alla pianificazione, mentre il budget impiegato ha registrato un lieve scostamento in eccesso rispetto alle previsioni.
A seguito della ripianificazione effettuata al termine dello Sprint 7, si nota nei periodi successivi una riduzione e stabilizzazione del budget speso (CPI), perdendo però qualche punto nell'SPI che rimane comunque accettabile.
Allo Sprint 11 il progetto si è concluso con CPI e SPI entrambi prossimi a 1.00, confermando il rispetto dei tempi e dei costi preventivati.

\subsection{MPC6 - Estimate At Completion}
\begin{figure}[H]
    \centering
    \includegraphics[width=1\textwidth]{./grafici/MPC6_7.pdf}
    \caption{Andamento EAC in base al periodo}
    \label{fig:mpc6}
\end{figure}

Il grafico rappresenta quanto riportato nella figura \ref{fig:mpc1-mpc2} (MPC1 e MPC2), ovvero una buona gestione del budget iniziale, con una previsione del costo al completamento in linea con il preventivo iniziale.
Al termine dello Sprint 5, momento in cui lo sforamento di budget ha raggiunto il valore massimo, è stata effettuata una lieve ripianificazione che ha contribuito alla riduzione dello scostamento dai costi previsti.
Dallo Sprint 7 allo Sprint 8 il costo stimato al completamento è rimasto al di sotto del preventivo iniziale, subendo un leggero aumento durante lo Sprint 9.
Allo Sprint 10 l'EAC si è attestato sensibilmente al di sotto del BAC, confermando una gestione del budget più efficiente precedentemente prevista. Allo Sprint 11 l'EAC si è ulteriormente ridotto, chiudendo il progetto con un costo finale inferiore al preventivo iniziale e senza sforamenti significativi.

\subsection{MPC8 - Requirements Stability Index}
\begin{figure}[H]
    \centering
    \includegraphics[width=1\textwidth]{./grafici/MPC8.pdf}
    \caption{Andamento RSI in base al periodo}
    \label{fig:mpc8}
\end{figure}

Il crollo della stabilità nello Sprint 3 è dovuto ad un iniziale errore nella raccolta dei requisiti che ha portato inevitabilmente a rivedere quasi tutti quelli individuati fino a quel momento.
Dallo Sprint 4 in poi i requisiti hanno avuto una stabilità pressoché ottima. 
Allo Sprint 11 l'RSI ha confermato il valore ottimo del 100\%, attestando che la base dei requisiti si è stabilizzata completamente senza ulteriori variazioni significative.

\subsection{MPC9 - Indice di Gulpease}
\begin{figure}[H]
    \centering
    \includegraphics[width=1\textwidth]{./grafici/MPC9.pdf}
    \caption{Andamento IG dei documenti principali, in base al periodo}
    \label{fig:mpc9}
\end{figure}

Durante tutti gli sprint la leggibilità dei documenti è rimasta buona e al di sopra della soglia di accettabilità. 
Allo Sprint 10 si è registrato un calo dell'indice di Gulpease per il Glossario e le Norme di Progetto, dovuto all'aggiunta di contenuti correlati alla progettazione e alla codifica.
Entrambi i documenti rimangono comunque al di sopra della soglia di accettabilità. Allo Sprint 11 il Piano di Qualifica e il Manuale Utente hanno registrato un incremento, raggiungendo entrambi valori superiori al valore ottimo, mentre la Specifica Tecnica ha anch'essa raggiunto un indice prossimo al valore ottimo.

\subsection{MPC10 - Indice di Frammentazione}
\begin{figure}[H]
    \centering
    \includegraphics[width=1\textwidth]{./grafici/MPC10.pdf}
    \caption{Andamento IF dei documenti principali, in base al periodo}
    \label{fig:mpc10}
\end{figure}

Durante tutti gli sprint la struttura dei documenti è rimasta buona, al di sopra della soglia di accettabilità.
Si rende evidente il fatto che il documento ``Glossario'' non è stato incluso nel controllo della metrica MPC10: è stata fatta questa scelta poiché quest'ultimo, essendo un documento per definizione molto frammentato (ad ogni termine viene associata una breve spiegazione), avrebbe ottenuto un valore superiore al 20\%. 
La verifica avrebbe dunque evidenziato un problema nel Glossario quando, in realtà, non è presente.
Allo Sprint 11 tutti i documenti monitorati hanno confermato valori stabili e coerenti con il range ottimo, attestandosi tra il 12\% e il 15\%.

\subsection{MPC12 - Percentuale Metriche Soddisfatte}
\begin{figure}[H]
    \centering
    \includegraphics[width=1\textwidth]{./grafici/MPC12.pdf}
    \caption{Andamento PMS in base al periodo}
    \label{fig:mpc12}
\end{figure}

Dal grafico è evidente l'impegno del gruppo nei confronti delle metriche di qualità di processo e del loro soddisfacimento.
Un lieve peggioramento è stato notato durante lo Sprint 4, soprattutto a causa dei rischi inattesi e dell'instabilità dei requisiti durante tale periodo.
Dopo averne preso consapevolezza, il gruppo ha adottato opportune azioni correttive, riuscendo a migliorare progressivamente le performance negli sprint successivi.
Il progetto si è concluso allo Sprint 11 con il 100\% delle metriche soddisfatte, attestando il pieno raggiungimento degli obiettivi di qualità di processo.

\subsection{MPC13 - Percentuale Rischi Inattesi}
\begin{figure}[H]
    \centering
    \includegraphics[width=1\textwidth]{./grafici/MPC13.pdf}
    \caption{Andamento PRI in base al periodo}
    \label{fig:mpc13}
\end{figure}

La percentuale rispetto al totale dei rischi inattesi è superiore al valore massimo accettabile nella maggior parte dei periodi.
Testimonia una capacità ancora non matura di previsione dei rischi da parte del gruppo.
Si può notare però che agli Sprint 6, 7 e 8 è diminuita, subendo un rialzo allo Sprint 9. 
Allo Sprint 10 la percentuale di rischi inattesi si è mantenuta al di sotto del valore massimo accettabile, confermando il miglioramento nella capacità di previsione dei rischi da parte del gruppo.
Allo Sprint 11 la percentuale si è azzerata, raggiungendo il valore ottimo e confermando la piena maturità del gruppo nella gestione e previsione dei rischi al termine del progetto.

\subsection{MPC14 - Labor Efficiency}
\begin{figure}[H]
    \centering
    \includegraphics[width=1\textwidth]{./grafici/MPC14.pdf}
    \caption{Andamento LE in base al periodo}
    \label{fig:mpc14}
\end{figure}

Dal grafico è possibile notare, fin dal primo Sprint, un'eccellente efficienza nel lavoro svolto, con un picco importante durante lo Sprint 6 e 8.
È opportuno sottolineare che tale andamento non indica un'eccessiva efficienza operativa del gruppo, bensì evidenzia che le ore preventivate sono risultate superiori rispetto a quelle effettivamente impiegate. 
Allo Sprint 10 si è registrato un calo della Labor Efficiency al di sotto del valore ottimo, dovuto a una stima delle ore più accurata rispetto ai periodi precedenti: le ore effettivamente impiegate si sono avvicinate a quelle preventivate, riducendo il surplus che aveva caratterizzato gli sprint precedenti.
Allo Sprint 11 la Labor Efficiency si è mantenuta al di sopra della soglia minima accettabile, confermando che il gruppo ha concluso il progetto con un utilizzo delle risorse coerente con quanto preventivato.

\end{document}
